% Canzoniere di Clan - template. I segnaposto tra doppi angoli (maiuscoli)
% sono sostituiti da scripts/generate_latex.py. Compilare con pdflatex (due passate) o latexmk.
\documentclass[a5paper,landscape,10pt,twoside]{article}

\usepackage{iftex}
\ifPDFTeX
  \usepackage[utf8]{inputenc}
  \usepackage[T1]{fontenc}
  \usepackage{lmodern}
\else % XeLaTeX/LuaLaTeX (es. Tectonic in locale): stesso carattere
  \usepackage{fontspec} % carattere predefinito: Latin Modern
\fi
\usepackage[italian]{babel}
\usepackage[a5paper,landscape,twoside,
  left=15mm,right=15mm,top=12mm,bottom=12mm,bindingoffset=15mm,
  includefoot,nohead,footskip=6mm]{geometry}
\usepackage{multicol}
\usepackage{graphicx}
\usepackage{xcolor}
\usepackage{url}
\usepackage{microtype}
\usepackage{needspace}

\setlength{\parindent}{0pt}
\setlength{\parskip}{0pt}
\setlength{\columnsep}{8mm}
\raggedbottom
\sloppy
\hyphenpenalty=1000

\definecolor{accordo}{gray}{0.25}
\definecolor{brLC}{HTML}{F2C230}
\definecolor{brEG}{HTML}{3F9A47}
\definecolor{brRS}{HTML}{C0392B}
\definecolor{fiamma}{HTML}{D9822B}

% Interlinea ampia: lascia spazio agli accordi sopra il testo.
\makeatletter
\renewcommand{\normalsize}{\@setfontsize\normalsize{10}{17}}

% Stile pagina: numero in basso sul lato esterno, nome della sezione sul lato interno.
\def\sezionecorrente{}
\def\ps@canz{%
  \def\@oddhead{}\def\@evenhead{}%
  \def\@oddfoot{{\footnotesize\scshape\color{gray}\sezionecorrente}\hfill\small\thepage}%
  \def\@evenfoot{\small\thepage\hfill{\footnotesize\scshape\color{gray}\sezionecorrente}}}

% \cantotesta{id}{titolo}{categoria}{autore}: nuova pagina + etichetta per gli indici
% \fiamme{n}: intensita' 1-5 come cinque barrette piene/vuote
\newcommand{\barra}{\textcolor{fiamma}{\rule{1.1mm}{2.6mm}}\,}
\newcommand{\barraoff}{\textcolor{black!20}{\rule{1.1mm}{2.6mm}}\,}
\newcommand{\fiamme}[1]{{\@tempcnta=\z@
  \@whilenum\@tempcnta<5\do{\advance\@tempcnta\@ne
    \ifnum\@tempcnta>#1\barraoff\else\barra\fi}}}
% Tag di branca: solo un quadratino colorato (giallo LC, verde EG, rosso RS)
\newcommand{\tagbranca}[1]{\textcolor{#1}{\rule{2.4mm}{2.4mm}}\,}
\newcommand{\bLC}{\tagbranca{brLC}}
\newcommand{\bEG}{\tagbranca{brEG}}
\newcommand{\bRS}{\tagbranca{brRS}}

% I canti scorrono di seguito su due colonne. \cantotesta{id}{titolo}{autore}{tag branche}{altezza}:
% se nella colonna non resta l'altezza stimata del canto si passa alla colonna successiva: i canti
% corti non si spezzano. Quelli oltre \sogliaspezza di colonna possono continuare nella colonna dopo
% (basta che ci stia 1/3 di colonna).
\newcommand{\sogliaspezza}{.6}
\newcommand{\cantotesta}[5]{%
  \par\addvspace{4.5mm}%
  \ifdim #5>\sogliaspezza\textheight \needspace{.33\textheight}\else\needspace{#5}\fi
  \def\@currentlabel{#1}\label{canto:#1}%
  \noindent\hangindent=12.5mm\hangafter=1
  \hbox to 12.5mm{\fontsize{17}{19}\selectfont\bfseries #1\hss}%
  {\fontsize{13}{15.5}\selectfont\bfseries #2\par}\nopagebreak
  \noindent{\small\itshape #3}\hfill{#4}\par\nopagebreak
  \vspace{.6mm}\hrule height 0.4pt\vspace{2mm}\nopagebreak}
% Canto senza accordi: interlinea normale (quella ampia serve solo a far stare gli accordi sopra il testo)
\newcommand{\senzaaccordi}{\fontsize{10}{12.6}\selectfont}

% Pagina divisoria di sezione: titolo grande + testo facoltativo (latex/sezioni/<sezione>.tex)
\newcommand{\sezione}[2]{%
  \clearpage\gdef\sezionecorrente{#1}%
  \thispagestyle{empty}
  \vspace*{\fill}
  \begin{center}
    {\fontsize{30}{34}\selectfont\bfseries #1\par}
    \vspace{3mm}\rule{40mm}{0.6pt}\par
  \end{center}
  \vspace{6mm}{\small #2\par}
  \vspace*{\fill}
  \clearpage}
% Titoletto di fase del fuoco / momento della messa, dentro la colonna
\newcommand{\sottosezione}[1]{%
  \par\addvspace{5mm}\needspace{.25\textheight}%
  {\centering\color{fiamma}\rule[.6ex]{8mm}{.6pt}\ \textcolor{black}{\scshape\bfseries #1}\ \rule[.6ex]{8mm}{.6pt}\par}
  \nopagebreak\vspace{1mm}}
\makeatother
\pagestyle{canz}

% Riga di testo con accordi: l'accordo sta sopra la sillaba che lo segue.
\newcommand{\ch}[2]{\leavevmode\vbox{\offinterlineskip
  \hbox{\fontsize{8}{9}\selectfont\bfseries\color{accordo}#1\kern.3em}%
  \kern.45ex
  \hbox{#2}}}
\newcommand{\rigo}[1]{{\hangindent=3mm\hangafter=1\noindent #1\par}}
\newcommand{\strofa}{\par\vspace{2.4mm}}
\newcommand{\commento}[1]{{\par\small\itshape\color{gray}#1\par\vspace{1mm}}}
\newenvironment{ritornello}{\par\begingroup\itshape\leftskip=3mm\relax}{\par\endgroup}

% Indici
\newcommand{\titolosezione}[1]{\clearpage\gdef\sezionecorrente{#1}%
  {\noindent\fontsize{16}{18}\selectfont\bfseries #1\par}
  \vspace{1mm}\hrule height 0.4pt\vspace{3mm}}
\newcommand{\lettera}[1]{\par\vspace{1.5mm}{\bfseries\color{accordo}#1}\par\nopagebreak}
\newcommand{\categoriaindice}[1]{\par\vspace{2mm}{\scshape\bfseries #1}\par\nopagebreak}
\newcommand{\voce}[2]{\noindent\hangindent=2.4em\hangafter=1
  \hbox to 2em{\hfil\bfseries #1}\hskip .4em #2\nobreak
  \leaders\hbox to .7em{\hss.\hss}\hfill\nobreak
  \hbox to 2em{\hfill\pageref{canto:#1}}\par}

\newcommand{\sottocategoriaindice}[1]{\par\vspace{.8mm}{\small\itshape #1}\par\nopagebreak}

\begin{document}

% ---- Frontespizio (numerazione romana, nessun numero stampato)
\pagenumbering{roman}
\thispagestyle{empty}
\begin{center}
  \vspace*{\fill}
  {\fontsize{40}{44}\selectfont\bfseries Canzoniere\par}
  \vspace{4mm}
  {\Large Clan - Gruppo Scout\par}
  \vspace{10mm}
  {\large\itshape Un regalo di Partenza, per i miei compagni di Clan.\par}
  \vspace*{\fill}
  {\small 2026\par}
\end{center}
% Pagina di dedica. BOZZA da personalizzare: sostituisci il testo e la firma.
\clearpage
\thispagestyle{empty}
\vspace*{\fill}
\begin{center}
  \begin{minipage}{0.8\textwidth}
    \centering\itshape\large
    Ai miei compagni di Clan.\par
    \vspace{5mm}
    \normalsize
    Ho camminato con voi per tanta strada, e quasi ogni passo l'abbiamo cantato.
    Questo libro raccoglie le canzoni che ci hanno tenuti insieme: quelle gridate
    in marcia, quelle sussurrate attorno al fuoco, quelle che non abbiamo mai
    finito di imparare.\par
    \vspace{3mm}
    Portatelo con voi, lasciate che si consumi, e quando io non ci sar\`o
    cantatene una anche per me.\par
    \vspace{6mm}
    Buona strada.\par
    \vspace{3mm}
    \upshape [Nome e totem]
  \end{minipage}
\end{center}
\vspace*{\fill}


\titolosezione{Come usare questo canzoniere}
\begingroup\small
\begin{multicols}{2}
\textbf{Struttura.} I canti sono numerati e raggruppati in sezioni, nell'ordine:
Canti associativi, Canti cerimoniali, Canti di lupettismo, Fuoco di bivacco, Canti alpini, Canti liturgici, Canti del risveglio. Ogni sezione si apre con una pagina divisoria; il numero grande accanto
al titolo è quello da cercare, sia qui sia nel sito. In fondo trovi l'indice alfabetico
e l'indice per sezione.

\medskip
\textbf{Accordi.} Sono scritti sopra la sillaba in cui cambiano. Il testo in
\textit{corsivo} è il ritornello.

\medskip
\textbf{Online.} Sull'ultima pagina c'è un QR code: porta al sito, con gli spartiti in PDF
(dove disponibili) e la modalità notte per il bivacco.

\columnbreak
\textbf{Branche.} A destra del titolo, le branche per cui il canto nasce o è
più adatto:

\smallskip
\bLC\ lupetti \quad \bEG\ esploratori \quad \bRS\ rover

\medskip
\textbf{Fuoco di bivacco.} I canti sono divisi per fase, dall'apertura agli ultimi:
la pagina all'inizio della sezione spiega come si succedono.

\medskip
Un canto senza tag è per tutti.
\end{multicols}
\endgroup


\titolosezione{Cantare insieme}
\begingroup\small
\begin{multicols}{2}
\textbf{1. Il tono.} Chi attacca sceglie una nota che stia comoda a tutti: meglio un po' bassa che
troppo alta. Se la prima strofa è faticosa, si riparte da capo, più in basso.

\smallskip
\textbf{2. Un solo attacco.} Si dà il tempo con un respiro o con la mano e si parte insieme.
Cantare a ritmi diversi è il modo più rapido per spegnere un canto.

\smallskip
\textbf{3. Ascoltare.} Si canta in modo da sentire il vicino. Se non lo senti più, stai coprendo
la sua voce.

\smallskip
\textbf{4. Il respiro.} Si respira alla fine delle frasi, non a metà parola. Chi respira in
momenti diversi tiene su il canto senza buchi.

\columnbreak
\textbf{5. Canoni e voci.} Nei canoni ogni gruppo tiene la propria parte e non segue quello
accanto. Si comincia a due voci, poi si aggiungono le altre.

\smallskip
\textbf{6. Bans.} Si insegnano un pezzo alla volta: prima i gesti, poi le parole, poi tutto
insieme. Chi guida sta davanti e li fa per primo.

\smallskip
\textbf{7. Chitarra.} Sta sotto le voci, non sopra. Un giro semplice e regolare vale più di uno
elaborato; con il capotasto si abbassa il tono senza cambiare accordi.

\smallskip
\textbf{8. Senza strumenti.} Un canto deve reggere anche a cappella: se si perde l'intonazione,
ci si ferma, si ritrova la nota e si riparte.
\end{multicols}
\endgroup


% ---- Canti (la numerazione araba riparte da 1: pagina stampata = campo "pagina")
\pagenumbering{arabic}
\sezione{Canti associativi}{}

\begin{multicols}{2}\raggedcolumns

\begingroup
\cantotesta{1}{Inno d'Italia}{G. Mameli - M. Novaro}{}{256pt}
\rigo{\ch{C}{Fratelli} d'Italia,}
\rigo{l'Italia s'è \ch{G}{desta,}}
\rigo{dell'elmo di \ch{C}{Scipio}}
\rigo{s'è cinta la \ch{G}{testa.}}
\rigo{Dov'è la Vit\ch{C}{toria?}}
\rigo{Le porga la \ch{G}{chioma,}}
\rigo{ché schiava di \ch{C}{Roma}}
\rigo{Iddio la cre\ch{G}{ò.}}
\strofa
\begin{ritornello}
\rigo{\ch{C}{Stringiamci} a coorte,}
\rigo{siam \ch{G}{pronti} alla morte,}
\rigo{siam \ch{C}{pronti} alla morte,}
\rigo{l'Italia chia\ch{G}{mò.} \ch{C}{Sì!}}
\end{ritornello}
\par\endgroup

\begingroup
\cantotesta{2}{Inno FSE}{}{}{580pt}
\begin{ritornello}
\rigo{Sulla \ch{Re}{strada} la \ch{Sol}{gran} \ch{Re}{strada}}
\rigo{che si snoda \ch{La}{intorno} al nostro vecchio \ch{Re}{mondo,}}
\rigo{io ti ascolto, \ch{Sol}{tu} mi a\ch{Re}{scolti,}}
\rigo{Scout d'Europa e cam\ch{La}{mino} con \ch{Re}{te.}}
\strofa
\end{ritornello}
\rigo{\ch{Sol}{Come} te \ch{Re}{presi} an\ch{Mi}{ch'io} del de\ch{La}{serto} il cammin,}
\rigo{quando in \ch{Mi}{cielo} spuntavan le \ch{La}{stelle:}}
\rigo{ai \ch{Mi}{bivacchi} \ch{La}{del} campo scoprimmo}
\rigo{\ch{Mi}{un} altro universo in pieno \ch{La}{splendor.}}
\strofa
\begin{ritornello}
\rigo{Sulla \ch{Re}{strada} la \ch{Sol}{gran} \ch{Re}{strada}}
\rigo{che si snoda \ch{La}{intorno} al nostro vecchio \ch{Re}{mondo,}}
\rigo{io ti ascolto, \ch{Sol}{tu} mi a\ch{Re}{scolti,}}
\rigo{Scout d'Europa e cam\ch{La}{mino} con \ch{Re}{te.}}
\strofa
\end{ritornello}
\rigo{Non ti chiedo il paese né ti chiedo chi sei tu}
\rigo{ma soltanto ti chiamo "fratello";}
\rigo{L'ideale e la legge che ci hanno uniti qui}
\rigo{son più forti dell'avversità.}
\strofa
\begin{ritornello}
\rigo{Sulla \ch{Re}{strada} la \ch{Sol}{gran} \ch{Re}{strada}}
\rigo{che si snoda \ch{La}{intorno} al nostro vecchio \ch{Re}{mondo,}}
\rigo{io ti ascolto, \ch{Sol}{tu} mi a\ch{Re}{scolti,}}
\rigo{Scout d'Europa e cam\ch{La}{mino} con \ch{Re}{te.}}
\strofa
\end{ritornello}
\rigo{L'avvenir può cambiar ma noi, no, non cambierem:}
\rigo{la promessa è per tutta la vita;}
\rigo{il ricordo di tanti bei giorni sarà}
\rigo{la certezza di questo ideal.}
\strofa
\begin{ritornello}
\rigo{Sulla \ch{Re}{strada} la \ch{Sol}{gran} \ch{Re}{strada}}
\rigo{che si snoda \ch{La}{intorno} al nostro vecchio \ch{Re}{mondo,}}
\rigo{io ti ascolto, \ch{Sol}{tu} mi a\ch{Re}{scolti,}}
\rigo{Scout d'Europa e cam\ch{La}{mino} con \ch{Re}{te.}}
\end{ritornello}
\par\endgroup

\begingroup\senzaaccordi
\cantotesta{3}{Canto di gruppo}{}{}{368pt}
\rigo{Ecco levarsi all'alba di ogni dì}
\rigo{fiero e gioioso il nostro canto,}
\rigo{misurando il passo del fratel}
\rigo{camminiamo fianco a fianco,}
\rigo{noi farem di ogni caduta}
\rigo{uno sprone per la meta,}
\rigo{se quel seme che è nei nostri cuor}
\rigo{non s'accresce frutta e amore,}
\rigo{non ha vita l'acqua del ruscel}
\rigo{non riscalda e asciuga il sole.}
\strofa
\begin{ritornello}
\rigo{Or che siamo qui dall'alto di una vetta,}
\rigo{comprendiamo assieme il valor della salita.}
\strofa
\end{ritornello}
\rigo{Il sorriso lacrime e il sudor}
\rigo{noi gli offriamo nel silenzio,}
\rigo{tutti i doni delle nostre man}
\rigo{li poniamo nel servire,}
\rigo{siamo giovani e bambini}
\rigo{luce e sale della terra,}
\rigo{testimoni di fraternità}
\rigo{costruiamo un buon domani,}
\rigo{cittadini per l'umanità}
\rigo{ci battiamo per la pace.}
\strofa
\begin{ritornello}
\rigo{Or che siamo qui dall'alto di una vetta,}
\rigo{ci affidiamo a lui per vivere la promessa.}
\end{ritornello}
\par\endgroup

\begingroup\senzaaccordi
\cantotesta{4}{Canto Riparto TV7}{Francesco Piazza}{}{268pt}
\rigo{Là sui monti sereni}
\rigo{nelle valli silenti}
\rigo{presso i bianchi torrenti}
\rigo{vive l'esplorator.}
\strofa
\rigo{Lui solo sa}
\rigo{del bosco la parola}
\rigo{dei prati in fior, dell'alba d'or}
\rigo{sa la serenità.}
\strofa
\rigo{Così seguir}
\rigo{la giusta traccia vuole}
\rigo{sotto la pioggia e il sole}
\rigo{la sua meta poter servir.}
\strofa
\rigo{Un richiamo risuona}
\rigo{su rispondi fratello}
\rigo{presta ascolto all'appello:}
\rigo{la stessa traccia tu puoi seguir.}
\par\endgroup

\end{multicols}

\sezione{Canti cerimoniali}{}

\begin{multicols}{2}\raggedcolumns

\begingroup\senzaaccordi
\cantotesta{5}{Addio lupetto}{}{\bLC}{268pt}
\rigo{Formiamo un grande cerchio}
\rigo{con le mani nelle man;}
\rigo{fratelli c'è un lupetto}
\rigo{che oggi il branco lascerà.}
\strofa
\rigo{Rit. Ma noi vogliamo dirgli ancor}
\rigo{che un dì ci rivedrem.}
\rigo{Arrivederci, fratellin,}
\rigo{presto ti seguirem.}
\strofa
\rigo{Tu pure come Mowgli}
\rigo{lasci il branco e vai lontan,}
\rigo{tu pure come Mowgli}
\rigo{sarai uomo nel doman.}
\strofa
\rigo{Ricordati del branco}
\rigo{noi di te fieri sarem}
\rigo{nessun ti scorderà,}
\rigo{fratello e un dì ti seguirem.}
\par\endgroup

\begingroup\senzaaccordi
\cantotesta{6}{Canto della promessa del lupetto}{}{\bLC}{382pt}
\rigo{Davanti a Te m'impegno}
\rigo{con tutto il cuor}
\rigo{diventerò tuo paggio}
\rigo{o mio Signor.}
\strofa
\rigo{Io voglio amarti sempre}
\rigo{sempre più}
\rigo{e la promessa mia}
\rigo{proteggi Tu.}
\strofa
\rigo{Voglio esser Tuo lupetto}
\rigo{per divenir}
\rigo{più tardi esploratore}
\rigo{per servir.}
\strofa
\rigo{Io voglio amarti sempre}
\rigo{sempre più}
\rigo{e la promessa mia}
\rigo{proteggi Tu.}
\strofa
\rigo{A Te Gesù che amo}
\rigo{obbedirò}
\rigo{a me dimenticando}
\rigo{ti ascolterò}
\strofa
\rigo{Io voglio amarti sempre}
\rigo{sempre di più}
\rigo{e la promessa mia}
\rigo{proteggi Tu.}
\par\endgroup

\begingroup
\cantotesta{7}{Canto della promessa}{J. Sevin - tradizionale}{\bEG}{366pt}
\rigo{\ch{Re}{Dinanzi} a voi m'impegno sul mio onor   \ch{La}{}}
\rigo{e vog\ch{Re}{lio} esserne degno per te o \ch{La}{Signor}         \ch{Re}{}}
\strofa
\commento{Rit.}
\rigo{La g\ch{Sol}{iusta} e retta vi\ch{Re}{a,} mostrami\ch{La7}{} Tu             \ch{Re7}{}}
\rigo{\ch{Sol}{e} la promessa m\ch{Re}{ia} acco\ch{La7}{gli} o Gesù.         \ch{Re}{}}
\strofa
\rigo{Fedele al tuo volere sempre sarò,}
\rigo{di Patria il mio dovere adempirò.}
\strofa
\commento{Rit.}
\strofa
\rigo{Apostolo tuo sono per tuo amor,}
\rigo{agli altri di me dono vo’ fare ognor.}
\strofa
\commento{Rit.}
\strofa
\rigo{Leale alla mia Legge sempre sarò,}
\rigo{se la tua man mi regge io manterrò!}
\strofa
\commento{Rit.}
\par\endgroup

\begingroup
\cantotesta{8}{Estote parati}{}{\bEG}{985pt}
\rigo{\ch{Re}{}\ch{La}{}\ch{Mi}{}\ch{La}{L’avventura} è una sfida nata dentro di te}
\rigo{\ch{Re}{}\ch{La}{}\ch{Mi}{}\ch{La}{tu} la senti spuntare, controllare non puoi}
\rigo{\ch{Fa\#m}{}\ch{Do\#m}{}\ch{Re}{}\ch{La}{senti} solo la voglia di mollar tutto e partir}
\rigo{\ch{Fa\#m}{}\ch{Do\#m}{}\ch{Re}{}\ch{Mi}{poi} le grida di amici che chiamano te.}
\strofa
\begin{ritornello}
\rigo{\ch{Fa\#m}{}\ch{Do\#m}{}\ch{Re}{}\ch{Mi}{Estote} Parati un grido s’alzerà}
\rigo{\ch{Fa\#m}{}\ch{Do\#m}{}\ch{Re}{}\ch{Mi}{e} mille voci a far da eco ad una voce fioca ormai}
\rigo{\ch{Fa\#m}{}\ch{Do\#m}{}\ch{Re}{}\ch{Mi}{e} allora dai, vieni con noi, è un’avventura in mare aperto}
\rigo{\ch{Fa\#m}{}\ch{Do\#m}{}\ch{Re}{}\ch{Mi}{e} viaggerai insieme a noi nella natura controvento.}
\strofa
\end{ritornello}
\rigo{Da mille paesi siamo giunti fino a qui,}
\rigo{anche se quegli amici mi dicevano di}
\rigo{non dar retta a quella gente che s’inventa fantasie,}
\rigo{lascia tutti vieni via, io ti venderò la mia.}
\strofa
\begin{ritornello}
\rigo{\ch{Fa\#m}{}\ch{Do\#m}{}\ch{Re}{}\ch{Mi}{Estote} Parati un grido s’alzerà}
\rigo{\ch{Fa\#m}{}\ch{Do\#m}{}\ch{Re}{}\ch{Mi}{e} mille voci a far da eco ad una voce fioca ormai}
\rigo{\ch{Fa\#m}{}\ch{Do\#m}{}\ch{Re}{}\ch{Mi}{e} allora dai, vieni con noi, è un’avventura in mare aperto}
\rigo{\ch{Fa\#m}{}\ch{Do\#m}{}\ch{Re}{}\ch{Mi}{e} viaggerai insieme a noi nella natura controvento.}
\strofa
\end{ritornello}
\rigo{Tutti insieme intorno al fuoco abbiam scoperto l’amicizia,}
\rigo{la pioggia sferza il viso ma non scalfisce la tua forza,}
\rigo{no non devi dare retta a chi deride l’avventura}
\rigo{lo scoutismo è la scoperta di una vita più matura.}
\strofa
\begin{ritornello}
\rigo{\ch{Fa\#m}{}\ch{Do\#m}{}\ch{Re}{}\ch{Mi}{Estote} Parati un grido s’alzerà}
\rigo{\ch{Fa\#m}{}\ch{Do\#m}{}\ch{Re}{}\ch{Mi}{e} mille voci a far da eco ad una voce fioca ormai}
\rigo{\ch{Fa\#m}{}\ch{Do\#m}{}\ch{Re}{}\ch{Mi}{e} allora dai, vieni con noi, è un’avventura in mare aperto}
\rigo{\ch{Fa\#m}{}\ch{Do\#m}{}\ch{Re}{}\ch{Mi}{e} viaggerai insieme a noi nella natura controvento.}
\strofa
\end{ritornello}
\rigo{Questo sole che dà luce e scalda il volo dei gabbiani}
\rigo{ti dice non dormire, oggi hai il mondo nelle mani}
\rigo{non lasciarti trascinare dalle facili correnti}
\rigo{guida tu la tua canoa fra gli scogli dei torrenti.}
\strofa
\begin{ritornello}
\rigo{\ch{Fa\#m}{}\ch{Do\#m}{}\ch{Re}{}\ch{Mi}{Estote} Parati un grido s’alzerà}
\rigo{\ch{Fa\#m}{}\ch{Do\#m}{}\ch{Re}{}\ch{Mi}{e} mille voci a far da eco ad una voce fioca ormai}
\rigo{\ch{Fa\#m}{}\ch{Do\#m}{}\ch{Re}{}\ch{Mi}{e} allora dai, vieni con noi, è un’avventura in mare aperto}
\rigo{\ch{Fa\#m}{}\ch{Do\#m}{}\ch{Re}{}\ch{Mi}{e} viaggerai insieme a noi nella natura controvento.}
\end{ritornello}
\par\endgroup

\begingroup
\cantotesta{9}{Un mondo, una promessa}{}{\bEG}{1131pt}
\rigo{Ed \ch{Re}{io} son qui che \ch{Sol}{aspetto} e anche \ch{Sim}{tu} sei insieme \ch{La}{a} \ch{La4}{me}}
\rigo{100 \ch{Re}{lune,} 100 \ch{Sol}{soli,} luce a \ch{Sim}{tutto} quel che \ch{La}{}\ch{La4}{c'è}}
\rigo{"Guarda \ch{Sol}{avanti} e \ch{Re}{sorridi",} sulle \ch{Fa\#}{tracce} di \ch{Sim}{B.P.}}
\rigo{L'avventura \ch{Sol}{ha} avuto \ch{Re}{inizio,} prima \ch{Mi}{volta} su a \ch{La4}{Brownsea}}
\rigo{Un \ch{Sol}{orizzonte} \ch{Re}{nuovo,} chi se \ch{La}{lo} immaginava \ch{Sim}{che}}
\rigo{Questo \ch{Sol}{tempo} costruito \ch{Re}{insieme,} di \ch{Mi}{noia} non \ch{La}{ce} n'è... e di \ch{Sol}{noia} \ch{La7}{non} ce \ch{Re}{n'è}}
\strofa
\begin{ritornello}
\rigo{\ch{Sol}{Un} \ch{Re}{grande} \ch{Sol}{}\ch{Re}{cerchio} che \ch{La}{abbraccia} un mondo \ch{Sim}{intero}}
\rigo{Pronti \ch{Sol}{a} fare il nostro \ch{Re}{meglio,} a \ch{La}{servire} per \ch{Re}{}\ch{Re7}{davvero}}
\rigo{\ch{Sol}{Un} \ch{Re}{grande} \ch{Sol}{}\ch{Re}{cerchio} e \ch{Fa\#}{quest'alba} da \ch{Sim}{scoprire}}
\rigo{Ecco il \ch{Sol}{sole} sorge \ch{Re}{ancora,} è \ch{Mi}{ora} di \ch{La}{partire}}
\rigo{\ch{Sol}{Un} \ch{Re}{grande} \ch{Sol}{}\ch{Re}{cerchio} con la \ch{Fa\#}{legge} che è la \ch{Sim}{stessa}}
\rigo{La mia \ch{Sol}{storia} nel \ch{Re}{futuro,} un \ch{Mi}{mondo} una \ch{La}{promessa...} un \ch{Sol}{mondo} \ch{La}{una} \ch{Re}{promessa}}
\strofa
\end{ritornello}
\rigo{E la gente che si incontra nella pista e nel sentiero}
\rigo{quando in route, che è la mia strada, si percorre il mondo intero}
\rigo{buon cammino verso il sole per la stessa direzione}
\rigo{mai uguale nel tracciato, tante vie tante persone}
\rigo{con un passo dopo l'altro, una bussola che non mi inganni}
\rigo{siamo in tanti a camminare, ho percorso 100 anni... ho percorso 100 anni}
\strofa
\begin{ritornello}
\rigo{\ch{Sol}{Un} \ch{Re}{grande} \ch{Sol}{}\ch{Re}{cerchio} che \ch{La}{abbraccia} un mondo \ch{Sim}{intero}}
\rigo{Pronti \ch{Sol}{a} fare il nostro \ch{Re}{meglio,} a \ch{La}{servire} per \ch{Re}{}\ch{Re7}{davvero}}
\rigo{\ch{Sol}{Un} \ch{Re}{grande} \ch{Sol}{}\ch{Re}{cerchio} e \ch{Fa\#}{quest'alba} da \ch{Sim}{scoprire}}
\rigo{Ecco il \ch{Sol}{sole} sorge \ch{Re}{ancora,} è \ch{Mi}{ora} di \ch{La}{partire}}
\rigo{\ch{Sol}{Un} \ch{Re}{grande} \ch{Sol}{}\ch{Re}{cerchio} con la \ch{Fa\#}{legge} che è la \ch{Sim}{stessa}}
\rigo{La mia \ch{Sol}{storia} nel \ch{Re}{futuro,} un \ch{Mi}{mondo} una \ch{La}{promessa...} un \ch{Sol}{mondo} \ch{La}{una} \ch{Re}{promessa}}
\strofa
\end{ritornello}
\rigo{Questa pace che respiri già alla luce dell'aurora}
\rigo{Costruita con le azioni, anche adesso, proprio ora,}
\rigo{la mia mano con la tua, oggi insieme come ieri}
\rigo{nell'impegno di ogni giorno, fanno veri i tuoi pensieri}
\rigo{questa voglia che ci unisce, che da forza alle nostre ali,}
\rigo{noi diversi, noi lontani, noi fratelli, noi uguali... noi fratelli, noi uguali}
\strofa
\begin{ritornello}
\rigo{\ch{Sol}{Un} \ch{Re}{grande} \ch{Sol}{}\ch{Re}{cerchio} che \ch{La}{abbraccia} un mondo \ch{Sim}{intero}}
\rigo{Pronti \ch{Sol}{a} fare il nostro \ch{Re}{meglio,} a \ch{La}{servire} per \ch{Re}{}\ch{Re7}{davvero}}
\rigo{\ch{Sol}{Un} \ch{Re}{grande} \ch{Sol}{}\ch{Re}{cerchio} e \ch{Fa\#}{quest'alba} da \ch{Sim}{scoprire}}
\rigo{Ecco il \ch{Sol}{sole} sorge \ch{Re}{ancora,} è \ch{Mi}{ora} di \ch{La}{partire}}
\rigo{\ch{Sol}{Un} \ch{Re}{grande} \ch{Sol}{}\ch{Re}{cerchio} con la \ch{Fa\#}{legge} che è la \ch{Sim}{stessa}}
\rigo{La mia \ch{Sol}{storia} nel \ch{Re}{futuro,} un \ch{Mi}{mondo} una \ch{La}{promessa...} un \ch{Sol}{mondo} \ch{La}{una} \ch{Re}{promessa}}
\strofa
\end{ritornello}
\rigo{\ch{Sol}{un} mondo \ch{Re}{una} \ch{Sol}{promessa...} un \ch{La}{mondo} \ch{Re}{una} promessa}
\par\endgroup

\begingroup\senzaaccordi
\cantotesta{10}{Canto delle Scolte d'Assisi}{}{\bRS}{319pt}
\rigo{Squilla!}
\rigo{Squilla la tromba che già il giorno finì}
\rigo{già del coprifuoco la canzone salì}
\rigo{Su, scolte, alle torri guardie armate ola'}
\rigo{attente in silenzio vigilate.}
\strofa
\rigo{O nostri santi che in cielo esultate}
\rigo{Vergini sante gloriose e beate}
\rigo{noi v'invochiamo.}
\strofa
\rigo{Questa città con il vostro amore}
\rigo{guardate contro il nemico che l'anima}
\rigo{tiene contro la morte che subito viene}
\rigo{in ogni cuor sia pace e bene}
\rigo{sia tregua ad ogni dolor. Pace.}
\strofa
\rigo{Squilla squilla la tromba che già il giorno finì}
\rigo{già del coprifuoco la canzone salì.}
\strofa
\rigo{Su, scolte, alle torri guardie armate ola'}
\rigo{attente in silenzio vigilate.}
\strofa
\rigo{Attente, scolte, vigilate.}
\par\endgroup

\begingroup
\cantotesta{11}{Se incontrerò}{}{\bRS}{389pt}
\rigo{\ch{Do}{Se} incontrerò \ch{Sol}{chi} \ch{Do}{con} me ha speranza}
\rigo{\ch{Sol}{nel} domani che splenda il sole,}
\rigo{\ch{Do}{se} incontrerò \ch{Sol}{chi} \ch{Do}{con} me ci crede}
\rigo{\ch{Fa}{io} gli stringerò \ch{Sol}{}\ch{Do}{la} mano.}
\strofa
\rigo{\ch{Do}{Ma} \ch{Rem}{se} incontrerò \ch{Sol7}{chi} stanco più \ch{Do}{non} vuole continuare,}
\rigo{\ch{Do}{io} \ch{Rem}{mi}  \ch{Do}{fermerò} per dirgli \ch{Fa}{andiamo} io ti \ch{Sol}{aiuterò.}}
\strofa
\rigo{Se incontrerò chi la stessa legge ha scelto di libertà e di amore}
\rigo{gli narrerò le avventure liete che mi hanno portato qui.}
\rigo{Ma se non saprà la gioia di cantare e di servire,}
\rigo{io mi fermerò per dargli un poco di felicità.}
\strofa
\rigo{Se incontrerò chi si sente solo sarà facile dirgli vieni,}
\rigo{siamo già in tanti che insieme credono nella fraternità.}
\par\endgroup

\begingroup\senzaaccordi
\cantotesta{12}{Servire}{}{\bRS}{384pt}
\rigo{Rit: Servire è parola che parla col cuore}
\rigo{servire è parola lucente d'affetto}
\rigo{servire è un concetto fraterno d'amore}
\rigo{che unisce, stupisce, che parla da se}
\strofa
\rigo{La strada fraterna è cosparsa d'amore}
\rigo{l'amore è paziente, è benigno l'amore}
\rigo{se dessi tutti i miei averi di cuore}
\rigo{e non dessi me stesso a che serve l'amore}
\strofa
\commento{Rit.}
\strofa
\rigo{L'amore non invidia, non riesce a vantare}
\rigo{non gonfia, non manca rispetto di dare}
\rigo{non cerca interesse, l'amore non si adira}
\rigo{ma cerca ed aspira soltanto a servire.}
\strofa
\commento{Rit.}
\strofa
\rigo{Se ami la feste, il bosco, il ruscello}
\rigo{se cerchi l'amore che cerca un fanciullo}
\rigo{sel male, sai, il conto non devi tenere}
\rigo{ma devi soltanto pensare a servire.}
\strofa
\commento{Rit.}
\strofa
\rigo{Fratello, fratello, servire è parola}
\rigo{pronunciata una volta è ricordata per sempre}
\rigo{sui campi, nei boschi da rover e da anziani}
\rigo{in mezzo alla folla da oggi al domani.}
\par\endgroup

\begingroup
\cantotesta{13}{Strade delle mie routes}{}{\bRS}{692pt}
\rigo{\ch{Sol}{}\ch{Re4}{}\ch{Mim7}{}\ch{Re4}{Un} fuoco per imparare}
\rigo{\ch{Sol}{}\ch{Re4}{}\ch{Mim7}{}\ch{Re4}{la} tua canoa a guidare}
\rigo{\ch{Mim}{}\ch{Re}{}\ch{Do}{}\ch{Re}{ti} guardi intorno a cercare}
\rigo{\ch{Mim}{}\ch{Re}{}\ch{Do}{}\ch{Re}{la} direzione giusta da seguir.}
\strofa
\begin{ritornello}
\rigo{\ch{Sol}{}\ch{Re4}{}\ch{Re}{Strade} delle mie routes,}
\rigo{\ch{Do9}{}\ch{Re}{è} il tempo di partire ancora}
\rigo{\ch{Sol}{}\ch{Re4}{}\ch{Re}{strade} delle mie routes,}
\rigo{\ch{Do9}{}\ch{Re}{e} la fatica è già l’aurora}
\rigo{\ch{Mim}{}\ch{Do}{di} un nuovo giorno che}
\rigo{\ch{Re}{le} mie mani aperte trova}
\rigo{\ch{Mim}{}\ch{Do}{e} amare allor sarà}
\rigo{\ch{Re}{la} mia vera libertà}
\strofa
\end{ritornello}
\rigo{\ch{Sol}{}\ch{Re4}{}\ch{Mim7}{}\ch{Re4}{Tempo} di scelte, viandante,}
\rigo{\ch{Sol}{}\ch{Re4}{}\ch{Mim7}{}\ch{Re4}{scopri} te stessa ed il mondo}
\rigo{\ch{Mim}{}\ch{Re}{}\ch{Do}{}\ch{Re}{impari} a servire chi è accanto a te}
\rigo{\ch{Mim}{}\ch{Re}{}\ch{Do}{}\ch{Re}{in} nome di un bene profondo.}
\strofa
\begin{ritornello}
\rigo{\ch{Sol}{}\ch{Re4}{}\ch{Re}{Strade} delle mie routes,}
\rigo{\ch{Do9}{}\ch{Re}{è} il tempo di partire ancora}
\rigo{\ch{Sol}{}\ch{Re4}{}\ch{Re}{strade} delle mie routes,}
\rigo{\ch{Do9}{}\ch{Re}{e} la fatica è già l’aurora}
\rigo{\ch{Mim}{}\ch{Do}{di} un nuovo giorno che}
\rigo{\ch{Re}{le} mie mani aperte trova}
\rigo{\ch{Mim}{}\ch{Do}{e} amare allor sarà}
\rigo{\ch{Re}{la} mia vera libertà}
\strofa
\end{ritornello}
\rigo{\ch{Sol}{}\ch{Re4}{}\ch{Mim7}{}\ch{Re4}{Essere} testimone}
\rigo{\ch{Sol}{}\ch{Re4}{}\ch{Mim7}{}\ch{Re4}{del} tuo amore e della fede}
\rigo{\ch{Mim}{}\ch{Re}{}\ch{Do}{}\ch{Re}{lasci} il fuoco per dare di più}
\rigo{\ch{Mim}{}\ch{Re}{}\ch{Do}{}\ch{Re}{fatta} sia la Sua volontà}
\strofa
\begin{ritornello}
\rigo{\ch{Sol}{}\ch{Re4}{}\ch{Re}{Strade} delle mie routes,}
\rigo{\ch{Do9}{}\ch{Re}{è} il tempo di partire ancora}
\rigo{\ch{Sol}{}\ch{Re4}{}\ch{Re}{strade} delle mie routes,}
\rigo{\ch{Do9}{}\ch{Re}{e} la fatica è già l’aurora}
\rigo{\ch{Mim}{}\ch{Do}{di} un nuovo giorno che}
\rigo{\ch{Re}{le} mie mani aperte trova}
\rigo{\ch{Mim}{}\ch{Do}{e} amare allor sarà}
\rigo{\ch{Re}{la} mia vera libertà}
\end{ritornello}
\par\endgroup

\end{multicols}

\sezione{Canti di lupettismo}{}

\begin{multicols}{2}\raggedcolumns

\begingroup\senzaaccordi
\cantotesta{14}{A caccia andiamo}{}{\bLC}{319pt}
\rigo{O lupi a caccia andiamo,}
\rigo{corriamo e pronti siam}
\rigo{sorriso lieto abbiamo}
\rigo{"del nostro meglio" urliam}
\rigo{e meglio urliam.}
\strofa
\rigo{Sì, bisognerà, bisognerà che vada.}
\rigo{Sì, bisognerà, bisognerà partir.}
\rigo{Sì, bisognerà, bisognerà che vada.}
\rigo{Sì, bisognerà, bisognerà partir.}
\strofa
\rigo{In mezzo alla natura}
\rigo{fratelli noi sarem.}
\rigo{Lupetti all'aria pura}
\rigo{felici correrem e correrem.}
\strofa
\rigo{Sì, bisognerà...}
\strofa
\rigo{Rientrati nella tana}
\rigo{ricorderemo, sì,}
\rigo{la caccia ormai lontana,}
\rigo{cantando ognor così, ognor così.}
\strofa
\rigo{Sì, bisognerà...}
\par\endgroup

\begingroup\senzaaccordi
\cantotesta{15}{Antica traccia}{}{\bLC}{273pt}
\rigo{Oh fratel mio gentile racconta tu (bis)}
\rigo{come Santo Francesco volle seguire Gesù.}
\rigo{Né feste, balli e canti, né guerreggiar (bis)}
\rigo{né ricchezze il suo cuor sapevano più}
\rigo{accontentar.}
\rigo{Ma povertà e preghiera volle per sé; (bis)}
\rigo{per portare l’amor divenne l’araldo del re!}
\strofa
\rigo{O fratel mio gentile, chi l’ascoltò (bis)}
\rigo{chi ascoltò il suo messaggio lieto}
\rigo{che il mondo incantò?}
\rigo{Lebbrosi e mendicanti, ricchi signor, (bis)}
\rigo{cittadini, viandanti e giovani puri di cuor.}
\rigo{Il lupo gli ubbidiva e l’usignol (bis)}
\rigo{con lui laudavan Dio}
\rigo{e il creato cantava il Signor.}
\strofa
\rigo{O fratel mio gentile, se vuoi puoi tu (bis)}
\rigo{come Santo Francesco puoi seguire Gesù!}
\par\endgroup

\begingroup\senzaaccordi
\cantotesta{16}{Attorno alla rupe}{}{\bLC}{307pt}
\rigo{Attorno alla rupe orsù lupi andiam}
\rigo{d'Akela e Baloo or le voci ascoltiam}
\rigo{del branco la forza in ciascun lupo sta}
\rigo{del lupo la forza del branco sarà.}
\strofa
\rigo{Rit.: uulla ullalala uullalala}
\rigo{del lupo la forza nel branco sarà (2v.)}
\strofa
\rigo{Or quando si sente un richiamo lontan}
\rigo{risponde un vicino più forte ulular}
\rigo{a quel che la legge fedel seguirà}
\rigo{a lui buona caccia si ripeterà}
\strofa
\rigo{Rit.: uulla ullalala uullalala}
\rigo{a lui buona caccia si ripeterà}
\strofa
\rigo{Il debole cucciolo via via crescerà}
\rigo{e forte ed ardito col clan caccerà}
\rigo{finché sulla rupe un dì porterà}
\rigo{la pelle striata del vile Shere Khan.}
\strofa
\rigo{Rit.: uulla ullalala uullalala}
\rigo{la pelle striata del vile Shere Khan (2v.)}
\par\endgroup

\begingroup\senzaaccordi
\cantotesta{17}{Branco di Akela}{}{\bLC}{210pt}
\rigo{Akela aveva un branco felice e fier,}
\rigo{sestiglie rossi e fulvi, pezzati e ner,}
\rigo{furbi i rossi in verità, con la loro abilità}
\rigo{nel seguir la traccia di Akela!.}
\strofa
\rigo{Ridenti anche i pezzati felici ognor}
\rigo{perché sestiglia in gamba è quella lor.}
\rigo{Ecco adesso venir giù dalla caccia di Baloo}
\rigo{vincitori i neri pien di virtù.}
\strofa
\rigo{Audaci e agili sono i fulvi ancor}
\rigo{i primi della corsa e allegri in cor.}
\rigo{La miglior quale sarà che per prima arriverà}
\rigo{a conquistare la preda di Akela.}
\par\endgroup

\begingroup\senzaaccordi
\cantotesta{18}{Buon lupetto}{}{\bLC}{350pt}
\rigo{Buon lupetto esser voglio.}
\rigo{Buon lupetto esser voglio.}
\rigo{Buon lupetto esser voglio.}
\rigo{Canterò, giocherò, servirò.}
\strofa
\rigo{Buon lupetto esser voglio.}
\rigo{Canterò, giocherò, servirò.}
\rigo{Buon lupetto esser voglio.}
\rigo{Canterò, giocherò, servirò.}
\strofa
\rigo{Penso alle due stelle,}
\rigo{penso, penso come saran belle,}
\rigo{come saran belle.}
\rigo{Penso alle due stelle,}
\rigo{penso, penso come saran belle}
\rigo{sul mio cuor.}
\strofa
\rigo{Oh presto sì, patapim patapom,}
\rigo{io le avrò, patapim patapom,}
\rigo{sopra il mio cuore,}
\rigo{sopra il mio cuore.}
\strofa
\rigo{Oh presto sì, patapim patapom,}
\rigo{io le avrò, patapim patapom,}
\rigo{sopra il mio cuore,}
\rigo{sì, sì le avrò. Le avrò.}
\par\endgroup

\begingroup\senzaaccordi
\cantotesta{19}{Buona caccia}{}{\bLC}{300pt}
\rigo{Partiamo col branco in caccia buona caccia, buona caccia;}
\rigo{partiamo col branco in caccia buona caccia fratellin.}
\strofa
\rigo{Cuor leale, lingua cortese fanno strada nella giungla;}
\rigo{cuor leale, lingua cortese nella giunglastrada fan.}
\strofa
\rigo{Siamo d'uno stesso sangue fratellino tu ed io;}
\rigo{siamo d'uno stesso sangue fratellino io e te.}
\strofa
\rigo{E' le legge della giungla vecchia e vera come il cielo}
\rigo{senza legge non c'e' vero non c'e' lupo che vivra'.}
\strofa
\rigo{Grande e' il bosco e tenebroso ed il lupo e' piccolino;}
\rigo{Baloo, l'orso coraggioso, viene a caccia insieme a te.}
\par\endgroup

\begingroup
\cantotesta{20}{Caccia al sambhur}{}{\bLC}{256pt}
\rigo{\ch{Re}{Nella} valletta chiara s'aggirano \ch{La7}{i} lupi,}
\rigo{\ch{Re}{annusano} le tracce \ch{La7}{del} sambur \ch{Re}{veloce,}}
\rigo{seguono Akela a caccia tra \ch{La7}{aspri} dirupi,}
\rigo{e da \ch{La7}{lontano} l'eco \ch{Re}{ripete} la voce:}
\strofa
\rigo{"E salta, lupo, \ch{La7}{salta} \ch{Re}{salta}}
\rigo{e salta \ch{La7}{sulla} preda.}
\rigo{Shere Kan che non ti veda,}
\rigo{\ch{Re}{che} te la vuol levar!}
\rigo{E salta, lupo, \ch{La7}{salta} \ch{Re}{salta}}
\rigo{e salta \ch{La7}{sulla} preda.}
\rigo{Shere Kan che non ti veda,}
\rigo{\ch{Re}{che} te la vuol levar".}
\par\endgroup

\begingroup
\cantotesta{21}{Canto del risveglio}{}{\bLC}{205pt}
\rigo{\ch{Mib}{Lupo} salta su che già l'or\ch{Sib}{a} è suonata !    \ch{Mib}{}}
\rigo{Lupo salta su, cominciato è \ch{Sib}{il} dì.   \ch{Mib}{}}
\rigo{Tra poco al ruscelletto \ch{Dom}{fra} \ch{Sib}{le} rocce corr\ch{Mib}{eremo}}
\rigo{e lieti all'acqu\ch{Dom}{a} fresca tutti\ch{Sib}{} insiem ci la\ch{Mib}{veremo.}}
\strofa
\rigo{La giungla già ci attende}
\rigo{presto in caccia partiremo,}
\rigo{la pista in mezzo al bosco}
\rigo{pronti e svelti seguiremo.}
\par\endgroup

\begingroup\senzaaccordi
\cantotesta{22}{Canto della sera}{}{\bLC}{216pt}
\rigo{Or scende la sera e noi siamo quaggiù}
\rigo{tutto il Branco è in preghiera}
\rigo{e Ti invoca, o Gesù.}
\rigo{Dall'alto del cielo proteggici Tu}
\rigo{ed il nostro branco benedici, o Gesù.}
\rigo{Quel che oggi abbiam fatto e le nostre B.A.}
\rigo{ecco a Te presentiamo con tutta lealtà.}
\strofa
\rigo{Proteggici a casa la mamma e il papà}
\rigo{fa' che noi meritiamo la loro bontà.}
\rigo{Silente è la giungla nel giorno che muor}
\rigo{su tutti discenda la gioia e l'amor.}
\rigo{Dall'alto del cielo proteggici Tu}
\rigo{ed il nostro branco benedici o Gesù.}
\par\endgroup

\begingroup\senzaaccordi
\cantotesta{23}{Danza dei cani rossi}{}{\bLC}{294pt}
\rigo{Rah, rah, rah! il cane rosso è qua.}
\rigo{Arriva giù, non lascia più la traccia di Wontolla.}
\strofa
\rigo{Rah, rah, rah! Il cane rosso è qua.}
\rigo{La rabbia sì lo inferocì, Wontolla lo annunciò.}
\strofa
\rigo{Or Bagheera, Kaa e Baloo alla lotta pronti son;}
\rigo{al comando di Akela coi suoi lupi Mowgli vien.}
\strofa
\rigo{Rah, rah, rah! Il cane rosso è qua.}
\rigo{Di peso fu portato giù dall'acqua del Waingunga.}
\strofa
\rigo{Rah, rah, rah! Il cane rosso è qua.}
\rigo{E quando là approderà il branco pronto è già.}
\strofa
\rigo{Tutti i lupi di Seeonee stan lottando con ardor}
\rigo{ed il Dhole invasor alla fin sconfiggeran.}
\par\endgroup

\begingroup\senzaaccordi
\cantotesta{24}{Danza del fiore rosso}{}{\bLC}{279pt}
\rigo{C'era un tempo un tigre perfida e sleal}
\rigo{nella giungla faceva tanto tanto mal}
\rigo{ma poi Mowgli il fiore rosso}
\rigo{alla rupe ci portò}
\rigo{ed il pelo ai suoi nemici}
\rigo{con un ramo bruciacchiò.}
\rigo{Fiore rosso avvampa ancor}
\rigo{tutto il branco ne ha terror.}
\strofa
\rigo{Mowgli prese pel mento}
\rigo{il muso di Shere-Khan}
\rigo{mentre il fior sbocciava}
\rigo{nella salda man;}
\rigo{tigre zoppa della giungla}
\rigo{la tua pelle brucerò}
\rigo{sulla rupe del Consiglio}
\rigo{se non fuggi stenderò.}
\rigo{Fiore rosso avvampa ancor}
\rigo{anche Shere-Khan ne ha terror.}
\par\endgroup

\begingroup\senzaaccordi
\cantotesta{25}{Danza della tregua dell'acqua}{}{\bLC}{256pt}
\rigo{Lupo la parola d'Hathi ascolterà}
\rigo{quando nella giungla tregua annunzierà.}
\strofa
\rigo{Rit. Tra la la la la ra la la la}
\rigo{tra la la la la ra la la la}
\rigo{tra la la la la ra la la la}
\rigo{tra la la la la ra la la la.}
\strofa
\rigo{Tigre col cerbiatto all'acqua insieme andrà}
\rigo{se nella Waingunga la roccia si vedrà.}
\strofa
\rigo{Acqua fresca e viva Baloo berrà con te...}
\rigo{ma sulla sorgente una macchia scura c'é.}
\strofa
\rigo{Gridano d'orrore Bagheera, Kaa, Baloo:}
\rigo{via da questa fonte, non ci tornare più!!!.}
\strofa
\rigo{E Shere-khan che arriva... d'umano sangue ha}
\rigo{zanne rosse orrende, ed ora a bere va.}
\par\endgroup

\begingroup\senzaaccordi
\cantotesta{26}{Fratelli alla candida luna}{}{\bLC}{286pt}
\rigo{RIT: Fratelli alla candida luna cantiamo la nostra canzone più bella:}
\rigo{tra faggi ed abeti, che gioia veniamo la nostra canzone a cantar.}
\strofa
\rigo{1- "Una via sola è vera" ci disse il lupo anziano}
\rigo{e dietro a lui noi siamo la traccia a seguitar.}
\rigo{Costanti alla ricerca, gioiosi nel cammino,}
\rigo{erriam fino al mattino seguendo il lupo anzian.}
\strofa
\rigo{2- La via è troppo scoscesa la traccia già si perde.}
\rigo{Solo non si disperde chi segue il lupo anzian.}
\rigo{Lupetto non lasciare oh, non lasciare il branco}
\rigo{il pelo è ancora bianco seguiamo il lupo anzian.}
\par\endgroup

\begingroup\senzaaccordi
\cantotesta{27}{Gli allegri folletti}{}{\bLC}{406pt}
\rigo{Dietro il muro d'un vecchio convento,}
\rigo{ogni sera, al calare del sole,}
\rigo{in un prato di rose e viole}
\rigo{cento folletti convegno si dan.}
\rigo{Vengon qui d'ogni parte del mondo}
\rigo{d'ogni razza e color essi son;}
\rigo{e facendo un gran girotondo}
\rigo{lieti cantano questa canzon:}
\strofa
\begin{ritornello}
\rigo{Siamo gli allegri, vispi folletti;}
\rigo{mille dispetti spesso facciam:}
\rigo{spargiamo il pepe sopra il guanciale,}
\rigo{versiamo il sale dentro il caffè;}
\rigo{mettiam gli spilli sulle poltrone}
\rigo{e dal balcone l'acqua spruzziam.}
\rigo{Ma se qualcuno ci chiede aiuto,}
\rigo{in un minuto noi accorriam!}
\rigo{Ma se qualcuno ci chiede aiuto,}
\rigo{in un minuto noi accorriam!}
\strofa
\end{ritornello}
\rigo{Un folletto si stacca dal gruppo}
\rigo{e bisbiglia ad ogni bambino:}
\rigo{Hai pensato, o mio fratellino,}
\rigo{che anche tu il folletto puoi far?}
\rigo{Per portare un po' d'allegria}
\rigo{non ci vuole, in fondo, gran che,}
\rigo{se anche tu vuoi destar simpatia,}
\rigo{canta, allora, insieme con me:}
\strofa
\rigo{Siamo gli allegri...}
\par\endgroup

\begingroup\senzaaccordi
\cantotesta{28}{Il canto del cuculo}{}{\bLC}{142pt}
\rigo{Sentiam nella foresta il cuculo cantar}
\rigo{ai piedi d'una quercia lo stiamo ad ascoltar.}
\strofa
\rigo{Cucu cucu...}
\strofa
\rigo{La notte è tenebrosa non c'è chiaror lunar}
\rigo{sentiam nel fitto bosco i lupi ulular.}
\strofa
\rigo{Ahu, ahu...}
\par\endgroup

\begingroup\senzaaccordi
\cantotesta{29}{Il lupo ferito}{}{\bLC}{312pt}
\rigo{Bagheera che vide narrò}
\rigo{d'un lupo dal giovane ardir}
\rigo{che il branco alla caccia lasciò}
\rigo{e volle l'istinto seguir.}
\strofa
\rigo{Rit. Lupo lupo torna nel branco}
\rigo{a cacciar con noi vieni vieni}
\rigo{vieni il grand'urlo a lanciar.}
\strofa
\rigo{La legge che insegna il dover}
\rigo{la legge che Akela ululò}
\rigo{la legge che un branco fa fier}
\rigo{il lupo randagio scordò.}
\strofa
\rigo{Ma quando dall'ombra sbucò}
\rigo{il toro selvaggio e guerrier}
\rigo{invano quel lupo saltò}
\rigo{ferito ululò sul sentier.}
\strofa
\rigo{E il lupo l'errore capì}
\rigo{tornò più veloce nel pian}
\rigo{e corse alla rupe e finì}
\rigo{fra le zampe di un lupo anzian.}
\par\endgroup

\begingroup\senzaaccordi
\cantotesta{30}{Il risveglio della giungla}{}{\bLC}{179pt}
\rigo{La luna è bianca e scuro è l'orizzonte}
\rigo{già tutto tace, sol canta la fonte.}
\strofa
\rigo{Lassù nel cielo bruno tutto solo}
\rigo{già Mang, il pipistrello, intreccia il volo.}
\strofa
\rigo{La giungla si ridesta piano piano}
\rigo{ed echeggia di lontano già il richiamo:}
\rigo{Buona caccia. Buona caccia.}
\strofa
\rigo{a tutti quelli che rispettano la Legge;}
\rigo{Buona caccia Buona caccia Buona caccia!}
\par\endgroup

\begingroup\senzaaccordi
\cantotesta{31}{L'ultimo canto di Akela}{}{\bLC}{242pt}
\rigo{Il lupo Wontolla non caccerà più}
\rigo{Akela da un pezzo è disteso laggiù}
\rigo{e s'ode d'intorno stormire di fronde}
\rigo{sommessa sussurra la giungla un sospir.}
\strofa
\rigo{Rit. E muoio contento fratello a sentir}
\rigo{che sei qui vicino a vedermi morir.}
\strofa
\rigo{E Mowgli il ranocchio lo venne a cercar}
\rigo{tra i lupi caduti l'udì rantolar}
\rigo{alzata la testa e le spalle di lì}
\rigo{a Mowgli in ginocchio diceva così.}
\strofa
\rigo{Già l'ultimo colpo il mio dente vibrò}
\rigo{tra poco la giungla lasciare dovrò.}
\rigo{Ma muoio contento, fratello, a sentir}
\rigo{che sei qui vicino a vedermi morir!}
\par\endgroup

\begingroup
\cantotesta{32}{La cassa misteriosa}{}{\bLC}{879pt}
\rigo{Un \ch{Do}{giorno} un mio lupetto passeggiava lungo il mar}
\rigo{\ch{Fa}{quand'ecco} a un tratto \ch{Do}{scorse} là una \ch{Sol7}{cassa} galleggiar.}
\rigo{E giù \ch{Do}{si} tuffò e presala su a \ch{Fa}{riva} la portò,}
\rigo{\ch{Sol7}{ma} apertala \ch{Do}{vide} la ... \ch{Sol7}{proprio} dinnanzi a sé,}
\rigo{\ch{Do}{ma} apertala vide la ... \ch{Sol7}{proprio} dinnanzi a \ch{Do}{sé.}}
\strofa
\rigo{Orbene imbarazzato su che fare e che non far}
\rigo{pensò: qui un buon consiglio forse Akela potrà dar}
\rigo{e in collo, perciò, se la carico e in Tana la portò,}
\rigo{ma appena io vidi la ... via, gli gridai di qui,}
\rigo{ma appena io vidi la ... via, gli gridai di qui.}
\strofa
\rigo{Pazienza disse e con la cassa prese a camminar}
\rigo{cercando un qualche amico a cui poterla regalar,}
\rigo{ma invano cercò e astuzie tentò che sempre a lui restò}
\rigo{la cassa con tutta la ... oh che disperazion!}
\rigo{La cassa con tutta la ... oh che disperazion!}
\strofa
\rigo{Via via che lui cresceva da Lupetto a Esplorator}
\rigo{la cassa inseparabile restava, oh che dolor,}
\rigo{ma un giorno morì e al cielo salì contento di lasciar}
\rigo{in terra la cassa e la ... senza più afflizion}
\rigo{in terra la cassa e la ... senza più afflizion.}
\strofa
\rigo{Ma quando al Gran Portone sbigottito si fermò}
\rigo{s'accorse che volata era la cassa fin lassù}
\rigo{allor s'infuriò, in cuor borbottò, ma alfin dove' capir}
\rigo{che anche la cassa e la ... in cielo volevano andar}
\rigo{che anche la cassa e la ... in cielo volevano andar.}
\strofa
\rigo{Moral di questa storia è che andando lungo il mar}
\rigo{se mai vedrete a un tratto là una cassa galleggiar}
\rigo{non ditevi: su, tuffiamoci giù, vediamo cosa c'è}
\rigo{perché quella perfida ... mai più vi lascerà,}
\rigo{perché quella perfida ... mai più vi lascerà.}
\par\endgroup

\begingroup
\cantotesta{33}{Lupo salta su}{}{\bLC}{205pt}
\rigo{\ch{Mib}{Lupo} salta su che già l’ora è \ch{Sib}{suonata} !     \ch{Mib}{}}
\rigo{Lupo salta su, cominciato è i\ch{Sib}{l} dì.    \ch{Mib}{}}
\rigo{Tra poco al ruscelletto fr\ch{Dom}{a} le\ch{Sib}{} rocce corre\ch{Mib}{remo}}
\rigo{e lieti all’acqua \ch{Dom}{fresca} tutti\ch{Sib}{} insiem ci la\ch{Mib}{veremo.}}
\strofa
\rigo{La giungla già ci attende}
\rigo{presto in caccia partiremo,}
\rigo{la pista in mezzo al bosco}
\rigo{pronti e svelti seguiremo.}
\par\endgroup

\begingroup
\cantotesta{34}{Nella giungla misteriosa}{}{\bLC}{406pt}
\rigo{\ch{Rem}{Nella} giungla misteriosa,}
\rigo{sotto un \ch{Solm}{albero} di \ch{Rem}{cocco}}
\rigo{un leone \ch{La}{se} ne \ch{Rem}{andava}}
\rigo{annusando \ch{Solm}{lo} scirocco.}
\rigo{\ch{La}{occo,} \ch{Rem}{occo}}
\strofa
\rigo{Annusava il venticello}
\rigo{per veder se il vicino}
\rigo{sulla sponda del ruscello}
\rigo{fosse nato un rosmarino.}
\rigo{ino ino}
\strofa
\rigo{Quella pianta prelibata}
\rigo{la delizia delle fiere}
\rigo{un bel giorno fu trovata}
\rigo{nella pancia del barbiere}
\rigo{ere ere}
\strofa
\rigo{Il barbiere era venuto}
\rigo{per tosare i Tumba Tumba,}
\rigo{ma il sentiero avea perduto}
\rigo{nel mistero della Giungla.}
\rigo{ungla, ungla.}
\par\endgroup

\begingroup\senzaaccordi
\cantotesta{35}{Salta o lupo}{}{\bLC}{325pt}
\rigo{Il sentiero di caccia è aperto innanzi a te}
\rigo{ascolta l'urlo, s'ode da lontan.}
\rigo{Il favor della giungla ti accompagnerà}
\rigo{sarà il tuo passo forte ardito e fier.}
\strofa
\rigo{Rit. Salta o lupo se tu vuoi ogni preda}
\rigo{sarà presto ai piedi tuoi}
\rigo{giuste tracce seguirai}
\rigo{e i segreti della giungla scoprirai.}
\strofa
\rigo{Se un buon lupo presto tu vuoi diventar}
\rigo{le Massime e la Legge non scordare,}
\rigo{la parola di Akela sia vicino a te}
\rigo{osserva sempre i detti di Baloo.}
\strofa
\rigo{Occhio aperto alle stelle dritto sa guardar}
\rigo{Bagheera a caccia ti può consigliar.}
\rigo{Del tuo branco la forza in te stesso sta}
\rigo{e la sua forza il branco a te darà.}
\strofa
\rigo{Quando il branco alla rupe si radunerà}
\rigo{di nuove imprese ancor potrem parlar.}
\rigo{Quella Legge che uniti sempre manterrà}
\rigo{il nostro vanto e il tuo onor sarà.}
\par\endgroup

\begingroup\senzaaccordi
\cantotesta{36}{Siam nella vasta giungla}{}{\bLC}{95pt}
\rigo{Siam nella vasta giungla}
\rigo{siam sulla pista di Akela}
\rigo{lieti cantando, sempre cantando}
\rigo{Lupi siam.}
\par\endgroup

\begingroup
\cantotesta{37}{Ula, ula, ula}{}{\bLC}{324pt}
\rigo{\ch{Do}{Ula,} ula, ula}
\rigo{\ch{Fa}{è} sera e \ch{Do}{stanchi} siam}
\rigo{\ch{Do}{Ula,} ula, ula}
\rigo{\ch{Re7}{doman} \ch{Sol}{lieti} sarem}
\rigo{\ch{Do}{Ula,} ula, ula}
\rigo{\ch{Fa}{tra} \ch{Do}{poco} dormirem}
\rigo{\ch{Fa}{per} \ch{Do}{tutta} la notte}
\rigo{\ch{Sol}{noi} \ch{Do}{ci} riposerem.}
\strofa
\rigo{Ula ula ula}
\rigo{il branco dorme già}
\rigo{ula ula ula}
\rigo{la giungla tacerà}
\rigo{ula ula ula}
\rigo{Gesù benedirà}
\rigo{questa nostra tana}
\rigo{e ci sorriderà...}
\par\endgroup

\begingroup\senzaaccordi
\cantotesta{38}{Un bravo lupo}{}{\bLC}{237pt}
\rigo{Un bravo lupo io voglio diventar}
\rigo{e la promessa per sempre rispettar}
\rigo{gentile e più cortese con tutti io sarò}
\rigo{la Buona Azion sempre farò.}
\strofa
\rigo{Akela oh Akela oh, Te lo prometto,}
\rigo{più in gamba io sarò.}
\strofa
\rigo{Caro Baloo, caro Baloo}
\rigo{io la legge osserverò di più. ( 2 volte )}
\strofa
\rigo{Bagheera oh, Bagheera oh}
\rigo{te lo prometto, in caccia io verrò. ( 2 volte )}
\strofa
\rigo{Mio caro Kaa, mio caro Kaa}
\rigo{te lo prometto, farò tante B.A. ( 2 volte )}
\par\endgroup

\begingroup\senzaaccordi
\cantotesta{39}{Zampe agili e silenti}{}{\bLC}{121pt}
\rigo{Zampe agili e silenti}
\rigo{zanne aguzze e rilucenti}
\rigo{acuto l'occhio scorge nell'oscurità}
\rigo{orecchie tese odono il vento che va.}
\rigo{D'un lupo sono i segni}
\rigo{d'un forte cacciator.}
\par\endgroup

\begingroup
\cantotesta{40}{Zanna bianca}{}{\bLC}{610pt}
\rigo{Un bel \ch{Re}{giorno} nella Giungla}
\rigo{Zanna \ch{La7}{Bianca} a caccia andò,}
\rigo{ma \ch{Re}{arrivato} giù al ruscello}
\rigo{sulle \ch{La7}{zampe} si fermò.}
\rigo{\ch{Re}{Guarda} un poco, \ch{La7}{guarda} un po',}
\rigo{guarda un poco, \ch{Re}{guarda} un po',}
\rigo{poveretto a quel che vide}
\rigo{tutto il \ch{La7}{pel} gli si \ch{Re}{drizzò.}}
\strofa
\rigo{Là riflessa dentro l’acqua}
\rigo{la sua immagine trovò}
\rigo{e nel folto della Giungla}
\rigo{col fiatone ritornò.}
\rigo{Guarda un poco, guarda un po',}
\rigo{guarda un poco, guarda un po',}
\rigo{e il Bandar-log da un rame}
\rigo{lo guardò e si sbellicò.}
\strofa
\rigo{L’altro giorno in un cespuglio}
\rigo{con il riccio s’incontrò}
\rigo{e per farne un bel boccone}
\rigo{ecco che gli capitò.}
\rigo{Guarda un poco, guarda un po',}
\rigo{guarda un poco, guarda un po',}
\rigo{con il naso pien di spilli}
\rigo{nella tana ritornò.}
\strofa
\rigo{Ma cacciando con Bagheera}
\rigo{e ascoltando il buon Baloo,}
\rigo{cose brutte come queste}
\rigo{non gli accaddero mai più.}
\rigo{Proprio più, proprio più,}
\rigo{proprio più, proprio più,}
\rigo{cose brutte come queste}
\rigo{non gli \ch{La7}{accaddero} mai \ch{Re}{più.}}
\par\endgroup

\end{multicols}

\sezione{Fuoco di bivacco}{% Testo della pagina divisoria "Fuoco di bivacco" (inserito da generate_latex.py).
\begin{center}\begin{minipage}{0.86\textwidth}
Un buon fuoco ha un andamento a onda: sale, tocca il punto più alto e poi scende fino al silenzio.
I canti sono raccolti per fase, nell'ordine in cui si cantano.

\medskip
\begin{tabular}{@{}l@{\hspace{4mm}}p{0.72\textwidth}@{}}
\textbf{Apertura} & Si accende il fuoco: i canti che lo introducono.\\[1mm]
\textbf{Bans} & Brevi, con gesti: sciolgono il gruppo e danno ritmo.\\[1mm]
\textbf{Crescendo} & Canti noti e ritmati, la voce sale.\\[1mm]
\textbf{Picco} & Il momento più acceso, breve.\\[1mm]
\textbf{Calando} & Ritmo e voce scendono: canti di strada e di ricordo.\\[1mm]
\textbf{Chiusura} & Canti raccolti, quasi sottovoce.\\[1mm]
\textbf{Ultimi} & \textit{Scende la notte} e \textit{Signor, fra le tende schierati}, poi silenzio.
\end{tabular}

\medskip
\textit{Regola d'oro:} chi guida legge il fuoco, non la scaletta. Se il gruppo è stanco, si scende prima.
\end{minipage}\end{center}}

\begin{multicols}{2}\raggedcolumns

\sottosezione{Apertura}

\begingroup\senzaaccordi
\cantotesta{41}{Andiamo al falò}{}{\bEG}{422pt}
\rigo{La fiamma è ammainata}
\rigo{e chiuse son le tende}
\rigo{finita è la giornata}
\rigo{e in ciel la luna ascende.}
\strofa
\begin{ritornello}
\rigo{Alò, alò, andiamo al falò.}
\rigo{Alò, alò, andiamo al falò.}
\rigo{Tralla lalla la la la}
\rigo{tralla lalla la la la.}
\rigo{Alò, alò, andiamo al falò.}
\rigo{Alò, alò, andiamo al falò.}
\strofa
\end{ritornello}
\rigo{Nel cielo costellato}
\rigo{Altair è già salita}
\rigo{il fuoco è preparato}
\rigo{la fiamma già c'invita.}
\strofa
\rigo{Alò, alò, andiamo...}
\strofa
\rigo{I boschi misteriosi}
\rigo{ripeton echeggianti}
\rigo{gli accenti melodiosi}
\rigo{dei nostri lieti canti.}
\strofa
\rigo{Alò, alò, andiamo...}
\strofa
\rigo{Seduti qui sul monte}
\rigo{al caldo della fiamma}
\rigo{sentiam sfiorar la fronte}
\rigo{dal bacio della mamma.}
\strofa
\rigo{Alò, alò, andiamo al falò.}
\rigo{Alò, alò, andiamo al falò.}
\par\endgroup

\begingroup
\cantotesta{42}{La leggenda del fuoco}{}{\bEG}{547pt}
\rigo{\ch{Do}{Splende} il fuoco nel cerchio dell’esplorato\ch{Sol}{r}   \ch{Do}{}}
\rigo{ascoltate la voce della f\ch{Sol}{iamma} d’or. \ch{Do}{}}
\strofa
\commento{Rit.}
\rigo{Sali al c\ch{Sol}{iel} fiamma leggera  \ch{Do}{}}
\rigo{del gran fuoc\ch{Fa}{o} caldo e buon     \ch{Sol}{}}
\rigo{sotto i pini alla brughie\ch{Do}{ra}}
\rigo{sali in a\ch{Fa}{lto} sali ancor \ch{Sol}{}}
\rigo{sali in alto sali anc\ch{Do}{or,}}
\rigo{fuoco dell’es\ch{Fa}{plorator.}  \ch{Sol}{}}
\strofa
\rigo{Ero un principe un giorno perfido e sleal}
\rigo{e spargevo d’intorno il terrore e il mal.}
\strofa
\commento{Rit.}
\strofa
\rigo{D’un gran mago l’incanto tosto mi punì}
\rigo{e nei tronchi del bosco mi rinchiuse un dì.}
\strofa
\commento{Rit.}
\strofa
\rigo{Da quel giorno nei tronchi prigioniero son}
\rigo{e costretto a soffrire freddo e solleon.}
\strofa
\commento{Rit.}
\strofa
\rigo{Nell’ardor della fiamma mi consumo qui}
\rigo{e col ceppo che arde brucio anch’io ogni dì.}
\strofa
\commento{Rit.}
\strofa
\rigo{Dal tremendo supplizio convertito son}
\rigo{e per l’uomo divento fuoco caldo e buon.}
\strofa
\commento{Rit.}
\par\endgroup

\begingroup\senzaaccordi
\cantotesta{43}{Salite o fiamme}{}{\bEG}{70pt}
\rigo{Salite o fiamme in questa sera,}
\rigo{portate in cielo la nostra preghiera.}
\par\endgroup

\sottosezione{Bans}

\begingroup
\cantotesta{44}{Alla fiera dell'est}{A. Branduardi, L. Zappa}{\bLC}{482pt}
\rigo{\ch{Mim}{Alla}\ch{Sol}{} fiera dell'est,}
\rigo{p\ch{Re}{er} \ch{Sol}{due} soldi}
\rigo{u\ch{Re}{n} t\ch{Sol}{opolino}}
\rigo{m\ch{Si7}{io} p\ch{Mim}{adre} comprò.}
\strofa
\rigo{E\ch{Re}{} ve\ch{Sol}{nne} il gatto}
\rigo{c\ch{Re}{he} \ch{Sol}{si} mangiò il topo}
\rigo{che al mercato}
\rigo{mio padre comprò.}
\strofa
\rigo{E venne il cane}
\rigo{che morse il gatto...}
\strofa
\rigo{E venne il bastone}
\rigo{che picchiò il cane...}
\strofa
\rigo{E venne il fuoco}
\rigo{che bruciò il bastone...}
\strofa
\rigo{E venne l'acqua}
\rigo{che spense il fuoco...}
\strofa
\rigo{E venne il toro}
\rigo{che bevve l'acqua...}
\strofa
\rigo{E venne il macellaio}
\rigo{che uccise il toro...}
\strofa
\rigo{E l'angelo della morte...}
\strofa
\rigo{E infine il Signore....}
\par\endgroup

\begingroup\senzaaccordi
\cantotesta{45}{Arturo e Gertrude}{}{\bEG}{70pt}
\rigo{C'è un buco nel secchio Arturo, Arturo (bis)}
\rigo{E stroppa chel buso Gertrude, Gertrude (bis)}
\par\endgroup

\begingroup\senzaaccordi
\cantotesta{46}{Banana cocco baobab}{}{\bEG}{298pt}
\rigo{Ero stanco di essere un uomo della città,}
\rigo{in questa gabbia di cemento a vivere come si fa?}
\rigo{Quanti problemi inutili, troppe preoccupazioni,}
\rigo{è meglio star sugli alberi, mettersi a penzoloni e fare...}
\strofa
\begin{ritornello}
\rigo{Uh uh ah ah, uh uh ah ah, uh uh ah ah}
\rigo{banana cocco baobab.}
\rigo{Banana banana banana banana}
\rigo{banana banana banana banana}
\rigo{banana banana banana banana}
\rigo{banana cocco baobab.}
\strofa
\end{ritornello}
\rigo{Andare con le liane invece che col tram}
\rigo{al posto del vestito una pelliccia di Pierre-Tarzan.}
\rigo{E poi stare sugli alberi a tirare le banane}
\rigo{in testa ai coccodrilli, ai bufali e ai caproni e fare...}
\par\endgroup

\begingroup\senzaaccordi
\cantotesta{47}{Bella però sta macchina}{}{\bLC}{153pt}
\rigo{Bella però sta macchina}
\rigo{con guida cibernetica}
\rigo{schiaccia il pulsante, gira il volante}
\rigo{e si muove una lancetta}
\rigo{e si accende una lucetta}
\rigo{e si sente un campanello}
\rigo{din don din don}
\strofa
\rigo{Bella però sta macchina...}
\par\endgroup

\begingroup\senzaaccordi
\cantotesta{48}{Colla dei francobolli}{}{\bLC}{146pt}
\rigo{Chissà perché}
\rigo{la colla dei francobolli}
\rigo{la fanno amara}
\rigo{al sapore del sale.}
\rigo{Avanti, chi inventa}
\rigo{i francobolli alla menta,}
\rigo{e quelli al limone,}
\rigo{e quelli al lampone?}
\par\endgroup

\begingroup\senzaaccordi
\cantotesta{49}{Fli flai}{}{\bEG}{196pt}
\rigo{A: Fly}
\rigo{T: Fly}
\rigo{A: Fly flay}
\rigo{T: Fly flay}
\rigo{A: Fly flay flo}
\rigo{T: Fly flay flo}
\rigo{A: Cummala cummala cummala de vista}
\rigo{T: Cummala cummala cummala de vista}
\rigo{A: Oh no, no, no de vista}
\rigo{T: Oh no, no, no de vista}
\rigo{A: Minni esta, minni esta, minni esta oh!}
\rigo{T: Minni esta, minni esta, minni esta oh!}
\par\endgroup

\begingroup\senzaaccordi
\cantotesta{50}{Il duca di Barnabò}{}{\bLC}{146pt}
\rigo{Il conte di Barnabò}
\rigo{aveva una compagnia}
\rigo{che andava su e giù}
\rigo{oppure a mezza via.}
\rigo{Saliti ch'eran su}
\rigo{discesi ch'eran giù}
\rigo{chi stava mezza via}
\rigo{non stava ne su e ne giù.}
\par\endgroup

\begingroup\senzaaccordi
\cantotesta{51}{Il ranocchio ciclista}{}{\bEG}{513pt}
\rigo{Un bel giorno un ranocchietto}
\rigo{che viveva nel pantan,}
\rigo{si stancò di quel laghetto,}
\rigo{volle andare via lontan.}
\rigo{Inforcò la bicicletta}
\rigo{da perfetto corridor:}
\rigo{verso il bosco in tutta fretta}
\rigo{s'avviò con grande ardor.}
\rigo{Oh com'è bello andar}
\rigo{verso quel bosco in fior.}
\rigo{Oh com'è bello andar}
\rigo{verso quel bosco in fior.}
\strofa
\rigo{Pedalava con passione}
\rigo{il ranocchio corridor,}
\rigo{quando a un tratto, dannazione,}
\rigo{fece un grosso ruzzolon.}
\rigo{Che disastro! Che macello!}
\rigo{La sua bici fracassò,}
\rigo{perciò sotto un alberello}
\rigo{tutto pesto si fermò.}
\rigo{Povero ranocchion,}
\rigo{dopo quel cascaton.}
\rigo{Povero ranocchion,}
\rigo{dopo quel cascaton.}
\strofa
\rigo{Col passare delle ore}
\rigo{tutto buio diventò}
\rigo{e il ranocchio dal terrore}
\rigo{a tremare incominciò.}
\rigo{E correndo a più-non-posso,}
\rigo{nello stagno ritornò.}
\rigo{Le altre rane di quel fosso}
\rigo{gli cantarono così:}
\rigo{“Chi te l’ha fatto far,}
\rigo{chi te l’ha fatto far,}
\rigo{chi te l’ha fatto far,}
\rigo{chi te l’ha fatto far”.}
\par\endgroup

\begingroup\senzaaccordi
\cantotesta{52}{Il tortellino}{}{\bLC}{133pt}
\rigo{Io sono un tortellino,}
\rigo{mi muovo un po' di qua,}
\rigo{mi muovo un po' di là.}
\rigo{Tutto ben condito,}
\rigo{con un poco di sale,}
\rigo{con un poco di pepe,}
\rigo{mi vuoi mangiare tu?}
\par\endgroup

\begingroup\senzaaccordi
\cantotesta{53}{Jack è in cucina}{}{\bLC}{108pt}
\rigo{Jack va in cucina con Tina (3 volte)}
\rigo{e suona il suo vecchio banjo}
\rigo{bom bom bom}
\rigo{fli flai fli flai flo (3 volte)}
\rigo{e suona il suo vecchio banjo.}
\par\endgroup

\begingroup\senzaaccordi
\cantotesta{54}{Jona nella balena}{}{\bLC}{184pt}
\rigo{Giona nella balena}
\rigo{felice fu, felice fu,}
\rigo{perché in prigion,}
\rigo{perché in prigion,}
\rigo{perché s'egli era in pena}
\rigo{potea mangiar, potea mangiar,}
\rigo{dello storion, dello storion,}
\rigo{ma al terzo dì, pirimpimpì,}
\rigo{quel gran bestion, porompompò,}
\rigo{lo rigettò, porompompò,}
\rigo{senza pietà!}
\par\endgroup

\begingroup\senzaaccordi
\cantotesta{55}{Jump down}{}{\bEG}{261pt}
\rigo{Jump down, turn around, pick a bale of cotton}
\rigo{Jump down, turn around, pick a bale a day}
\rigo{Jump down, turn around, pick a bale of cotton}
\rigo{Jump down, turn around, pick a bale a day}
\strofa
\rigo{Oh Lordy, pick a bale of cotton}
\rigo{Oh Lordy, pick a bale a day}
\rigo{Oh Lordy, pick a bale of cotton}
\rigo{Oh Lordy, pick a bale a day}
\strofa
\rigo{Pick a, pick a, pick a, pick a, pick a bale of cotton}
\rigo{Pick a, pick a, pick a, pick a, pick a bale a day}
\rigo{Pick a, pick a, pick a, pick a, pick a bale of cotton}
\rigo{Pick a, pick a, pick a, pick a, pick a bale a day}
\par\endgroup

\begingroup\senzaaccordi
\cantotesta{56}{L'albero piantato}{}{\bLC}{472pt}
\rigo{C'è un albero piantato in mezzo al prato (bis)}
\rigo{e sopra l'albero indovina cosa c'è?}
\rigo{ci sono i rami, i rami sopra l'albero}
\strofa
\rigo{C'è un albero...}
\rigo{e sopra i rami indovina cosa c'è?}
\rigo{ci sono le foglie, le foglie sopra i rami, i rami sopra l'albero}
\strofa
\rigo{e sopra le foglie indovina cosa c'è?}
\rigo{c'è un nido, un nido fra le foglie, le foglie sopra i rami, i rami sopra l'albero}
\strofa
\rigo{e dentro il nido indovina cosa c'è?}
\rigo{c'è un uovo, un uovo dentro il nido, un nido fra le foglie, le foglie sopra i rami, i rami sopra l'albero}
\strofa
\rigo{e dentro l’uovo indovina cosa c'è?}
\rigo{c'è un pulcino, un pulcino dentro l’uovo, un uovo dentro il nido, un nido fra le foglie, le foglie sopra i rami, i rami sopra l'albero}
\strofa
\rigo{e dentro il pulcino indovina cosa c'è?}
\rigo{c'è un cuore, un cuore nel pulcino, un pulcino dentro l’uovo, un uovo dentro il nido, un nido fra le foglie, le foglie sopra i rami, i rami sopra l'albero}
\strofa
\rigo{e dentro il cuore sapete cosa c'è?}
\rigo{c'è un mondo, un mondo dentro il cuore, un cuore nel pulcino, un pulcino dentro il nido, un nido fra le foglie, le foglie sopra i rami, i rami sopra l'albero}
\strofa
\rigo{e nel mondo indovina cosa c'è?}
\rigo{C'è un albero piantato in mezzo al prato (bis)}
\par\endgroup

\begingroup\senzaaccordi
\cantotesta{57}{La battaglia di Magenta}{}{\bLC}{102pt}
\rigo{Era una sera Battaglia di Magenta}
\rigo{o che piacere caricare i cavalieri}
\rigo{cavalieri caricate con una mano.}
\strofa
\rigo{(per finire:) Cavalieri carica!}
\par\endgroup

\begingroup\senzaaccordi
\cantotesta{58}{La bella polenta}{}{\bEG}{198pt}
\rigo{Come sea pianta la bella polenta,}
\rigo{la bella polenta sea pianta cussì.}
\strofa
\rigo{Oh, oh, oh, la bella polenta cussì.}
\strofa
\rigo{Come sea innaffia..}
\rigo{Come chea cresse...}
\rigo{Come sea taja...}
\rigo{Come sea sbatte...}
\rigo{Come sea masena...}
\rigo{Come sea cusina...}
\rigo{Come sea magna...}
\rigo{Come sea...}
\par\endgroup

\begingroup\senzaaccordi
\cantotesta{59}{La macchina del capo}{}{\bLC}{115pt}
\rigo{-La macchina del capo}
\rigo{ha un buco nella gomma,- (3 volte)}
\rigo{ripariamola col chewing-gum.}
\strofa
\rigo{(si ripete sostituendo ogni volta con i rumori)}
\par\endgroup

\begingroup\senzaaccordi
\cantotesta{60}{Napuletana mia}{}{\bEG}{192pt}
\rigo{Napuletana mia, napuletana.}
\rigo{Quant' si bella assaje ma nun 'o saje}
\rigo{Napuletana mia, napuletana.}
\rigo{Quant' si bella assaje ma nun 'o saje}
\strofa
\rigo{La voce di Persona1,}
\rigo{la voce di Persona1,}
\rigo{la voce di Persona1}
\rigo{da Persona1 vogliam sentir}
\strofa
\rigo{Persona1: Vegetale1!}
\strofa
\rigo{Ad libitum.}
\par\endgroup

\begingroup\senzaaccordi
\cantotesta{61}{Onci onci}{}{\bEG}{248pt}
\rigo{Noi siamo gli onci onci,}
\rigo{veniamo giù dai monti,}
\rigo{se vuoi ballar con noi,}
\rigo{tu devi far così:}
\strofa
\rigo{Onci onci onci, fai un saltino ino ino,}
\rigo{non far cadere ere ere il tuo vicino,}
\rigo{vai verso destra estra estra,}
\rigo{muovi la testa esta esta,}
\rigo{vai a sinistra istra istra,}
\rigo{alza la vista ista ista,}
\rigo{fai un sorrisetto etto etto,}
\rigo{muovi il culetto etto etto,}
\rigo{fai l’aeroplano...}
\rigo{dammi la mano...}
\strofa
\rigo{rip.}
\par\endgroup

\begingroup\senzaaccordi
\cantotesta{62}{One banana}{}{\bEG}{90pt}
\rigo{One banana, two banana, three banana, four...}
\rigo{Five banana, six banana, seven and anymore.}
\strofa
\rigo{Ad libitum.}
\par\endgroup

\begingroup\senzaaccordi
\cantotesta{63}{Quattro pirati}{}{\bEG}{146pt}
\rigo{Quattro pirati del mar dei Sargassi}
\rigo{hanno una zattera fatta di assi}
\rigo{vanno remando, dicono loro}
\rigo{alla ricerca di un grande tesoro}
\rigo{però uno è alto, uno è basso, uno è zoppo}
\rigo{e il quarto ha la benda sull'occhio}
\rigo{vanno remando, dicono loro}
\rigo{alla ricerca di un grande tesoro}
\par\endgroup

\begingroup\senzaaccordi
\cantotesta{64}{San Michele}{}{\bEG}{95pt}
\rigo{San Michele aveva un gallo}
\rigo{bianco, rosso, verde e giallo.}
\rigo{E per farlo crescere bene}
\rigo{gli dava latte e miele.}
\par\endgroup

\begingroup
\cantotesta{65}{Siamo andati alla caccia del leon}{}{\bLC}{451pt}
\rigo{\ch{Sol}{Sono} andato alla cacci\ch{Do}{a} del leon...\ch{Sol}{} pem pem}
\rigo{Sono andato alla caccia del leon... pem p\ch{Re}{em}}
\rigo{Sono andato \ch{Sol}{alla} caccia\ch{Sol7}{} (3v)\ch{Do}{} de\ch{Lam}{l} le\ch{Sol}{on..}\ch{Re}{} pe\ch{Sol}{m} pem}
\strofa
\commento{Rit.}
\rigo{SOL               DO                 SOL (RE)}
\rigo{Singing ja ja yuppi yuppi ja (2v)}
\rigo{Singing ja j\ch{Sol}{a} yu\ch{Sol7}{ppi} ja j\ch{Do}{a} yupp\ch{Lam}{i} ja ja y\ch{Sol}{uppi} yuppi \ch{Re}{ja}      \ch{Sol}{}}
\strofa
\rigo{Taglia l'ancora la nave per partir... tu tu  RIT.}
\strofa
\rigo{Sulla riva son gli amici a salutar... bye bye  RIT.}
\strofa
\rigo{Siamo ormai nella foresta equatorial... brr brr  RIT.}
\strofa
\rigo{Un leone sta dormendo non lontan ... ron ron  RIT.}
\strofa
\rigo{Se ci vede ci divora in un boccon... gnam gnam  RIT.}
\strofa
\rigo{Cari amici non facciamolo svegliar... pss pss  RIT.}
\par\endgroup

\begingroup\senzaaccordi
\cantotesta{66}{Sul lago Tanganika}{}{\bEG}{308pt}
\rigo{Sul lago Tanganika zum baibai zum baibai}
\rigo{Ci stava una piroga zum baibai zum baibai}
\rigo{Ci stava una piroga zum baibai zum baibai}
\rigo{E con il capo degli Zulù}
\strofa
\rigo{Ela mazù, ela mazù, evviva il capo degli Zulù}
\strofa
\rigo{E in quella piroga}
\rigo{Ci stavano i guerrieri (2v)}
\rigo{E con il capo degli Zulù}
\strofa
\rigo{Ela mazù, ela mazù, evviva il capo degli Zulù}
\strofa
\rigo{E i forti guerrieri}
\rigo{Andavano alla guerra (2v)}
\rigo{E con il capo degli Zulù}
\strofa
\rigo{Ela mazù, ela mazù, evviva il capo degli Zulù}
\strofa
\rigo{Andavano alla guerra}
\rigo{Perché avevano rapito (2v)}
\rigo{E la regina degli Zulù}
\strofa
\rigo{Ela mazù, ela mazù, evviva il capo degli Zulù.}
\par\endgroup

\begingroup\senzaaccordi
\cantotesta{67}{Un ciodo de fero vecio}{}{\bEG}{153pt}
\rigo{Un ciodo de fero vecio}
\rigo{dela mecanica dela mecanica.}
\rigo{Un ciodo de fero vecio}
\rigo{dela mecanica de precision.}
\strofa
\rigo{Ohi bela mechimecomeca}
\rigo{mechimecomeca, mechimecomeca.}
\rigo{Ohi bela mechimecomeca}
\rigo{mechimecomeca, mechimecò.}
\par\endgroup

\begingroup\senzaaccordi
\cantotesta{68}{Un elefante}{}{\bEG}{128pt}
\rigo{Un elefante si dondolava}
\rigo{sopra il filo di una ragnatela}
\rigo{e ritenendo la cosa interessante}
\rigo{andò a chiamare un altro elefante.}
\strofa
\rigo{Due elefanti si dondolavano...}
\rigo{Tre elefanti...}
\par\endgroup

\begingroup\senzaaccordi
\cantotesta{69}{Uomini della filibusta}{}{\bEG}{184pt}
\rigo{Capitano urla: "Uomini della filibusta, Parados?"}
\rigo{Soldati: "Parados!"}
\rigo{Capitano: "Al puente?"}
\rigo{Soldati: "Al puente!"}
\rigo{Capitano: "Fuego!"}
\rigo{Soldati: "BOOM!"}
\rigo{Capitano: "Fuego!"}
\rigo{Soldati: "BOOM!"}
\rigo{Capitano: "Fuego!"}
\rigo{Soldati: "BOOM! BOOM! BOOM!"}
\par\endgroup

\sottosezione{Crescendo}

\begingroup
\cantotesta{70}{Al passo del guidon}{}{\bEG}{726pt}
\rigo{\ch{Re}{Al} passo del Guidon,}
\rigo{\ch{La}{fratello} scout, t'attende l'avventura}
\rigo{\ch{Mim}{tra} \ch{La}{il} verde delle macchie e sotto il sol.}
\rigo{\ch{Re}{Al} passo del Guidon,}
\rigo{\ch{La}{avanti} ad esplorar la natura.}
\rigo{\ch{Mim}{un} \ch{La}{nido,} un'erba, un fior t'aspetta ed è}
\rigo{\ch{Re}{}\ch{La}{tutto} per te!}
\strofa
\begin{ritornello}
\rigo{\ch{Re}{}\ch{La}{}\ch{Re}{Apri} gli occhi fratello scout}
\rigo{\ch{La}{tutto} il mondo che è intorno a te}
\rigo{\ch{Sim}{è} una \ch{Mim}{cosa} \ch{La}{meravigliosa}}
\rigo{\ch{Re}{Apri} \ch{La}{}\ch{Re}{gli} occhi fratello scout}
\rigo{\ch{La}{tutto} il mondo che è intorno a te}
\rigo{\ch{Sim}{è} una \ch{Mim}{cosa} \ch{La}{meravigliosa} da scoprir.}
\strofa
\end{ritornello}
\rigo{Al lato del sentier}
\rigo{la pista ancor fratel, non è battuta,}
\rigo{la bussola ti guida senza error...}
\rigo{Al lato del sentier}
\rigo{il mondo è tutta terra sconosciuta;}
\rigo{ma certo c'è un amico che di là}
\rigo{ti aspetterà.}
\strofa
\begin{ritornello}
\rigo{\ch{Re}{}\ch{La}{}\ch{Re}{Apri} gli occhi fratello scout}
\rigo{\ch{La}{tutto} il mondo che è intorno a te}
\rigo{\ch{Sim}{è} una \ch{Mim}{cosa} \ch{La}{meravigliosa}}
\rigo{\ch{Re}{Apri} \ch{La}{}\ch{Re}{gli} occhi fratello scout}
\rigo{\ch{La}{tutto} il mondo che è intorno a te}
\rigo{\ch{Sim}{è} una \ch{Mim}{cosa} \ch{La}{meravigliosa} da scoprir.}
\strofa
\end{ritornello}
\rigo{Al fuoco del falò}
\rigo{la gioia dei fratelli è la più pura;}
\rigo{fa un'unica gran tenda il vasto ciel...}
\rigo{Al fuoco del falò}
\rigo{si sente ancor più limpida e sicura}
\rigo{la Voce che ci vuol esplorator.}
\strofa
\begin{ritornello}
\rigo{\ch{Re}{}\ch{La}{}\ch{Re}{Apri} gli occhi fratello scout}
\rigo{\ch{La}{tutto} il mondo che è intorno a te}
\rigo{\ch{Sim}{è} una \ch{Mim}{cosa} \ch{La}{meravigliosa}}
\rigo{\ch{Re}{Apri} \ch{La}{}\ch{Re}{gli} occhi fratello scout}
\rigo{\ch{La}{tutto} il mondo che è intorno a te}
\rigo{\ch{Sim}{è} una \ch{Mim}{cosa} \ch{La}{meravigliosa} da scoprir.}
\end{ritornello}
\par\endgroup

\begingroup\senzaaccordi
\cantotesta{71}{El pescador da Ciosa}{}{\bRS}{427pt}
\rigo{El pescator de Ciosa}
\rigo{frisendo el pesse…}
\rigo{ghe xe cascà 'na iossa… soe braghesse!}
\rigo{E me nona co a scoresa a fa un}
\rigo{subioon… che se sente fino… a}
\rigo{Maeamoco!}
\strofa
\rigo{E me morosa vien}
\rigo{da Villafranca}
\rigo{la g'ha 'na tetta nera}
\rigo{e una bianca.}
\strofa
\rigo{E me morosa vien}
\rigo{da Carbonera}
\rigo{la g'ha 'na tetta bianca}
\rigo{e una nera.}
\strofa
\rigo{E me morosa vien}
\rigo{da Castelfranco}
\rigo{la g'ha do tette piene}
\rigo{de vin bianco.}
\strofa
\rigo{E me morosa g'à}
\rigo{un difetto solo}
\rigo{la g'à na tetta}
\rigo{sensa bichignolo.}
\strofa
\rigo{E me morosa vien}
\rigo{da Santa Bona}
\rigo{co' solo cento franchi}
\rigo{a te dà la m…}
\strofa
\rigo{E me morosa vien..}
\par\endgroup

\begingroup
\cantotesta{72}{Fare well}{}{\bEG}{765pt}
\rigo{\ch{Re}{Là} sul fondo val un coniglio sta}
\rigo{Singing Polly Wolly Doodle all the \ch{La}{day}}
\rigo{che ben presto in ciel se ne vuol volar}
\rigo{Singing \ch{La7}{Polly} Wolly Doodle all the \ch{Re}{day}}
\strofa
\begin{ritornello}
\rigo{Farewell, farewell, farewell, my fairy fay}
\rigo{I'm off to Louisiana for to see my Sousy Anna}
\rigo{Singing Polly Wolly Doodle all the day}
\strofa
\end{ritornello}
\rigo{Quando in alto è già sulla nube là}
\rigo{Singing Polly Wolly Doodle all the day}
\rigo{l'apparecchio allor prende a dondolar}
\rigo{Singing Polly Wolly Doodle all the day}
\strofa
\rigo{Farewell, farewell...}
\strofa
\rigo{Il coniglio allor tutto pien d'ardor}
\rigo{Singing Polly Wolly Doodle all the day}
\rigo{prende l'ombrellon e con emozion}
\rigo{Singing Polly Wolly Doodle all the day}
\strofa
\rigo{Farewell, farewell...}
\strofa
\rigo{Quasi a terra è già quando vede là}
\rigo{Singing Polly Wolly Doodle all the day}
\rigo{Scritto sul cartel "Non si può sostar"}
\rigo{Singing Polly Wolly Doodle all the day}
\strofa
\rigo{Farewell, farewell...}
\strofa
\rigo{Quando a terra ahimè ha già messo il piè}
\rigo{Singing Polly Wolly Doodle all the day}
\rigo{un sever guardian vuol la coda in man}
\rigo{Singing Polly Wolly Doodle all the day}
\strofa
\rigo{Farewell, farewell...}
\strofa
\rigo{Così sta in prigion per la presunzion}
\rigo{Singing Polly Wolly Doodle all the day}
\rigo{per voler volar sta senza mangiar}
\rigo{Singing Polly Wolly Doodle all the day}
\strofa
\rigo{Farewell, farewell...}
\strofa
\rigo{A grasshopper sittin' on a railroad track}
\rigo{Singing Polly Wolly Doodle all the day}
\rigo{and pickin' his teeth with a carpet tack}
\rigo{Singing Polly Wolly Doodle all the day}
\strofa
\rigo{Farewell, farewell...}
\par\endgroup

\begingroup\senzaaccordi
\cantotesta{73}{Masafacia'}{}{\bEG}{95pt}
\rigo{Masafacia, masafacia}
\rigo{can che l'abbaia, can che l'abbaia}
\rigo{ultimo grido, ultimo grido}
\rigo{gingibagi, gingibagi}
\par\endgroup

\begingroup\senzaaccordi
\cantotesta{74}{Milan brucia}{}{\bEG}{95pt}
\rigo{Milan brucia, Milan brucia}
\rigo{guarda guarda (bis)}
\rigo{al fuoco, al fuoco}
\rigo{ma l'acqua non c'è!}
\par\endgroup

\begingroup\senzaaccordi
\cantotesta{75}{Sul colle}{}{\bEG}{235pt}
\rigo{Rit. Sul colle scorre il ruscello ancor,}
\rigo{nel bosco canta il cucù,}
\rigo{è sorto il sole, esplorator,}
\rigo{non indugiare più, non indugiar.}
\strofa
\rigo{È sorto il sole esplorator,}
\rigo{non indugiare più!}
\rigo{Rit.rna al tuo lavoro,}
\rigo{che grande gioia ti dà}
\rigo{diffondi attorno a te così}
\rigo{letizia e gran bontà.}
\strofa
\rigo{La buona azione non tralasciar}
\rigo{sii sempre pronto a servir;}
\rigo{ricorda sempre: nel tuo lavor}
\rigo{lealtà devi seguir.}
\par\endgroup

\begingroup\senzaaccordi
\cantotesta{76}{Tinchete Ritinchete}{}{\bLC}{83pt}
\rigo{Tinchete ritenchete, tinchete ritinchete}
\rigo{avete lavoro per lo stagnaro}
\rigo{pentole e coperchioli da riparare}
\par\endgroup

\begingroup\senzaaccordi
\cantotesta{77}{Vecchio scarpone}{}{\bRS}{349pt}
\rigo{Quando la neve scende giù lieve non esitar si deve}
\rigo{prendi il tuo sacco e lesto va, dove più fioccherà.}
\rigo{Vecchio scarpone d' ogni passione vinci la tentazione}
\rigo{non abbracciar la tua Mimì, ma i tuoi fedeli sci...}
\strofa
\rigo{Si va sulla montagna}
\rigo{dove la neve il volto ci abbronzerà.}
\rigo{L'ardor che ci accompagna,}
\rigo{come una fiamma il cuor ci riscalderá.}
\rigo{Salir, sermpre salir,}
\rigo{mentre ogni valle ripeterà così... SCI SCI}
\rigo{Sciator, ripete il vento}
\rigo{con ardimento íl tuo motto sarà SCIAR}
\strofa
\rigo{Signorinella pallida e snella, lascia la tua gonnella e la pelliccia di vison,}
\rigo{mettiti i pantalon. Bon Bo Bon Bon...}
\rigo{Vecchio Scarpone d'ogni passione, vinci la tentazione,}
\rigo{non abbracciar la Mimi, ma i tuoi fedeli sci...}
\par\endgroup

\sottosezione{Picco}

\begingroup\senzaaccordi
\cantotesta{78}{Bevi bevi}{}{\bRS}{823pt}
\rigo{Se ti viene il mal di pancia, non far uso di laudano;}
\rigo{bevi, bevi del Cinzano - ed il mal ti passerà!}
\strofa
\begin{ritornello}
\rigo{Per gli uomini dabbene - come noi}
\rigo{generoso il vin si versi - fallo pieno:}
\rigo{l'acqua è fatta pei perversi,}
\rigo{e il diluvio il dimostrò!}
\strofa
\end{ritornello}
\rigo{Se ti viene il mal di denti, non far uso di calmanti,}
\rigo{bevi, bevi del vin Chianti - ed il mal ti passerà!}
\strofa
\begin{ritornello}
\rigo{Per gli uomini dabbene - come noi}
\rigo{generoso il vin si versi - fallo pieno:}
\rigo{l'acqua è fatta pei perversi,}
\rigo{e il diluvio il dimostrò!}
\strofa
\end{ritornello}
\rigo{Se ti viene il mal di gola, non far uso di clorato,}
\rigo{bevi, bevi del moscato - ed il mal ti passerà.}
\strofa
\begin{ritornello}
\rigo{Per gli uomini dabbene - come noi}
\rigo{generoso il vin si versi - fallo pieno:}
\rigo{l'acqua è fatta pei perversi,}
\rigo{e il diluvio il dimostrò!}
\strofa
\end{ritornello}
\rigo{Se ti viene il mal di cuore, non far uso di strofanto,}
\rigo{bevi, bevi, del vin santo - ed il mal ti passerà.}
\strofa
\begin{ritornello}
\rigo{Per gli uomini dabbene - come noi}
\rigo{generoso il vin si versi - fallo pieno:}
\rigo{l'acqua è fatta pei perversi,}
\rigo{e il diluvio il dimostrò!}
\strofa
\end{ritornello}
\rigo{Se ti viene il mal di calli, non usare il cicarelli,}
\rigo{bevi, bevi il vino dei castelli - ed il mal ti passerà.}
\strofa
\begin{ritornello}
\rigo{Per gli uomini dabbene - come noi}
\rigo{generoso il vin si versi - fallo pieno:}
\rigo{l'acqua è fatta pei perversi,}
\rigo{e il diluvio il dimostrò!}
\strofa
\end{ritornello}
\rigo{Se ti vien l'insolazione, non far uso di rinfresco,}
\rigo{bevi, bevi il Barbaresco - ed il mal ti passerà.}
\strofa
\begin{ritornello}
\rigo{Per gli uomini dabbene - come noi}
\rigo{generoso il vin si versi - fallo pieno:}
\rigo{l'acqua è fatta pei perversi,}
\rigo{e il diluvio il dimostrò!}
\strofa
\end{ritornello}
\rigo{Se ti viene il mal d'amore, non far uso di stricnina,}
\rigo{corri subito in cantina - ed il mal ti passerà.}
\strofa
\begin{ritornello}
\rigo{Per gli uomini dabbene - come noi}
\rigo{generoso il vin si versi - fallo pieno:}
\rigo{l'acqua è fatta pei perversi,}
\rigo{e il diluvio il dimostrò!}
\end{ritornello}
\par\endgroup

\begingroup
\cantotesta{79}{Canzone dell'amicizia}{}{\bEG}{437pt}
\rigo{\ch{Do}{In} \ch{Sol}{un} mondo di maschere}
\rigo{d\ch{Lam}{ove} \ch{Mim}{sembra} impossibile}
\rigo{r\ch{Fa}{ius}\ch{Do}{cire} a sconfiggere}
\rigo{t\ch{Re}{utt}\ch{Sol}{o} ciò che annienta l'uomo.}
\strofa
\rigo{Il potere, la falsità,}
\rigo{la ricchezza, l'avidità,}
\rigo{sono mostri da abbattere,}
\rigo{noi però non siamo soli.}
\strofa
\rigo{C\ch{Do}{ant}\ch{Sol}{a} con noi, batti le mani,}
\rigo{a\ch{Lam}{lzal}\ch{Mim}{e} in\ch{Fa}{} al\ch{Sol}{to,} muovile al ritmo del canto.}
\rigo{S\ch{Do}{tri}\ch{Sol}{ngi} la mano del tuo vicino}
\rigo{e\ch{Lam}{} sco\ch{Mim}{prir}\ch{Fa}{ai} \ch{Sol}{che} \ch{Do}{è} meno duro il cammino così.}
\strofa
\rigo{Ci hai promesso il tuo spirito,}
\rigo{lo sentiamo in mezzo a noi}
\rigo{e perciò possiam credere}
\rigo{che ogni cosa può cambiare.}
\strofa
\rigo{Non possiamo più assistere}
\rigo{impotenti ed attoniti}
\rigo{perché siam responsabili}
\rigo{della vita intorno a noi.}
\strofa
\commento{Rit.}
\par\endgroup

\begingroup\senzaaccordi
\cantotesta{80}{Caramba}{}{\bLC}{268pt}
\rigo{Lassù in cima al monte Nero,}
\rigo{c'è una piccola caverna}
\rigo{ci son dodici briganti,}
\rigo{al chiaror d'una lanterna.}
\strofa
\rigo{Rit. Caramba beviamo del whisky}
\rigo{caramba beviamo del gin.}
\strofa
\rigo{E tu non dar retta al cuore}
\rigo{che tutto passerà.}
\strofa
\rigo{Mentre tutti son festanti}
\rigo{uno solo resta muto:}
\rigo{ancor pieno è il suo bicchiere:}
\rigo{come mai non ha bevuto?}
\rigo{Ma non può dimenticare}
\rigo{il brigante la sua bella}
\rigo{i suoi occhi come il mare,}
\rigo{la sua bocca tanto bella.}
\par\endgroup

\begingroup
\cantotesta{81}{Drizza la tenda}{}{\bEG}{256pt}
\rigo{\ch{Re}{Drizza} la tenda, \ch{La}{pianta} il tuo guidone}
\rigo{\ch{Re}{Questo} è il posto che fa per \ch{La}{te.}}
\rigo{La \ch{Re}{gialla} fiamma sale \ch{La}{verso} il sole}
\rigo{\ch{Re}{E} il suo giglio ha il color \ch{La}{del} ciel.}
\rigo{\ch{Re7}{Coraggio} guida}
\rigo{\ch{Sol}{Canta} il tuo cuore la \ch{La}{tua} \ch{Re}{felicità.}}
\strofa
\rigo{Ed alla notte sai c’è una stella}
\rigo{Su cui punta il tuo giglio d’or}
\rigo{Orienta il passo la strada giusta è quella}
\rigo{Ed ora corri guida ognor.}
\rigo{È questo il canto}
\rigo{Del primo campo della stella polar.}
\par\endgroup

\begingroup\senzaaccordi
\cantotesta{82}{El tram de Opcina}{}{\bEG}{690pt}
\rigo{E anche el tram de Opcina xe nato disgrazià}
\rigo{Vignindo zo de Scorcola 'na casa el ga ribaltà.}
\rigo{Bona de Dio che jera giorno de lavor,}
\rigo{e drento no ghe jera che el povero frenador!}
\strofa
\begin{ritornello}
\rigo{E come la bora che vien e che va,}
\rigo{i dixi che 'l mondo se ga ribaltà!}
\rigo{E come la bora che vien e che va,}
\rigo{i dixi che 'l mondo se ga ribaltà!}
\strofa
\end{ritornello}
\rigo{E anche el tram de Servola xe nato disgrazià.}
\rigo{Corendo in Galeria in piazza el xe sbrissà.}
\rigo{Drento ghe jera diverso personl}
\rigo{che se ga ribaltado e se ga fato mal!}
\strofa
\begin{ritornello}
\rigo{E come la bora che vien e che va,}
\rigo{i dixi che 'l mondo se ga ribaltà!}
\rigo{E come la bora che vien e che va,}
\rigo{i dixi che 'l mondo se ga ribaltà!}
\strofa
\end{ritornello}
\rigo{Col "due" se va a Servola, col "quatro" in Arsenal,}
\rigo{col "sie" se va a Barcola, col "cinque" in Ospedal}
\rigo{col "uno" in Zimitero, col "sette" alla Stazion,}
\rigo{col "nove" in Manicomio, col "diese" in Canon!}
\strofa
\begin{ritornello}
\rigo{E come la bora che vien e che va,}
\rigo{i dixi che 'l mondo se ga ribaltà!}
\rigo{E come la bora che vien e che va,}
\rigo{i dixi che 'l mondo se ga ribaltà!}
\strofa
\end{ritornello}
\rigo{E anche 'ste mulete tute mate pe' l' capèl,}
\rigo{le zerca de compagnarse a qualche bel putel;}
\rigo{ma co le 'riva a casa se senti un gran bordel,}
\rigo{e pare, mare e fia copa i zimisi col martel!}
\strofa
\begin{ritornello}
\rigo{E come la bora che vien e che va,}
\rigo{i dixi che 'l mondo se ga ribaltà!}
\rigo{E come la bora che vien e che va,}
\rigo{i dixi che 'l mondo se ga ribaltà!}
\strofa
\end{ritornello}
\rigo{L'Italia ga pan bianco, la Francia ga bon vin,}
\rigo{Trieste ga putele, tute carighe de morbin,}
\rigo{carbon ga l'Inghiltera, la Rusia ga cavial}
\rigo{e l'Austria ga capuzi che no se pol magnar!}
\strofa
\begin{ritornello}
\rigo{E come la bora che vien e che va,}
\rigo{i dixi che 'l mondo se ga ribaltà!}
\rigo{E come la bora che vien e che va,}
\rigo{i dixi che 'l mondo se ga ribaltà!}
\end{ritornello}
\par\endgroup

\begingroup\senzaaccordi
\cantotesta{83}{I do gobeti}{}{\bEG}{357pt}
\rigo{Una sera, una sera di notte}
\rigo{do gobetti si davan le bòtte}
\rigo{do gobetti si davan le bòtte}
\rigo{se sté zitti ve digo 'l perché.}
\strofa
\begin{ritornello}
\rigo{Gobi, gobi}
\strofa
\end{ritornello}
\rigo{Do gobetti di mezza statura}
\rigo{si parlavan di cose amorose,}
\rigo{ma i gaveva 'na matta paura}
\rigo{che i passanti li stesse a sentir.}
\strofa
\rigo{Uno gera 'l famoso Mattia,}
\rigo{l'altro gera el "fabrica inchiostro",}
\rigo{che 'mbriago de sgnapa, ('sto mostro)}
\rigo{insultava l'amico fedel.}
\strofa
\rigo{A ragion del vero 'l gha dito,}
\rigo{ma va' là, ma va' là ti xe gobo!}
\rigo{E quell'altro el ga dato risposta:}
\rigo{se son gobo ti no te si dritto,}
\rigo{drio la schena ti ga 'n botesèl}
\strofa
\rigo{I s'ha dito parole da ciodi,}
\rigo{i s'ha dato carèghe sul muso,}
\rigo{po' i xe 'ndà a finire in quel buso}
\rigo{'ndò i se beve 'n bicer de quel bon.}
\par\endgroup

\begingroup\senzaaccordi
\cantotesta{84}{Il mio galletto}{I girasoli}{\bRS}{357pt}
\rigo{El mio galletto gà un grosso becco}
\rigo{E le galline ghe portan rispetto}
\rigo{Galletto quì, galletto là}
\rigo{El mio galletto l'è mai a ca'}
\strofa
\begin{ritornello}
\rigo{Muoio, muoio}
\rigo{Muoio per te, mio galletto}
\rigo{Muoio, muoio}
\rigo{Muoio, per te morirò}
\rigo{No no no no...}
\strofa
\end{ritornello}
\rigo{El mio galletto gà una gran prua}
\rigo{E le galline le porta a ca' sua}
\rigo{Galletto quì, galletto là}
\rigo{E le galline non son mai a ca'}
\strofa
\rigo{El mio galletto varda le gambe}
\rigo{Alle galline qu'en senza mutande}
\rigo{Galletto quì, galletto là}
\rigo{El mio galletto l'è mai a ca'}
\strofa
\rigo{El mio galletto gà una gran cresta}
\rigo{Alle galline lui ghe fa' la festa}
\rigo{Galletto quì, galletto là}
\rigo{El mio galletto l'è quasi stufà}
\strofa
\rigo{Muoio...}
\par\endgroup

\begingroup\senzaaccordi
\cantotesta{85}{John Brown}{tradizionale (USA)}{\bEG}{459pt}
\rigo{John Brown giace nella tomba là nel pian}
\rigo{dopo una lunga lotta contro l'oppressor;}
\rigo{John Brown giace nella tomba là nel pian,}
\rigo{ma l'anima vive ancor.}
\strofa
\rigo{Rit. Glory, glory hallelujah}
\rigo{glory, glory hallelujah,}
\rigo{glory, glory hallelujah}
\rigo{ma l'anima vive ancor.}
\strofa
\rigo{Con diciannove suoi compagni di valor,}
\rigo{dall'Est all'Ovest la Virginia conquistò;}
\rigo{con diciannove suoi compagni di valor,}
\rigo{ma l'anima vive ancor.}
\strofa
\rigo{Poi l'hanno ucciso come fosse un traditor}
\rigo{ma il traditor fu quell'uomo che l'impiccò;}
\rigo{poi l'hanno ucciso come fosse un traditor,}
\rigo{ma l'anma vive ancor.}
\strofa
\rigo{John Brown è morto ma lo schiavo è in libertà,}
\rigo{tutti fratelli, bianchi e neri siamo già.}
\strofa
\rigo{John Brown è morto, ma lo schiavo è in libertà,}
\rigo{ma l'anima vive ancor.}
\strofa
\rigo{Oh, non temere colui che il corpo ucciderà}
\rigo{se la tua anima rapir non ti potrà;}
\rigo{Oh, non temere colui che il corpo ucciderà}
\rigo{se l'anima vive ancor.}
\strofa
\rigo{Stelle dei cieli non piangete su John Brown,}
\rigo{stelle dei cieli sorridete con John Brown;}
\rigo{stelle dei cieli non piangete su John Brown,}
\rigo{la sua anima vive ancor.}
\par\endgroup

\begingroup\senzaaccordi
\cantotesta{86}{L'uomo di Cromagnon}{}{\bLC}{427pt}
\rigo{Era l'alba della preistoria}
\rigo{mille e mille anni fa}
\rigo{appariva senza gloria}
\rigo{un essere simile all'orangutan;}
\rigo{aveva il corpo coperto di peli}
\rigo{l'ascia di pietra ed il baston...}
\strofa
\rigo{e volea conquistare la terra}
\rigo{era l'uomo di Cromagnon.}
\strofa
\rigo{Rit. L'uomo di CRO}
\rigo{l'uomo di MA}
\rigo{l'uomo di GNON...}
\strofa
\rigo{l'uomo di Cromagnon;}
\rigo{l'uomo di Cromagnon}
\rigo{l'era proprio un bestion}
\rigo{l'era proprio un bestion.}
\strofa
\rigo{L'uomo di Cromagnon}
\rigo{l'era proprio un bestion}
\rigo{l'era proprio un bestion.}
\strofa
\rigo{Pon pon pon para pam pon pon.}
\strofa
\rigo{Scavalcando i secoli e i mesi}
\rigo{chè la storia la sappiam:}
\rigo{lasciamo perdere gli Esquimesi}
\rigo{gli antichi Assiri e gli Egizian:}
\rigo{non parliam delle armi da guerra:}
\rigo{la spingarda ed il cannon:}
\rigo{chi vuol mai conquistare la terra}
\rigo{è ancor l'uomo di Cromagnon.}
\par\endgroup

\begingroup\senzaaccordi
\cantotesta{87}{La bella la va al fosso}{tradizionale}{\bRS}{221pt}
\rigo{La bella la va al fosso.}
\rigo{Ravanel, remolass, barbabietole e spinass}
\rigo{Tre palanche al mass}
\rigo{La bella la va la fosso, al fosso a resentar}
\rigo{Sì, al fosso a resentar}
\rigo{Intant che la resenta e gh'è burlà giò l'anel}
\rigo{La alza gli occhi al cielo, la vide il ciel seren}
\rigo{La sbassa gli occhi all'onda, la vide un pescator}
\rigo{Ohi pescator dell'onda, pescatemi l'anel}
\rigo{E qul'avrai pescato un regalo ti farò}
\rigo{Andrem lassù sui monti, sui monti a far l'amor.}
\par\endgroup

\begingroup\senzaaccordi
\cantotesta{88}{La bomba imbriaga}{B. De Marzi}{\bRS}{546pt}
\rigo{Quarantatre giorni ca semo in trincea}
\rigo{Magnar pane smarso, dormire par tera.}
\rigo{Nissùn se ricorda, nissùni che scrive,}
\rigo{nissùn che tien nota chi more e chi vive….}
\strofa
\begin{ritornello}
\rigo{Ossst/regheta}
\rigo{Chissà che ‘ne riva}
\rigo{Doman pi domani}
\rigo{‘na botta de vin'}
\rigo{Ossst/regheta}
\rigo{Ciapemo ‘na bala}
\rigo{Coi morti e coi vivi}
\rigo{Strucà li vizin.}
\strofa
\end{ritornello}
\rigo{Silenzio sul fronte. Qualcun ne prepara}
\rigo{Un bel funerale, con banda e con bara…}
\rigo{Silenzio, ecco il fis-cio: l'ariva, la viene}
\rigo{Doman sarà festa: vestive par bene…}
\strofa
\begin{ritornello}
\rigo{Ossst/regheta}
\rigo{Chissà che ‘ne riva}
\rigo{Doman pi domani}
\rigo{‘na botta de vin'}
\rigo{Ossst/regheta}
\rigo{Ciapemo ‘na bala}
\rigo{Coi morti e coi vivi}
\rigo{Strucà li vizin.}
\strofa
\end{ritornello}
\rigo{Se snànara un merlo nel ziel de cobalto}
\rigo{Un boto, una sberla, s-giantizi su in alto.}
\rigo{Se mai sarà ver che ghe xe il padreterno,}
\rigo{‘sta raza de bechi va drita a l'inferno.}
\strofa
\rigo{Ossst/regheta}
\rigo{Se fusse ‘na bota}
\rigo{Se fusse ‘na bota}
\rigo{Ripiena de vin.}
\rigo{Ossst/regheta}
\rigo{Ma l'era una bomba,}
\rigo{ma l'era una bomba}
\rigo{s-ciopà lì vizin.}
\rigo{Bombà.}
\par\endgroup

\begingroup\senzaaccordi
\cantotesta{89}{La damigiana}{Gianni Dego}{\bRS}{490pt}
\rigo{Per combattere la nebbia in Valpadana}
\rigo{L'unico rimedio sai qual è? La damigiana!}
\rigo{Con Albano e Sangiovese,}
\rigo{Malvasia a più riprese,}
\rigo{Poi Barbera e Gambellara,}
\rigo{Sai, la vita... è meno amara.}
\strofa
\rigo{Per combattere il maltempo in Valpadana}
\rigo{Tu sfogati con la tua damigiana,}
\rigo{E poi vedrai che tutto ti andrà bene,}
\rigo{Senza bisogno di antinebbia e di catene.}
\strofa
\rigo{Per combatter la calura in Valpadana}
\rigo{L'unico rimedio sai qual è? La damigiana!}
\rigo{Con Trebbiano e Moscatello,}
\rigo{Botticino e Rosatello,}
\rigo{Poi Lambrusco di Sorbara,}
\rigo{Sai, la vita... è meno amara.}
\strofa
\rigo{Quando arriva il Ferragosto in Valpadana}
\rigo{Rinfrescati con la tua damigiana,}
\rigo{E poi vedrai che tutto ti andrà bene,}
\rigo{Saran lontani gli antinebbia e le catene.}
\strofa
\rigo{Perché vada tutto bene in Valpadana}
\rigo{L'unico rimedio sai qual è? La damigiana!}
\rigo{Con Albano e Sangiovese,}
\rigo{Malvasia a più riprese,}
\rigo{Poi Barbera e Gambellara,}
\rigo{Sai, la vita... è meno amara.}
\strofa
\rigo{Per risolvere i problemi in Valpadana}
\rigo{Tu sfogati con la tua damigiana,}
\rigo{E poi vedrai che tutto ti andrà bene,}
\rigo{Senza bisogno di antinebbia e di catene.}
\strofa
\rigo{E poi vedrai che tutto ti andrà bene,}
\rigo{Saran lontani gli antinebbia e le catene.}
\par\endgroup

\begingroup\senzaaccordi
\cantotesta{90}{La Gigia l'è malada}{tradizionale}{\bRS}{382pt}
\rigo{La Gigia l’è malada,}
\rigo{Ohimè ohimè ohimè ohimè!}
\rigo{La Gigia l’è malada,}
\rigo{Ohimè ohimè ohimè ohimè!}
\strofa
\rigo{Se l’è malada, ciamée ciamée 'l dottore,}
\rigo{ciamée ‘l dottore che la guarirà.}
\rigo{Se l’è malada, ciamée ciamée 'l dottore,}
\rigo{ciamée ‘l dottore che la guarirà.}
\strofa
\rigo{Dottore l'entra in camera,}
\rigo{Ohimè ohimè ohimè ohimè!}
\rigo{Dottore l'entra in camera,}
\rigo{Ohimè ohimè ohimè ohimè!}
\strofa
\rigo{E l'entra in camera, ghe palpa palpa il polso:}
\rigo{la bella Gigia l'è malà d'amor.}
\rigo{E l'entra in camera, ghe palpa palpa il polso:}
\rigo{la bella Gigia l'è malà d'amor.}
\strofa
\rigo{O Gigia, o bella Gigia.}
\rigo{Ohimè ohimè ohimè ohimè!}
\rigo{O Gigia, bella Gigia.}
\rigo{Ohimè ohimè ohimè ohimè!}
\strofa
\rigo{Ma dimmi pure la verità:}
\rigo{chi l'è 'l papà di questo bel bambin?}
\rigo{Ma dimmi pure la verità:}
\rigo{chi l'è 'l papà di questo bel bambin?}
\par\endgroup

\begingroup
\cantotesta{91}{Lazy boy}{}{\bEG}{968pt}
\rigo{Lazy Boy}
\strofa
\rigo{C\ch{Lam}{'è} una stella bianca,}
\rigo{c\ch{Do}{he} splende su nel ciel,}
\rigo{g\ch{Lam}{uardando} solitaria}
\rigo{i\ch{Do}{} pascoli del Far West,}
\rigo{è\ch{Rem}{} la \ch{Mim}{gran} stella}
\rigo{d\ch{Fa}{el} \ch{Mim}{vecchio} Texas,}
\rigo{l\ch{Rem}{a} st\ch{Mim}{ella}\ch{Lam}{} del cow boy.}
\strofa
\rigo{Nell'Ottantasette,}
\rigo{brillando su Fort Joy,}
\rigo{vide venire al mondo}
\rigo{il piccolo Lazy Boy.}
\rigo{E qui comincia la lunga storia}
\rigo{del pallido cow-boy.}
\strofa
\rigo{RIT: DO LA-}
\rigo{Un dì la mamma gli disse va,}
\rigo{m\ch{Do}{a} r\ch{Sol}{esta}\ch{Lam}{} un bravo cow boy,}
\rigo{u\ch{Do}{n} d\ch{Lam}{ì} la tromba ti chiamerà,}
\rigo{D\ch{Do}{ani}\ch{Sol}{ele} \ch{Lam}{la} suonerà.}
\strofa
\rigo{Lazy vuol dire pigro,}
\rigo{ma il nostro Lazy Boy,}
\rigo{non era affatto pigro}
\rigo{lontano da fort Joy.}
\rigo{Ma un brutto giorno}
\rigo{conobbe Jessie,}
\rigo{conobbe la sua colt.}
\strofa
\rigo{Cominciò a sparare}
\rigo{ai sassi della via,}
\rigo{poi preferì i cavalli}
\rigo{e infine la ferrovia,}
\rigo{ai passeggieri,}
\rigo{bucò sombreri,}
\rigo{rubò pepite d'or.}
\strofa
\rigo{Tutti gli sceriffi, cercavan Lazy Boy,}
\rigo{la taglia era grossa,}
\rigo{piaceva al sergente Roy,}
\rigo{ma ad una spanna, trovò la canna}
\rigo{del pallido cow boy.}
\strofa
\rigo{Stava per sparare}
\rigo{il nostro Lazy Boy,}
\rigo{ma udì da una missione,}
\rigo{un coro di Little Boy:}
\rigo{"Oh lode, oh glory, oh alleluia}
\rigo{e il colpo non partì.}
\strofa
\rigo{C'è una stella bianca,}
\rigo{che splende su nel ciel,}
\rigo{guardando solitaria}
\rigo{i pascoli del Far West}
\rigo{e qui finisce}
\rigo{la lunga storia}
\rigo{del pallido cow boy.}
\par\endgroup

\begingroup\senzaaccordi
\cantotesta{92}{Me compare Giacometo}{tradizionale}{\bEG}{210pt}
\rigo{Me compare Giacometo}
\rigo{El gh'aveva un bel galeto}
\rigo{Quando canta al verdail bevo}
\rigo{Ch'el fa propi innamorar}
\strofa
\rigo{Un bel giorno la padrona}
\rigo{Per far festa agli invitati}
\rigo{La g'ha tira il collo al gallo}
\rigo{E lo mette a cucinar}
\strofa
\rigo{Le galline tutte matte}
\rigo{Per la perdita del gallo}
\rigo{Le g'ha rotto el caponaro}
\rigo{De la rabbia che le g'ha.}
\par\endgroup

\begingroup
\cantotesta{93}{Non si va in cielo}{}{\bLC}{417pt}
\rigo{\ch{Sol}{Non} si va in cielo in bicicl\ch{Re}{etta}}
\rigo{perché in cielo si va un po’ più in frett\ch{Sol}{a...}}
\rigo{Non si va in ciel\ch{Do}{o} in bicicletta}
\rigo{perché in cie\ch{Sol}{lo} si va un po’ più in frett\ch{Re}{a..}\ch{Sol}{.}   \ch{Do}{}   \ch{Sol}{}}
\strofa
\rigo{Non si va in cielo in Pininfarina}
\rigo{perché in cielo non c’è la benzina.}
\strofa
\rigo{Non si va in cielo con la vestaglia}
\rigo{perché in cielo si lavora a maglia.}
\strofa
\rigo{Non si va in cielo con il portafoglio}
\rigo{perché in cielo non c’è imbroglio.}
\strofa
\rigo{Non si va in cielo con i cardinali}
\rigo{perché in cielo sono tutti uguali.}
\strofa
\rigo{Non si va in cielo con i capitalisti}
\rigo{perché in cielo sono comunisti.}
\strofa
\rigo{Se tu vai in cielo prima di me}
\rigo{fai un buchino e tira su anche me.}
\strofa
\rigo{Se tu vai all’inferno prima di me}
\rigo{tappa tutti i buchi che io non passi più.}
\par\endgroup

\begingroup
\cantotesta{94}{Oleanna}{tradizionale (Norvegia)}{\bLC}{451pt}
\rigo{\ch{Re}{Ole-Oleanna} vecchia\ch{Sol}{} terra del Far-W\ch{Re}{est}}
\rigo{f\ch{Sol}{ra} le terre\ch{Re}{} la più strana tutt\ch{La}{a} ti vorrei per me.\ch{Re}{}}
\strofa
\commento{Rit.}
\rigo{O\ch{Re}{le-Oleanna,} Ole-Ole\ch{Sol}{anna-na,}        \ch{Re}{}}
\rigo{O\ch{Sol}{le-Ole,} Ole\ch{Re}{-Ole,} Ole-O\ch{La}{leanna-na.}\ch{Re}{}}
\strofa
\rigo{Sheeee...nengo il caballero}
\rigo{quando ad Oleanna fu}
\rigo{mise in piedi un grattacielo}
\rigo{con le fondamenta in su.}
\strofa
\commento{Rit.}
\strofa
\rigo{Le mucche campagnole}
\rigo{tutte migrano in città}
\rigo{vi frequentano le scuole}
\rigo{dove imparano a ballar.}
\strofa
\commento{Rit.}
\strofa
\rigo{Se vino tu vuoi bere}
\rigo{che ti faccia inebriar}
\rigo{dalle pompe di un pompiere}
\rigo{non ce l’hai che da succhiar.}
\strofa
\commento{Rit.}
\par\endgroup

\begingroup\senzaaccordi
\cantotesta{95}{Si va}{}{\bEG}{541pt}
\rigo{Quando la neve scende giù lieve}
\rigo{non esitar si deve,}
\rigo{prendi il tuo sacco e lesto va}
\rigo{dove ti porterà.}
\strofa
\begin{ritornello}
\rigo{Si va, sulla montagna}
\rigo{dove la neve il volto ci abbronzerà.}
\rigo{L'ardor, che ci accompagna,}
\rigo{sarà la fiamma, che ci riscalderà.}
\rigo{Salir, sempre salir,}
\rigo{finché ogni valle canterà così: sì, sì!}
\rigo{Sciator, ripete il vento,}
\rigo{e ardimento il tuo motto sarà!}
\strofa
\end{ritornello}
\rigo{Signorinella, pallida e snella,}
\rigo{getta la tua gonnella}
\rigo{la tua pelliccia di vison,}
\rigo{mettiti i pantalon!}
\strofa
\begin{ritornello}
\rigo{Si va, sulla montagna}
\rigo{dove la neve il volto ci abbronzerà.}
\rigo{L'ardor, che ci accompagna,}
\rigo{sarà la fiamma, che ci riscalderà.}
\rigo{Salir, sempre salir,}
\rigo{finché ogni valle canterà così: sì, sì!}
\rigo{Sciator, ripete il vento,}
\rigo{e ardimento il tuo motto sarà!}
\strofa
\end{ritornello}
\rigo{Vecchio “scarpone” d'ogni emozione,}
\rigo{vinci la tentazione,}
\rigo{non riabbracciar più la mimi,}
\rigo{solo i tuoi vecchi sci!}
\strofa
\begin{ritornello}
\rigo{Si va, sulla montagna}
\rigo{dove la neve il volto ci abbronzerà.}
\rigo{L'ardor, che ci accompagna,}
\rigo{sarà la fiamma, che ci riscalderà.}
\rigo{Salir, sempre salir,}
\rigo{finché ogni valle canterà così: sì, sì!}
\rigo{Sciator, ripete il vento,}
\rigo{e ardimento il tuo motto sarà!}
\end{ritornello}
\par\endgroup

\begingroup\senzaaccordi
\cantotesta{96}{Sul cappello un bel fior}{}{\bEG}{350pt}
\rigo{Voi che passare ci vedete}
\rigo{sotto la pioggia o del sol al raggio}
\rigo{qualche volta penserete}
\rigo{che ci vuole un bel coraggio}
\rigo{per andare sempre così}
\rigo{pur se ripida è la via:}
\rigo{voi V.P. non ci capite}
\strofa
\rigo{perché certo non avete...}
\strofa
\rigo{Rit. Sul cappello un bel fior}
\rigo{sulle labbra una canzon}
\rigo{un cuor gioioso e sincero}
\rigo{questo solo ci vuol}
\rigo{per un bravo Esplorator}
\rigo{per girare il mondo intero.}
\strofa
\rigo{Oh potessimo nei boschi}
\rigo{ogni nostro dì passare}
\rigo{lieti poi sotto le tende}
\rigo{alla sera riposare.}
\strofa
\rigo{Dalle vette delle Alpi}
\rigo{alle sponde del gran mare}
\rigo{con il nostro lieto canto}
\rigo{noi diremo a tutto il mondo...}
\par\endgroup

\begingroup\senzaaccordi
\cantotesta{97}{Vinassa, vinassa}{tradizionale}{\bRS}{440pt}
\rigo{Là nella valle, c'è un'osteria}
\rigo{è l'allegria, è l'allegria}
\rigo{là nella valle, c'è un'osteria}
\rigo{è l'allegria di noi alpin!}
\strofa
\rigo{E se son pallida né miei colori}
\rigo{no' vo' dotori, no' vo' dotori}
\rigo{e se son pallida come 'na strassa}
\rigo{vinassa, vinassa e fiaschi de vin!}
\strofa
\rigo{Là sul Cervino c'è una colonna}
\rigo{è la Madonna, è la Madonna}
\rigo{là sul Cervino c'è una colonna}
\rigo{è la Madonna di noi alpin!}
\strofa
\rigo{Là nella valle c'è un filo d'erba}
\rigo{l'è la riserva, l'è la riserva.}
\rigo{Là nella valle c'è un filo d'erba}
\rigo{l'è la riserva di noi alpin.}
\strofa
\rigo{Là nella valle c'è la Rosina}
\rigo{l'è la rovina, l'è la rovina.}
\rigo{Là nella valle c'è la Rosina}
\rigo{l'è la rovina di noi alpin.}
\strofa
\rigo{Là nella valle c'è un buco nero}
\rigo{l'è il cimitero, l'è il cimitero}
\rigo{Là nella valle c'è un buco nero}
\rigo{l'è il cimitero di noi alpin.}
\strofa
\rigo{Là su quel monte c'è una lanterna}
\rigo{requiem aeternam, requiem aeternam,}
\rigo{Là su quel monte c'è una lanterna}
\rigo{requiem aeternam per noi alpin.}
\par\endgroup

\sottosezione{Calando}

\begingroup
\cantotesta{98}{Ah io vorrei tornare}{}{\bRS}{403pt}
\rigo{A\ch{Sol}{h,} i\ch{Do}{o} v\ch{Lam}{orrei} tornare anche solo per un dì}
\rigo{L\ch{Re}{ass}\ch{Sol}{ù} nella valle alpina.}
\rigo{Lassù tra gli alti abeti ed i rododendri in fior}
\rigo{Distendermi a terra e sognar.}
\strofa
\rigo{P\ch{Sol}{orta}\ch{Do}{mi} \ch{Sol}{tu} lassù, o Signor}
\rigo{D\ch{Re}{ove} meglio ti veda…}
\strofa
\rigo{P\ch{Sol}{orta}\ch{Do}{mi} \ch{Lam}{tra} il verde dei tuoi pascoli, lassù}
\rigo{P\ch{Re}{er} \ch{Sol}{non} farmi scender mai più (2 vv)}
\strofa
\rigo{Là, sotto il pino antico noi lasciammo nel partir}
\rigo{La croce del nostro altare}
\rigo{Ma sotto il pino antico con la croce là restò}
\rigo{Un poco del nostro cuor.}
\strofa
\rigo{Portami tu …}
\strofa
\rigo{E quando quest'inverno qui la neve scenderà}
\rigo{Bianca sarà la valle}
\rigo{Ma sotto il pino antico un bel giglio fiorirà}
\rigo{Il giglio dell'esplorator.}
\par\endgroup

\begingroup
\cantotesta{99}{Bocca di rosa}{Fabrizio De André}{\bRS}{1355pt}
\rigo{La \ch{Lam}{chiamavano} Bocca di Rosa, metteva \ch{Mi}{l'amore,} metteva \ch{Lam}{l'amore}}
\rigo{la \ch{Lam}{chiamavano} Bocca di Rosa, metteva \ch{Mi}{l'amore} sopra ogni \ch{Lam}{cosa}}
\rigo{appena \ch{Lam}{scese} alla stazione del \ch{Mi}{paesino} di \ch{Lam}{sant'Ilario}}
\rigo{tutti si \ch{Lam}{accorsero} con uno sguardo che non si \ch{Mi}{trattava} di un \ch{Lam}{missionario.}}
\rigo{C'è chi \ch{La7}{l'amore} lo fa per \ch{Rem}{noia,} chi se lo \ch{Sol7}{sceglie} per \ch{Do}{professione,}}
\rigo{Bocca di \ch{Rem}{Rosa} né l'uno né \ch{Lam}{l'altro,} lei lo \ch{Mi}{faceva} per \ch{Lam}{passione.}}
\strofa
\rigo{Ma la \ch{Lam}{passione} spesso conduce a soddisfare \ch{Mi}{le} proprie \ch{Lam}{voglie}}
\rigo{senza \ch{Lam}{indagare} se il concupito ha il cuore \ch{Mi}{libero} oppure ha \ch{Lam}{moglie.}}
\rigo{E fu \ch{Lam}{così} che da un giorno all'altro Bocca di \ch{Mi}{Rosa} si tirò \ch{Lam}{addosso}}
\rigo{l'ira \ch{Lam}{funesta} delle cagnette a cui \ch{Mi}{aveva} sottratto \ch{Lam}{l'osso.}}
\rigo{Ma le \ch{La7}{comari} di un \ch{Rem}{paesino} non brillano \ch{Sol7}{certo} in \ch{Do}{iniziativa,}}
\rigo{le contromisure \ch{Rem}{fino} a quel \ch{Lam}{punto} si \ch{Mi}{limitavano} \ch{Lam}{all'invettiva.}}
\strofa
\rigo{Si sa che la \ch{Lam}{gente} dà buoni consigli, sentendosi \ch{Mi}{come} Gesù nel \ch{Lam}{tempio,}}
\rigo{si sa che la \ch{Lam}{gente} dà buoni consigli se non può più \ch{Mi}{dare} cattivo \ch{Lam}{esempio.}}
\rigo{Così una \ch{Lam}{vecchia} mai stata moglie, senza mai \ch{Mi}{figli} senza più \ch{Lam}{voglie,}}
\rigo{si prese la \ch{Lam}{briga} e di certo il gusto di dare \ch{Mi}{a} tutte il consiglio \ch{Lam}{giusto.}}
\rigo{E \ch{La7}{rivolgendosi} alle \ch{Rem}{cornute} le apostrofò \ch{Sol7}{con} parole \ch{Do}{argute:}}
\rigo{“Il furto \ch{Rem}{d'amore} sarà \ch{Lam}{punito”} disse \ch{Mi}{“dall'ordine} \ch{Lam}{costituito”.}}
\strofa
\rigo{E quelle \ch{Lam}{andarono} dal commissario e dissero \ch{Mi}{senza} \ch{Lam}{parafrasare:}}
\rigo{“Quella \ch{Lam}{schifosa} ha già troppi clienti più di un \ch{Mi}{consorzio} \ch{Lam}{alimentare”.}}
\rigo{Ed \ch{Lam}{arrivarono} quattro gendarmi con i \ch{Mi}{pennacchi,} con i \ch{Lam}{pennacchi,}}
\rigo{ed \ch{Lam}{arrivarono} quattro gendarmi con i \ch{Mi}{pennacchi} e con le \ch{Lam}{armi.}}
\rigo{Spesso gli \ch{La7}{sbirri} e i carabinieri \ch{Rem}{al} proprio \ch{Sol7}{dovere} vengono \ch{Do}{meno}}
\rigo{ma non quando \ch{Rem}{sono} in alta \ch{Lam}{uniforme} e l'accompagnarono \ch{Mi}{al} primo \ch{Lam}{treno.}}
\strofa
\rigo{Alla \ch{Lam}{stazione} c'erano tutti, dal \ch{Mi}{commissario} al \ch{Lam}{sagrestano,}}
\rigo{alla \ch{Lam}{stazione} c'erano tutti, con gli occhi \ch{Mi}{rossi} e il cappello in \ch{Lam}{mano.}}
\rigo{A \ch{Lam}{salutare} chi per un poco, senza \ch{Mi}{pretese,} senza \ch{Lam}{pretese,}}
\rigo{a \ch{Lam}{salutare} chi per un poco portò \ch{Mi}{l'amore} nel \ch{Lam}{paese}}
\rigo{c'era un \ch{La7}{cartello} \ch{Rem}{giallo} con una \ch{Sol7}{scritta} \ch{Do}{nera,}}
\rigo{diceva: \ch{Rem}{“Addio,} Bocca di \ch{Lam}{Rosa,} con te se ne \ch{Mi}{parte} la \ch{Lam}{primavera”.}}
\strofa
\rigo{Ma una \ch{Lam}{notizia} un po' originale non ha \ch{Mi}{bisogno} di alcun \ch{Lam}{giornale,}}
\rigo{come una \ch{Lam}{freccia} dall'arco scocca, vola \ch{Mi}{veloce} di bocca in \ch{Lam}{bocca.}}
\rigo{E alla \ch{Lam}{stazione} successiva molta più \ch{Mi}{gente} di quando \ch{Lam}{partiva,}}
\rigo{chi manda un \ch{Lam}{bacio,} chi getta un fiore, chi si \ch{Mi}{prenota} per due \ch{Lam}{ore.}}
\rigo{Persino il \ch{La7}{parroco} che non \ch{Rem}{disprezza} tra un \ch{Sol7}{Miserere} e un'estrema \ch{Do}{unzione}}
\rigo{il bene \ch{Rem}{effimero} della \ch{Lam}{bellezza,} la vuole \ch{Mi}{accanto} in \ch{Lam}{processione:}}
\rigo{e con la \ch{Lam}{vergine} in prima fila, e Bocca di \ch{Mi}{Rosa} poco \ch{Lam}{lontano,}}
\rigo{si porta a \ch{Lam}{spasso} per il paese l'amore \ch{Mi}{sacro} e l'amor \ch{Lam}{}\ch{Mi}{}\ch{Lam}{profano.}}
\par\endgroup

\begingroup\senzaaccordi
\cantotesta{100}{Canto del ricordo}{}{\bEG}{301pt}
\rigo{Quando tramonta il sol B.P. tu torni tra noi,}
\rigo{quando riflette la fiamma le ombre di noi}
\rigo{che cantiamo in cor; tu torni B.P.,}
\strofa
\rigo{tu torni tra noi ci guardi non visto}
\rigo{da un angolo d'ombra cantando con noi.}
\strofa
\rigo{E tu rivedi a sera i fuochi dei militar,}
\rigo{rivedi i fuochi, le tende, i cavalli, le veglie,}
\rigo{i deserti, il mar. Tu torni B.P.,}
\strofa
\rigo{tu torni a cantar; confuso nel vento}
\rigo{si perde il tuo canto col nostro che va.}
\strofa
\rigo{(... fischio... )}
\rigo{Quando tramonta il sol B.P., tu torni fra noi,}
\strofa
\rigo{quando riflette la fiamma le ombre di noi}
\rigo{che cantiamo in cor; tu torni B.P.}
\strofa
\rigo{tu torni fra noi, ritorni ogni sera}
\rigo{cantando sull'eco, dei canti d'allor.}
\par\endgroup

\begingroup
\cantotesta{101}{Canto delle Aquile Randagie (La luna delle vette)}{}{\bEG}{280pt}
\rigo{\ch{Mi}{}\ch{Si}{}\ch{Do\#m}{La} luna che risplende inonda di luce}
\rigo{\ch{Mi}{}\ch{Si}{}\ch{Mi}{le} vette che scintillano lassù}
\rigo{\ch{Mi}{}\ch{Si}{}\ch{Do\#m}{La} nenia che cantiamo sull’ali del sogno}
\rigo{\ch{Mi}{}\ch{Si}{}\ch{Mi}{lontano} porta i cuori e fa sognar}
\rigo{\ch{Mi}{}\ch{Sol\#}{}\ch{Do\#m}{}\ch{Fa\#m}{}\ch{Si}{}\ch{Mi}{Sognar} lontani dì l’antica libertà}
\rigo{\ch{Fa\#m}{}\ch{Si}{}\ch{Mi}{}\ch{Fa\#m}{}\ch{Si}{}\ch{Mi}{del} tempo che già fu del tempo che sarà.}
\strofa
\rigo{Lontano ci risponde lo scroscio del fiume}
\rigo{che scorre tra le rocce con fragor}
\rigo{Sotto un manto di stelle la fiamma s’innalza}
\rigo{guizzando verso il ciel finché muor}
\rigo{Ma mai non può morir non morirà mai più}
\rigo{la fiamma che ravviva la nostra gioventù}
\strofa
\rigo{Ma non morirà mai più.}
\par\endgroup

\begingroup
\cantotesta{102}{Canto di Montegemoli}{}{\bEG}{474pt}
\rigo{\ch{Sol}{È} segnata la mia strada}
\rigo{\ch{Re}{da} chi prima la passò}
\rigo{\ch{Sol}{cercherò} in ogni contrada}
\rigo{\ch{Re}{ed} \ch{Sol}{il} messaggio scoprirò.}
\strofa
\begin{ritornello}
\rigo{\ch{Sol}{A} chi mi accoglierà}
\rigo{\ch{Re}{la} \ch{Sol}{pace} porterò}
\rigo{\ch{Sol}{a} chi mi ascolterà}
\rigo{\ch{Re}{di} \ch{Sol}{Lui} io parlerò}
\rigo{\ch{Sol}{la} gioia che ho con me}
\rigo{\ch{Re}{la} \ch{Sol}{voglio} far sentir}
\rigo{\ch{Sol}{a} tutti e a quello che}
\rigo{\ch{Re}{con} \ch{Sol}{me} vorrà venir.}
\strofa
\end{ritornello}
\rigo{Per servire ho camminato}
\rigo{so il messaggio dell’amor}
\rigo{coi fratelli che ho trovato}
\rigo{sarà lieve anche il dolor.}
\strofa
\begin{ritornello}
\rigo{\ch{Sol}{A} chi mi accoglierà}
\rigo{\ch{Re}{la} \ch{Sol}{pace} porterò}
\rigo{\ch{Sol}{a} chi mi ascolterà}
\rigo{\ch{Re}{di} \ch{Sol}{Lui} io parlerò}
\rigo{\ch{Sol}{la} gioia che ho con me}
\rigo{\ch{Re}{la} \ch{Sol}{voglio} far sentir}
\rigo{\ch{Sol}{a} tutti e a quello che}
\rigo{\ch{Re}{con} \ch{Sol}{me} vorrà venir.}
\end{ritornello}
\par\endgroup

\begingroup\senzaaccordi
\cantotesta{103}{Canto Eurojam 94}{F. Piazza}{\bEG}{300pt}
\rigo{Canta sonoro il fringuello}
\rigo{La luna fa crescere i mari}
\rigo{Umile scorre il ruscello}
\rigo{Svettano picchi solitari}
\strofa
\rigo{Ma il mio posto dov’è}
\rigo{In questa immensa armonia del creato}
\rigo{Solo ascoltando te}
\rigo{io capirò perché sono nato}
\strofa
\rigo{Cerco il tuo volto signore}
\rigo{È il volto d’ogni fratello}
\rigo{Solo così nell’amore}
\rigo{lascerò il mondo più bello}
\strofa
\rigo{E anche il fuoco sarà}
\rigo{La nostra lode la nostra preghiera}
\rigo{Resta con noi Signore}
\rigo{Perché si fa sera}
\strofa
\rigo{Resta con noi Signore}
\rigo{Perché si fa sera}
\par\endgroup

\begingroup\senzaaccordi
\cantotesta{104}{Carnet di marcia}{}{\bEG}{493pt}
\rigo{Guardo nei campi brulli le stoppie aride}
\rigo{e nel canneto osservo levarsi un vol.}
\rigo{Mi chiedo che fanno queste cose intorno}
\rigo{è un sogno o un inganno}
\rigo{questa vita accanto a me.}
\rigo{Sei tu Signor che ti nascondi cercano Te.}
\rigo{La mia tendina chiara spicca tra gli alberi}
\rigo{nella radura erbosa declina il dì.}
\rigo{trattiene il respiro ogni cosa intorno}
\rigo{il fuoco che miro raccoglie tutto in sè.}
\rigo{Marcio con zaino in spalla per valli insolite}
\rigo{divido il pane e l'acqua con il fratel.}
\rigo{La gente che vedo mi ridà il saluto}
\rigo{le cose in cui credo son concrete accanto a me.}
\rigo{Dolci ricordi tornano}
\rigo{Ah io vorrei tornare anche solo per un dì}
\rigo{lassù nella valle alpina,}
\rigo{lassù tra gli alti abeti ed i rododendri in fior}
\rigo{distendermi a terra e sognar.}
\strofa
\rigo{Portami tu lassù, Signor}
\rigo{dove meglio ti veda!}
\rigo{Oh portami tra il verde}
\rigo{dei tuoi pascoli lassù}
\rigo{per non farmi scendere mai più.}
\rigo{Là, sotto il pino antico noi lasciammo nel partir}
\rigo{la croce del nostro altare}
\rigo{là, sotto il pino antico con la croce là restò}
\rigo{un poco del nostro cuor.}
\rigo{E quando quest'inverno qui la neve scenderà}
\rigo{bianca sarà la valle,}
\rigo{ma sopra quella croce un bel giglio fiorirà}
\rigo{il giglio dell'esplorator.}
\par\endgroup

\begingroup\senzaaccordi
\cantotesta{105}{Col cappellone}{}{\bEG}{368pt}
\rigo{Col cappellone e un giglio d'or}
\rigo{sempre restiamo esplorator.}
\rigo{Se l'Asci sciolta non morirem}
\rigo{con voce franca cantiamo assiem.}
\strofa
\begin{ritornello}
\rigo{L'esploratore tenace resterà}
\rigo{e la promessa sua non tradirà,}
\rigo{ma forte ognor, fedele ancor}
\rigo{all'ideale che non muor.}
\end{ritornello}
\rigo{L'Asci un bel giorno il cuore ci rapì}
\rigo{e tanta gioia ci portò così}
\rigo{che se un bel dì l'Asci risorgerà,}
\rigo{tutti compatti ci troverà.}
\rigo{E se un bel dì l'Asci risorgerà,}
\rigo{tutti compatti ci troverà.}
\strofa
\rigo{Su per i monti lieti saliam}
\rigo{e in faccia al sole l'urlo lanciam.}
\rigo{Dall'alte cime si leverà}
\rigo{il nostro canto che ancor dirà:}
\rigo{L'esploratore tenace resterà...}
\strofa
\rigo{Quando quell'ora udrem suonar}
\rigo{e l'Asci ancora potrà marciar,}
\rigo{gigliate fiamme, vecchi guidon}
\rigo{sventoleranno tra le canzon.}
\rigo{L'esploratore tenace resterà...}
\par\endgroup

\begingroup\senzaaccordi
\cantotesta{106}{Colore del sole}{}{\bRS}{494pt}
\rigo{Colore del sol, più giallo dell'oro fino}
\rigo{la gioia che brilla negli occhi tuoi fratellino;}
\rigo{colore di un fior, il primo della mattina,}
\rigo{in te sorellina vedo rispecchiar;}
\rigo{nel gioco che fai non sai ma già stringi forte}
\rigo{segreti che sempre avrai nelle mani in ogni tua sorte;}
\rigo{nel canto che sai la voce tua chiara dice:}
\rigo{famiglia felice oggi nel mondo poi su nel ciel}
\rigo{per sempre sarà.}
\strofa
\rigo{Colore del grano che presto avrà d'or splendore,}
\rigo{la verde tua tenda fra gli alberi esploratore;}
\rigo{colore lontano di un prato che al cielo grida}
\rigo{sorella mia guida con me scoprirai...}
\strofa
\rigo{Dove il tuo sentiero ti porta in un mondo ignoto}
\rigo{che aspetta il tuo passo come di un cavaliere fidato;}
\rigo{dove il tuo pensiero dilata questa natura}
\rigo{verso l'avventura che ti farà capire te stessa}
\rigo{e amare il Signor.}
\strofa
\rigo{Colore rubino del sangue di un testimone,}
\rigo{il sangue che sgorga dal cuor tuo fratello rover;}
\rigo{color vespertino di nubi sul sol raccolte,}
\rigo{l'amore che Scolta tu impari a donar:}
\rigo{conosci quel bene che prima tu hai ricevuto}
\rigo{sai che non potrai tenerlo per te neppure un minuto;}
\rigo{conosci la gioia di spenderti in sacrificio}
\rigo{prometti servizio ad ogni fratello ovunque}
\rigo{la strada ti porterà.}
\par\endgroup

\begingroup\senzaaccordi
\cantotesta{107}{E mi la dona bionda}{Nanni Svampa}{\bRS}{738pt}
\rigo{E mi la donna bionda la vôri no,}
\rigo{E mi la donna bionda la vôri no:}
\rigo{Tutti i omem ghe fan la ronda}
\rigo{Tutti...}
\rigo{E mi la donna bionda la vôri no.}
\strofa
\rigo{E mi la donna nera la vôri no,}
\rigo{E mi...}
\rigo{Donna nera non è sincera}
\rigo{Donna ...}
\rigo{E mi...}
\strofa
\rigo{E mi la donna mora la vôri no,}
\rigo{E mi...}
\rigo{Donna mora l'è traditora}
\rigo{Donna...}
\rigo{E mi...}
\strofa
\rigo{E mi la donna rossa la vôri no,}
\rigo{E mi...}
\rigo{Se te tuca te dà la scossa}
\rigo{Se te...}
\rigo{E mi ...}
\strofa
\rigo{E mi la donna granda la vôri no,}
\rigo{E mi ...}
\rigo{La camina, la par'na stanga,}
\rigo{La cam...}
\rigo{E mi...}
\strofa
\rigo{E mi la donna riccia la vôri no,}
\rigo{E mi...}
\rigo{Suta ai ricci la g'à i capricci}
\rigo{Suta...}
\rigo{E mi...}
\strofa
\rigo{E mi la donna picula la vôri no,}
\rigo{E mi...}
\rigo{La camina, la par'na chichera,}
\rigo{La cam...}
\rigo{E mi...}
\strofa
\rigo{E mi la donna grassa la vôri no,}
\rigo{E mi...}
\rigo{Tutta quanta la te sconquassa}
\rigo{Tutta...}
\rigo{E mi...}
\strofa
\rigo{E mi la donna seca la vôri no,}
\rigo{E mi...}
\rigo{E la pesa comè 'na bisteca}
\rigo{E la pesa...}
\rigo{E mi...}
\strofa
\rigo{E mi la donna bella la vôri sì,}
\rigo{E mi...}
\rigo{La ga i oci che par'na stella}
\rigo{La gh'ha i oci che par'na stella}
\rigo{E mi...}
\par\endgroup

\begingroup\senzaaccordi
\cantotesta{108}{Faria faria oh}{Paul Doncoeur}{\bEG}{286pt}
\rigo{Per il mondo ha camminato, faria, faria oh!}
\rigo{Monti e mari ha già varcato faria, faria oh!}
\rigo{Ma gli è dolce continuar}
\rigo{rischi e incognite a sfidar}
\rigo{faria, faria, faria, faria, faria, faria oh!}
\rigo{Non gli pesa il borsellino, faria, faria oh!}
\rigo{Ma è signor del suo destino faria, faria oh!}
\rigo{Sempre lieto e pien d'ardor canta e va l'esplorator.}
\strofa
\rigo{Faria, faria, faria, faria, faria, faria oh!}
\rigo{Col fagotto e col bastone, faria, faria oh!}
\rigo{Parte senza colazione, faria, faria oh!}
\rigo{Lunga strada far dovrà molti amici troverà.}
\strofa
\rigo{Faria, faria, faria, faria, faria, faria oh!}
\rigo{Parte e sogna cose belle, faria, faria oh!}
\rigo{Albe, fiori, prati e stelle, faria, faria oh!}
\rigo{Quando Iddio lo chiamerà sorridendo tornerà}
\rigo{Faria, faria, faria, faria, faria, faria oh!}
\par\endgroup

\begingroup
\cantotesta{109}{Fiume rosso}{}{\bEG}{567pt}
\begin{ritornello}
\rigo{Il \ch{Re}{sole} tramonta a Ovest,}
\rigo{il \ch{Fa\#m}{bestiame} \ch{La}{discende} \ch{Sol}{al} \ch{Re}{fiume,}}
\rigo{Il tordo ritorna al suo \ch{Fa\#m}{nido,}}
\rigo{per un \ch{La}{cowboy} è \ch{Sol}{l'ora} di \ch{Re}{sognar.}}
\strofa
\end{ritornello}
\rigo{Luce rossa là nel \ch{Sol}{canyon}}
\rigo{ecco \ch{Re}{dove} devo an\ch{Sol}{dar;}}
\rigo{\ch{Re}{tornerò} \ch{Sol}{dal} mio amore}
\rigo{with \ch{Re}{my} rifle, \ch{Sol}{my} pony and \ch{Re}{me.}}
\strofa
\rigo{Il sole tramonta a Ovest...}
\strofa
\rigo{Non ruberò più cavalli,}
\rigo{caccerò i traditor,}
\rigo{tornerò dal mio amore}
\rigo{with my rifle, my pony and me.}
\strofa
\rigo{Il sole tramonta a Ovest...}
\strofa
\rigo{Un uccellin là sul salce}
\rigo{fa una dolce melodia,}
\rigo{tornerò dal mio amore}
\rigo{with my rifle, my pony and me.}
\strofa
\rigo{Il sole tramonta a Ovest...}
\strofa
\rigo{Appenderò il mio sombrero}
\rigo{quando il sole scenderà,}
\rigo{tornerò dal mio amore}
\rigo{with my rifle, my pony and me.}
\strofa
\rigo{Il sole tramonta a Ovest...}
\strofa
\rigo{(CODA)}
\rigo{Tornerò dal mio amore}
\rigo{with my rifle, my pony and me.}
\par\endgroup

\begingroup\senzaaccordi
\cantotesta{110}{Guarda lontano}{}{\bEG}{335pt}
\rigo{Sono gesti ispirati}
\rigo{Trasparenti colori:}
\rigo{C' è un pittore invisibile}
\rigo{che dipinge un nuovo mondo che}
\rigo{prende forma già davanti a noi;}
\rigo{la sua tela è la terra}
\rigo{la cornice il cielo blu.}
\rigo{Sono note di musica}
\rigo{nuovi suonì un pò magici:}
\rigo{un musicista invisibile}
\rigo{sta orchestrando una canzone che}
\rigo{si fa strada proprio grazie a lui;}
\rigo{toccherà nuove corde che}
\rigo{suoneranno dentro noi.}
\rigo{Uomini uomini d' oggi}
\rigo{che guardano lontano uomini}
\rigo{presi da un vento e portati in una mano.}
\rigo{E’ disegno insondabile}
\rigo{che da un senso di vertigine:}
\rigo{chi di noi è ripetibile?}
\rigo{Siamo tutti strumenti}
\rigo{che lascian chiaro un segno dietrodi sé}
\rigo{siamo fatti per vivere.}
\par\endgroup

\begingroup
\cantotesta{111}{Guardati intorno scout}{}{\bEG}{488pt}
\rigo{Guardati \ch{Fa}{intorno,} scout}
\rigo{il bosco si \ch{Do7}{ridesta}}
\rigo{salta di ramo in ramo}
\rigo{lo scoiattolo in \ch{Fa}{festa}}
\rigo{e par che ti ripeta:}
\strofa
\begin{ritornello}
\rigo{\ch{Sib}{Sei} sulla giusta \ch{Fa}{traccia,}}
\rigo{\ch{Sib}{prosegui} il tuo cammino.}
\rigo{\ch{Fa}{Fratello,} \ch{Do7}{buona} \ch{Fa}{caccia.}}
\strofa
\end{ritornello}
\rigo{Il sol splende radioso}
\rigo{là nell'azzurro cielo,}
\rigo{si libra in alto un'aquila}
\rigo{nell'aria senza velo.}
\rigo{Anch'essa par che dica:}
\strofa
\begin{ritornello}
\rigo{sei sulla giusta traccia,}
\rigo{prosegui il tuo cammino.}
\rigo{Fratello, buona caccia.}
\strofa
\end{ritornello}
\rigo{Quando tramonta il sole}
\rigo{e discende la notte}
\rigo{escono dalle tane}
\rigo{a caccia i lupi a frotte.}
\rigo{E anch'essi ti salutano:}
\strofa
\begin{ritornello}
\rigo{sei sulla giusta traccia,}
\rigo{prosegui il tuo cammino.}
\rigo{Fratello, buona caccia.}
\end{ritornello}
\par\endgroup

\begingroup
\cantotesta{112}{Il pescatore}{Fabrizio De André}{\bRS}{1164pt}
\rigo{All'ombra \ch{Do}{dell'ultimo} \ch{Sol}{}\ch{Do}{sole}}
\rigo{si era assopito \ch{Fa}{un} \ch{Do}{pescatore}}
\rigo{e aveva un \ch{Fa}{solco} \ch{Sol}{lungo} il \ch{Do}{viso}}
\rigo{\ch{Fa}{come} una \ch{Do}{specie} \ch{Sol}{di} \ch{Do}{sorriso.} \ch{Fa}{}\ch{Do}{}\ch{Sol}{}}
\strofa
\rigo{\ch{Do}{}\ch{Fa}{}\ch{Do}{}\ch{Sol}{}}
\strofa
\rigo{Venne alla \ch{Do}{spiaggia} \ch{Sol}{un} \ch{Do}{assassino}}
\rigo{due occhi \ch{Fa}{grandi} da \ch{Do}{bambino}}
\rigo{due occhi \ch{Fa}{enormi} \ch{Sol}{di} \ch{Do}{paura}}
\rigo{\ch{Fa}{eran} gli \ch{Do}{specchi} di \ch{Sol}{}\ch{Do}{un'avventura.}}
\strofa
\begin{ritornello}
\rigo{La La La \ch{Fa}{La} La La La \ch{Do}{La}}
\rigo{La La La \ch{Sol}{La} La La La \ch{Do}{La}}
\rigo{La La La \ch{Fa}{La} La \ch{Sol}{La} La \ch{Lam}{La}}
\rigo{\ch{Fa}{La} La La \ch{Do}{La} La \ch{Sol}{La} La \ch{Lam}{La}}
\strofa
\end{ritornello}
\rigo{E chiese al \ch{Do}{vecchio,} \ch{Sol}{"Dammi} il \ch{Do}{pane}}
\rigo{ho poco \ch{Fa}{tempo} e troppa \ch{Do}{fame"}}
\rigo{e chiese al \ch{Fa}{vecchio,} \ch{Sol}{"Dammi} il \ch{Do}{vino}}
\rigo{\ch{Fa}{ho} sete e \ch{Do}{sono} un \ch{Sol}{}\ch{Do}{assassino".} \ch{Fa}{}\ch{Do}{}\ch{Sol}{}}
\strofa
\rigo{\ch{Do}{}\ch{Fa}{}\ch{Do}{}\ch{Sol}{}}
\strofa
\rigo{Gli occhi \ch{Do}{dischiuse} il \ch{Sol}{vecchio} al \ch{Do}{giorno}}
\rigo{non si guardò \ch{Fa}{neppure} \ch{Do}{intorno}}
\rigo{ma versò il \ch{Fa}{vino} e \ch{Sol}{spezzò} il \ch{Do}{pane}}
\rigo{\ch{Fa}{per} chi \ch{Do}{diceva} ho \ch{Sol}{sete} e ho \ch{Do}{fame.}}
\strofa
\begin{ritornello}
\rigo{La La La \ch{Fa}{La} La La La \ch{Do}{La}}
\rigo{La La La \ch{Sol}{La} La La La \ch{Do}{La}}
\rigo{La La La \ch{Fa}{La} La \ch{Sol}{La} La \ch{Lam}{La}}
\rigo{\ch{Fa}{La} La La \ch{Do}{La} La \ch{Sol}{La} La \ch{Lam}{La}}
\strofa
\end{ritornello}
\rigo{E fu il \ch{Do}{calore} \ch{Sol}{di} un \ch{Do}{momento}}
\rigo{poi via di \ch{Fa}{nuovo} verso il \ch{Do}{vento}}
\rigo{poi via di \ch{Fa}{nuovo} \ch{Sol}{verso} il \ch{Do}{sole}}
\rigo{\ch{Fa}{dietro} alle \ch{Do}{spalle} un \ch{Sol}{}\ch{Do}{pescatore.} \ch{Fa}{}\ch{Do}{}\ch{Sol}{}}
\strofa
\rigo{\ch{Do}{}\ch{Fa}{}\ch{Do}{}\ch{Sol}{}}
\strofa
\rigo{Dietro alle \ch{Do}{spalle} \ch{Sol}{un} \ch{Do}{pescatore}}
\rigo{e la \ch{Fa}{memoria} è già \ch{Do}{dolore}}
\rigo{è già il \ch{Fa}{rimpianto} \ch{Sol}{di} un \ch{Do}{aprile}}
\rigo{\ch{Fa}{giocato} \ch{Do}{all'ombra} di \ch{Sol}{un} \ch{Do}{cortile.}}
\strofa
\begin{ritornello}
\rigo{La La La \ch{Fa}{La} La La La \ch{Do}{La}}
\rigo{La La La \ch{Sol}{La} La La La \ch{Do}{La}}
\rigo{La La La \ch{Fa}{La} La \ch{Sol}{La} La \ch{Lam}{La}}
\rigo{\ch{Fa}{La} La La \ch{Do}{La} La \ch{Sol}{La} La \ch{Lam}{La}}
\strofa
\end{ritornello}
\rigo{Vennero in \ch{Do}{sella} \ch{Sol}{due} \ch{Do}{gendarmi}}
\rigo{vennero in \ch{Fa}{sella} con le \ch{Do}{armi}}
\rigo{chiesero al \ch{Fa}{vecchio} \ch{Sol}{se} lì \ch{Do}{vicino}}
\rigo{\ch{Fa}{fosse} \ch{Do}{passato} un \ch{Sol}{}\ch{Do}{assassino.} \ch{Fa}{}\ch{Do}{}\ch{Sol}{}}
\strofa
\rigo{\ch{Do}{}\ch{Fa}{}\ch{Do}{}\ch{Sol}{}}
\strofa
\rigo{Ma all'ombra \ch{Do}{dell'ultimo} \ch{Sol}{}\ch{Do}{sole}}
\rigo{si era assopito \ch{Fa}{il} \ch{Do}{pescatore}}
\rigo{e aveva un \ch{Fa}{solco} \ch{Sol}{lungo} il \ch{Do}{viso}}
\rigo{\ch{Fa}{come} una \ch{Do}{specie} \ch{Sol}{di} \ch{Do}{sorriso.}}
\rigo{e aveva un \ch{Fa}{solco} \ch{Sol}{lungo} il \ch{Do}{viso}}
\rigo{\ch{Fa}{come} una \ch{Do}{specie} \ch{Sol}{di} \ch{Do}{sorriso.}}
\strofa
\begin{ritornello}
\rigo{La La La \ch{Fa}{La} La La La \ch{Do}{La}}
\rigo{La La La \ch{Sol}{La} La La La \ch{Do}{La}}
\rigo{La La La \ch{Fa}{La} La \ch{Sol}{La} La \ch{Lam}{La}}
\rigo{\ch{Fa}{La} La La \ch{Do}{La} La \ch{Sol}{La} La \ch{Lam}{La}}
\strofa
\end{ritornello}
\rigo{\ch{La}{} \ch{La}{} \ch{La}{} \ch{La}{} \ch{La}{} \ch{La}{} \ch{La}{} \ch{La}{}}
\rigo{\ch{La}{} \ch{La}{} \ch{La}{} \ch{La}{} \ch{La}{} \ch{La}{} \ch{La}{} \ch{La}{}}
\rigo{\ch{La}{} \ch{La}{} \ch{La}{} \ch{La}{} \ch{La}{} \ch{La}{} \ch{La}{} \ch{La}{}}
\rigo{\ch{La}{} \ch{La}{} \ch{La}{} \ch{La}{} \ch{La}{} \ch{La}{} \ch{La}{} \ch{La}{}}
\par\endgroup

\begingroup\senzaaccordi
\cantotesta{113}{Joska la rossa}{B. De Marzi}{\bRS}{888pt}
\rigo{El muro bianco, drio de la to casa}
\rigo{Ti te saltavi come un oseleto}
\rigo{(Jo) Joska la rossa, pèle de bombasa}
\rigo{Tute le sere prima de 'nà in leto}
\rigo{Te stavi lì, co le to scarpe rote}
\rigo{Te ne vardavi drio da j' oci mori}
\rigo{E te balavi alegra tuta note}
\rigo{E i baldi alpini te cantava i cori, oh}
\strofa
\begin{ritornello}
\rigo{Joska, Joska, Joska}
\rigo{Salta la mura!}
\rigo{Fin che la dura, oh}
\rigo{Joska, Joska, Joska}
\rigo{Salta la mura!}
\rigo{Balla con mi! Oh}
\rigo{Joska, Joska, Joska}
\rigo{Salta la mura!}
\rigo{Fin che la dura, oh}
\rigo{Joska, Joska, Joska}
\rigo{Salta la mura!}
\rigo{Balla con mi!}
\strofa
\end{ritornello}
\rigo{Ti te portavi el sole ogni matina}
\rigo{E de j 'Alpini te geri la morosa}
\rigo{(So) Sorela, mama, boca canterina}
\rigo{Oci del sol, meravigliosa rosa}
\rigo{Che tanto e tanto nù, ca te zerchémo}
\rigo{Joska la rossa, amor, rosa spanja}
\rigo{Ma dove sito andà? Ma dove andemo?}
\rigo{Semo ramenghi, o morti, e così sia, oh}
\strofa
\begin{ritornello}
\rigo{Joska, Joska, Joska}
\rigo{Salta la mura!}
\rigo{Fin che la dura, oh}
\rigo{Joska, Joska, Joska}
\rigo{Salta la mura!}
\rigo{Balla con mi! Oh}
\rigo{Joska, Joska, Joska}
\rigo{Salta la mura!}
\rigo{Fin che la dura, oh}
\rigo{Joska, Joska, Joska}
\rigo{Salta la mura!}
\rigo{Balla con mi!}
\strofa
\end{ritornello}
\rigo{Busa con crose, sarà sta i putèi?}
\rigo{La par 'na bara e invece xe 'na cuna}
\rigo{E dentro dorme tuti i to fradei}
\rigo{Fermi, impalà, co i oci ne la luna, oh}
\strofa
\begin{ritornello}
\rigo{Joska, Joska, Joska}
\rigo{Salta la mura!}
\rigo{Fin che la dura, oh}
\rigo{Joska, Joska, Joska}
\rigo{Salta la mura!}
\rigo{Fermate là! Oh}
\rigo{Joska, Joska, Joska}
\rigo{Salta la mura!}
\rigo{Fin che la dura, oh}
\rigo{Joska, Joska, Joska}
\rigo{Salta la mura!}
\rigo{Fermate là!}
\rigo{Joska, Joska, Joska}
\rigo{Salta la mura!}
\strofa
\end{ritornello}
\rigo{Finchè la dura! Oh}
\strofa
\begin{ritornello}
\rigo{Joska, Joska, Joska}
\rigo{Salta la mura!}
\rigo{Fermate là!}
\rigo{Fermate là!}
\end{ritornello}
\par\endgroup

\begingroup\senzaaccordi
\cantotesta{114}{Kamaludu}{}{\bEG}{158pt}
\rigo{Kamaludu, kamaludu,}
\rigo{ludu kamaludu, uandaisii}
\rigo{kamaludu, kamaludu uandaisii}
\rigo{Behbel ina, behbel ina,}
\rigo{ina behbel ina, uandaisi}
\rigo{behbel ina, behbel ina uandaisii}
\rigo{Matulefa, matulefa,}
\rigo{efa matulefa, uandaisi}
\rigo{matulefa, matulefa uandaisi.}
\par\endgroup

\begingroup\senzaaccordi
\cantotesta{115}{La ceseta de Transaqua}{}{\bRS}{331pt}
\rigo{As tu vist la ceseta}
\rigo{de Transacqua, de Transacqua}
\rigo{col Cimon de la Pala sóra i copi? Oh!}
\rigo{Te g'ha i oci ciari}
\rigo{come l'aqua, come l'aqua,}
\rigo{e i 'cavei tuti rizi e senza sgropi, oh!}
\strofa
\rigo{Mi g'ho vist la ceseta}
\rigo{de Transacqua, de Transacqua,}
\rigo{ma 'l Cimon de la Pala no ghe xera, oh!}
\rigo{Sóra i copi lustri}
\rigo{de tant'aqua, de tant'aqua,}
\rigo{se sgropava 'na nuvolona nera, oh!}
\strofa
\rigo{Nella cesa cjanta}
\rigo{messa il prete, messa il prete,}
\rigo{sul Cimon de la Pala fis'cia 'l vento, oh!}
\rigo{Cosa importa se g'ho}
\rigo{le scarpe rote, scarpe rote?}
\rigo{Mi te vardo e me sento il cor contento, oh!}
\strofa
\rigo{Cosa importa se g'ho}
\rigo{le scarpe rote, scarpe rote?}
\rigo{Mi te vardo e me sento il cor contento, oh!}
\par\endgroup

\begingroup
\cantotesta{116}{La guadalupa}{}{\bLC}{407pt}
\rigo{Coro 1: zim bum, zim bum bum, zim bum, zim bum bum (ad libitum)}
\strofa
\rigo{\ch{Do}{}\ch{Sol}{}\ch{Fa}{}\ch{Do}{Ho} lasciato da sol, sopra una scialuppa}
\rigo{\ch{Do}{}\ch{Sol}{}\ch{Fa}{}\ch{Do}{il} paese del sol della Guadalupa.}
\strofa
\begin{ritornello}
\rigo{\ch{Do}{}\ch{Sol}{}\ch{Fa}{}\ch{Sol}{}\ch{Do}{Ah!} Piango ancor la mia capanna e la savana}
\rigo{\ch{Do}{}\ch{Sol}{}\ch{Fa}{}\ch{Sol}{}\ch{Do}{ed} il paese dove sogno di tornar.}
\strofa
\end{ritornello}
\rigo{Se volevo mangiar, c’eran le banane;}
\rigo{se volevo dormir, c’erano le palme.}
\strofa
\begin{ritornello}
\rigo{\ch{Do}{}\ch{Sol}{}\ch{Fa}{}\ch{Sol}{}\ch{Do}{Ah!} Piango ancor la mia capanna e la savana}
\rigo{\ch{Do}{}\ch{Sol}{}\ch{Fa}{}\ch{Sol}{}\ch{Do}{ed} il paese dove sogno di tornar.}
\strofa
\end{ritornello}
\rigo{Sempre triste quaggiù sogno le savane;}
\rigo{il bel sol e l’amor e tante banane.}
\strofa
\begin{ritornello}
\rigo{\ch{Do}{}\ch{Sol}{}\ch{Fa}{}\ch{Sol}{}\ch{Do}{Ah!} Piango ancor la mia capanna e la savana}
\rigo{\ch{Do}{}\ch{Sol}{}\ch{Fa}{}\ch{Sol}{}\ch{Do}{ed} il paese dove sogno di tornar.}
\strofa
\end{ritornello}
\rigo{Si diffonde il “tam tam” lungo la pianura,}
\rigo{mentre il sole tramonta nella radura.}
\strofa
\begin{ritornello}
\rigo{\ch{Do}{}\ch{Sol}{}\ch{Fa}{}\ch{Sol}{}\ch{Do}{Ah!} Piango ancor la mia capanna e la savana}
\rigo{\ch{Do}{}\ch{Sol}{}\ch{Fa}{}\ch{Sol}{}\ch{Do}{ed} il paese dove sogno di tornar.}
\end{ritornello}
\par\endgroup

\begingroup
\cantotesta{117}{La guerra di Piero}{Fabrizio De André}{\bRS}{1088pt}
\rigo{\ch{Mi}{Dormi} sepolto in un \ch{Lam}{campo} di grano,}
\rigo{non e' la \ch{Sol}{rosa,} non e' il \ch{Do}{tulipano}}
\rigo{che ti fan \ch{Mi}{veglia} dall'ombra dei \ch{Lam}{fossi}}
\rigo{ma sono \ch{Rem}{mille} \ch{Mi}{papaveri} \ch{Lam}{rossi.}}
\strofa
\rigo{Lungo le \ch{Mi}{sponde} del mio \ch{Lam}{torrente}}
\rigo{voglio che \ch{Rem}{scendano} i lucci \ch{Lam}{argentati,}}
\rigo{non più i \ch{Sol}{cadaveri} dei \ch{Do}{soldati}}
\rigo{portati in \ch{Mi}{braccio} dalla \ch{Lam}{corrente.}}
\strofa
\rigo{Così \ch{Mi}{dicevi} ed era \ch{Lam}{d'Inverno} e come \ch{Rem}{gli} altri,}
\rigo{verso \ch{Lam}{l'inferno} te ne vai \ch{Sol}{triste}}
\rigo{come chi \ch{Do}{deve} ed il vento ti \ch{Mi}{sputa} in faccia \ch{Lam}{la} neve.}
\rigo{Fermati \ch{Mi}{Piero,} fermati \ch{Lam}{adesso,}}
\strofa
\rigo{lascia che il \ch{Sol}{vento} ti passi un po' \ch{Do}{addosso,}}
\rigo{dei morti in \ch{Mi}{battaglia} ti porti la \ch{Lam}{voce,}}
\rigo{chi diede la \ch{Rem}{vita} ebbe in \ch{Mi}{cambio} \ch{Lam}{una} croce.}
\rigo{Ma tu non lo \ch{Mi}{udisti} ed il tempo \ch{Lam}{passava}}
\strofa
\rigo{con le \ch{Rem}{stagioni} a passo di \ch{Lam}{giava}}
\rigo{ed \ch{Sol}{arrivasti} a passar la \ch{Do}{frontiera}}
\rigo{in un bel \ch{Mi}{giorno} di \ch{Lam}{Primavera.}}
\rigo{E mentre \ch{Mi}{marciavi} con l'anima in \ch{Lam}{spalle}}
\strofa
\rigo{vedesti un \ch{Rem}{uomo} in fondo alla \ch{Lam}{valle}}
\rigo{che aveva il tuo \ch{Sol}{stesso} identico \ch{Do}{umore}}
\rigo{ma la \ch{Mi}{divisa} di un altro \ch{Lam}{colore.}}
\rigo{Sparagli \ch{Mi}{Piero,} sparagli \ch{Lam}{ora} e dopo \ch{Sol}{un} colpo sparagli \ch{Do}{ancora,}}
\strofa
\rigo{fino a che \ch{Mi}{tu} non lo vedrai \ch{Lam}{esangue}}
\rigo{cadere in \ch{Rem}{terra} a \ch{Mi}{coprire} il \ch{Lam}{suo} sangue.}
\rigo{E se gli \ch{Mi}{sparo} in fronte o nel \ch{Lam}{cuore}}
\rigo{soltanto il \ch{Rem}{tempo} avrà per \ch{Lam}{morire,}}
\strofa
\rigo{ma il tempo a \ch{Sol}{me} resterà per \ch{Do}{vedere,}}
\rigo{vedere gli \ch{Mi}{occhi} d'un uomo che \ch{Lam}{muore.}}
\rigo{E mentre gli \ch{Mi}{usi} questa \ch{Lam}{premura} quello \ch{Rem}{si} volta,}
\rigo{ti vede, ha \ch{Lam}{paura} ed \ch{Sol}{imbracciata} \ch{Do}{l'artiglieria}}
\strofa
\rigo{non ti \ch{Mi}{ricambia} la \ch{Lam}{cortesia.}}
\rigo{Cadesti a \ch{Mi}{terra,} senza un \ch{Lam}{lamento}}
\rigo{e ti \ch{Rem}{accorgesti} in un \ch{Sol}{solo} \ch{Do}{momento}}
\rigo{che il tempo \ch{Mi}{non} ti sarebbe \ch{Lam}{bastato}}
\strofa
\rigo{a chieder \ch{Rem}{perdono} per \ch{Mi}{ogni} \ch{Lam}{peccato.}}
\rigo{Cadesti a \ch{Mi}{terra,} senza un \ch{Lam}{lamento}}
\rigo{e ti \ch{Rem}{accorgesti} in un \ch{Sol}{solo} \ch{Do}{momento}}
\rigo{che la tua \ch{Mi}{vita} finiva quel \ch{Lam}{giorno}}
\strofa
\rigo{e non ci \ch{Rem}{sarebbe} \ch{Mi}{stato} \ch{Lam}{ritorno.}}
\rigo{Ninetta \ch{Mi}{mia,} crepare di \ch{Lam}{Maggio}}
\rigo{ci vuole \ch{Rem}{tanto,} troppo \ch{Lam}{coraggio.}}
\rigo{Ninetta \ch{Sol}{bella} dritto \ch{Do}{all'Inferno}}
\strofa
\rigo{avrei \ch{Mi}{preferito} andarci in \ch{Lam}{Inverno.}}
\rigo{E mentre il \ch{Mi}{grano} ti stava a \ch{Lam}{sentire}}
\rigo{dentro le \ch{Rem}{mani} stringevi il \ch{Lam}{fucile,}}
\rigo{dentro la \ch{Sol}{bocca} stringevi \ch{Do}{parole}}
\strofa
\rigo{troppo \ch{Mi}{gelate} per sciogliersi al \ch{Lam}{sole.}}
\rigo{Dormi \ch{Mi}{sepolto} in un campo di \ch{Lam}{grano,}}
\rigo{non e' la \ch{Sol}{rosa,} non e' il \ch{Do}{tulipano}}
\rigo{che ti fan \ch{Mi}{veglia} dall'ombra dei \ch{Lam}{fossi}}
\strofa
\rigo{ma sono \ch{Rem}{mille} \ch{Mi}{papaveri} \ch{Lam}{rossi.}}
\par\endgroup

\begingroup\senzaaccordi
\cantotesta{118}{La montanara}{T. Ortelli}{\bRS}{361pt}
\rigo{La su per le montagne,}
\rigo{fra boschi e valli d'or,}
\rigo{tra l'aspre rupi echeggia}
\rigo{un cantico d'amor.}
\rigo{La su per le montagne}
\rigo{fra boschi e valli d'or,}
\rigo{Tra l'aspre rupi echeggia}
\rigo{un cantico d'amor.}
\strofa
\begin{ritornello}
\rigo{“La montanara, ohè!”}
\rigo{si sente cantare,}
\rigo{cantiam la montanara}
\rigo{e chi non la sa?}
\rigo{La montanara ohe}
\rigo{si sente cantare.}
\rigo{Cantiam la montanara}
\rigo{e chi non la sa.}
\strofa
\end{ritornello}
\rigo{Là su sui monti}
\rigo{dai rivi d'argento}
\rigo{una capanna cosparsa di fior.}
\rigo{Era la piccola}
\rigo{dolce dimora}
\rigo{di Soreghina,}
\rigo{la figlia del Sol,}
\rigo{la figlia del Sol!}
\par\endgroup

\begingroup\senzaaccordi
\cantotesta{119}{La mula de Parenzo}{}{\bEG}{402pt}
\rigo{La mula de Parenzo}
\rigo{L'ha messo su bottega}
\rigo{E tutto la vendeva}
\rigo{E tutto la vendeva}
\strofa
\rigo{La mula de Parenzo}
\rigo{L'ha messo su bottega}
\rigo{E tutto la vendeva}
\rigo{Föra che 'l baccalà}
\strofa
\rigo{Perchè non m'ami più}
\strofa
\rigo{La me morosa vegia}
\rigo{La tengo di riserva}
\rigo{Ma quando spunta l'erba}
\rigo{La mando a pascolar}
\strofa
\rigo{La mando a pascolare}
\rigo{Insieme alle caprette}
\rigo{L'amor con le servette}
\rigo{Non lo farò mai più.}
\strofa
\rigo{Tutti mi dicono bionda}
\rigo{Ma io bionda non sono}
\rigo{Porto i capelli neri}
\rigo{Neri come il carbon.}
\strofa
\rigo{Se il mare fosse cacio}
\rigo{E i monti di polenta}
\rigo{Ohi, mamma che pociate}
\rigo{Polenta e baccalà.}
\par\endgroup

\begingroup\senzaaccordi
\cantotesta{120}{Pim pam}{}{\bEG}{287pt}
\rigo{Nella casa sulla montagna}
\rigo{un camino grande grande sta,}
\rigo{nel camino grande grande grande}
\rigo{un gran fuoco fuoco fuoco va.}
\strofa
\rigo{RIT: Percio' pim pam, le scarpe pim pam,}
\rigo{di notte fan sul sentiero di pietre grosse,}
\rigo{pim pam, le scarpe pim pam,}
\rigo{di notte fan sul sentiero cosi'.}
\strofa
\rigo{Nella casa sulla montagna}
\rigo{un signore grande grande sta,}
\rigo{nella casa verde viola e bianca}
\rigo{tante sedie rosse e gialle ha.}
\strofa
\rigo{RIT:...}
\strofa
\rigo{Nella casa sulla montagna}
\rigo{una sedia a tutti tutti da,}
\rigo{e ciascuno togliera' le scarpe}
\rigo{tutti insieme poi si danzera'.}
\par\endgroup

\begingroup\senzaaccordi
\cantotesta{121}{Sul pajon}{}{\bEG}{431pt}
\rigo{Il battaglione Aosta}
\rigo{sul pajon}
\rigo{è sempre sulle cime,}
\rigo{sul pajon}
\rigo{ma quando scende a valle}
\rigo{sul pajon}
\rigo{attente ragazzine!}
\strofa
\begin{ritornello}
\rigo{Sul pajon de la caserma}
\rigo{requiem aeternam e così sia,}
\rigo{va' a ramengo ti,}
\rigo{to pare, ta mare e ta zia}
\rigo{e la naja e compagnia,}
\rigo{sul pajon, sul pajon!}
\strofa
\end{ritornello}
\rigo{Il parroco di Aosta}
\rigo{sul pajon}
\rigo{l'ha detto predicando:}
\rigo{sul pajon}
\rigo{«Attente ragazzine}
\rigo{sul pajon}
\rigo{che il "Quarto" sta arrivando.»}
\rigo{Sul pajon...}
\strofa
\rigo{Ma una delle più belle}
\rigo{sul pajon}
\rigo{ha detto piano piano:}
\rigo{sul pajon}
\rigo{«Se il "Quarto" sta arrivando}
\rigo{sul pajon}
\rigo{è quello che vogliamo!»}
\rigo{Sul pajon...}
\par\endgroup

\begingroup
\cantotesta{122}{Terra della betulla}{}{\bEG}{369pt}
\rigo{\ch{Rem}{Terra} di betull\ch{Fa}{a} casa \ch{Rem}{del} castoro,    \ch{Fa}{}}
\rigo{l\ch{Sib}{à} dove errando \ch{Fa}{va} il lupo \ch{Solm}{ancora.}      \ch{Rem}{}}
\rigo{V\ch{Sib}{oglio} tornare a\ch{Fa}{ncor} al mio bel\ch{Solm}{} lago blu\ch{Sib}{.}               \ch{Rem}{}}
\strofa
\commento{Rit.}
\rigo{B\ch{Rem}{umbadiaidi,} bumbadia\ch{Lam}{idi,} bumbadiaidi bum\ch{Solm}{}                 \ch{Rem}{}}
\rigo{B\ch{Rem}{umbadiaidi,} bumbadia\ch{Lam}{idi,} bumbadiaidi bum\ch{Solm}{.}                \ch{Rem}{}}
\strofa
\rigo{La mia canoa scivola leggera}
\rigo{sulle lucenti vie del grande fiume.}
\rigo{Voglio tornare ancor al mio bel lago blu.}
\strofa
\commento{Rit.}
\strofa
\rigo{La tra gli abeti la luna appare,}
\rigo{il tuo bel viso , mamma, vedo ancora.}
\rigo{Voglio tornare ancora al mio bel lago blu.}
\strofa
\commento{Rit.}
\par\endgroup

\begingroup\senzaaccordi
\cantotesta{123}{Vento della sera}{}{\bEG}{217pt}
\rigo{Rit. Vento della sera tiepida e leggera,}
\rigo{in quest'atmosfera di serenità;}
\rigo{tutti attorno al fuoco, riposiamo un poco,}
\rigo{e cantiamo assieme le vecchie canzon.}
\strofa
\rigo{Cantiamo anche per voi fratelli lontan,}
\rigo{un giorno voi partendo ci deste la man;}
\rigo{qui attorno vi vediamo, voi siete ancora qui,}
\rigo{insieme a voi cantiamo, come cantammo un dì.}
\strofa
\rigo{Un giorno voi partiste lasciandoci qua,}
\rigo{il fuoco si spegneva nell'oscurità.}
\strofa
\rigo{Sentite: noi vi amiamo come vi amammo un dì:}
\rigo{tornate vi preghiamo, noi vi aspettiamo qui.}
\par\endgroup

\begingroup\senzaaccordi
\cantotesta{124}{Vento sottile}{}{\bLC}{216pt}
\rigo{Vento sottile, vento del mattino}
\rigo{Vento che scuoti la cima del mio pino}
\rigo{Vento che danzi, che balli la gioia tu ci porti}
\rigo{Vento sottile}
\rigo{Oh-oh-oh-oh}
\rigo{Oh-oh-oh-oh}
\rigo{La la la la la la la la la lalalala lala}
\rigo{Vento sottile}
\strofa
\rigo{Vent fin, vent du matin,}
\rigo{vent qui souffles au bout des sapins,}
\rigo{vent qui chantes, vent qui danses,}
\rigo{vent, vent fin.}
\par\endgroup

\begingroup\senzaaccordi
\cantotesta{125}{Voglio girare il mondo}{I girasoli}{\bRS}{305pt}
\rigo{Voglio girare il mondo}
\rigo{anche se la gente ride,}
\rigo{anche se tutti quanti}
\rigo{mi consigliano di restare.}
\strofa
\rigo{Rit: E voglio andarmene via}
\rigo{proprio perchè secondo loro}
\rigo{dormire sotto le stelle,}
\rigo{vuol dire solo sentir freddo.}
\strofa
\rigo{Voglio girare il mondo anche se sul mio cammino}
\rigo{incontrerò la pioggia che scendendo mi bagnerà}
\rigo{anche con l’acqua negli occhi}
\rigo{potrò veder tutte le cose}
\rigo{che in una stanza all’asciutto}
\rigo{io non potrei sognare mai.}
\strofa
\rigo{Rit: E voglio andarmene via}
\rigo{proprio perchè secondo loro}
\rigo{dormire sotto le stelle,}
\rigo{vuol dire solo sentir freddo.}
\par\endgroup

\sottosezione{Chiusura}

\begingroup
\cantotesta{126}{Airone}{}{\bEG}{324pt}
\rigo{Airone \ch{Fa}{dove} vai airone}
\rigo{dalle ali \ch{Do7}{forti} contro il vento;}
\rigo{\ch{Fa}{il} tuo volo è come una canzone}
\rigo{\ch{Do7}{per} il cielo libero e contento.}
\rigo{La \ch{Fa}{tua} rotta \ch{Sib}{segui} da \ch{Fa}{solo,}}
\rigo{hai \ch{Rem}{lasciato} gli altri fratelli;}
\rigo{\ch{Do7}{forse} stanco \ch{Sib}{sei} del tuo volo;}
\rigo{\ch{Fa}{oo} cerchi \ch{Rem}{cieli} più \ch{Do7}{belli.}}
\strofa
\rigo{Se divisi siam nell'avventura,}
\rigo{uno stesso azzurro cielo abbiamo}
\rigo{nubi d'oro e rivi d'acqua pura}
\rigo{lungo i quali stanchi riposiamo.}
\rigo{Al di là dell'arco baleno,}
\rigo{al di là dei monti e dei pini,}
\rigo{noi ancora ci ritroveremo}
\rigo{e saremo ancora vicini.}
\par\endgroup

\begingroup
\cantotesta{127}{Geordie}{tradizionale (versione di F. De André)}{\bEG}{563pt}
\rigo{\ch{Lam}{Mentre} attraversavo \ch{Sol}{}\ch{Fa}{London} \ch{Mi7}{Bridge}}
\rigo{\ch{Lam}{un} giorno \ch{Do}{senza} \ch{Sol}{sole}}
\rigo{\ch{Lam}{vidi} una \ch{Do}{donna} \ch{Sol}{pianger} \ch{Mim}{d'amore,}}
\rigo{\ch{Lam}{piangeva} per \ch{Mim}{il} suo \ch{Lam}{Geordie.}}
\strofa
\rigo{\ch{Lam}{Impiccheranno} \ch{Sol}{Geordie} con \ch{Fa}{una} corda \ch{Mi7}{d'oro,}}
\rigo{\ch{Lam}{è} un \ch{Do}{privilegio} \ch{Sol}{raro.}}
\rigo{\ch{Lam}{Rubò} sei \ch{Do}{cervi} nel \ch{Sol}{parco} del \ch{Mim}{Re}}
\rigo{\ch{Lam}{vendendoli} \ch{Mim}{per} \ch{Lam}{denaro.}}
\strofa
\rigo{Sellate il suo cavallo dalla bianca criniera}
\rigo{Sellatele il suo pony}
\rigo{Cavalcherà fino a Londra stasera}
\rigo{Ad implorare per Geordie}
\strofa
\rigo{Geordie non rubò mai neppure per me}
\rigo{Un frutto o un fiore raro}
\rigo{Rubò sei cervi nel parco del re}
\rigo{Vendendoli per denaro}
\strofa
\rigo{Salvate le sue labbra, salvate il suo sorriso}
\rigo{Non ha vent'anni ancora}
\rigo{Cadrà l'inverno anche sopra il suo viso}
\rigo{Potrete impiccarlo allora}
\strofa
\rigo{Né il cuore degli inglesi né lo scettro del re}
\rigo{Geordie potran salvare}
\rigo{Anche se piangeranno con te}
\rigo{La legge non può cambiare}
\strofa
\rigo{\ch{Lam}{Così} lo \ch{Sol}{impiccheranno} con \ch{Fa}{una} corda \ch{Mi7}{d'oro}}
\rigo{\ch{Lam}{È} un \ch{Do}{privilegio} \ch{Sol}{raro}}
\rigo{\ch{Lam}{Rubò} sei \ch{Do}{cervi} nel \ch{Sol}{parco} del \ch{Mim}{re}}
\rigo{\ch{Lam}{Vendendoli} \ch{Mim}{per} \ch{Lam}{denaro}}
\par\endgroup

\begingroup
\cantotesta{128}{Il falco}{}{\bEG}{584pt}
\rigo{Un fal\ch{Re}{co} volava nel cielo\ch{Sol}{} un mattino}
\rigo{r\ch{La}{icordo} quel tempo quando ero bambin\ch{Re}{o}}
\rigo{io lo seguivo nel ros\ch{Sol}{so} tramonto}
\rigo{dall’\ch{La}{alto} di un monte vedevo il \ch{Re}{suo} mondo.}
\strofa
\commento{Rit.}
\rigo{E allora \ch{Sol}{eha,} eha eha ...}
\rigo{e allora \ch{Re}{eha,} eha}
\rigo{eha, eha, eha, eha, eha, \ch{Sol}{eha} eh.     \ch{Re}{}}
\strofa
\rigo{Fiumi mari e boschi senza confine}
\rigo{i chiari orizzonti e le verdi colline}
\rigo{e quando partivo per un lungo sentiero}
\rigo{partivo ragazzo e tornavo guerriero.}
\strofa
\commento{Rit.}
\strofa
\rigo{Le tende rosse vicino al torrente}
\rigo{la vita felice tra la mia gente}
\rigo{e quando il mio arco colpiva lontano}
\rigo{sentivo l’orgoglio di essere indiano.}
\strofa
\commento{Rit.}
\strofa
\rigo{Fiumi mari e boschi mossi dal vento}
\rigo{luna su luna i miei capelli d’argento}
\rigo{e quando era l’ora dell’ultimo sonno}
\rigo{partivo dal campo per non farvi ritorno.}
\strofa
\commento{Rit.}
\strofa
\rigo{Un falco volava nel cielo un mattino}
\rigo{e verso il sole mi indicava il cammino}
\rigo{un falco che un giorno era stato colpito}
\rigo{ma no, non è morto, era solo ferito.}
\strofa
\commento{Rit.}
\par\endgroup

\begingroup
\cantotesta{129}{Il testamento del capitano}{}{\bRS}{454pt}
\rigo{Il capitan de la \ch{Fa}{compagnia}}
\rigo{E l'è ferito sta per \ch{Do}{morir}}
\rigo{E manda a dire ai suoi alpini}
\rigo{\ch{Fa}{Perché} lo \ch{Do}{vengano} a \ch{Fa}{ritrovar}}
\strofa
\rigo{I suoi alpini ghe manda a dire}
\rigo{Che non han scarpe per camminar}
\rigo{(O con le scarpe, o senza scarpe}
\rigo{I miei alpini li voglio qua)}
\strofa
\rigo{Cosa comanda siòr Capitano}
\rigo{Che noi adesso semo arriva'}
\rigo{(Ed io comando che il mio corpo}
\rigo{In cinque pezzi sia taglia)}
\strofa
\rigo{Il primo pezzo alla mia patria}
\rigo{Secondo pezzo al battaglion}
\rigo{(Il terzo pezzo alla mia mamma}
\rigo{Che si ricordi del suo figliuol)}
\strofa
\rigo{Il quarto pezzo alla mia bella}
\rigo{Che si ricordi del suo primo amor}
\rigo{L'ultimo pezzo alle montagne}
\rigo{Ché lo fioriscano di rose e fior}
\strofa
\rigo{L'ultimo pezzo alle montagne}
\rigo{Ché lo fioriscano di rose e fior}
\par\endgroup

\begingroup
\cantotesta{130}{Il vascello fantasma}{}{\bEG}{741pt}
\rigo{\ch{Lam}{Pende} un uomo dal pennone,}
\rigo{\ch{Rem}{tutto} \ch{Lam}{nero} di catrame}
\rigo{\ch{Rem}{non} \ch{Lam}{è} certo un buon boccone,}
\rigo{\ch{Mi}{per} \ch{Lam}{i} corvi che hanno fame.}
\strofa
\begin{ritornello}
\rigo{\ch{Lam}{Pa} zum! Pa pa zum!}
\rigo{\ch{Mi}{pa} \ch{Lam}{pa} pa pa pa zach!}
\strofa
\end{ritornello}
\rigo{Cinque teschi tutti neri}
\rigo{stan sul cassero di prua}
\rigo{son dei cinque bucanieri}
\rigo{che son morti alla tortura.}
\strofa
\begin{ritornello}
\rigo{\ch{Lam}{Pa} zum! Pa pa zum!}
\rigo{\ch{Mi}{pa} \ch{Lam}{pa} pa pa pa zach!}
\strofa
\end{ritornello}
\rigo{Venti ombre tutte nere}
\rigo{vengon su dal boccaporto}
\rigo{sono venti schiavi neri}
\rigo{che son morti nel trasporto.}
\strofa
\begin{ritornello}
\rigo{\ch{Lam}{Pa} zum! Pa pa zum!}
\rigo{\ch{Mi}{pa} \ch{Lam}{pa} pa pa pa zach!}
\strofa
\end{ritornello}
\rigo{Sulla tolda biancheggianti}
\rigo{stan tre scheletri a ballare}
\rigo{sono i resti dei briganti}
\rigo{giustiziati in alto mare.}
\strofa
\begin{ritornello}
\rigo{\ch{Lam}{Pa} zum! Pa pa zum!}
\rigo{\ch{Mi}{pa} \ch{Lam}{pa} pa pa pa zach!}
\strofa
\end{ritornello}
\rigo{Se una notte tutta scura}
\rigo{sentirete un gran lamento}
\rigo{è la voce di Tortuga}
\rigo{morto in ammutinamento.}
\strofa
\begin{ritornello}
\rigo{\ch{Lam}{Pa} zum! Pa pa zum!}
\rigo{\ch{Mi}{pa} \ch{Lam}{pa} pa pa pa zach!}
\strofa
\end{ritornello}
\rigo{Coricandovi stanotte}
\rigo{sentirete un sordo tonfo}
\rigo{è il fantasma del vascello}
\rigo{che reclama il suo trionfo.}
\strofa
\begin{ritornello}
\rigo{\ch{Lam}{Pa} zum! Pa pa zum!}
\rigo{\ch{Mi}{pa} \ch{Lam}{pa} pa pa pa zach!}
\end{ritornello}
\par\endgroup

\begingroup\senzaaccordi
\cantotesta{131}{La piroga}{}{\bEG}{396pt}
\rigo{La notte è piena di stelle}
\rigo{ci fa sognare le cose più belle,}
\rigo{più belle, più belle.}
\strofa
\rigo{Tu sogni, guardi lontano}
\rigo{vedi un gran fiume che scorre pian piano,}
\rigo{pian piano, pian piano.}
\strofa
\rigo{Sul fiume c’è una piroga}
\rigo{e dentro questa c’è un uomo che voga,}
\rigo{che voga, che voga.}
\strofa
\rigo{Intanto dietro la duna}
\rigo{vedi pian piano spuntare la luna,}
\rigo{la luna, la luna.}
\strofa
\rigo{L’uomo lascia il vogare}
\rigo{guarda la luna e si mette a cantare,}
\rigo{a cantare, a cantare.}
\strofa
\rigo{"Ti prego, o madre luna,}
\rigo{dona al mio popol ricchezza e fortuna,}
\rigo{fortuna, fortuna.}
\strofa
\rigo{Proteggi tutte le greggi,}
\rigo{fa’ che il mio popol rispetti le leggi,}
\rigo{le leggi, le leggi.}
\strofa
\rigo{Proteggi l’acqua del fonte,}
\rigo{l’erba del piano e la pianta del monte,}
\rigo{del monte, del monte.}
\par\endgroup

\sottosezione{Ultimi}

\begingroup
\cantotesta{132}{Scende la notte (esploratori)}{}{\bEG}{516pt}
\rigo{\ch{Rem}{Scende} \ch{La}{la} \ch{Rem}{notte} e \ch{La}{distende} il \ch{Rem}{suo} mantello \ch{La}{di} vel}
\rigo{\ch{Rem}{e} il campo \ch{Do}{calmo} e \ch{Fa}{silente} si \ch{Sib}{raccoglie} \ch{Do}{nel} \ch{Fa}{mister}}
\strofa
\begin{ritornello}
\rigo{\ch{Rem}{O} \ch{Sib}{Vergine} \ch{La}{di} luce, \ch{Rem}{stella} dei nostri \ch{La}{cuor,}}
\rigo{\ch{Rem}{ascolta} la \ch{Sib}{preghiera,} \ch{La}{}\ch{Rem}{madre} degli \ch{La}{}\ch{Re}{esplorator}}
\strofa
\end{ritornello}
\rigo{O delle stelle Signora volgi lo sguardo quaggiù}
\rigo{dove i tuoi figli sotto le tende ti ameranno sempre più}
\strofa
\begin{ritornello}
\rigo{\ch{Rem}{O} \ch{Sib}{Vergine} \ch{La}{di} luce, \ch{Rem}{stella} dei nostri \ch{La}{cuor,}}
\rigo{\ch{Rem}{ascolta} la \ch{Sib}{preghiera,} \ch{La}{}\ch{Rem}{madre} degli \ch{La}{}\ch{Re}{esplorator}}
\strofa
\end{ritornello}
\rigo{O Tu più bianca che neve nel tuo mantel verginal}
\rigo{con la dolcezza tua lievi ci proteggi contro ogni mal}
\strofa
\begin{ritornello}
\rigo{\ch{Rem}{O} \ch{Sib}{Vergine} \ch{La}{di} luce, \ch{Rem}{stella} dei nostri \ch{La}{cuor,}}
\rigo{\ch{Rem}{ascolta} la \ch{Sib}{preghiera,} \ch{La}{}\ch{Rem}{madre} degli \ch{La}{}\ch{Re}{esplorator}}
\strofa
\end{ritornello}
\rigo{Come le tende a noi care s’avvolgono sul partir}
\rigo{così avvolgici col tuo pregare quando saremo per morir}
\strofa
\begin{ritornello}
\rigo{\ch{Rem}{O} \ch{Sib}{Vergine} \ch{La}{di} luce, \ch{Rem}{stella} dei nostri \ch{La}{cuor,}}
\rigo{\ch{Rem}{ascolta} la \ch{Sib}{preghiera,} \ch{La}{}\ch{Rem}{madre} degli \ch{La}{}\ch{Re}{esplorator}}
\end{ritornello}
\par\endgroup

\begingroup
\cantotesta{133}{Signor, fra le tende schierati}{}{\bEG}{273pt}
\rigo{\ch{Sol}{Sign}\ch{Do}{or} tra le tende sch\ch{Sol}{ierati}}
\rigo{per s\ch{Do}{alutar}\ch{Sol}{} il dì che muor \ch{Re}{}       \ch{Re7}{}}
\rigo{l\ch{Sol}{e} no\ch{Do}{te} di canti accorat\ch{Sol}{i}}
\rigo{levia\ch{Do}{mo} \ch{Sol}{a} te, calde d'am\ch{Re}{or.}    \ch{Re7}{}}
\strofa
\rigo{\ch{Sol}{Ascolta} tu \ch{Do}{l'umil} preghier\ch{Sol}{a}}
\rigo{che d'aspro s\ch{Do}{uol} s'ode innalzar \ch{Re}{}}
\rigo{a\ch{Sol}{} te cui\ch{Mim}{} mancava, la sera,  \ch{Re}{}}
\rigo{un te\ch{Do}{tto} ancor per ripos\ch{Re}{ar.}        \ch{Sol}{}}
\rigo{C\ch{Sol}{hiedon} sol tutti i nostri cu\ch{Re}{ori}}
\rigo{a\ch{Re7}{} te sempre meglio servir\ch{Sol}{,}}
\rigo{g\ch{Do}{enufle}\ch{Sol}{tton} qu\ch{Re}{i} nel pian \ch{Sol}{}}
\rigo{i tuoi esplor\ch{Do}{ator.}      \ch{Dom}{}}
\rigo{Tu dal ci\ch{Sol}{el} bene\ch{Re}{dici} o Signor. \ch{Do}{}       \ch{Dom}{}       \ch{Sol}{}}
\par\endgroup

\end{multicols}

\sezione{Canti alpini}{}

\begin{multicols}{2}\raggedcolumns

\begingroup
\cantotesta{134}{Bersagliere ha cento penne}{tradizionale}{\bRS}{345pt}
\rigo{\ch{Sol}{Il} bersagliere ha \ch{Re7}{cento} \ch{Sol}{penne}}
\rigo{e l'alpino ne ha una \ch{Re7}{sola,}}
\rigo{il partigiano ne ha \ch{Sol}{nessuna}}
\rigo{e sta sui \ch{Re}{monti} a \ch{Sol}{guerreggiar.}}
\strofa
\rigo{Là sui monti vien giù la neve,}
\rigo{la bufera dell'inverno,}
\rigo{ma se venisse anche l'inferno}
\rigo{il partigiano riman lassù.}
\strofa
\rigo{Quando viene la notte scura}
\rigo{tutti dormono alla pieve,}
\rigo{ma camminando sopra la neve}
\rigo{il partigiano scende in azion.}
\strofa
\rigo{Quando poi ferito cade}
\rigo{non piangetelo dentro al cuore,}
\rigo{perchè se libero un uomo muore}
\rigo{che cosa importa di morir.}
\par\endgroup

\begingroup
\cantotesta{135}{Dove sei stato mio bell'alpino}{tradizionale}{\bRS}{185pt}
\rigo{\ch{Fa}{Dove} sei \ch{Solm}{stato} mio bel\ch{Do7}{l'alpi}\ch{Fa}{no,} (bis)}
\rigo{\ch{Do7}{Dove} sei stato bell'alpino \ch{Sib}{che} ti \ch{Fa}{ha} \ch{Re7}{cangià} \ch{Solm}{co}\ch{Do7}{lo}\ch{Fa}{re?} (bis)}
\strofa
\rigo{L'è stata l'aria del Trentino (bis)}
\rigo{che mi ha cangià colore (bis)}
\strofa
\rigo{I tuoi colori ritorneranno (bis)}
\rigo{questa sera a fare l'amore (bis)}
\par\endgroup

\begingroup\senzaaccordi
\cantotesta{136}{Era una notte che pioveva}{tradizionale}{\bRS}{268pt}
\rigo{Era una notte che pioveva}
\rigo{E che tirava un forte vento}
\rigo{Immaginatevi che grande tormento}
\rigo{Per un alpino che stava a vegliar (bis)}
\strofa
\rigo{A mezzanotte arriva il cambio}
\rigo{Accompagnato dal capo posto}
\rigo{Ohi sentinella ritorna al tuo posto}
\rigo{Sotto la tenda a ri a riposar (bis)}
\strofa
\rigo{Quando fui stato sotto la tenda}
\rigo{Sentii un rumore giù per la valle}
\rigo{Sentivo l'acqua giù per le spalle}
\rigo{Sentivo i sassi a ro a rotolar (bis)}
\strofa
\rigo{Quando fui stato addormentato}
\rigo{Sognavo di esser con la mia bella}
\rigo{Invece ero di sentinella}
\rigo{Fare la guardia allo, allo stranier (bis)}
\par\endgroup

\begingroup
\cantotesta{137}{Gerusalemme}{B. De Marzi}{\bRS}{219pt}
\rigo{\ch{Do}{Gerusalemme,} città di \ch{Sol}{pietre} \ch{Do}{bianche,}}
\rigo{\ch{Lam}{cuore} inaridito, \ch{Mim}{anima} \ch{Fa}{perduta}\ch{Sol}{,}}
\rigo{\ch{Do}{cos'hai} fatto \ch{Fa}{sul} Monte \ch{Do}{Calvario,} \ch{Sol}{dimmi,} \ch{Do}{cos'hai} fatto?}
\strofa
\rigo{\ch{Do}{Gerusalemme,} città di \ch{Sol}{rose} rosse\ch{Do}{,}}
\rigo{\ch{Lam}{cuore} di \ch{Mim}{pietra,} \ch{Fa}{anima} \ch{Sol}{perduta.}}
\strofa
\rigo{\ch{Do}{Gerusalemme,} città di \ch{Fa}{pietre} \ch{Do}{bianche,}}
\rigo{\ch{Fa}{guarda} chi piange \ch{Do}{ai} \ch{Sol7}{piedi} della \ch{Do}{croce:}}
\rigo{\ch{Fa}{è} la sua \ch{Do}{mamma,} che \ch{Sol7}{ti} perdona\ch{Do}{!}}
\par\endgroup

\begingroup
\cantotesta{138}{La leggenda del Piave}{E. A. Mario}{\bRS}{872pt}
\rigo{Il \ch{Sol}{Piave} \ch{Re7/Fa\#}{mormorava} calmo e \ch{Sol7/Fa}{placido} al \ch{Do/Mi}{passaggio}}
\rigo{Dei \ch{Sol/Re}{primi} fanti il \ch{Re7}{ventiquattro} \ch{Sol}{Maggio}}
\rigo{L'esercito \ch{Re7/Fa\#}{marciava} per \ch{Sol7/Fa}{raggiunger} la \ch{Do/Mi}{frontiera}}
\rigo{Per \ch{Sol/Re}{far} contro il \ch{La7/Do\#}{nemico} una \ch{Re7}{barriera}}
\rigo{\ch{Re}{Muti} \ch{Re/Mi}{passaron} \ch{Re}{quella} \ch{Re/Mi}{notte} i \ch{Re}{fanti}}
\rigo{\ch{Re}{Tacere} \ch{Re/Mi}{bisognava} \ch{Re}{e} \ch{Re/La}{andare} \ch{Re}{avanti}}
\rigo{\ch{Sol}{S'udiva} intanto dalle amate \ch{Re}{sponde,}}
\rigo{Sommesso e lieve il tripudiar \ch{Sol}{dell'onde.}}
\rigo{Era un presagio dolce e \ch{Re}{lusinghiero}}
\rigo{Il \ch{Do}{Piave} mormorò: \ch{Sol/Si}{«Non} \ch{Re7/Fa\#}{passa} lo \ch{Sol}{straniero!»}}
\strofa
\rigo{Ma in una notte \ch{Re7/Fa\#}{trista} si \ch{Sol7/Fa}{parlò} di un fosco \ch{Do/Mi}{evento}}
\rigo{E il \ch{Sol/Re}{Piave} udiva \ch{Re7}{l'ira} e lo \ch{Sol}{sgomento}}
\rigo{Ahi, quanta gente ha \ch{Re7/Fa\#}{vista} venir \ch{Sol7/Fa}{giù,} lasciare il \ch{Do/Mi}{tetto}}
\rigo{Per \ch{Sol/Re}{l'onta} \ch{La7/Do\#}{consumata} a \ch{Re7}{Caporetto}}
\rigo{\ch{Re}{Profughi} \ch{Re/Mi}{ovunque,} \ch{Re}{dai} \ch{Re/Mi}{lontani} \ch{Re}{monti}}
\rigo{\ch{Re}{Venivan} \ch{Re/Mi}{a} gremir \ch{Re/La}{}\ch{Re}{tutti} i suoi \ch{Re}{ponti}}
\rigo{\ch{Sol}{S'udiva} allor dalle violate \ch{Re}{sponde}}
\rigo{Sommesso e triste il mormorio de \ch{Sol}{l'onde}}
\rigo{Come un singhiozzo in quell'autunno \ch{Re}{nero}}
\rigo{Il \ch{Do}{Piave} mormorò: \ch{Re7/Fa\#}{}\ch{Sol/Si}{«Ritorna} lo \ch{Sol}{straniero!»}}
\strofa
\rigo{E ritornò il \ch{Re7/Fa\#}{nemico} per \ch{Sol7/Fa}{l'orgoglio} e per la \ch{Do/Mi}{fame}}
\rigo{\ch{Sol/Re}{Volea} sfogare \ch{Re7}{tutte} le sue \ch{Sol}{brame}}
\rigo{Vedeva il piano \ch{Re7/Fa\#}{aprico,} di \ch{Sol7/Fa}{lassù} voleva \ch{Do/Mi}{ancora}}
\rigo{\ch{Sol/Re}{Sfamarsi} e \ch{La7/Do\#}{tripudiare} come \ch{Re7}{allora}}
\rigo{\ch{Re}{«No!»,} disse il \ch{Re/Mi}{Piave.} \ch{Re}{«No!»,} \ch{Re/Mi}{dissero} i \ch{Re}{fanti}}
\rigo{\ch{Re}{Mai} più il \ch{Re/Mi}{nemico} \ch{Re}{faccia} un \ch{Re/La}{passo} \ch{Re}{avanti}}
\rigo{Si \ch{Sol}{vide} il Piave rigonfiar le \ch{Re}{sponde}}
\rigo{E come i fanti combatteron \ch{Sol}{l'onde}}
\rigo{Rosso di sangue del nemico \ch{Re}{altero}}
\rigo{Il \ch{Do}{Piave} \ch{Sol/Si}{comandò:} \ch{Re7/Fa\#}{«Indietro} va', \ch{Sol}{straniero!»}}
\strofa
\rigo{Indietreggiò il \ch{Re7/Fa\#}{nemico} fino a \ch{Sol7/Fa}{Trieste,} fino a \ch{Do/Mi}{Trento}}
\rigo{E \ch{Sol/Re}{la} vittoria \ch{Re7}{sciolse} le ali al \ch{Sol}{vento}}
\rigo{Fu sacro il patto \ch{Re7/Fa\#}{antico,} tra le \ch{Sol7/Fa}{schiere} furon \ch{Do/Mi}{visti}}
\rigo{\ch{Sol/Re}{Risorgere} \ch{La7/Do\#}{Oberdan,} Sauro e \ch{Re7}{Battisti}}
\rigo{\ch{Re}{Infranse} \ch{Re/Mi}{alfin} \ch{Re}{l'italico} \ch{Re}{}\ch{Re/Mi}{valore}}
\rigo{\ch{Re}{Le} forche e \ch{Re/Mi}{l'armi} \ch{Re}{}\ch{Re/La}{}\ch{Re}{dell'impiccatore}}
\rigo{\ch{Sol}{Sicure} l'Alpi, libere le \ch{Re}{sponde}}
\rigo{E tacque il Piave, si placaron \ch{Sol}{l'onde}}
\rigo{Sul patrio suolo, vinti i torvi \ch{Re}{Imperi}}
\rigo{La \ch{Do}{Pace} non \ch{Sol/Si}{trovò} né \ch{Re7/Fa\#}{oppressi,} né \ch{Sol}{stranieri}}
\par\endgroup

\begingroup
\cantotesta{139}{La tradotta}{tradizionale}{\bRS}{651pt}
\rigo{\ch{Sol}{La} tradotta che \ch{Re}{parte} da Torino}
\rigo{\ch{Re7}{a} Milano non si ferma più\ch{Sol}{}}
\rigo{\ch{Sol}{ma} la va diretta al \ch{Re}{Piave,}}
\rigo{\ch{Re7}{ma} la va diretta al Piave\ch{Sol}{,}}
\rigo{\ch{Sol}{La} tradotta che \ch{Re}{parte} da Torino}
\rigo{\ch{Re7}{a} Milano non si ferma più\ch{Sol}{}}
\rigo{\ch{Sol}{ma} la va diretta al \ch{Re}{Piave,}}
\rigo{\ch{Re7}{cimitero} della gioventù\ch{Sol}{.}}
\strofa
\rigo{Siam partiti, siam partiti}
\rigo{in ventisette,}
\rigo{solo in cinque siam tornati qua,}
\rigo{e gli altri ventidue?}
\rigo{e gli altri ventidue?}
\rigo{Siam partiti, siam partiti}
\rigo{in ventisette,}
\rigo{solo in cinque siam tornati qua,}
\rigo{e gli altri ventidue?}
\rigo{sono morti tutti a San Donà.}
\strofa
\rigo{A Nervesa, a Nervesa c'è una croce}
\rigo{mio fratello è sepolto là,}
\rigo{io c'ho scritto su Ninetto,}
\rigo{io c'ho scritto su Ninetto.}
\rigo{A Nervesa, a Nervesa c'è una croce}
\rigo{mio fratello è sepolto là,}
\rigo{io c'ho scritto su Ninetto,}
\rigo{che la mamma lo ritroverà.}
\strofa
\rigo{Cara suora, cara suora son ferito}
\rigo{a domani non c'arrivo più,}
\rigo{se non c'è qui la mia mamma,}
\rigo{se non c'è qui la mia mamma.}
\rigo{Cara suora, cara suora son ferito}
\rigo{a domani non c'arrivo più,}
\rigo{se non c'è qui la mia mamma,}
\rigo{un bel fiore me lo porti tu.}
\par\endgroup

\begingroup
\cantotesta{140}{La vien giù dalle montagne}{}{\bRS}{345pt}
\rigo{\ch{Re}{La} vien giù da le montagne,}
\rigo{l'è vesti\ch{La7}{ta} a la fran\ch{Re}{cese,}}
\rigo{da un bel \ch{Sol}{giovane} cortese}
\rigo{gli fu \ch{La7}{chiesto} far l'a\ch{Re}{mor.}}
\strofa
\rigo{«Lo ringra\ch{Re}{zio,} o giovanotto,}
\rigo{lo ringrazio \ch{La7}{del} buon cuore,}
\rigo{appartengo a un altro amore}
\rigo{che mi \ch{La7}{ama} e mi vuol \ch{Re}{ben.}}
\strofa
\rigo{Vatten via, o sciagurata,}
\rigo{vatten via su le montagne}
\rigo{a raccoglier le castagne}
\rigo{con gli agnelli a pascolar.}
\strofa
\rigo{Sono nata in mezzo ai fiori,}
\rigo{in mezzo ai fiori di Vermiglio,}
\rigo{sono pura come un giglio,}
\rigo{come un giglio vòi morir»}
\par\endgroup

\begingroup\senzaaccordi
\cantotesta{141}{Monte Canino}{tradizionale}{\bRS}{541pt}
\rigo{Non ti ricordi quel mese d'Aprile,}
\rigo{quel lungo treno che andava al confine.}
\rigo{Che trasportavano migliaia degli alpini:}
\rigo{sù, sù correte: è l'ora di partir!}
\rigo{Che trasportavano migliaia degli alpini:}
\rigo{sù, sù correte: è l'ora di partir!}
\strofa
\rigo{Dopo tre giorni di strada ferrata,}
\rigo{ed altri due di lungo cammino,}
\rigo{siamo arrivati sul Monte Canino}
\rigo{e a ciel sereno ci tocca riposar...}
\rigo{siamo arrivati sul Monte Canino}
\rigo{e a ciel sereno ci tocca riposar...}
\strofa
\rigo{Se avete fame guardate lontano,}
\rigo{se avete sete la tazza alla mano.}
\rigo{Se avete sete la tazza alla mano}
\rigo{che ci rinfresca la neve ci sarà.}
\rigo{Se avete sete la tazza alla mano}
\rigo{che ci rinfresca la neve ci sarà.}
\strofa
\rigo{Non più coperte lenzuola pulite.}
\rigo{Non più il sapore dei caldi tuoi baci.}
\rigo{Solo si sentono gli uccelli rapaci,}
\rigo{tra la tormenta e il rombo del cannon.}
\rigo{Solo si sentono gli uccelli rapaci}
\rigo{tra la tormenta e il rombo del cannon.}
\strofa
\rigo{Alla mattina il tenente fa sveglia}
\rigo{e tutt'a un tratto riunisce i plotoni}
\rigo{e sulle cime degli alti burroni}
\rigo{tutti insieme il fucile si sparò}
\rigo{e sulle cime degli alti burroni}
\rigo{tutti insieme il fucile si sparò}
\strofa
\rigo{E più di venti li ho visti a morire}
\rigo{e tutti gli altri li ho visti a scappare}
\rigo{e si sentivano tra loro a gridare}
\rigo{"se ci rendiamo, saremo prigionier"}
\rigo{e si sentivano tra loro a gridare}
\rigo{"se ci rendiamo, saremo prigionier"}
\par\endgroup

\begingroup\senzaaccordi
\cantotesta{142}{Monte Pasubio}{C. Geminiani, G. De Marzi}{\bRS}{389pt}
\rigo{Su la strada del Monte Pasubio}
\rigo{bomborombom bom bomborombom}
\rigo{lenta sale una lunga colonna}
\rigo{bomborombom bom bomborombom.}
\rigo{L'è la marcia de chi non torna}
\rigo{de chi se ferma a morir lassù.}
\strofa
\rigo{Ma gli alpini non hanno paura}
\rigo{bomborombom bom bomborombom.}
\strofa
\rigo{Su la cima del Monte Pasubio}
\rigo{bomborombom bom bomborombom,}
\rigo{soto i denti ghe ze 'na miniera}
\rigo{bomborombom bom bomborombom.}
\rigo{Ze i alpini che scava e spera}
\rigo{de ritornare a trovar l'amor.}
\strofa
\rigo{Ma gli alpini non hanno paura,}
\rigo{bomborombom bom bomborombom.}
\strofa
\rigo{Su la strada del Monte Pasubio}
\rigo{bomborombom bom bomborombom}
\rigo{ze rimasta soltanto 'na crose}
\rigo{bomborombom bom bomborombom.}
\rigo{No se sente ma più 'na vose,}
\rigo{ma solo el vento che basa i fior.}
\strofa
\rigo{Ma gli alpini non hanno paura}
\rigo{bomborombom bom bomborombà.}
\par\endgroup

\begingroup\senzaaccordi
\cantotesta{143}{Quel mazzolin di fiori}{tradizionale}{\bRS}{332pt}
\rigo{Quel mazzolin di fiori che vien dalla montagna …}
\rigo{E bada ben che non si bagna}
\rigo{Che lo voglio regalar}
\strofa
\rigo{Lo voglio regalare perché l'è un bel mazzetto …}
\rigo{Lo voglio dare al mio moretto}
\rigo{Questa sera quando vien}
\strofa
\rigo{Stasera quando viene sarà una brutta sera …}
\rigo{E lui sabato sera,}
\rigo{Lui non è venuto a me}
\strofa
\rigo{Non è venuto a me, l'è andai dalla Rosina …}
\rigo{E perché mi son poverina}
\rigo{Mi fa pianger, sospirar.}
\strofa
\rigo{Fa pianger, sospirare sul letto dei lamenti …}
\rigo{E cosa mai dirà la gente}
\rigo{Cosa mai dirà di me}
\strofa
\rigo{Dirà che son tradita, tradita nell' amore …}
\rigo{E perché mì, me piange il cuore}
\rigo{E per sempre piangerà}
\par\endgroup

\begingroup\senzaaccordi
\cantotesta{144}{Se non ci conoscete}{tradizionale (canto degli Arditi)}{\bRS}{634pt}
\rigo{Noi siamo gli arditi della Terza Armata:}
\rigo{dove attacchiamo è sicura l'avanzata.}
\strofa
\rigo{Bombe a man e colpi di pugnal!}
\strofa
\rigo{Se non ci conoscete guardateci dall'alto:}
\rigo{noi siamo gli arditi del battaglion d'assalto.}
\rigo{Bombe a man...}
\strofa
\rigo{Se non ci conoscete, guardateci nel viso:}
\rigo{veniamo dall'inferno e andiamo in paradiso.}
\rigo{Bombe a man...}
\strofa
\rigo{Se non ci conoscete guardateci il maglione,}
\rigo{noi siamo gli arditissimi del colonnel Pavone.}
\rigo{Bombe a man...}
\strofa
\rigo{Se giri tutta Italia non trovi un lazzarone;}
\rigo{li ha requisiti tutti il colonnel Pavone.}
\rigo{Bombe a man...}
\strofa
\rigo{Gli arditi son soldati creduti farabutti,}
\rigo{ma loro non s'arrendon se non son morti tutti.}
\rigo{Bombe a man...}
\strofa
\rigo{Nelle patrie galere non ci sono più banditi}
\rigo{perchè andaron tutti ai battaglion d'arditi.}
\rigo{Bombe a man...}
\strofa
\rigo{Se volete far la guerra con dei bravi soldati,}
\rigo{andate al cellulare, prendete i carcerati.}
\rigo{Bombe a man...}
\strofa
\rigo{Se ha più di sei anni puoi farlo caporale:}
\rigo{se è condannato a vita puoi farlo generale.}
\rigo{Bombe a man...}
\strofa
\rigo{Se non ci conoscete guardate le mostrine,}
\rigo{noi siamo gli arditi e rubiamo le galline.}
\rigo{Bombe a man...}
\strofa
\rigo{La figlia di.... non vuol prender marito,}
\rigo{perchè dopo la guerra si sposerà un ardito.}
\rigo{Bombe a man...}
\strofa
\rigo{Perchè noi siamo arditi sfidiamo la mitraglia:}
\rigo{sarà 'l nostro pugnale che salverà l'Italia.}
\rigo{Bombe a man...}
\strofa
\rigo{Ragazze di Trieste apriteci le porte:}
\rigo{ripasseremo il Piave a costo della morte.}
\rigo{Bombe a man...}
\par\endgroup

\begingroup\senzaaccordi
\cantotesta{145}{Signore delle cime}{B. De Marzi}{\bRS}{297pt}
\rigo{Dio del cielo,}
\rigo{signore delle cime}
\rigo{un nostro amico}
\rigo{,hai chiesto alla montagna.}
\rigo{Ma ti preghiamo,}
\rigo{ma ti preghiamo}
\rigo{su nel paradiso,}
\rigo{su nel paradiso}
\rigo{lascialo andare}
\rigo{,per le tue montagne.}
\rigo{Santa Maria,}
\rigo{signora della neve}
\rigo{copri col bianco,}
\rigo{soffice mantello}
\rigo{il nostro amico}
\rigo{il nostro fratello.}
\rigo{Su nel paradiso,}
\rigo{su nel paradiso}
\rigo{lascialo andare,}
\rigo{per le tue montagne.}
\par\endgroup

\begingroup
\cantotesta{146}{Stelutis alpinis}{A. Zardini}{\bRS}{669pt}
\rigo{Se tu vens ca \ch{Fa}{su} ta' \ch{Do7}{cretis,}}
\rigo{là \ch{Fa7}{che} lôr \ch{Sib}{mi} àn soterât,}
\rigo{\ch{Fa}{al} è un \ch{Do7}{splaz} plen di \ch{Fa}{stelutis:}}
\rigo{dal gno \ch{Do}{sanc} \ch{Sol7}{'l} è stât \ch{Do}{bagnât.}}
\strofa
\rigo{Par se\ch{Solm}{gnâl} \ch{Do7}{une} \ch{Fa}{crosute}}
\rigo{je scul\ch{Sib}{pide} lì tal \ch{Fa}{cret:}}
\rigo{fra \ch{Fa7}{chês} \ch{Sib}{stelis} \ch{Do7}{nas} \ch{Fa}{l'arbute,}}
\rigo{sot di \ch{Do7}{lôr} jo duâr \ch{Fa}{cuièt.}}
\strofa
\rigo{Cjol sù, cjol une stelute:}
\rigo{je a ricuart dal nestri ben,}
\rigo{tu i darâs 'ne bussadute,}
\rigo{e po platile tal sen.}
\strofa
\rigo{Cuant che a cjase tu sês sole}
\rigo{e di cûr tu preis par me,}
\rigo{il gno spirt atôr ti svole:}
\rigo{jo e la stele o sin cun te.}
\strofa
\rigo{Traduzione italiana:}
\strofa
\rigo{Se tu vieni quassù tra le rocce,}
\rigo{laddove mi hanno sepolto,}
\rigo{c'è uno spiazzo pieno di stelle alpine:}
\rigo{dal mio sangue è stato bagnato.}
\strofa
\rigo{Come segno una piccola croce}
\rigo{è scolpita lì nella roccia:}
\rigo{fra quelle stelle nasce l'erbetta,}
\rigo{sotto di loro io dormo sereno.}
\strofa
\rigo{Cogli, cogli una piccola stella:}
\rigo{a ricordo del nostro amore.}
\rigo{Dalle un bacio,}
\rigo{e poi nascondila in seno.}
\strofa
\rigo{Quando a casa tu sei sola}
\rigo{e di cuore preghi per me}
\rigo{il mio spirito ti aleggia intorno}
\rigo{io e la stella siamo con te.}
\par\endgroup

\begingroup
\cantotesta{147}{Sul cappello}{tradizionale}{\bEG}{481pt}
\rigo{\ch{Fa}{Sul} cappello, sul cappello che noi por\ch{Do7}{tiamo}}
\rigo{c'è una lunga, c'è una lunga penna \ch{Fa}{nera}}
\rigo{che a noi serve, che a noi serve da \ch{Do7}{bandiera}}
\rigo{su pei monti, su pei monti a guerreg\ch{Fa}{giar,} oilalà.}
\strofa
\rigo{Su pei monti, su pei monti}
\rigo{che noi saremo}
\rigo{coglieremo, coglieremo}
\rigo{le stelle alpine}
\rigo{per donarle, per donarle}
\rigo{alle bambine,}
\rigo{farle pianger, farle pianger}
\rigo{e sospirar, oilalà.}
\strofa
\rigo{Su pei monti, su pei monti}
\rigo{che noi saremo}
\rigo{pianteremo, pianteremo}
\rigo{l'accampamento}
\rigo{brinderemo, brinderemo}
\rigo{al reggimento,}
\rigo{viva il Corpo, viva il Corpo}
\rigo{degli Alpin, oilalà.}
\strofa
\begin{ritornello}
\rigo{Evviva \ch{Sib}{viva} il reggi\ch{Fa}{mento}}
\rigo{Evviva \ch{Do7}{viva} il Corpo degli Alpin}
\rigo{\ch{Fa}{Oilalà.} (bis)}
\end{ritornello}
\par\endgroup

\begingroup\senzaaccordi
\cantotesta{148}{Sul Monte Nero}{}{\bRS}{461pt}
\rigo{Lassù sul monte nero}
\rigo{c'è una piccola taverna;}
\rigo{ci son dodici briganti}
\rigo{al chiaror di una lanterna.}
\strofa
\rigo{Rit.: Caramba beviamo del whisky}
\rigo{caramba beviamo del gin !}
\rigo{E tu non dar retta al cuore}
\rigo{che tutto passa e và...}
\strofa
\rigo{Mentre gli altri son festanti}
\rigo{uno solo resta muto}
\rigo{ha il bicchiere ancora pieno}
\rigo{come mai non ha bevuto?}
\strofa
\commento{Rit.}
\strofa
\rigo{Ma non può dimenticare}
\rigo{il brigante la sua bella}
\rigo{i sui occhi luminosi}
\rigo{la sua bocca tanto bella.}
\strofa
\commento{Rit.}
\strofa
\rigo{Mentre stava spigolando}
\rigo{in un campo di frumento}
\rigo{è passato un aeroplano}
\rigo{che compì un mitragliamento.}
\strofa
\commento{Rit.}
\strofa
\rigo{Or la bella par che dorma}
\rigo{con la spiga ancora in mano}
\rigo{par che pensi al suo brigante}
\rigo{che la piange da lontano.}
\strofa
\commento{Rit.}
\par\endgroup

\begingroup\senzaaccordi
\cantotesta{149}{Sul ponte di Bassano}{tradizionale}{\bRS}{555pt}
\rigo{Eccole che le riva}
\rigo{Ste quattro moscardine...}
\rigo{Eccole che le riva}
\rigo{Ste quattro moscardine...}
\rigo{L'è tutte verdoline}
\rigo{Colori no le ghe n'ha...}
\rigo{Color no le ghe n'ha...}
\rigo{Color no le ghe n'ha!...}
\strofa
\rigo{Colori non ghe n'avemo,}
\rigo{Ne manco ghe cichemo...}
\rigo{Ma un bacio ci daremo}
\rigo{Sul ponte di Bassan...}
\strofa
\rigo{Sul ponte di Bassano}
\rigo{Là ci darem la mano...}
\rigo{Noi ci darem la mano}
\rigo{Ed un bacin d'amor.}
\strofa
\rigo{Per un bacin d'amore}
\rigo{Successer tanti guai!...}
\rigo{No, non credevo mai}
\rigo{Doverti abbandonar.}
\strofa
\rigo{Doverti bandonare}
\rigo{Volerti tanto bene!...}
\rigo{E' un giro di catene}
\rigo{Che m'incantena il cor.}
\strofa
\rigo{Che m'incantena il core}
\rigo{Che m'incatena a vita!...}
\rigo{Ohi mè io son tradita}
\rigo{Per un bacin d'amor.}
\strofa
\rigo{Chi m'incatena il core}
\rigo{Sarà la mia morosa!...}
\rigo{A maggio la và sposa}
\rigo{E mi vo fa el soldà.}
\strofa
\rigo{E mi farò 'l soldato}
\rigo{Nel 6° Reggimento...}
\rigo{Non partirò contento}
\rigo{Se non l'avarò sposà.}
\par\endgroup

\begingroup
\cantotesta{150}{Sul ponte di Perati}{tradizionale}{\bRS}{570pt}
\rigo{\ch{Sol}{Sul} ponte di \ch{Re7}{Perati,}}
\rigo{\ch{Do}{bandiera} \ch{Sol}{nera,}}
\rigo{È il lutto \ch{Do}{degli} \ch{Sol}{Alpini}}
\rigo{che \ch{Re7}{fan} la \ch{Sol}{guerra.}}
\strofa
\rigo{L'è il lutto della Julia}
\rigo{che va a la guerra}
\rigo{la meglio gioventù}
\rigo{che va sottoterra.}
\strofa
\rigo{Sull'ultimo vagone}
\rigo{c'è l'amor mio}
\rigo{col fazzoletto in mano}
\rigo{mi dà l'addio.}
\strofa
\rigo{Col fazzoletto in mano}
\rigo{mi salutava}
\rigo{e con la bocca i baci}
\rigo{la mi mandava.}
\strofa
\rigo{Quelli che son partiti}
\rigo{non son tornati,}
\rigo{sui monti della Grecia}
\rigo{sono restati.}
\strofa
\rigo{Sui monti della Grecia}
\rigo{c'è la Vojussa}
\rigo{col sangue degli Alpini}
\rigo{s'è fatta rossa.}
\strofa
\rigo{Alpini della Julia,}
\rigo{in alto il cuore:}
\rigo{sul ponte di Perati}
\rigo{c'è il tricolore!}
\par\endgroup

\begingroup\senzaaccordi
\cantotesta{151}{Ta-pum}{tradizionale}{\bRS}{308pt}
\rigo{Ho lasciato la mamma mia,}
\rigo{l'ho lasciata per fare il soldà.}
\strofa
\rigo{Rit. Tapum, tapum, tapum !}
\rigo{Venti giorni su l'Ortigara,}
\rigo{senza cambio per dismontar.}
\strofa
\rigo{Quando portano la pagnotta,}
\rigo{il cecchino comincia a sparar.}
\strofa
\rigo{E domani si andrà all'assalto:}
\rigo{soldatino, non farti ammazzar.}
\strofa
\rigo{Quando poi si discende al piano,}
\rigo{battaglione non ha più soldà.}
\strofa
\rigo{Nella valle c'è un cimitero,}
\rigo{cimitero dei nostri soldà.}
\strofa
\rigo{Cimitero di noi soldà,}
\rigo{forse un giorno ti vengo a trovar.}
\strofa
\rigo{Battaglion di tutti i morti,}
\rigo{noi giuriamo l'Italia salvar.}
\par\endgroup

\begingroup
\cantotesta{152}{Valsugana}{tradizionale}{\bRS}{481pt}
\rigo{\ch{Re}{Quando} seremo fora, \ch{La7}{fora} de la \ch{Re}{Valsugana}}
\rigo{Quando seremo fora, \ch{La7}{fora} de la \ch{Re}{Valsugana}}
\rigo{E a ritrovar la \ch{La7}{mama} a veder come la \ch{Re}{sta}}
\rigo{E a ritrovar la \ch{La7}{mama} a veder come, come la \ch{Re}{sta}}
\strofa
\rigo{La mama la sta \ch{La7}{bene,} il papà l'è \ch{Re}{ammalato}}
\rigo{La mama la sta \ch{La7}{bene,} il papà l'è \ch{Re}{ammalato}}
\rigo{Il mio bel partì \ch{La7}{soldato} chi sa quando tornerà\ch{Re}{}}
\rigo{Il mio bel partì \ch{La7}{soldato} chi sa quando, mai \ch{Re}{tornerà}}
\strofa
\rigo{Tuti i me dis che \ch{La7}{lu} 'l s'è zercà n'altra \ch{Re}{morosa}}
\rigo{Tuti i me dis che \ch{La7}{lu} 'l s'è zercà n'altra \ch{Re}{morosa}}
\rigo{L'è 'na storia \ch{La7}{dolorosa} che mi credere non \ch{Re}{so}}
\rigo{L'è 'na storia \ch{La7}{dolorosa} che mi cre, mi credere non \ch{Re}{so}}
\strofa
\rigo{Mi no la credo, \ch{La7}{ma} se 'l fussa propi, propi \ch{Re}{vera}}
\rigo{Mi no la credo, \ch{La7}{ma} se 'l fussa propi, propi \ch{Re}{vera}}
\rigo{Biondo o moro ancor'sta \ch{La7}{sera} 'n altro merlo \ch{Re}{troverò}}
\rigo{Biondo o moro ancor'sta \ch{La7}{sera} 'n altro merlo, mi \ch{Re}{retroverò}}
\par\endgroup

\end{multicols}

\sezione{Canti liturgici}{}

\begin{multicols}{2}\raggedcolumns

\sottosezione{Ingresso}

\begingroup\senzaaccordi
\cantotesta{153}{Acqua siamo noi}{}{}{45pt}
\par\endgroup

\begingroup
\cantotesta{154}{Camminerò}{}{}{526pt}
\commento{Rit.}
\rigo{C\ch{Re}{amminerò,} camminerò\ch{Sim}{}}
\rigo{s\ch{Sol}{ulla} Tua strada, Signor.\ch{La}{}}
\rigo{D\ch{Re}{ammi} la mano, voglio restar\ch{Sim}{}}
\rigo{per s\ch{Sol}{empre} insie\ch{La}{me} a Te.       \ch{Re}{}}
\strofa
\rigo{Q\ch{Re}{uando} ero solo, solo e stanco d\ch{Sim}{el} mondo,}
\rigo{q\ch{Sol}{uando} non c'era l'Amor, \ch{La}{}}
\rigo{t\ch{Re}{ante} persone vidi i\ch{Sim}{ntorno} a me,}
\rigo{s\ch{Sol}{entivo} cant\ch{La}{are} così.  \ch{Re}{}}
\strofa
\commento{Rit.}
\strofa
\rigo{Io non capivo ma rimasi a sentire}
\rigo{quando il Signore mi parlò:}
\rigo{Lui mi chiamava, chiamava anche me}
\rigo{e la mia risposta s'alzò.}
\strofa
\commento{Rit.}
\strofa
\rigo{Or non importa se uno ride di me,}
\rigo{lui certamente non sa}
\rigo{del gran regalo che ebbi quel dì}
\rigo{che dissi al Signore così.}
\strofa
\commento{Rit.}
\strofa
\rigo{A volte sono triste, ma mi guardo intorno,}
\rigo{scopro il mondo e l'amor;}
\rigo{son questi i doni che Lui fa a me}
\rigo{felice ritorno a cantar.}
\strofa
\commento{Rit.}
\par\endgroup

\begingroup\senzaaccordi
\cantotesta{155}{Cantiamo a te}{}{}{1773pt}
\rigo{Cantiamo a Te, Signore della vita:}
\rigo{il nome tuo è grande sulla terra}
\rigo{tutto parla di Te e canta la tua gloria.}
\rigo{Grande Tu sei e compi meraviglie: Tu sei Dio.}
\strofa
\rigo{Cantiamo Te, Signore Gesù Cristo:}
\rigo{Figlio di Dio venuto sulla terra,}
\rigo{fatto uomo per noi nel grembo di Maria.}
\rigo{Dolce Gesù risorto dalla Morte sei con noi.}
\strofa
\rigo{Cantiamo Te, Amore senza fine:}
\rigo{Tu che sei Dio lo Spirito del Padre}
\rigo{Brani e testi contenuti : 1145}
\rigo{vivi dentro di noi e guidi i nostri passi.}
\rigo{Accendi in noi il fuoco dell’eterna carità.}
\strofa
\rigo{Cantiamo la nostra lode}
\rigo{A Cristo Signore,}
\rigo{vivente nella Chiesa.}
\strofa
\begin{ritornello}
\rigo{Cantiamo la nostra lode,}
\rigo{cantiamo il nostro amen!}
\rigo{Alleluia! Alleluia!}
\strofa
\end{ritornello}
\rigo{A Cristo pastore,}
\rigo{presente nella Chiesa.}
\strofa
\rigo{Cantiamo le meraviglie}
\rigo{Ho visto una folla immensa:}
\rigo{cantavano un canto ineffabile. Alleluia!}
\strofa
\begin{ritornello}
\rigo{Cantiamo le meraviglie}
\rigo{del Signore dell’universo.}
\rigo{Le genti innanzi a lui verranno}
\rigo{inneggiando per l’eternità.}
\strofa
\end{ritornello}
\rigo{Ho visto volti di gioia,}
\rigo{ho udito arpe celesti. Alleluia!}
\strofa
\rigo{L’agnello che vinto la morte,}
\rigo{nel sangue ha portato vittoria. Alleluia!}
\strofa
\rigo{Cantiamo l’inizio}
\rigo{= E non videro che Gesù}
\strofa
\begin{ritornello}
\rigo{Cantico dei redenti}
\rigo{Il Signore è la mia salvezza}
\rigo{e con lui non temo più,}
\rigo{perché ho nel cuore la certezza,}
\rigo{la salvezza è qui con me.}
\strofa
\end{ritornello}
\rigo{Ti lodo Signore perché}
\rigo{un giorno eri lontano da me:}
\rigo{ora invece sei tornato}
\rigo{e mi hai preso con te.}
\strofa
\rigo{Berrete con gioia alle fonti,}
\rigo{alle fonti della salvezza}
\rigo{e quel giorno voi direte:}
\rigo{lodate il Signore invocate il suo nome.}
\strofa
\rigo{Fate conoscere ai popoli}
\rigo{tutto quello che lui ha compiuto}
\rigo{e ricordino per sempre,}
\rigo{ricordino sempre che il suo nome è grande.}
\strofa
\rigo{Cantate a chi ha fatto grandezze}
\rigo{e si è fatto sapere nel mondo;}
\rigo{grida forte la tua gioia, abitante di Sion,}
\rigo{perché grande con te è il Signore.}
\strofa
\rigo{Brani e testi contenuti : 1145}
\rigo{Cantico della vergine}
\rigo{= Magnificat Magnifica il Signore anima mia}
\strofa
\rigo{Cantico delle creature Benediciamo il Signore}
\rigo{= Benediciamo il Signore}
\strofa
\rigo{Cantico delle Creature laudato si’}
\rigo{Laudato si’ mi’ Signore}
\rigo{per frate sole, sora luna,}
\rigo{frate vento, il cielo, le stelle}
\rigo{per sora acqua, frate focu.}
\strofa
\begin{ritornello}
\rigo{Laudato si’ mi’ Signore}
\rigo{per la terra e le tue creature}
\strofa
\end{ritornello}
\rigo{Laudato si’ mi’ Signore}
\rigo{quello che porta la tua pace:}
\rigo{e saprà perdonare,}
\rigo{per il tuo amore saprà amare.}
\strofa
\rigo{Laudato si’ mi’ Signore}
\rigo{per sora morte corporale}
\rigo{dalla quale homo vivente}
\rigo{non potrà mai, mai scappare}
\strofa
\rigo{Cantico di liberazione}
\rigo{= Voglio cantare}
\strofa
\rigo{Cantico di nozze}
\rigo{Tu che inventi per ognuno di noi}
\rigo{un cammino che porta verso te.}
\rigo{Tu che chiami ognuno per il suo nome, Tu,}
\rigo{oggi ci chiami insieme.}
\rigo{Tu che hai in mente per ognuno di noi}
\rigo{una strada bellissima con te,}
\rigo{Tu stavolta parli di un’avventura che,}
\rigo{che si percorre insieme.}
\strofa
\begin{ritornello}
\rigo{Rendi agile il vascello dell’amore nostro}
\rigo{docile al soffio del tuo spirito.}
\rigo{E per gli immaginabili mari conducilo,}
\rigo{quelli che tu sai.}
\rigo{Rendi agile il vascello dell’amore nostro}
\rigo{che tocchi lidi che tu solo sai}
\rigo{e la scia della rotta disegni un’onda:}
\rigo{quella che tu vuoi.}
\strofa
\end{ritornello}
\rigo{Dacci forza perché ognuno di noi}
\rigo{con la vita dispieghi la tua idea.}
\rigo{Dacci forza per svelare il pensiero tuo}
\rigo{su questa vita insieme.}
\strofa
\begin{ritornello}
\rigo{Cantico di Zaccaria}
\rigo{Sia benedetto il Signore, Dio d’Abramo,}
\rigo{perché ha visitato e redento il suo popolo,}
\rigo{come aveva promesso, come aveva promesso.}
\strofa
\end{ritornello}
\rigo{È la salvezza da tutti i nostri nemici}
\rigo{e dalla mani di quanti ci odiano.}
\rigo{Ha avuto misericordia di noi,}
\rigo{si è ricordato del suo patto santo.}
\rigo{Gioisci popolo fedele, perché da quel giorno}
\rigo{tu non sei più schiavo}
\rigo{un’alba nuova è sorta per te,}
\rigo{un’alba nuova è sorta per te.}
\rigo{Brani e testi contenuti : 1145}
\strofa
\rigo{Quando sei stanco, quando non sai più cosa fare,}
\rigo{quando sei triste e quando sei tanto solo}
\rigo{e perfino quando la vita ti lascia}
\rigo{ricordati che è stata detta una parola di vita.}
\strofa
\rigo{Gioisci popolo fedele, perché da quel giorno}
\rigo{tu non sei più schiavo}
\rigo{un’alba nuova è sorta per te,}
\rigo{un’alba nuova è sorta per te.}
\par\endgroup

\begingroup\senzaaccordi
\cantotesta{156}{Che gioia ci hai dato}{}{}{356pt}
\rigo{Che gioia ci hai dato, Signore del cielo,}
\rigo{Signore del grande universo!}
\rigo{Che gioia ci hai dato vestito di luce,}
\rigo{vestito di gloria infinita,}
\rigo{vestito di gloria infinita!}
\strofa
\rigo{Vederti risorto, vederti Signore,}
\rigo{il cuore sta per impazzire!}
\rigo{Tu sei ritornato, tu sei qui tra noi,}
\rigo{e adesso ti avremo per sempre,}
\rigo{e adesso ti avremo per sempre.}
\strofa
\rigo{Chi cercate, donne, quaggiù?}
\rigo{Chi cercate, donne, quaggiù?}
\rigo{Quello che era morto non è qui:}
\rigo{è risorto sì come aveva detto anche a voi,}
\rigo{voi gridate a tutti che è risorto lui,}
\rigo{a tutti che è risorto lui.}
\strofa
\rigo{Tu hai vinto il mondo, Gesù!}
\rigo{Tu hai vinto il mondo, Gesù!}
\rigo{Liberiamo la felicità}
\rigo{e la morte, no, non esiste più, l’hai vinta tu,}
\rigo{hai salvato tutti noi, uomini con te,}
\rigo{tutti noi, uomini con te.}
\par\endgroup

\begingroup
\cantotesta{157}{Dall'aurora al tramonto}{}{}{495pt}
\begin{ritornello}
\rigo{\ch{Do\#m}{Dall'aurora} io \ch{Mi}{cerco} te,\ch{Fa\#m}{}\ch{Si}{}}
\rigo{\ch{Do\#m}{fino} al tramonto \ch{Mi}{ti} \ch{Fa\#m}{chiamo.}\ch{Si}{}}
\rigo{\ch{Do\#m}{Ha} sete solo \ch{Sol\#m}{di} te\ch{La}{}\ch{Si}{}}
\rigo{\ch{Do\#m}{l'anima} mia \ch{Sol\#m}{come} terra \ch{La}{deserta.}\ch{Si}{}}
\strofa
\end{ritornello}
\rigo{Non mi fermerò un solo istante}
\rigo{sempre canterò la tua lode}
\rigo{perchè sei il mio Dio}
\rigo{il mio riparo}
\rigo{mi proteggerai all'ombra delle tue ali.}
\strofa
\begin{ritornello}
\rigo{\ch{Do\#m}{Dall'aurora} io \ch{Mi}{cerco} te,\ch{Fa\#m}{}\ch{Si}{}}
\rigo{\ch{Do\#m}{fino} al tramonto \ch{Mi}{ti} \ch{Fa\#m}{chiamo.}\ch{Si}{}}
\rigo{\ch{Do\#m}{Ha} sete solo \ch{Sol\#m}{di} te\ch{La}{}\ch{Si}{}}
\rigo{\ch{Do\#m}{l'anima} mia \ch{Sol\#m}{come} terra \ch{La}{deserta.}\ch{Si}{}}
\strofa
\end{ritornello}
\rigo{Non mi fermerò un solo istante}
\rigo{io racconterò le tue opere}
\rigo{perchè sei il mio Dio}
\rigo{unico bene}
\rigo{Nulla mai potrà la notte contro di me}
\strofa
\begin{ritornello}
\rigo{\ch{Do\#m}{Dall'aurora} io \ch{Mi}{cerco} te,\ch{Fa\#m}{}\ch{Si}{}}
\rigo{\ch{Do\#m}{fino} al tramonto \ch{Mi}{ti} \ch{Fa\#m}{chiamo.}\ch{Si}{}}
\rigo{\ch{Do\#m}{Ha} sete solo \ch{Sol\#m}{di} te\ch{La}{}\ch{Si}{}}
\rigo{\ch{Do\#m}{l'anima} mia \ch{Sol\#m}{come} terra \ch{La}{deserta.}\ch{Si}{}}
\strofa
\rigo{\ch{Do\#m}{Ha} sete solo \ch{Sol\#m}{di} te\ch{La}{}\ch{Si}{}}
\rigo{\ch{Do\#m}{l'anima} mia \ch{Sol\#m}{come} terra \ch{La}{deserta.}\ch{Si}{}}
\end{ritornello}
\par\endgroup

\begingroup
\cantotesta{158}{Emmanuel}{}{}{498pt}
\rigo{\ch{Mim}{Come} in mezzo al mare brilla un \ch{Re}{riflesso} del so\ch{Mim}{le}}
\rigo{come tra le nuvole un\ch{Re}{} buco mostra il\ch{Mim}{} cielo azzurro}
\rigo{c\ch{Sol}{ome} sopra un pi\ch{Re}{cco} c'è una stella \ch{Do}{alpina} che     \ch{Mim}{}}
\rigo{è velluto bia\ch{Re}{nco} tra le rocc\ch{Do}{e} di lassù.    \ch{Mim}{}}
\strofa
\rigo{Tu vicino alla mia anima e i miei figli in te}
\rigo{la mia strada verso la luce è l'avventura in Te,}
\rigo{nuova vita, nuovo amore in questo mondo che}
\rigo{ha dimenticato il tempo della libertà.}
\strofa
\rigo{Emmanuel, Emmanuel, Emmanuel, Dio con noi}
\rigo{nel mio amore, nel nostro amore}
\rigo{c'è l'amore di Dio con noi}
\rigo{nella nostra vita Lui ha preso il posto suo,}
\rigo{Primo posto al primo posto è venuto Lui.}
\strofa
\rigo{Una Chiesa non è soltanto una casa di mattoni}
\rigo{una Chiesa non è soltanto un altare colorato,}
\rigo{una Chiesa è la realtà che nasce tra noi}
\rigo{se ci amiamo se ci amiamo con l'amore di Dio.}
\strofa
\rigo{Cambierà cambierà cambierà l'umanità}
\rigo{cambierà cambierà cambierà l'umanità}
\rigo{nella terra è seminato e un mondo nascerà}
\rigo{nascerà nascerà una nuova umanità}
\par\endgroup

\begingroup
\cantotesta{159}{Gioia che invade l'anima}{}{\bRS}{413pt}
\rigo{\ch{La}{Gioia} che invade l'\ch{Re}{anima} e canta, \ch{La}{gioia} di avere \ch{Re}{Te,}}
\rigo{\ch{La}{resurrezione} e \ch{Re}{vita} infinita, \ch{Mi}{vita} dell'uni\ch{Re}{tà.}}
\rigo{\ch{La}{La} porteremo al \ch{Re}{mondo} che attende, \ch{La}{la} porteremo \ch{Re}{là}}
\rigo{\ch{La}{dove} si sta spe\ch{Re}{gnendo} la vita, \ch{La}{vita} si accende\ch{Re}{rà,}}
\strofa
\rigo{\ch{La}{Perché} la tua casa è an\ch{Re}{cora} più grande, \ch{La}{grande} come sei \ch{Re}{tu,}}
\rigo{\ch{Fa\#m}{grande} come la \ch{Re}{terra} nell'uni\ch{Fa\#m}{verso} che vive in \ch{Mi}{Te;}}
\rigo{\ch{Re}{continueremo} il canto \ch{La}{delle} tue \ch{Mi}{lodi,}}
\rigo{\ch{Re}{noi} con la nostra \ch{La}{vita} con \ch{Mi}{Te.}}
\strofa
\rigo{\ch{La}{Ed} ora, \ch{Re}{via!} A por\ch{La}{tare} l'a\ch{Re}{more} nel \ch{La}{mondo,}}
\rigo{\ch{Re}{carità}  \ch{Mi}{nelle} case, nei \ch{Re}{campi,} nella cit\ch{La}{tà.}}
\rigo{\ch{Re}{Liberi} a por\ch{La}{tare} l'a\ch{Re}{more} nel \ch{La}{mondo,}}
\rigo{\ch{Re}{verità} \ch{Mi}{nelle} scuole, in uf\ch{Re}{ficio,} dove sa\ch{La}{rà:}}
\strofa
\rigo{\ch{Re}{e} sa\ch{La}{rà} \ch{Re}{vita} nuova! \ch{Fa\#m}{Fuori} il \ch{Re}{mondo} \ch{Fa\#m}{chiama}\ch{Mi}{}}
\rigo{\ch{Re}{anche} noi con il canto \ch{La}{delle} tue \ch{Mi}{lodi,} \ch{Re}{nella} \ch{La}{vita} con \ch{Mi}{Te.}}
\par\endgroup

\begingroup
\cantotesta{160}{Grandi cose}{}{}{290pt}
\rigo{\ch{Do}{Grandi} cos\ch{Sol}{e} ha fatto il Si\ch{Rem}{gnore} per no\ch{Lam}{i}}
\rigo{h\ch{Do}{a} fatto germog\ch{Fa}{liare} i fiori tra l\ch{Do}{e} rocce.       \ch{Sol}{}}
\rigo{G\ch{Do}{randi} cose\ch{Sol}{} ha fatto il Sig\ch{Rem}{nore} per noi\ch{Lam}{}}
\rigo{c\ch{Do}{i} ha riportati\ch{Fa}{} liberi all\ch{Do}{a} nostra terra.\ch{Sol}{}}
\rigo{LA-                     RE-     7                 MI-7}
\rigo{Ed ora possiamo cantare, possiamo gridare}
\rigo{l'amore che Dio h\ch{Fa}{a} versato su noi.      \ch{Sol}{}}
\strofa
\rigo{\ch{Do}{Tu} che sai\ch{Sol}{} strapp\ch{Rem}{are} dalla morte     \ch{Lam}{}}
\rigo{h\ch{Do}{ai} solleva\ch{Fa}{to} il nostro viso d\ch{Do}{alla} polvere.  \ch{Sol}{}}
\rigo{T\ch{Do}{u} che hai \ch{Sol}{sentito}\ch{Rem}{} il nostro pianto   \ch{Lam}{}}
\rigo{n\ch{Do}{el} nostro cuor\ch{Fa}{e} hai messo un seme di \ch{Do}{felicità.}  \ch{Sol}{}  RIT.}
\par\endgroup

\begingroup
\cantotesta{161}{Invochiamo la tua presenza}{}{}{478pt}
\rigo{\ch{Rem}{}\ch{Sib}{}\ch{Fa}{}\ch{Do}{Invochiamo} la tua presenza vieni Signor.}
\rigo{\ch{Rem}{}\ch{Sib}{}\ch{Fa}{}\ch{Do}{Invochiamo} la tua presenza scendi su di noi.}
\rigo{\ch{Solm}{}\ch{Rem}{}\ch{Sib}{}\ch{Do}{Vieni} Consolatore e dona pace e umiltà.}
\rigo{\ch{Solm}{}\ch{Rem}{}\ch{Sib}{}\ch{La4}{}\ch{La}{Acqua} viva d'amore questo cuore apriamo a te.}
\strofa
\begin{ritornello}
\rigo{\ch{Rem}{}\ch{Sib}{}\ch{Fa}{}\ch{Do}{Vieni} spirito, vieni spirito scendi su di noi}
\rigo{\ch{Rem}{}\ch{Sib}{}\ch{Fa}{}\ch{Do}{Vieni} spirito, vieni spirito scendi su di noi}
\rigo{\ch{Sib}{}\ch{Do}{}\ch{La}{}\ch{Rem}{}\ch{Sib}{}\ch{Do}{}\ch{Rem}{Vieni} su noi maranathà, vieni su noi spirito}
\rigo{\ch{Rem}{}\ch{Sib}{}\ch{Fa}{}\ch{Do}{Vieni} spirito, vieni spirito scendi su di noi}
\rigo{\ch{Rem}{}\ch{Sib}{}\ch{Fa}{}\ch{Do}{Vieni} spirito, vieni spirito scendi su di noi}
\strofa
\end{ritornello}
\rigo{\ch{Rem}{Scendi} su di noi (Intro)}
\strofa
\rigo{Invochiamo la tua presenza vieni Signor.}
\rigo{Invochiamo la tua presenza scendi su di noi.}
\rigo{Vieni luce dei cuori dona forza e fedeltà.}
\rigo{Fuoco eterno d'amore questa vita offriamo a te.}
\strofa
\begin{ritornello}
\rigo{\ch{Rem}{}\ch{Sib}{}\ch{Fa}{}\ch{Do}{Vieni} spirito, vieni spirito scendi su di noi}
\rigo{\ch{Rem}{}\ch{Sib}{}\ch{Fa}{}\ch{Do}{}\ch{Soldim}{Vieni} spirito, vieni spirito scendi su di noi}
\rigo{\ch{Fam}{}\ch{Do\#}{}\ch{Lab}{}\ch{Mib}{Vieni} spirito, vieni spirito scendi su di noi}
\rigo{\ch{Fam}{}\ch{Do\#}{}\ch{Lab}{}\ch{Mib}{Vieni} spirito, vieni spirito scendi su di noi}
\rigo{\ch{Do\#}{}\ch{Mib}{}\ch{Do}{}\ch{Fam}{}\ch{Do\#}{}\ch{Mib}{}\ch{Fam}{Vieni} su noi maranathà, vieni su noi spirito}
\rigo{\ch{Fam}{}\ch{Do\#}{}\ch{Lab}{}\ch{Mib}{Vieni} spirito, vieni spirito scendi su di noi}
\rigo{\ch{Fam}{}\ch{Do\#}{}\ch{Lab}{}\ch{Mib}{Vieni} spirito, vieni spirito scendi su di noi}
\strofa
\end{ritornello}
\rigo{\ch{Fam}{Scendi} su di noi}
\par\endgroup

\begingroup
\cantotesta{162}{Vieni spirito d'amore}{}{}{638pt}
\begin{ritornello}
\rigo{\ch{Mim}{}\ch{Lam}{Vieni,} vieni spirito d'amore}
\rigo{\ch{Mim}{}\ch{Sim}{}\ch{Sim7}{ad} insegnar le cose di Dio.}
\rigo{\ch{Mim}{}\ch{Lam}{Vieni,} vieni spirito di pace}
\rigo{\ch{Mim}{}\ch{Sim}{}\ch{Sim7}{a} suggerir le cose che Lui}
\rigo{\ch{Mim}{ha} detto a noi!}
\strofa
\end{ritornello}
\rigo{\ch{Mim}{}\ch{Lam}{Noi} ti invochiamo, spirito di Cristo}
\rigo{\ch{Mim}{}\ch{Sim}{}\ch{Sim7}{vieni} tu dentro di noi.}
\rigo{\ch{Mim}{}\ch{Lam}{Cambia} i nostri occhi, fa' che noi vediamo}
\rigo{\ch{Mim}{}\ch{Sim}{}\ch{Sim7}{la} bontà di Dio per noi.}
\strofa
\begin{ritornello}
\rigo{\ch{Mim}{}\ch{Lam}{Vieni,} vieni spirito d'amore}
\rigo{\ch{Mim}{}\ch{Sim}{}\ch{Sim7}{ad} insegnar le cose di Dio.}
\rigo{\ch{Mim}{}\ch{Lam}{Vieni,} vieni spirito di pace}
\rigo{\ch{Mim}{}\ch{Sim}{}\ch{Sim7}{a} suggerir le cose che Lui}
\rigo{\ch{Mim}{ha} detto a noi!}
\strofa
\end{ritornello}
\rigo{Vieni o spirito dai quattro venti}
\rigo{e soffia su chi non ha vita.}
\rigo{Vieni o spirito e soffia su di noi}
\rigo{perché anche noi riviviamo.}
\strofa
\begin{ritornello}
\rigo{\ch{Mim}{}\ch{Lam}{Vieni,} vieni spirito d'amore}
\rigo{\ch{Mim}{}\ch{Sim}{}\ch{Sim7}{ad} insegnar le cose di Dio.}
\rigo{\ch{Mim}{}\ch{Lam}{Vieni,} vieni spirito di pace}
\rigo{\ch{Mim}{}\ch{Sim}{}\ch{Sim7}{a} suggerir le cose che Lui}
\rigo{\ch{Mim}{ha} detto a noi!}
\strofa
\end{ritornello}
\rigo{Insegnaci a sperare, insegnaci ad amare}
\rigo{insegnaci a lodare Iddio.}
\rigo{Insegnaci a pregare, insegnaci la via}
\rigo{insegnaci Tu l'unità.}
\strofa
\begin{ritornello}
\rigo{\ch{Mim}{}\ch{Lam}{Vieni,} vieni spirito d'amore}
\rigo{\ch{Mim}{}\ch{Sim}{}\ch{Sim7}{ad} insegnar le cose di Dio.}
\rigo{\ch{Mim}{}\ch{Lam}{Vieni,} vieni spirito di pace}
\rigo{\ch{Mim}{}\ch{Sim}{}\ch{Sim7}{a} suggerir le cose che Lui}
\rigo{\ch{Mim}{ha} detto a noi!}
\end{ritornello}
\par\endgroup

\begingroup\senzaaccordi
\cantotesta{163}{Vivere la vita}{}{}{422pt}
\rigo{Vivere la vita con le gioie e coi dolori di ogni giorno}
\rigo{è quello che Dio vuole da te.}
\strofa
\rigo{Vivere la vita e inabissarti nell’amore è il tuo destino}
\rigo{è quello che Dio vuole da te.}
\strofa
\rigo{Fare insieme agli altri la tua strada verso Lui correre con i fratelli tuoi...}
\rigo{Scoprirai allora il cielo dentro di te, una scia di luce lascerai.}
\strofa
\rigo{Vivere la vita è l’avventura più stupenda dell’amore,}
\rigo{è quello che Dio vuole da te.}
\strofa
\rigo{Vivere la vita e generare ogni momento il Paradiso}
\rigo{è quello che Dio vuole da te.}
\strofa
\rigo{Vivere perché ritorni al mondo l’unità perché dio sta nei fratelli tuoi...}
\rigo{Scoprirai allora il cielo dentro di te,}
\rigo{una scia di luce lascerai.}
\strofa
\rigo{(...musica)}
\strofa
\rigo{Vivere perché ritorni al mondo l’unità perché dio sta nei fratelli tuoi...}
\rigo{Scoprirai allora il cielo dentro di te, una scia di luce lascerai.}
\rigo{una scia di luce lascerai.}
\par\endgroup

\begingroup
\cantotesta{164}{Vocazione}{P. Sequeri}{}{508pt}
\rigo{\ch{Do}{Era} un giorno come \ch{Sol}{tanti} altri \ch{Fa}{}}
\rigo{e quel gi\ch{Sol}{orno} Lui pas\ch{Do}{sò.}       \ch{Fa}{}   \ch{Do}{}   \ch{Sol}{}}
\rigo{E\ch{Do}{ra} un uomo come tutt\ch{Sol}{i} gli altri \ch{Fa}{}}
\rigo{e passand\ch{Sol}{o} mi chiamò.\ch{Do}{}          \ch{Fa}{}   \ch{Do}{}   \ch{Mi}{}}
\rigo{C\ch{Lam}{ome} lo sapesse \ch{Mim}{che} il mio nome         \ch{Fa}{}}
\rigo{era p\ch{Sol}{roprio} quello,}
\rigo{c\ch{Do}{ome} mai vedess\ch{Sol}{e} proprio me        \ch{Fa}{}}
\rigo{nella Sua vit\ch{Sol}{a,} non lo so\ch{Do}{.}         \ch{Fa}{}   \ch{Do}{}   \ch{Sol}{}}
\rigo{E\ch{Do}{ra} un giorno come t\ch{Sol}{anti} altri  \ch{Fa}{}}
\rigo{e quel gi\ch{Sol}{orno} mi chiamò. \ch{Do}{}           \ch{Fa}{}   \ch{Do}{}   \ch{Mi}{}}
\strofa
\commento{Rit.}
\rigo{T\ch{Lam}{u,} Dio,\ch{Mim}{} che co\ch{Fa}{nosci} il nome mio  \ch{Sol}{}}
\rigo{f\ch{Lam}{a'} c\ch{Mim}{he} asco\ch{Fa}{ltando} la Tua v\ch{Sol}{oce}}
\rigo{i\ch{Do}{o} ricordi \ch{Sol}{dove} porta l\ch{Fa}{a} mia strada       \ch{Sol}{}}
\rigo{n\ch{Do}{ella} vita \ch{Sol}{all'incontro}\ch{Fa}{} con Te.       \ch{Do}{}}
\strofa
\rigo{Era un'alba triste e senza vita}
\rigo{e qualcuno mi chiamò.}
\rigo{Era un uomo come tanti altri}
\rigo{ma la voce, quella no;}
\rigo{quante volte un uomo}
\rigo{con il nome giusto mi ha chiamato,}
\rigo{una volta sola l'ho sentito}
\rigo{pronunciare con amore.}
\rigo{Era un uomo come nessun altro}
\rigo{e quel giorno mi chiamò.}
\strofa
\commento{Rit.}
\par\endgroup

\sottosezione{Alleluia}

\begingroup
\cantotesta{165}{Alleluia la tua festa}{}{}{263pt}
\commento{Rit.}
\rigo{A\ch{Do}{llelu}\ch{Fa}{ia} alleluia}
\rigo{A\ch{Do}{llelu}\ch{Sol}{ia} alleluia}
\rigo{A\ch{Do}{llelu}\ch{Fa}{ia} alleluia al\ch{Do}{lel}\ch{Sol}{uia} (bis\ch{Fa}{)}  \ch{Sol}{}}
\rigo{L\ch{Do}{a} tua fes\ch{Fa}{ta} non deve f\ch{Sol}{inire,}\ch{Do}{}}
\rigo{non deve f\ch{Lam}{inire} e \ch{Rem}{non} fi\ch{Sol7}{nirà} (\ch{Do}{bis)}      \ch{Do7}{}}
\strofa
\rigo{Perché\ch{Fa}{} la fest\ch{Sol}{a} siamo \ch{Do}{noi}    \ch{Lam}{}}
\rigo{che cammi\ch{Fa}{niamo} ve\ch{Sol}{rso} te          \ch{Do}{}   \ch{Do7}{}}
\rigo{perché la\ch{Fa}{} festa\ch{Sol}{} siamo \ch{Do}{noi}     \ch{Lam}{}}
\rigo{canta\ch{Re7}{ndo} insieme così…      \ch{Sol4}{}       \ch{Sol}{}}
\rigo{\ch{La}{} \ch{La}{} \ch{La}{} \ch{La}{} \ch{La}{} \ch{La}{} \ch{La}{}}
\strofa
\commento{Rit.}
\par\endgroup

\begingroup\senzaaccordi
\cantotesta{166}{Alleluia! È il tempo della gioia}{}{}{497pt}
\rigo{Alleluia, alleluia!}
\rigo{E' il tempo della gioia, è il tempo della gioia.}
\rigo{Alleluia, alleluia!}
\rigo{E' il tempo della gioia, è il tempo della gioia.}
\strofa
\rigo{Il Signore viene sulla terra!}
\rigo{Il Signore discende}
\rigo{a rivelare l'amore di Dio,}
\rigo{a rivelare il Cielo, a rivelare il Cielo!}
\rigo{Alleluia.}
\strofa
\rigo{Alleluia, Alleluia!}
\rigo{Con lui cammineremo, con lui cammineremo.}
\rigo{Alleluia, Alleluia!}
\rigo{Con lui cammineremo, con lui cammineremo.}
\strofa
\rigo{Il Signore viene sulla terra!}
\rigo{Il Signore discende}
\rigo{a illuminare la strada per noi,}
\rigo{a illuminare il viaggio, a illuminare il viaggio!}
\rigo{Alleluia.}
\strofa
\rigo{Alleluia, Alleluia!}
\rigo{In lui dimoreremo, in lui dimoreremo.}
\rigo{Alleluia, Alleluia!}
\rigo{In lui dimoreremo, in lui dimoreremo.}
\strofa
\rigo{Il Signore viene sulla terra!}
\rigo{Il Signore discende}
\rigo{a rinnovare la pace tra noi,}
\rigo{a rinnovare il cuore, a rinnovare il cuore!}
\rigo{Alleluia.}
\strofa
\rigo{Alleluia. Alleluia. Alleluia. Alleluia.}
\rigo{Alleluia. Alleluia. Alleluia.}
\par\endgroup

\begingroup
\cantotesta{167}{Alleluia, ed oggi ancora}{}{}{253pt}
\begin{ritornello}
\rigo{\ch{Rem}{}\ch{Solm}{}\ch{Do7}{Alleluia,} alleluia}
\rigo{\ch{Fa}{}\ch{Sib}{alleluia,} alleluia}
\rigo{\ch{Solm}{}\ch{La7}{alleluia,} alleluia}
\rigo{\ch{Rem}{alleluia,} alleluia...}
\strofa
\end{ritornello}
\rigo{\ch{Rem}{}\ch{Solm}{Ed} oggi ancora, mio Signore,}
\rigo{\ch{Do}{}\ch{Fa}{ascolterò} la Tua parola}
\rigo{\ch{Sib}{}\ch{Solm}{}\ch{La4}{}\ch{La}{che} mi guida nel cammino della vita.}
\strofa
\begin{ritornello}
\rigo{\ch{Rem}{}\ch{Solm}{}\ch{Do7}{Alleluia,} alleluia}
\rigo{\ch{Fa}{}\ch{Sib}{alleluia,} alleluia}
\rigo{\ch{Solm}{}\ch{La7}{alleluia,} alleluia}
\rigo{\ch{Rem}{alleluia,} alleluia...}
\end{ritornello}
\par\endgroup

\begingroup
\cantotesta{168}{Alleluia, la parola tua}{}{}{287pt}
\begin{ritornello}
\rigo{\ch{Re}{}\ch{La}{}\ch{Sim}{}\ch{Fa\#m}{Alleluia,} alleluia, alleluia,}
\rigo{\ch{Sol}{}\ch{Re}{}\ch{Mim}{}\ch{Sol}{}\ch{La}{alleluia,} alleluia, alleluia,}
\rigo{\ch{Re}{}\ch{La}{}\ch{Sim}{}\ch{Fa\#m}{alleluia,} alleluia, alleluia,}
\rigo{\ch{Sol}{}\ch{Re}{}\ch{Mim}{}\ch{Sol}{}\ch{La}{}\ch{Re}{alleluia,} alleluia, alleluia, alleluia.}
\strofa
\end{ritornello}
\rigo{\ch{Re}{}\ch{Mim}{}\ch{Re}{La} parola Tua ci accompagnerà}
\rigo{\ch{Mim}{}\ch{Fa\#m}{nel} cammino che porta verso Te,}
\rigo{\ch{Sol}{}\ch{Re}{solo} non sarò, Tu sarai con me}
\rigo{\ch{Mim}{}\ch{La}{e} la vita mia, e la vita mia}
\rigo{\ch{Re}{danzerà} con Te.}
\strofa
\begin{ritornello}
\rigo{\ch{Re}{}\ch{La}{}\ch{Sim}{}\ch{Fa\#m}{Alleluia,} alleluia, alleluia,}
\rigo{\ch{Sol}{}\ch{Re}{}\ch{Mim}{}\ch{Sol}{}\ch{La}{alleluia,} alleluia, alleluia,}
\rigo{\ch{Re}{}\ch{La}{}\ch{Sim}{}\ch{Fa\#m}{alleluia,} alleluia, alleluia,}
\rigo{\ch{Sol}{}\ch{Re}{}\ch{Mim}{}\ch{Sol}{}\ch{La}{}\ch{Re}{alleluia,} alleluia, alleluia, alleluia.}
\end{ritornello}
\par\endgroup

\begingroup\senzaaccordi
\cantotesta{169}{Alleluia, passeranno i cieli}{}{}{172pt}
\rigo{Alleluia alleluia alleluia,}
\rigo{alleluia alleluia alleluia,}
\rigo{alleluia alleluia.}
\strofa
\rigo{Passeranno i cieli e passerà la terra,}
\rigo{la sua Parola non passerà...}
\rigo{alleluia alleluia.}
\strofa
\rigo{Alleluia alleluia alleluia,}
\rigo{alleluia alleluia alleluia,}
\rigo{alleluia alleluia.}
\par\endgroup

\sottosezione{Offertorio}

\begingroup
\cantotesta{170}{Accogli i nostri doni}{}{}{420pt}
\rigo{\ch{Re}{}\ch{Fa\#m7}{}\ch{Mim7}{}\ch{La4}{}\ch{La}{Accogli} i nostri doni, Dio dell'universo,}
\rigo{\ch{Sim}{}\ch{Fa\#m7}{}\ch{Sol}{}\ch{La4}{}\ch{La}{in} questo misterioso incontro col tuo Figlio.}
\rigo{\ch{Sim}{}\ch{Fa\#m}{}\ch{Sol}{}\ch{Re}{Ti} offriamo il pane che Tu ci dai:}
\rigo{\ch{Sim}{}\ch{Mim7}{}\ch{Sol}{}\ch{La}{}\ch{Re}{}\ch{La}{trasformalo} in Te, Signor.}
\strofa
\begin{ritornello}
\rigo{\ch{Re}{}\ch{Mim}{}\ch{La4}{}\ch{La}{Benedetto} nei secoli il Signore}
\rigo{\ch{Sol}{}\ch{La}{}\ch{Fa\#m}{}\ch{Sim}{infinita} sorgente della vita.}
\rigo{\ch{Mim7}{}\ch{La6}{}\ch{Re}{}\ch{Re/Do\#}{}\ch{Sim}{Benedetto} nei secoli}
\rigo{\ch{Sol}{}\ch{La}{}\ch{Re}{}\ch{Sol}{}\ch{Re}{Benedetto} nei secoli.}
\strofa
\end{ritornello}
\rigo{\ch{Re}{}\ch{Fa\#m7}{}\ch{Mim7}{}\ch{La4}{}\ch{La}{Accogli} i nostri doni, Dio dell'universo,}
\rigo{\ch{Sim}{}\ch{Fa\#m7}{}\ch{Sol}{}\ch{La4}{}\ch{La}{in} questo misterioso incontro col tuo Figlio.}
\rigo{\ch{Sim}{}\ch{Fa\#m}{}\ch{Sol}{}\ch{Re}{Ti} offriamo il vino che Tu ci dai:}
\rigo{\ch{Sim}{}\ch{Mim7}{}\ch{Sol}{}\ch{La}{}\ch{Re}{}\ch{La}{trasformalo} in Te, Signor.}
\strofa
\begin{ritornello}
\rigo{\ch{Re}{}\ch{Mim}{}\ch{La4}{}\ch{La}{Benedetto} nei secoli il Signore}
\rigo{\ch{Sol}{}\ch{La}{}\ch{Fa\#m}{}\ch{Sim}{infinita} sorgente della vita.}
\rigo{\ch{Mim7}{}\ch{La6}{}\ch{Re}{}\ch{Re/Do\#}{}\ch{Sim}{Benedetto} nei secoli}
\rigo{\ch{Sol}{}\ch{La}{}\ch{Re}{}\ch{Sol}{}\ch{Re}{Benedetto} nei secoli.}
\strofa
\rigo{\ch{Mib}{}\ch{Fam}{}\ch{Sib4}{}\ch{Sib}{Benedetto} nei secoli il Signore}
\rigo{\ch{Lab}{}\ch{Sib}{}\ch{Solm}{}\ch{Dom}{infinita} sorgente della vita.}
\rigo{\ch{Fam7}{}\ch{Sib6}{}\ch{Mib}{}\ch{Mib/Re}{}\ch{Dom}{Benedetto} nei secoli}
\rigo{\ch{Lab}{}\ch{Sib}{}\ch{Lab}{}\ch{Mib}{Benedetto} nei secoli.}
\end{ritornello}
\par\endgroup

\begingroup
\cantotesta{171}{Antica, eterna danza}{Gen Verde}{}{577pt}
\rigo{\ch{Sol}{}\ch{Re}{}\ch{Mi}{Spighe} d'oro al vento,}
\rigo{\ch{Mim}{}\ch{Re}{}\ch{Do}{antica} eterna danza}
\rigo{\ch{Lam}{}\ch{Re}{}\ch{Mim}{}\ch{Re}{per} fare un solo pane}
\rigo{\ch{Do}{}\ch{La(no3)}{}\ch{Si}{spezzato} sulla mensa.}
\strofa
\rigo{\ch{Sol}{}\ch{Re}{}\ch{Mi}{Grappoli} dei colli,}
\rigo{\ch{Mim}{}\ch{Re}{}\ch{Do}{profumo} di letizia}
\rigo{\ch{Lam}{}\ch{Re}{}\ch{Mim}{}\ch{Re}{per} fare un solo vino,}
\rigo{\ch{Lam}{}\ch{Sim}{}\ch{Mi}{bevanda} della grazia.}
\strofa
\rigo{\ch{Sol7+}{}\ch{Mim7}{}\ch{Lam7}{}\ch{Mi}{}}
\rigo{\ch{Sol7+}{}\ch{Do}{}\ch{Sol}{}}
\strofa
\rigo{\ch{Sol}{}\ch{Re}{}\ch{Mi}{Con} il pane e il vino,}
\rigo{\ch{Mim}{}\ch{Re}{}\ch{Do}{Signore} ti doniamo}
\rigo{\ch{Lam}{}\ch{Re}{}\ch{Mim}{}\ch{Re}{Le} nostre gioie pure}
\rigo{\ch{Do}{}\ch{La(no3)}{}\ch{Si}{Le} attese le paure,}
\strofa
\rigo{\ch{Sol}{}\ch{Re}{}\ch{Mi}{frutti} del lavoro}
\rigo{\ch{Mim}{}\ch{Re}{}\ch{Do}{e} fede nel futuro}
\rigo{\ch{Lam}{}\ch{Re}{}\ch{Mim}{}\ch{Re}{la} voglia di cambiare}
\rigo{\ch{Lam}{}\ch{Sim}{}\ch{Mi}{e} di ricominciare}
\strofa
\rigo{\ch{Sol}{}\ch{Re}{} / \ch{Mim}{} / \ch{Mim}{}\ch{Re}{} / \ch{Do}{} / \ch{Lam}{}\ch{Re}{} /}
\rigo{\ch{Mim}{}\ch{Re}{} / \ch{Do}{}\ch{La}{} / \ch{Si4}{}\ch{Si}{}}
\strofa
\rigo{\ch{Sol}{}\ch{Re}{}\ch{Mi}{Dio} della speranza}
\rigo{\ch{Mim}{}\ch{Re}{}\ch{Do}{Sorgente} di ogni dono}
\rigo{\ch{Lam}{}\ch{Re}{}\ch{Mim}{}\ch{Re}{Accogli} questa offerta}
\rigo{\ch{Do}{}\ch{La(no3)}{}\ch{Si}{Che} insieme ti portiamo}
\strofa
\rigo{\ch{Sol}{}\ch{Re}{}\ch{Mi}{Dio} dell'universo}
\rigo{\ch{Mim}{}\ch{Re}{}\ch{Do}{Raccogli} chi è disperso}
\rigo{\ch{Lam}{}\ch{Re}{}\ch{Mim}{}\ch{Re}{E} facci tutti chiesa,}
\rigo{\ch{Lam}{}\ch{Sim}{}\ch{Mi}{una} cosa in Te.}
\par\endgroup

\begingroup
\cantotesta{172}{Benedetto sei tu}{M. Frisina}{}{331pt}
\rigo{B\ch{La7}{ened}\ch{Re}{ett}\ch{Mim}{o} se\ch{Fa\#m}{i} Tu \ch{Sol}{Dio} \ch{Re}{del}\ch{Sim}{l'universo!}}
\rigo{D\ch{Sol}{alla}\ch{La}{} Tu\ch{Re}{a} b\ch{Sim}{ontà}}
\rigo{a\ch{Sol}{bbia}\ch{La}{mo} \ch{Re}{ric}\ch{Sim}{evuto} questo pane,}
\rigo{f\ch{Sol}{rutt}\ch{La}{o} d\ch{Re}{ell}\ch{Sim}{a} terra e del nostro lavoro}
\rigo{l\ch{Sol}{o} pr\ch{Fa\#m}{esentiamo} a Te}
\rigo{p\ch{Mim}{erch}\ch{Re}{é} diventi per noi}
\rigo{c\ch{Mim}{ibo} \ch{La}{di} vita eterna.}
\strofa
\rigo{Benedetto sei Tu Dio dell'universo!}
\rigo{Dalla Tua bontà}
\rigo{abbiamo ricevuto questo vino,}
\rigo{frutto della vite e del nostro lavoro}
\rigo{lo presentiamo a Te}
\rigo{perché diventi per noi}
\rigo{bevanda di salvezza.}
\strofa
\rigo{B\ch{Re}{ene}\ch{Mim}{dett}\ch{Fa\#m}{o} sei Tu Signor!}
\rigo{B\ch{Sol}{ened}\ch{La}{ett}\ch{Re}{o} sei Tu Signor!}
\par\endgroup

\begingroup
\cantotesta{173}{Benedici, o Signore}{Gen Rosso}{}{481pt}
\rigo{\ch{Rem}{Nebbia} e freddo, giorni lunghi e amari      \ch{Do}{}}
\rigo{mentre il seme muore.\ch{Rem}{}}
\rigo{P\ch{Fa}{oi} il prodigio antico e sempre nuovo   \ch{Do}{}}
\rigo{del primo filo d'erba\ch{Sib}{.}}
\rigo{E nel ven\ch{Fa}{to} dell'estate \ch{Do}{ondeggiano} \ch{Rem}{le} spighe           \ch{Fa}{}}
\rigo{avrem\ch{Do}{o} ancora pane.     \ch{La}{}       \ch{Re}{}}
\strofa
\commento{Rit.}
\rigo{B\ch{Sol}{enedici}\ch{Re}{,} oh S\ch{Sol}{ignore,}\ch{Re}{}}
\rigo{q\ch{Do}{uesta} offerta \ch{Sol}{che} portiamo a T\ch{La}{e,}}
\rigo{f\ch{Sol}{acci} "u\ch{Re}{no"} come i\ch{Sim}{l} pane      \ch{Fa\#m}{}}
\rigo{c\ch{Mi}{he} anche oggi \ch{Sol}{hai} dato a n\ch{Re}{oi.}}
\strofa
\rigo{Nei filari dopo il lungo inverno}
\rigo{fremono le viti.}
\rigo{La rugiada avvolge nel silenzio}
\rigo{i primi tralci verdi.}
\rigo{Poi i colori dell'autunno}
\rigo{coi grappoli maturi}
\rigo{avremo ancora vino.}
\strofa
\commento{Rit.}
\strofa
\rigo{Benedici, oh Signore,}
\rigo{questa offerta che portiamo a Te,}
\rigo{facci "uno" come il vino}
\rigo{che anche oggi hai dato a noi.}
\par\endgroup

\begingroup
\cantotesta{174}{Ecco quel che abbiamo}{Gen Verde}{}{699pt}
\begin{ritornello}
\rigo{\ch{Sol}{}\ch{Re/Sol}{}\ch{Do/Sol}{}\ch{Sol}{Ecco} quel che abbiamo, nulla ci appartiene ormai}
\rigo{\ch{Mim}{}\ch{Sim7}{}\ch{Do}{}\ch{Re7}{ecco} i frutti della terra che Tu moltiplicherai.}
\rigo{\ch{Sol}{}\ch{Re/Sol}{}\ch{Do/Sol}{}\ch{Sol}{Ecco} queste mani, puoi usarle se lo vuoi}
\rigo{\ch{Mim}{}\ch{Sim7}{}\ch{Do}{}\ch{Re}{}\ch{Sol}{per} dividere nel mondo il pane che Tu hai dato a noi.}
\strofa
\end{ritornello}
\rigo{\ch{Mim}{}\ch{Sim}{Solo} una goccia hai messo fra le mani mie}
\rigo{\ch{Mi+}{}\ch{Lam7}{}\ch{Do}{}\ch{Re/Do}{}\ch{Sim}{}\ch{Mi}{solo} una goccia che Tu ora chiedi a me... (vocalizzo)}
\rigo{\ch{Lam}{}\ch{Sim}{}\ch{Do}{}\ch{Sim}{}\ch{Do}{}\ch{Re}{Una} goccia che, in mano a Te, una pioggia diventerà}
\rigo{\ch{Do7+}{e} la terra feconderà.}
\strofa
\begin{ritornello}
\rigo{\ch{Sol}{}\ch{Re/Sol}{}\ch{Do/Sol}{}\ch{Sol}{Ecco} quel che abbiamo, nulla ci appartiene ormai}
\rigo{\ch{Mim}{}\ch{Sim7}{}\ch{Do}{}\ch{Re7}{ecco} i frutti della terra che Tu moltiplicherai.}
\rigo{\ch{Sol}{}\ch{Re/Sol}{}\ch{Do/Sol}{}\ch{Sol}{Ecco} queste mani, puoi usarle se lo vuoi}
\rigo{\ch{Mim}{}\ch{Sim7}{}\ch{Do}{}\ch{Re}{}\ch{Sol}{per} dividere nel mondo il pane che Tu hai dato a noi.}
\strofa
\end{ritornello}
\rigo{\ch{Mim}{}\ch{Sim}{Le} nostre gocce, pioggia tra le mani Tue,}
\rigo{\ch{Mi+}{}\ch{Lam7}{}\ch{Do}{}\ch{Re/Do}{}\ch{Sim}{}\ch{Mi}{saranno} linfa di una nuova civiltà... (vocalizzo)}
\rigo{\ch{Lam}{}\ch{Sim}{}\ch{Do}{}\ch{Sim}{}\ch{Do}{}\ch{Re}{E} la terra preparerà la festa del pane che}
\rigo{\ch{Do7+}{}\ch{Mi}{ogni} uomo condividerà.}
\strofa
\begin{ritornello}
\rigo{\ch{Sol}{}\ch{Re/Sol}{}\ch{Do/Sol}{}\ch{Sol}{Ecco} quel che abbiamo, nulla ci appartiene ormai}
\rigo{\ch{Mim}{}\ch{Sim7}{}\ch{Do}{}\ch{Re7}{ecco} i frutti della terra che Tu moltiplicherai.}
\rigo{\ch{Sol}{}\ch{Re/Sol}{}\ch{Do/Sol}{}\ch{Sol}{Ecco} queste mani, puoi usarle se lo vuoi}
\rigo{\ch{Mim}{}\ch{Sim7}{}\ch{Do}{}\ch{Re}{}\ch{Sol}{per} dividere nel mondo il pane che Tu hai dato a noi.}
\strofa
\end{ritornello}
\rigo{\ch{Sol}{}\ch{Re/Sol}{}\ch{Do/Sol}{}\ch{Sol}{Sulle} strade il vento da lontano porterà}
\rigo{\ch{Mim}{}\ch{Sim7}{}\ch{Do}{}\ch{Re7}{il} profumo del frumento, che tutti avvolgerà.}
\rigo{\ch{Sol}{}\ch{Re/Sol}{}\ch{Do/Sol}{}\ch{Sol}{E} sarà l'amore che il raccolto spartirà}
\rigo{\ch{Mim}{}\ch{Sim7}{}\ch{Do}{}\ch{Re}{}\ch{Sol}{e} il miracolo del pane in terra si ripeterà}
\par\endgroup

\begingroup
\cantotesta{175}{Frutto della nostra terra}{}{}{515pt}
\rigo{\ch{Sol}{}\ch{Do}{}\ch{Sol}{}\ch{Re}{Frutto} della nostra terra, del lavoro di ogni uomo:}
\rigo{\ch{Mim}{}\ch{Sim}{}\ch{Do}{}\ch{Re}{pane} della nostra vita, cibo della quotidianità.}
\rigo{\ch{Sol}{}\ch{Do}{}\ch{Sol}{}\ch{Re}{Tu} che lo prendevi un giorno, lo spezzavi per i tuoi,}
\rigo{\ch{Mim}{}\ch{Sim}{}\ch{Do}{}\ch{Re4}{}\ch{Re}{oggi} vieni in questo pane, cibo vero dell'umanità.}
\strofa
\begin{ritornello}
\rigo{\ch{Sol}{}\ch{Re}{}\ch{Mim}{}\ch{Sim}{E} sarò pane, e sarò vino nella mia vita, nelle tue mani.}
\rigo{\ch{Do}{}\ch{Re}{Ti} accoglierò dentro di me,}
\rigo{\ch{Mim7}{}\ch{Do7+}{farò} di me un'offerta viva,}
\rigo{\ch{Lam}{}\ch{Re4}{}\ch{Sol}{}\ch{Do}{}\ch{Re}{}\ch{Sol}{un} sacrificio gradito a Te.}
\strofa
\end{ritornello}
\rigo{\ch{Sol}{}\ch{Do}{}\ch{Sol}{}\ch{Re}{Frutto} della nostra terra, del lavoro di ogni uomo:}
\rigo{\ch{Mim}{}\ch{Sim}{}\ch{Do}{}\ch{Re}{vino} delle nostre vigne, sulla mensa dei fratelli tuoi}
\rigo{\ch{Sol}{}\ch{Do}{}\ch{Sol}{}\ch{Re}{Tu} che lo prendevi un giorno, lo bevevi con i tuoi,}
\rigo{\ch{Mim}{}\ch{Sim}{}\ch{Do}{}\ch{Re4}{}\ch{Re}{oggi} vieni in questo vino e ti doni per la vita mia}
\strofa
\begin{ritornello}
\rigo{\ch{Sol}{}\ch{Re}{}\ch{Mim}{E} sarò pane, e sarò vino nella mia vita,}
\rigo{\ch{Sim}{}\ch{Do}{}\ch{Re}{nelle} tue mani. Ti accoglierò dentro di me,}
\rigo{\ch{Mim7}{}\ch{Do7+}{farò} di me un'offerta viva,}
\rigo{\ch{Lam}{}\ch{Re}{}\ch{Mim}{}\ch{Do}{un} sacrificio gradito a Te,}
\rigo{\ch{Lam}{}\ch{Re4}{}\ch{Sol}{un} sacrificio gradito a Te.}
\end{ritornello}
\par\endgroup

\begingroup\senzaaccordi
\cantotesta{176}{Ogni mia parola}{}{}{158pt}
\rigo{Come la pioggia e la neve, scendono giù dal cielo}
\rigo{e non vi ritornano senza irrigare e far germogliare la terra,}
\rigo{così ogni mia parola non ritornerà a me,}
\rigo{senza operare quanto desidero,}
\rigo{senza aver compiuto ciò per cui l'avevo mandata.}
\rigo{Ogni mia parola. Ogni mia parola.}
\par\endgroup

\begingroup
\cantotesta{177}{Se m'accogli}{P. Sequeri}{}{440pt}
\rigo{\ch{Do}{Tra} le mani non ho nien\ch{Fa}{te} so che Tu mi\ch{Do}{} accogliera\ch{Mi}{i}     \ch{Lam}{}}
\rigo{chiedo so\ch{Do}{lo} di restare a\ch{Fa}{ccanto} a Te.       \ch{Sol}{}}
\rigo{Sono \ch{Do}{ricco} sola\ch{Fa}{mente} dell'amor\ch{Do}{e} che \ch{Mi}{mi} dai:    \ch{Lam}{}}
\rigo{è per que\ch{Do}{lli} che non l'hanno\ch{Fa}{} avuto mai.        \ch{Sol}{}}
\strofa
\commento{Rit.}
\rigo{S\ch{Do}{e} m'accogli, mio Signor\ch{Mi}{e,} altro non Ti\ch{Lam}{} chiederò       \ch{Do}{}}
\rigo{e per sem\ch{Fa}{pre} la Tua stra\ch{Rem}{da} la mia strada\ch{Mi}{} resterà       \ch{Lam}{}}
\rigo{nella\ch{Fa}{} gioia e n\ch{Do}{el} dolore fino \ch{Fa}{a} quando Tu vor\ch{Mi}{rai}}
\rigo{con la ma\ch{Do}{no} nella Tu\ch{Sol}{a} camminerò.    \ch{Fa}{}      \ch{Do}{}}
\strofa
\rigo{Io Ti prego con il cuore, so che Tu mi ascolterai}
\rigo{rendi forte la mia fede più che mai.}
\rigo{Tieni accesa la mia luce fino al giorno che Tu sai,}
\rigo{con i miei fratelli incontro a Te verrò.}
\strofa
\commento{Rit.}
\par\endgroup

\begingroup
\cantotesta{178}{Servo per amore}{}{}{457pt}
\rigo{\ch{Sim}{Una} notte di sudore}
\rigo{sulla barca in mezzo al mare \ch{Re}{}}
\rigo{e mentre il c\ch{La}{ielo} s'imbianca già    \ch{Fa\#m}{}}
\rigo{tu guardi\ch{Sol}{} le tue reti vuo\ch{Sim}{te.}}
\rigo{Ma la voc\ch{Re}{e} che ti chiama}
\rigo{un altro \ch{La}{mare} ti mostrerà   \ch{Re}{}}
\rigo{e sulle r\ch{Sol}{ive} di ogni cuor\ch{Re}{e}}
\rigo{le tue re\ch{Mim}{ti} g\ch{Sol}{ette}\ch{Re}{rai.}}
\strofa
\commento{Rit.}
\rigo{O\ch{Sim}{ffri} la vita tu\ch{Re}{a}}
\rigo{come Mari\ch{La}{a} ai piedi \ch{Fa\#m}{della} croce  \ch{Sol}{}}
\rigo{e\ch{Sim}{} sarai \ch{Re}{servo} di ogni uomo,    \ch{La}{}}
\rigo{s\ch{Re}{ervo} per amore,    \ch{Sol}{}}
\rigo{s\ch{Re}{acerdote} dell'\ch{Mim}{umanità.}    \ch{Sol}{}       \ch{Re}{}}
\strofa
\rigo{Avanzavi nel silenzio}
\rigo{tra le lacrime e speravi}
\rigo{che il seme sparso davanti a te}
\rigo{cadesse sulla buona terra.}
\rigo{Ora il cuore tuo è in festa}
\rigo{perché il grano biondeggia ormai,}
\rigo{è maturato sotto il sole,}
\rigo{puoi riporlo nei granai.}
\strofa
\commento{Rit.}
\par\endgroup

\sottosezione{Santo}

\begingroup
\cantotesta{179}{Santo (Dio del cielo)}{}{}{597pt}
\rigo{\ch{Sol}{}\ch{Re}{}\ch{Mim}{}\ch{Re}{}\ch{Do}{}\ch{Sol}{}\ch{Re}{}\ch{Sol}{}\ch{Re}{Santo,} Santo, Santo il Signore, Dio del cielo, Dio del cielo.}
\rigo{\ch{Sol}{}\ch{Re}{}\ch{Mim}{}\ch{Re}{}\ch{Do}{}\ch{Sol}{}\ch{Re}{}\ch{Sol}{}\ch{Re}{Santo,} Santo, Santo il Signore, Dio del cielo, Dio del cielo.}
\strofa
\rigo{\ch{Do}{I} cieli e la terra con noi,}
\rigo{\ch{Re}{si} uniscono a cantare per Te,}
\rigo{\ch{Si7}{per} questa nostra vita che Tu,}
\rigo{\ch{Mim}{regali} ogni giorno a noi,}
\rigo{\ch{Do}{e} canta la mia vita Signore,}
\rigo{\ch{Re}{un} canto imparato da Te,}
\rigo{\ch{Si7}{un} canto imparato da Te.}
\strofa
\begin{ritornello}
\rigo{\ch{Mi}{}\ch{Sol}{Osanna,} osanna, osanna a Dio nell'alto dei cieli,}
\rigo{\ch{Re}{}\ch{Do}{Osanna} a Dio nell'alto dei cieli, osanna nei cieli e dentro di noi. (bis)}
\strofa
\end{ritornello}
\rigo{\ch{Sol}{}\ch{Re}{}\ch{Mim}{}\ch{Re}{}\ch{Do}{}\ch{Sol}{}\ch{Re}{}\ch{Sol}{}\ch{Re}{Benedetto} colui che viene, viene per noi, Dio del cielo.}
\rigo{\ch{Sol}{}\ch{Re}{}\ch{Mim}{}\ch{Re}{}\ch{Do}{}\ch{Sol}{}\ch{Re}{}\ch{Sol}{}\ch{Re}{Benedetto} colui che viene, viene per noi, Dio del cielo.}
\strofa
\rigo{\ch{Do}{I} cieli e la terra con noi}
\rigo{\ch{Re}{si} uniscono a cantare per Te,}
\rigo{\ch{Si7}{per} questa nostra vita che Tu,}
\rigo{\ch{Mim}{regali} ogni giorno a noi,}
\rigo{\ch{Do}{e} canta la mia vita Signore,}
\rigo{\ch{Re}{un} canto imparato da Te,}
\rigo{\ch{Si7}{un} canto imparato da Te.}
\strofa
\begin{ritornello}
\rigo{\ch{Mi}{}\ch{Sol}{Osanna,} osanna, osanna a Dio nell'alto dei cieli,}
\rigo{\ch{Re}{}\ch{Do}{Osanna} a Dio nell'alto dei cieli, osanna nei cieli e dentro di noi. (bis)}
\end{ritornello}
\par\endgroup

\sottosezione{Comunione}

\begingroup\senzaaccordi
\cantotesta{180}{Beati voi}{}{}{102pt}
\rigo{Beati, beati voi che mi seguite:}
\rigo{vostro sarà il regno dei cieli.}
\strofa
\rigo{Beati coloro che muoiono}
\rigo{= Alleluia beati coloro che muoiono.}
\par\endgroup

\begingroup
\cantotesta{181}{Benedetto Signore (cerco solo te)}{}{}{665pt}
\rigo{\ch{Re}{Cerco} solo Te mio Signor}
\rigo{\ch{Sol}{Perché} solo Tu \ch{La}{dai} gioia al mio cuo\ch{Re}{re,}}
\rigo{si rallegra l'anima mi\ch{Sol}{a}}
\rigo{solo in \ch{La}{Te,} solo in Te.\ch{Re}{}}
\strofa
\rigo{\ch{Re}{Cerco} solo Te mio Signor}
\rigo{\ch{Sol}{Perché} la Tua vi\ch{La}{a} conduce alla vi\ch{Re}{ta,}}
\rigo{si rallegra l'anima mi\ch{Sol}{a}}
\rigo{solo in \ch{La}{Te,} solo in Te.\ch{Re}{}}
\strofa
\begin{ritornello}
\rigo{\ch{Re4}{}\ch{Re}{Benedetto} il Si\ch{Sol}{gnore,}}
\rigo{\ch{La}{benedetto} il Tuo \ch{Re}{nome,}}
\rigo{\ch{Sim}{come} un tenero \ch{Sol}{padre} sei verso di \ch{La}{me,} mio Signor.}
\rigo{\ch{Re}{Bene}\ch{Re/Fa\#}{detto} il Si\ch{Sol}{gnore,}}
\rigo{\ch{La}{benedetto} il Tuo \ch{Sim9}{nome,}}
\rigo{\ch{Re/Fa\#}{dalle} Tue mani \ch{Sol}{questa} mia vita}
\rigo{ri\ch{La}{ceve} salvezza e a\ch{Re}{mor.}}
\end{ritornello}
\rigo{\ch{Re4}{}\ch{Re}{}}
\strofa
\rigo{\ch{Mi}{Cerco} solo Te mio Signor}
\rigo{\ch{La}{perché} la Tua gra\ch{Si}{zia} rimane in eter\ch{Mi}{no,}}
\rigo{si rallegra l'anima mi\ch{La}{a}}
\rigo{solo in \ch{Si}{Te,} solo in Te.\ch{Mi}{}}
\strofa
\rigo{\ch{Mi}{Cerco} solo Te mio Signor}
\rigo{\ch{La}{perché} mi coro\ch{Si}{ni} di misericor\ch{Mi}{dia,}}
\rigo{si rallegra l'anima mi\ch{La}{a}}
\rigo{solo in \ch{Si}{}\ch{Mi4}{Te,} solo in Te.\ch{Mi}{}}
\strofa
\begin{ritornello}
\rigo{\ch{Mi}{Benedetto} il Si\ch{La}{gnore,}}
\rigo{\ch{Si}{benedetto} il Tuo \ch{Mi}{nome,}}
\rigo{\ch{Do\#m}{come} un tenero \ch{La}{padre} sei verso di \ch{Si}{me,} mio Signor.}
\rigo{\ch{Mi}{Bene}\ch{Mi/Sol\#}{detto} il Si\ch{La}{gnore,}}
\rigo{\ch{Si}{benedetto} il Tuo \ch{Do\#m9}{nome,}}
\rigo{\ch{Mi/Sol\#}{dalle} Tue mani \ch{La}{questa} mia vita}
\rigo{ri\ch{Si}{ceve} salvezza e a\ch{Mi}{mor.}}
\end{ritornello}
\rigo{\ch{Mi4}{}\ch{Mi}{}}
\par\endgroup

\begingroup\senzaaccordi
\cantotesta{182}{C'è ancora un sole}{}{\bEG}{286pt}
\rigo{Ho solo un pane ma se spezzarlo tu vuoi con me}
\rigo{crescera' la letizia di marciare insieme fratel}
\rigo{Ho qui un po' d'acqua un sorso solo vuoi berlo tu ?}
\rigo{anche l'acqua di fonte a spartirlaa e' di piu'.}
\strofa
\rigo{Vecchie parole non han più senso ne voce qui,}
\rigo{sotto il fiato del vento ogni antico ricordo svanì.}
\rigo{Parole nuove sentiamo nascere in fondo al cuor,}
\rigo{sono fatte di passi di fatica, speranza e sudor.}
\strofa
\rigo{C'e' ancora un sole l'abbiamo ritrovato}
\rigo{seguiva le ombre mobili dei passi sul sentir}
\rigo{C'e' ancora un sole scaldava le tue spalle}
\rigo{quio toccai lo zaino che tu portavi per me.}
\par\endgroup

\begingroup
\cantotesta{183}{Canto di San Damiano}{}{}{318pt}
\rigo{\ch{Re}{Ogni} u\ch{La7}{omo} sem\ch{Sim}{plice} p\ch{La7}{orta} in \ch{Re}{cuore} un sogno,\ch{La7}{}            \ch{Sim7}{}        \ch{La7}{}}
\rigo{con amore ed umiltà potrà costruirlo.}
\rigo{Se davvero tu saprai vivere umilmente}
\rigo{più felice tu sarai anche senza niente.}
\strofa
\commento{Rit.}
\rigo{S\ch{Mi}{e} vorrai ogni \ch{Sol}{giorno} con il tu\ch{Re}{o} sudore       \ch{La7}{}}
\rigo{u\ch{Mi}{na} pietra dopo\ch{Sol}{} l'altra alto ar\ch{Re}{riverai.}       \ch{La7}{}}
\strofa
\rigo{E le gioie semplici sono le più belle,}
\rigo{sono quelle che alla fine sono le più grandi.}
\strofa
\commento{Rit.}
\strofa
\rigo{Nella vita semplice troverai la strada}
\rigo{che la calma donerà al tuo cuore puro.}
\strofa
\commento{Rit.}
\par\endgroup

\begingroup
\cantotesta{184}{Come fuoco vivo}{Gen Rosso - Gen Verde}{}{624pt}
\begin{ritornello}
\rigo{\ch{La}{}\ch{Mi/Sol\#}{}\ch{Fa\#m}{Come} fuoco vivo si accende in noi}
\rigo{\ch{Sim7}{}\ch{Mi}{}\ch{La}{un'immensa} felicità}
\rigo{\ch{Re}{}\ch{Mi}{}\ch{La/Do\#}{}\ch{Re}{che} mai più nessuno ci toglierà}
\rigo{\ch{Sim7}{}\ch{Mi4}{}\ch{Mi}{perché} tu sei ritornato.}
\rigo{\ch{La}{}\ch{Mi/Sol\#}{}\ch{Fa\#m}{Chi} potrà tacere da ora in poi}
\rigo{\ch{Sim7}{}\ch{Mi}{}\ch{La}{che} sei tu in cammino con noi,}
\rigo{\ch{Re}{}\ch{Mi}{}\ch{La}{}\ch{Re}{che} la morte è vinta per sempre, che}
\rigo{\ch{Sim}{}\ch{Mi4}{}\ch{Mi}{ci} hai ridonato la vita.}
\strofa
\end{ritornello}
\rigo{\ch{Fa\#m}{}\ch{Re}{}\ch{La}{Spezzi} il pane davanti a noi,}
\rigo{\ch{Mi}{mentre} il sole è al tramonto:}
\rigo{\ch{Mim/Sol}{}\ch{Fa\#}{}\ch{Sim}{ora} gli occhi ti vedono,}
\rigo{\ch{Re}{}\ch{Mi4}{}\ch{Mi}{sei} Tu! Resta con noi.}
\strofa
\begin{ritornello}
\rigo{\ch{La}{}\ch{Mi/Sol\#}{}\ch{Fa\#m}{Come} fuoco vivo si accende in noi}
\rigo{\ch{Sim7}{}\ch{Mi}{}\ch{La}{un'immensa} felicità}
\rigo{\ch{Re}{}\ch{Mi}{}\ch{La/Do\#}{}\ch{Re}{che} mai più nessuno ci toglierà}
\rigo{\ch{Sim7}{}\ch{Mi4}{}\ch{Mi}{perché} tu sei ritornato.}
\rigo{\ch{La}{}\ch{Mi/Sol\#}{}\ch{Fa\#m}{Chi} potrà tacere da ora in poi}
\rigo{\ch{Sim7}{}\ch{Mi}{}\ch{La}{che} sei tu in cammino con noi,}
\rigo{\ch{Re}{}\ch{Mi}{}\ch{La}{}\ch{Re}{che} la morte è vinta per sempre, che}
\rigo{\ch{Sim}{}\ch{Mi4}{}\ch{Mi}{ci} hai ridonato la vita.}
\strofa
\end{ritornello}
\rigo{E per sempre ti mostrerai}
\rigo{in quel gesto d'amore:}
\rigo{mani che ancora spezzano}
\rigo{pane d'eternità.}
\strofa
\begin{ritornello}
\rigo{\ch{La}{}\ch{Mi/Sol\#}{}\ch{Fa\#m}{Come} fuoco vivo si accende in noi}
\rigo{\ch{Sim7}{}\ch{Mi}{}\ch{La}{un'immensa} felicità}
\rigo{\ch{Re}{}\ch{Mi}{}\ch{La/Do\#}{}\ch{Re}{che} mai più nessuno ci toglierà}
\rigo{\ch{Sim7}{}\ch{Mi4}{}\ch{Mi}{perché} tu sei ritornato.}
\rigo{\ch{La}{}\ch{Mi/Sol\#}{}\ch{Fa\#m}{Chi} potrà tacere da ora in poi}
\rigo{\ch{Sim7}{}\ch{Mi}{}\ch{La}{che} sei tu in cammino con noi,}
\rigo{\ch{Re}{}\ch{Mi}{}\ch{La}{}\ch{Re}{che} la morte è vinta per sempre, che}
\rigo{\ch{Sim}{}\ch{Mi4}{}\ch{Mi}{ci} hai ridonato la vita.}
\end{ritornello}
\par\endgroup

\begingroup
\cantotesta{185}{Come ti ama Dio}{}{}{587pt}
\rigo{\ch{Do}{}\ch{Lam}{}\ch{Fa}{}\ch{Sol}{Io} vorrei saperti amare come Dio}
\rigo{\ch{Do}{}\ch{Lam}{}\ch{Fa}{che} ti prende per mano, ma ti lascia anche andare}
\rigo{\ch{Sol}{}\ch{Do}{}\ch{Lam}{}\ch{Fa}{vorrei} saperti amare senza farti mai domande}
\rigo{\ch{Sol}{}\ch{Do}{}\ch{Lam}{}\ch{Fa}{}\ch{Sol}{}\ch{Do}{felice} perché esisti e così io posso darti il meglio di me.}
\strofa
\begin{ritornello}
\rigo{\ch{Sol}{}\ch{Lam}{}\ch{Fa}{}\ch{Do}{Con} la forza del mare, l'eternità dei giorni,}
\rigo{\ch{Sol}{}\ch{Lam}{}\ch{Fa}{}\ch{Do}{la} gioia dei voli, la pace della sera,}
\rigo{\ch{Sol}{}\ch{Lam}{}\ch{Fa}{}\ch{Do}{l'immensità} del cielo, come ti ama Dio.}
\strofa
\end{ritornello}
\rigo{\ch{Do}{}\ch{Lam}{}\ch{Fa}{}\ch{Sol}{Io} vorrei saperti amare come ti ama Dio}
\rigo{\ch{Do}{}\ch{Lam}{}\ch{Fa}{che} ti conosce e ti accetta come sei}
\rigo{\ch{Sol}{}\ch{Do}{}\ch{Lam}{}\ch{Fa}{tenerti} fra le mani come i voli nell'azzurro}
\rigo{\ch{Sol}{}\ch{Do}{}\ch{Lam}{}\ch{Fa}{}\ch{Sol}{}\ch{Do}{felice} perché esisti e così io posso darti il meglio di me.}
\strofa
\begin{ritornello}
\rigo{\ch{Sol}{}\ch{Lam}{}\ch{Fa}{}\ch{Do}{Con} la forza del mare, l'eternità dei giorni,}
\rigo{\ch{Sol}{}\ch{Lam}{}\ch{Fa}{}\ch{Do}{la} gioia dei voli, la pace della sera,}
\rigo{\ch{Sol}{}\ch{Lam}{}\ch{Fa}{}\ch{Do}{l'immensità} del cielo, come ti ama Dio.}
\strofa
\end{ritornello}
\rigo{\ch{Do}{}\ch{Lam}{}\ch{Fa}{}\ch{Sol}{Io} vorrei saperti amare come Dio}
\rigo{\ch{Do}{}\ch{Lam}{}\ch{Fa}{che} ti fa migliorare con l'amore che ti dona}
\rigo{\ch{Sol}{}\ch{Do}{}\ch{Lam}{}\ch{Fa}{seguirti} fra la gente con la gioia che hai dentro}
\rigo{\ch{Sol}{}\ch{Do}{}\ch{Lam}{}\ch{Fa}{}\ch{Sol}{}\ch{Do}{felice} perché esisti e così io posso darti il meglio di me.}
\strofa
\begin{ritornello}
\rigo{\ch{Sol}{}\ch{Lam}{}\ch{Fa}{}\ch{Do}{Con} la forza del mare, l'eternità dei giorni,}
\rigo{\ch{Sol}{}\ch{Lam}{}\ch{Fa}{}\ch{Do}{la} gioia dei voli, la pace della sera,}
\rigo{\ch{Sol}{}\ch{Lam}{}\ch{Fa}{}\ch{Do}{l'immensità} del cielo, come ti ama Dio.}
\strofa
\rigo{\ch{Sol}{}\ch{Lam}{}\ch{Fa}{}\ch{Do}{Con} la forza del mare, l'eternità dei giorni,}
\rigo{\ch{Sol}{}\ch{Lam}{}\ch{Fa}{}\ch{Do}{la} gioia dei voli, la pace della sera,}
\rigo{\ch{Sol}{}\ch{Lam}{}\ch{Fa}{}\ch{Do}{l'immensità} del cielo, come ti ama Dio.}
\end{ritornello}
\par\endgroup

\begingroup
\cantotesta{186}{Come tu mi vuoi}{}{}{430pt}
\rigo{\ch{Sol}{}\ch{Lam}{}\ch{Mi}{}\ch{Sol}{}\ch{Re}{Eccomi} signor, vengo a te mio Re, che si compia in me la tua volontà.}
\rigo{\ch{Sol}{}\ch{Lam}{}\ch{Mim}{}\ch{Sim}{Eccomi} Signor, vengo a te mio Dio, plasma il cuore mio e di te vivrò.}
\rigo{\ch{Do}{}\ch{Re}{}\ch{Do}{}\ch{Lam}{}\ch{Sim}{}\ch{Re}{Se} tu lo vuoi, Signore manda me e il tuo nome annuncerò}
\strofa
\begin{ritornello}
\rigo{\ch{Re7}{}\ch{Sol}{}\ch{Re}{}\ch{Mim}{}\ch{Sim}{}\ch{Do}{COME} TU MI VUOI IO SARO', DOVE TU MI VUOI IO ANDRO'. QUESTA VITA}
\rigo{\ch{Sim}{}\ch{Mim}{}\ch{Fa}{}\ch{Re4}{IO} VOGLIO DONARLA A TE PER DAR GLORIA AL TUO NOME, MIO RE.}
\rigo{\ch{Re}{}\ch{Do}{}\ch{Re}{}\ch{Si}{}\ch{Mim}{}\ch{Sim}{}\ch{Do}{COME} TU MI VUOI IO SARO', DOVE TU MI VUOI IO ANDRO'. SE MI GUIDA IL}
\rigo{\ch{Re}{}\ch{Si}{}\ch{Mim}{}\ch{Lam}{}\ch{Do}{}\ch{Re4}{}\ch{Re}{}\ch{Sol}{}\ch{Do}{}\ch{Re7}{}\ch{Sol}{TUO} AMORE PAURA NON HO. PER SEMPRE IO SARO' COME TU MI VUOI}
\strofa
\end{ritornello}
\rigo{\ch{Sol}{}\ch{Lam}{}\ch{Mim}{}\ch{Sol}{}\ch{Re}{Eccomi,} signor, vengo a te mio Re, che si compia in me la tua volontà.}
\rigo{\ch{Sol}{}\ch{La}{}\ch{Mi}{}\ch{Sim}{Eccomi,} Signor, vengo a te mio Dio, plasma il cuore mio e di te vivrò.}
\rigo{\ch{Do}{}\ch{Re}{}\ch{Do}{}\ch{Lam}{}\ch{Sim}{}\ch{Re}{Tra} le tue mani mai più vacillerò e strumento tuo sarò}
\strofa
\rigo{\ch{Re7}{}\ch{Sol}{}\ch{Re}{}\ch{Mim}{}\ch{Sim}{}\ch{Do}{COME} TU MI VUOI IO SARO'}
\par\endgroup

\begingroup
\cantotesta{187}{Dimmelo tu}{}{}{607pt}
\rigo{Il nostro c\ch{Mi}{ane} non mi riconosce p\ch{La}{iù,}\ch{Do\#m}{}     \ch{Si}{}}
\rigo{MI                              LA \#DO\#-}
\rigo{altri profumi, altre valige da portare giù,}
\rigo{c\ch{Si}{erco} le chiavi di casa,}
\rigo{m\ch{Si7}{a} questa non è casa mi\ch{Mi}{a,} più mia\ch{La}{.}}
\strofa
\rigo{I libri, quelli li ho portati via,}
\rigo{anche l'ultimo sguardo, l'ultima fotografia,}
\rigo{quella di nonna bambina}
\rigo{che gioca con il cerchio in piazza, la piazza...}
\strofa
\rigo{Amore, dimmelo tu cos\ch{Re}{'è} \ch{Fa\#m}{}     \ch{Mi}{}}
\rigo{quello che anc\ch{La}{ora} ci manca,}
\rigo{quello che ancora non c\ch{Re}{'è} \ch{Fa\#m}{}}
\rigo{e\ch{Mi}{} che ti prende alle spalle}
\rigo{e\ch{Mi7}{} non ti fa tornare indi\ch{La}{etro,} più indiet\ch{Re}{ro.}}
\rigo{T\ch{Mi}{u,} dimmelo tu cos'è\ch{La}{,}  \ch{Do\#m}{}     \ch{Si}{}}
\rigo{oh t\ch{Mi}{u,} dimmelo tu cos'è\ch{La}{.}}
\strofa
\rigo{Scopare bene, scopare bene questa è la prima cosa,}
\rigo{cercare un'altra donna un'altra casa che non sia troppo vuota}
\rigo{per ritornare di sera}
\rigo{e non sentirsi ancora più soli, ancora più soli.}
\strofa
\rigo{In questa notte d'Agosto ogni stella al suo posto lassù,}
\rigo{in questa notte di strada quella che manca sei tu,}
\rigo{in questa estate romana}
\rigo{di musica e fotografia, ancora più mia.}
\strofa
\commento{Rit.}
\par\endgroup

\begingroup
\cantotesta{188}{E sono solo un uomo}{P. Sequeri}{}{525pt}
\rigo{\ch{Re}{Io} lo so, S\ch{Fa\#m}{ignore} che vengo d\ch{Sol}{a} lontano,      \ch{Re}{}}
\rigo{p\ch{Re}{rima} nel pe\ch{Fa\#m}{nsiero} e poi nell\ch{Sol}{a} tua mano;      \ch{La}{}}
\rigo{i\ch{Re}{o} mi rendo co\ch{La}{nto} che Tu sei la\ch{Sol}{} mia vita          \ch{Re}{}}
\rigo{e n\ch{Sol}{on} mi sembr\ch{Sol/Fa\#}{a} vero d\ch{Mim}{i} pregarTi \ch{Mi7}{così:}      \ch{La7}{}}
\rigo{Padre d’ogni uomo e non Ti ho visto mai,}
\rigo{Spirito di vita e nacqui da una donna}
\rigo{Figlio mio fratello e sono solo un uomo}
\rigo{Eppure io capisco che Tu sei verità.}
\strofa
\rigo{Rit. E imp\ch{Re}{arerò} a guardar\ch{Sol}{e} tutto il mondo\ch{La}{}   \ch{Re7}{}}
\rigo{con gli o\ch{Sol}{cchi} trasparenti d\ch{La}{i} un bambino      \ch{Re7}{}}
\rigo{e insegner\ch{Sol}{ò} a chiamarTi Pa\ch{La}{dre} nostro         \ch{Re}{}       \ch{Si7}{}}
\rigo{Mi-                Mi7 (La)        La (Re)}
\strofa
\rigo{ad ogni figlio che diventa uomo. (2 v.)}
\rigo{Io lo so Signore, che Tu mi sei vicino}
\rigo{Luce alla mia mente, guida il mio cammino,}
\rigo{mano che sorregge, sguardo che perdona,}
\rigo{e non mi sembra vero che Tu esisti così.}
\rigo{Dove nasce amore Tu sei la sorgente,}
\rigo{dove c’è una Croce Tu sei la speranza,}
\rigo{dove il tempo ha fine Tu sei la vita eterna:}
\rigo{e so che posso sempre contare su di Te.}
\strofa
\rigo{Rit. E accoglierò la vita come un dono}
\rigo{E avrò il coraggio di morire anch’io}
\rigo{E incontro a Te verrò con mio fratello}
\rigo{Che non si sente amato da nessuno.}
\par\endgroup

\begingroup
\cantotesta{189}{Fuoco per il mondo}{}{}{529pt}
\rigo{\ch{Do}{}\ch{Sol}{Il} mondo che nasce ogni giorno alla vita}
\rigo{\ch{Fa}{}\ch{Do}{Il} mio sguardo non lo conterrà}
\rigo{\ch{Do}{}\ch{Sol}{Il} mondo che cresce come un seme germoglia}
\rigo{\ch{Fa}{}\ch{Do}{Il} mio cuore non lo porterà}
\strofa
\rigo{\ch{Fa}{}\ch{Do}{Il} mondo che ha fame lotta e cerca giustizia}
\rigo{\ch{Sol}{}\ch{Do}{La} mia vita non lo salverà}
\rigo{\ch{Do}{}\ch{Sol}{Ma} se alzo lo sguardo ed allargo il mio cuore}
\rigo{\ch{Fa}{}\ch{Sol7}{Questo} mondo Dio mi affiderà}
\strofa
\begin{ritornello}
\rigo{\ch{Do}{}\ch{Sol}{}\ch{Fa}{}\ch{Do}{Dio} cerca te come fuoco per il mondo}
\rigo{\ch{Do}{}\ch{Sol}{}\ch{Rem}{}\ch{Sol}{Santo} come Lui per cambiare il male in bene}
\rigo{\ch{Do}{}\ch{Sol}{}\ch{Fa}{}\ch{Do}{Dio} vuole te appassionato fino in fondo}
\rigo{\ch{Rem}{}\ch{Lam}{}\ch{Sol}{}\ch{Do}{per} la vita di ogni uomo che con te Lui salverà}
\strofa
\end{ritornello}
\rigo{\ch{Do}{}\ch{Sol}{Pregare} col cuore è un gran gesto d'amore}
\rigo{\ch{Fa}{}\ch{Do}{Sono} uomo e ho bisogno di Te}
\rigo{\ch{Do}{}\ch{Sol}{Fatica} e dolore è costruire la storia}
\rigo{\ch{Fa}{}\ch{Do}{Il} mio posto nel mondo qual è}
\strofa
\rigo{\ch{Fa}{}\ch{Do}{Dare} la vita con tutti i miei doni}
\rigo{\ch{Sol}{}\ch{Do}{Fedele} nella carità}
\rigo{\ch{Do}{}\ch{Sol}{Gioco} la libertà obbediente alla vita}
\rigo{\ch{Fa}{}\ch{Sol7}{E} il suo Spirito mi accenderà}
\strofa
\begin{ritornello}
\rigo{\ch{Do}{}\ch{Sol}{}\ch{Fa}{}\ch{Do}{Dio} cerca te come fuoco per il mondo}
\rigo{\ch{Do}{}\ch{Sol}{}\ch{Rem}{}\ch{Sol}{Santo} come Lui per cambiare il male in bene}
\rigo{\ch{Do}{}\ch{Sol}{}\ch{Fa}{}\ch{Do}{Dio} vuole te appassionato fino in fondo}
\rigo{\ch{Rem}{}\ch{Lam}{}\ch{Sol}{}\ch{Do}{per} la vita di ogni uomo che con te Lui salverà}
\end{ritornello}
\par\endgroup

\begingroup
\cantotesta{190}{Gesù sei qui con me}{}{}{471pt}
\rigo{\ch{Do}{}\ch{Fa}{GESU'} SEI QUI CON ME}
\rigo{\ch{Do}{}\ch{Fa}{GESU'} SEI QUI CON ME}
\rigo{\ch{Do}{}\ch{Fa}{SEI} VENUTO NEL MIO CUORE}
\rigo{\ch{Do}{}\ch{Fa}{SEI} VENUTO NEL MIO CUORE}
\rigo{\ch{Do}{}\ch{Fa}{SUL} PECCATO HAI TRIONFATO}
\rigo{\ch{Do}{}\ch{Fa}{SUL} PECCATO HAI TRIONFATO}
\rigo{\ch{Do}{}\ch{Fa}{PACE} E AMORE MI HAI DONATO}
\rigo{\ch{Do}{}\ch{Fa}{PACE} E AMORE MI HAI DONATO}
\strofa
\rigo{\ch{Do}{}\ch{Fa}{GESU'} SEI IL RE DEI RE}
\rigo{VIA VERITA' E VITA}
\rigo{TU SEI IL SANTO SEI IL SIGNORE}
\rigo{IL GLORIOSO, IL SALVATORE (BIS)}
\strofa
\rigo{\ch{Do}{}\ch{Sib}{}\ch{Fa}{}\ch{Do}{} (3 volte)}
\rigo{\ch{Sol}{}}
\strofa
\rigo{\ch{La7}{}}
\rigo{\ch{Re}{}\ch{Sol}{GESU'} SEI IL REDENTOR}
\rigo{SULLA CROCE PER AMORE}
\rigo{DALLA MORTE CI HAI SALVATO}
\rigo{VITA NUOVA CI HAI DONATO}
\strofa
\rigo{\ch{Mi}{A} DIO PADRE CREATOR}
\rigo{ALL'AGNELLO AL SALVATORE}
\rigo{ALLO SPIRITO D'AMORE}
\rigo{GLORIA, POTENZA E ONORE !}
\par\endgroup

\begingroup
\cantotesta{191}{I.N.R.I. (io non ritorno indietro)}{}{}{750pt}
\rigo{\ch{Do}{Ha} senso solo co\ch{Fa}{sì,} Perdo\ch{Sol}{nando} proprio \ch{Do}{tutto}}
\rigo{\ch{Do}{Ha} senso solo co\ch{Fa}{sì,} Am\ch{Sol}{ando} fino in \ch{Do}{fondo}}
\rigo{\ch{Do}{Ha} senso solo co\ch{Fa}{sì,} Anche quando \ch{Sol}{hai} tutti \ch{Do}{contro}}
\rigo{\ch{Do}{Ha} senso solo co\ch{Fa}{sì,} E ora \ch{Sol}{io} lo difen\ch{Do}{do.}}
\strofa
\rigo{\ch{Do}{Io} Non Ritorno Indi\ch{Fa}{etro,} \ch{Sol}{Io} Non Ritorno Indi\ch{Do}{etro}}
\rigo{\ch{Do}{Da} un Amore gi\ch{Fa}{gante} così, Io non \ch{Sol}{posso} tornare \ch{Do}{indietro}}
\rigo{Io \ch{Fa}{Non} Ritorno Indi\ch{Do}{etro,} \ch{Sol}{Io} no che Non Ritorno Indi\ch{Lam}{etro}}
\rigo{\ch{Fa}{Da} un Amore gi\ch{Do}{gante} così Io resto \ch{Sol}{appesa} e non \ch{Lam}{scendo}}
\rigo{\ch{Fa}{Io} non sc\ch{Do}{endo}\ch{Sol}{}\ch{Lam}{}\ch{Fa}{}}
\rigo{\ch{Do}{}\ch{Sol}{}\ch{Lam}{}\ch{Fa}{}}
\strofa
\rigo{\ch{Do}{Ha} senso solo co\ch{Fa}{sì,} Risch\ch{Sol}{iando} proprio \ch{Do}{tutto}}
\rigo{\ch{Do}{Ha} senso solo co\ch{Fa}{sì,} Continu\ch{Sol}{ando} fino in \ch{Do}{fondo}}
\rigo{\ch{Do}{Ha} senso solo co\ch{Fa}{sì,} Soprattutto quando \ch{Sol}{hai} il buio ad\ch{Do}{dosso}}
\rigo{\ch{Do}{Ha} senso solo co\ch{Fa}{sì,} E ora \ch{Sol}{io} lo difen\ch{Do}{do.}}
\strofa
\rigo{\ch{Do}{Io} Non Ritorno Indi\ch{Fa}{etro,} \ch{Sol}{Io} Non Ritorno Indi\ch{Do}{etro}}
\rigo{\ch{Do}{Da} un Amore gi\ch{Fa}{gante} così, Io non \ch{Sol}{posso} tornare \ch{Do}{indietro}}
\rigo{Io \ch{Fa}{Non} Ritorno Indi\ch{Do}{etro,} \ch{Sol}{Io} no che Non Ritorno Indi\ch{Lam}{etro}}
\rigo{\ch{Fa}{Da} un Amore gi\ch{Do}{gante} così Io resto \ch{Sol}{appesa} e non \ch{Lam}{scendo}}
\rigo{\ch{Fa}{Io} non sc\ch{Do}{endo}\ch{Sol}{}\ch{Lam}{}\ch{Fa}{}}
\rigo{\ch{Do}{}\ch{Sol}{}\ch{Lam}{}\ch{Fa}{}}
\strofa
\rigo{\ch{Do}{È} questa la natura \ch{Sol}{del} vero Amore \ch{Lam}{l'operare} da \ch{Fa}{Dio}}
\rigo{\ch{Do}{la} fermezza e il non \ch{Sol}{ritirarsi} \ch{Lam}{mai} \ch{Fa}{}}
\rigo{\ch{Do}{È} questa la natura \ch{Sol}{del} vero Amore \ch{Lam}{l'operare} \ch{Fa}{da} Dio}
\rigo{\ch{Do}{la} fermezza e il non \ch{Sol}{ritirarsi} \ch{Lam}{mai}}
\rigo{\ch{Fa}{A} costo di qualunque Sacri\ch{Do}{ficio}}
\strofa
\rigo{\ch{Sol}{Io} Non Ritorno Ind\ch{Lam}{ietro} \ch{Fa}{}\ch{Do}{}}
\rigo{\ch{Sol}{Io} Non Ritorno Ind\ch{Lam}{ietro} \ch{Fa}{}\ch{Do}{}}
\par\endgroup

\begingroup
\cantotesta{192}{Il canto dell'amore}{}{}{590pt}
\rigo{\ch{Do}{}\ch{Lam}{}\ch{Fa}{}\ch{Do}{}}
\rigo{\ch{Do}{}\ch{Lam}{}\ch{Fa}{}\ch{Do}{}}
\strofa
\rigo{\ch{Do}{}\ch{Lam}{}\ch{Fa}{}\ch{Do}{Se} dovrai attraversare il deserto non temere io sarò con te}
\rigo{\ch{Sol4}{}\ch{Lam}{}\ch{Fa}{}\ch{Do}{se} dovrai camminare nel fuoco la sua fiamma non ti brucerà}
\rigo{\ch{Mim7}{}\ch{Fa}{}\ch{Do}{}\ch{Rem}{}\ch{Sol}{}\ch{La}{seguirai} la mia luce nella notte sentirai la mia forza nel cammino}
\rigo{\ch{Sib}{}\ch{Fa}{}\ch{Do}{}\ch{Lam}{}\ch{Fa}{}\ch{Do}{io} sono il tuo Dio, il Signore.}
\strofa
\rigo{\ch{Do}{}\ch{Lam}{}\ch{Fa}{}\ch{Do}{Sono} io che ti ho fatto e plasmato, ti ho chiamato per nome}
\rigo{\ch{Sol4}{}\ch{Lam}{}\ch{Fa}{}\ch{Do}{io} da sempre ti ho conosciuto e ti ho dato il mio amore}
\rigo{\ch{Mim7}{}\ch{Fa}{}\ch{Do}{}\ch{Rem}{}\ch{Sol}{}\ch{La}{perché} tu sei prezioso ai miei occhi, vali più del più grande dei tesori}
\rigo{\ch{Sib}{}\ch{Fa}{}\ch{Do}{}\ch{La4}{}\ch{La}{io} sarò con te dovunque andrai.}
\strofa
\rigo{\ch{Re}{}\ch{Sim}{}\ch{Sol}{}\ch{Re}{Non} pensare alle cose di ieri, cose nuove fioriscono già}
\rigo{\ch{La4}{}\ch{Sim}{}\ch{Sol}{}\ch{Re}{aprirò} nel deserto sentieri, darò acqua nell'aridità}
\rigo{\ch{Fa\#m}{}\ch{Sol}{}\ch{Re}{}\ch{Mim}{}\ch{La}{}\ch{Si}{perché} tu sei prezioso ai miei occhi, vali più del più grande dei tesori}
\rigo{\ch{Do}{}\ch{Sol}{}\ch{Re}{}\ch{Sim}{}\ch{Sol}{}\ch{Re}{io} sarò con te dovunque andrai}
\rigo{\ch{Fa\#m}{}\ch{Sol}{}\ch{Re}{}\ch{Mim}{}\ch{La}{}\ch{Si}{perché} tu sei prezioso ai miei occhi, vali più del più grande dei tesori}
\rigo{\ch{Do}{}\ch{Sol}{} (stop)}
\rigo{io sarò con te}
\rigo{dovunque andrai.}
\strofa
\rigo{\ch{Re}{}\ch{Sim}{}\ch{Sol}{}\ch{Mi}{Io} ti sarò accanto sarò con te}
\rigo{\ch{Re}{}\ch{Sim}{}\ch{Sol}{}\ch{Re}{per} tutto il tuo viaggio sarò con te (2 volte)}
\par\endgroup

\begingroup\senzaaccordi
\cantotesta{193}{Insieme}{}{\bEG}{294pt}
\rigo{Insieme abbiam marciato un dì}
\rigo{per strade non battute}
\rigo{insieme abbiam raccolto un fior}
\rigo{sull'orlo di una rupe.}
\strofa
\rigo{RIT: Insieme, insieme}
\rigo{è il motto di fraternità}
\rigo{insieme, nel bene crediam.}
\strofa
\rigo{Insieme abbiam portato un dì}
\rigo{lo zaino che ci spezza}
\rigo{insieme abbiam goduto alfin}
\rigo{del vento la carezza.}
\strofa
\rigo{RIT: Insieme, insieme...}
\strofa
\rigo{Insieme abbiamo appreso ciò}
\rigo{che il libro non addita}
\rigo{abbiamo appreso che l'amor}
\rigo{è il senso della vita.}
\strofa
\rigo{RIT: Insieme, insieme...}
\par\endgroup

\begingroup
\cantotesta{194}{Io sarò con te}{}{}{495pt}
\begin{ritornello}
\rigo{\ch{Re}{}\ch{Sim}{}\ch{Mim}{}\ch{La}{Io} sarò con te sulle strade della vita,}
\rigo{\ch{Re}{}\ch{Sim}{}\ch{Mim}{}\ch{La}{Io} sarò con te anche quando non lo sai.}
\rigo{\ch{Re}{}\ch{Fa\#m}{}\ch{Sol}{}\ch{La}{Io} sarò con te custodendoti per sempre}
\rigo{\ch{Sol}{}\ch{La}{}\ch{Re}{nella} fedeltà è il mio Amore.}
\strofa
\rigo{\ch{Re}{}\ch{Sim}{}\ch{Mim}{}\ch{La}{Tu} sarai con me sulle strade della vita,}
\rigo{\ch{Re}{}\ch{Sim}{}\ch{Mim}{}\ch{La}{tu} sarai con me anche quando non lo so.}
\rigo{\ch{Re}{}\ch{Fa\#m}{}\ch{Sol}{}\ch{La}{Tu} sarai con me custodendomi per sempre}
\rigo{\ch{Sol}{}\ch{La}{}\ch{Re}{nella} fedeltà è il Tuo amore.}
\strofa
\end{ritornello}
\rigo{\ch{Re}{}\ch{Fa\#m}{}\ch{Sim}{Se} forte e coraggioso tu sarai}
\rigo{\ch{Sol}{}\ch{La}{}\ch{Re}{scegliendo} la via della vita,}
\rigo{\ch{Sol}{}\ch{La}{}\ch{Re}{}\ch{Sim}{ascoltando} la mia Parola}
\rigo{\ch{Sol}{}\ch{La}{}\ch{Re}{}\ch{La}{custodendola} nel tuo cuore...}
\strofa
\begin{ritornello}
\rigo{\ch{Re}{}\ch{Sim}{}\ch{Mim}{}\ch{La}{Io} sarò con te sulle strade della vita,}
\rigo{\ch{Re}{}\ch{Sim}{}\ch{Mim}{}\ch{La}{Io} sarò con te anche quando non lo sai.}
\rigo{\ch{Re}{}\ch{Fa\#m}{}\ch{Sol}{}\ch{La}{Io} sarò con te custodendoti per sempre}
\rigo{\ch{Sol}{}\ch{La}{}\ch{Re}{nella} fedeltà è il mio Amore.}
\strofa
\end{ritornello}
\rigo{Se forte e coraggioso tu sarai}
\rigo{annunciando ad ogni uomo la salvezza}
\rigo{cantando la speranza che non muore}
\rigo{camminerò con te dovunque andrai...}
\strofa
\begin{ritornello}
\rigo{\ch{Re}{}\ch{Sim}{}\ch{Mim}{}\ch{La}{Io} sarò con te sulle strade della vita,}
\rigo{\ch{Re}{}\ch{Sim}{}\ch{Mim}{}\ch{La}{Io} sarò con te anche quando non lo sai.}
\rigo{\ch{Re}{}\ch{Fa\#m}{}\ch{Sol}{}\ch{La}{Io} sarò con te custodendoti per sempre}
\rigo{\ch{Sol}{}\ch{La}{}\ch{Re}{nella} fedeltà è il mio Amore.}
\end{ritornello}
\par\endgroup

\begingroup
\cantotesta{195}{Nel tuo volto la mia strada}{}{\bRS}{631pt}
\rigo{\ch{Do}{Ho} cercato \ch{Mim}{il} mio volto, \ch{Lam}{la} mia \ch{Sol}{identità}}
\rigo{\ch{Fa}{Con} la bussola \ch{Lam}{in} mano \ch{Fa}{diretta} a \ch{Sol}{nord}}
\rigo{\ch{Do}{Ho} \ch{Mim}{trovato} timori \ch{Fa}{}\ch{Sol}{e} fragilità}
\rigo{\ch{Fa}{Ed} un cuore \ch{Do}{che} cerca \ch{Re}{la} \ch{Sol}{felicità}}
\rigo{\ch{Lam}{Ma} qualcuno \ch{Mim}{m'ha} preso per \ch{Fa}{mano} lo \ch{La7}{so}}
\rigo{\ch{Fa}{Dietro} \ch{Fam}{a} Lui \ch{Sol7}{sicuro} sarò.}
\strofa
\begin{ritornello}
\rigo{\ch{Do}{Il} tuo volto io cercherò \ch{Rem}{sulla} strada che hai \ch{Mim}{disegnato} in me\ch{Fa}{}}
\rigo{\ch{Do/La}{}\ch{Fa}{}\ch{Fam}{}\ch{Do}{}\ch{Rem}{}\ch{Mim}{}\ch{Sol7}{E} sarò ciò che vorrai, Signor: vita donata!}
\end{ritornello}
\rigo{(la seconda volta: Passerà dai piedi questa nostra vita sulla strada che hai disegnato in me e sarò ciò che vorrai Signor: vita donata.)}
\strofa
\rigo{Noi non siamo mai soli ma insieme affrontiam}
\rigo{La fatica e le sfide che il mondo ci dà}
\rigo{Una sola la legge, promessa sarem}
\rigo{Di un mondo migliore che vive l'amor!}
\rigo{Nel guidar la canoa più forza ci dà}
\rigo{La certezze che Tu sei con noi.}
\strofa
\begin{ritornello}
\rigo{\ch{Do}{Il} tuo volto io cercherò \ch{Rem}{sulla} strada che hai \ch{Mim}{disegnato} in me\ch{Fa}{}}
\rigo{\ch{Do/La}{}\ch{Fa}{}\ch{Fam}{}\ch{Do}{}\ch{Rem}{}\ch{Mim}{}\ch{Sol7}{E} sarò ciò che vorrai, Signor: vita donata! (Passerà dai piedi…)}
\strofa
\end{ritornello}
\rigo{Orizzonti lontani si aprono allor}
\rigo{Scopri il dono dell'altro, riveli il tuo cuor}
\rigo{Hai trovato la strada, la felicità}
\rigo{E tendendo le braccia servire saprai.}
\rigo{Una tenda che accoglie tu diventerai}
\rigo{E pienezza di vita sarà.}
\strofa
\begin{ritornello}
\rigo{\ch{Do}{Il} tuo volto io cercherò \ch{Rem}{sulla} strada che hai \ch{Mim}{disegnato} in me\ch{Fa}{}}
\rigo{\ch{Do/La}{}\ch{Fa}{}\ch{Fam}{}\ch{Do}{}\ch{Rem}{}\ch{Mim}{}\ch{Sol7}{E} sarò ciò che vorrai, Signor: vita donata! (Passerà dai piedi…)}
\end{ritornello}
\par\endgroup

\begingroup\senzaaccordi
\cantotesta{196}{Pane del cielo}{}{}{343pt}
\begin{ritornello}
\rigo{Pane del cielo}
\rigo{sei tu Gesù,}
\rigo{via d’amore:}
\rigo{tu ci fai come te.}
\strofa
\end{ritornello}
\rigo{No, non è rimasta fredda la terra:}
\rigo{tu sei rimasto con noi}
\rigo{per nutrirci di te,}
\rigo{pane di vita;}
\rigo{ed infiammare col tuo amore}
\rigo{tutta l’umanità.}
\strofa
\rigo{Sì, il cielo è qui su questa terra:}
\rigo{tu sei rimasto con noi}
\rigo{ma ci porti con te}
\rigo{nella tua casa}
\rigo{dove vivremo insieme a te}
\rigo{tutta l’eternità.}
\strofa
\rigo{No, la morte non può farci paura:}
\rigo{tu sei rimasto con noi.}
\rigo{E chi vive di te}
\rigo{vive per sempre.}
\rigo{Sei Dio con noi, sei Dio per noi,}
\rigo{Dio in mezzo a noi.}
\par\endgroup

\begingroup\senzaaccordi
\cantotesta{197}{Preghiera semplice}{}{}{261pt}
\rigo{Se anche parlassi le lingue degli uomini e degli angeli,}
\rigo{ma non avessi la carità, sono come un bronzo che risuona,}
\rigo{o un cembalo che tintinna.}
\strofa
\rigo{E se avessi il dono della profezia,}
\rigo{e conoscessi tutti i misteri e tutta la scienza,}
\rigo{e possedessi la pienezza della Fede così da trasportare}
\rigo{le montagne, ma non avessi la carità non avrei nulla.}
\strofa
\rigo{E se anche distribuissi tutte le mie sostanze}
\rigo{e dessi il mio corpo per essere bruciato, ma non avessi la carità}
\rigo{niente mi giova.}
\par\endgroup

\begingroup
\cantotesta{198}{Sei fuoco e vento}{}{}{645pt}
\rigo{In un \ch{Lam}{mare} calmo e immobile}
\rigo{con un \ch{Do}{cielo} senza nuvole}
\rigo{non si \ch{Sol}{riesce} a na\ch{Rem}{vigare}}
\rigo{\ch{Fa}{proseguire} non si può.}
\strofa
\rigo{Una \ch{Lam}{brezza} lieve e debole}
\rigo{poi di\ch{Do}{venta} un vento a raffiche}
\rigo{soffia \ch{Sol}{forte} sulle \ch{Rem}{barche}}
\rigo{e ci \ch{Fa}{spinge} via da qua.}
\strofa
\rigo{Come il \ch{Do}{vento} da la \ch{Sol}{forza}}
\rigo{per \ch{Lam}{viaggiare} in un \ch{Mim}{oceano}}
\rigo{così \ch{Fa}{Tu} ci dai lo \ch{Do}{Spirito}}
\rigo{che ci \ch{Rem}{guiderà} da \ch{Sol}{Te...}}
\strofa
\begin{ritornello}
\rigo{Sei come \ch{Do}{vento} che \ch{Fa}{gonfia} le \ch{Do}{vele}}
\rigo{sei come \ch{Do}{fuoco} che \ch{Fa}{accende} l'\ch{Sol}{Amore}\ch{Mi}{}}
\rigo{sei come \ch{Lam}{l'aria} che \ch{Mim}{si} respira \ch{Fa}{libera}}
\rigo{chiara \ch{Do}{luce} che \ch{Sol}{il} cammino \ch{Fa}{indica.} (2 volte)}
\strofa
\end{ritornello}
\rigo{Nella \ch{Lam}{notte} impenetrabile}
\rigo{ogni \ch{Do}{cosa} è irraggiungibile}
\rigo{non puoi \ch{Sol}{scegliere} la \ch{Rem}{strada}}
\rigo{se non \ch{Fa}{vedi} davanti a te.}
\strofa
\rigo{Una \ch{Lam}{luce} fioca e debole}
\rigo{sembra \ch{Do}{sorgere} e poi crescere}
\rigo{come \ch{Sol}{fiamma} che \ch{Rem}{rigenera}}
\rigo{e che \ch{Fa}{illumina} la via.}
\strofa
\rigo{Come il \ch{Do}{fuoco} scioglie il \ch{Sol}{gelo}}
\rigo{e \ch{Lam}{rischiara} ogni \ch{Mim}{sentiero}}
\rigo{così \ch{Fa}{Tu} riscaldi \ch{Do}{il} cuore}
\rigo{di chi \ch{Rem}{Verbo} \ch{Sol}{annuncerà...}}
\strofa
\begin{ritornello}
\rigo{Sei come \ch{Do}{vento} che \ch{Fa}{gonfia} le \ch{Do}{vele}}
\rigo{sei come \ch{Do}{fuoco} che \ch{Fa}{accende} l'\ch{Sol}{Amore}\ch{Mi}{}}
\rigo{sei come \ch{Lam}{l'aria} che \ch{Mim}{si} respira \ch{Fa}{libera}}
\rigo{chiara \ch{Do}{luce} che \ch{Sol}{il} cammino \ch{Fa}{indica.} (3 volte)}
\end{ritornello}
\par\endgroup

\begingroup
\cantotesta{199}{Su ali d'aquila}{}{}{594pt}
\rigo{\ch{Sol7+}{Tu} che abiti al riparo\ch{Re7+}{} del Signore}
\rigo{e\ch{Sol7+}{} che dimori alla Sua ombra\ch{Re7+}{,}}
\rigo{d\ch{Fa}{i'} al Signore "Mio rifugio,\ch{Sib}{}}
\rigo{\ch{Solm7}{mia} roccia in cui confido"\ch{La4}{.}           \ch{La7}{}}
\strofa
\commento{Rit.}
\rigo{E ti \ch{Re}{rialzerà,} ti so\ch{Re7+}{lleverà}}
\rigo{su ali d'\ch{Mim}{aquila,} ti r\ch{La7}{eggerà}}
\rigo{sulla br\ch{Lam7}{ezza} dell'alb\ch{Re7}{a} ti fa\ch{Sol}{rà} bril\ch{Mim}{lar}}
\rigo{come il sole,\ch{Sim}{} cos\ch{Fa\#m}{ì} nelle S\ch{Sol}{ue} m\ch{Mim}{ani} vivrai.\ch{La7}{}            \ch{Re}{}}
\strofa
\rigo{Dal laccio del cacciatore ti libererà}
\rigo{e dalla carestia che distrugge,}
\rigo{poi ti coprirà con le Sue ali}
\rigo{e rifugio troverai.}
\strofa
\commento{Rit.}
\strofa
\rigo{Non devi temere i terrori della notte}
\rigo{né freccia che vola di giorno,}
\rigo{mille cadranno al tuo fianco}
\rigo{ma nulla ti colpirà.}
\strofa
\commento{Rit.}
\strofa
\rigo{Perché‚ ai Suoi angeli ha dato un comando}
\rigo{di preservarti in tutte le tue vie,}
\rigo{ti porteranno sulle loro mani}
\rigo{contro la pietra non inciamperai.}
\strofa
\commento{Rit.}
\strofa
\rigo{E ti rialzerò, ti solleverò}
\rigo{su ali d'aquila, ti reggerò}
\rigo{sulla brezza dell'alba ti farò brillar}
\rigo{come il sole, così nelle mie mani vivrai.}
\par\endgroup

\begingroup
\cantotesta{200}{Symbolum 77}{P. Sequeri}{}{508pt}
\rigo{\ch{Lam}{Tu} sei la mia vita, \ch{Fa}{altro} i\ch{Sol}{o} non h\ch{Do}{o;} \ch{Lam}{}}
\rigo{Tu sei la mia str\ch{Fa}{ada,} la mia\ch{Sol}{} verità,    \ch{Mi}{}}
\rigo{n\ch{Rem}{ella} Tua Parola\ch{Sol}{} io \ch{Do}{camminerò}      \ch{Lam}{}}
\rigo{f\ch{Fa}{inché‚} avrò respiro\ch{Sol}{,} fino a quando \ch{Do}{Tu} vorrai;     \ch{Mi}{}}
\rigo{n\ch{Rem}{on} avrò paura, \ch{Sol}{sai,} se Tu s\ch{Do}{ei} con me,         \ch{Lam}{}}
\rigo{i\ch{Fa}{o} Ti prego: re\ch{Sol}{sta} con me. \ch{Lam}{}}
\strofa
\rigo{Credo in Te, Signore, nato da Maria,}
\rigo{Figlio eterno e Santo, uomo come noi;}
\rigo{morto per amore, vivo in mezzo a noi,}
\rigo{una cosa sola con il Padre e con i tuoi,}
\rigo{fino a quando, io lo so, Tu ritornerai}
\rigo{per aprirci il Regno di Dio.}
\strofa
\rigo{Tu sei la mia vita, altro io non ho;}
\rigo{Tu sei la mia pace, la mia libertà,}
\rigo{niente nella vita ci separerà,}
\rigo{so che la Tua Mano forte non mi lascerà,}
\rigo{so che da ogni male Tu mi libererai}
\rigo{e nel Tuo perdono vivrò.}
\strofa
\rigo{Padre della vita, noi crediamo in Te;}
\rigo{Figlio Salvatore, noi speriamo in Te,}
\rigo{Spirito d'Amore vieni in mezzo a noi,}
\rigo{Tu da mille strade ci raduni in unità,}
\rigo{e per mille strade poi dove Tu vorrai}
\rigo{noi saremo il seme di Dio.}
\par\endgroup

\sottosezione{Finale}

\begingroup
\cantotesta{201}{Abbracciami}{}{}{519pt}
\rigo{\ch{Mi}{}\ch{La}{}\ch{Do\#m}{}\ch{La}{}}
\strofa
\rigo{\ch{Mi}{}\ch{La}{Gesù} parola viva e vera}
\rigo{\ch{Do\#m}{}\ch{La}{Sorgente} che disseta e cura ogni ferita}
\rigo{\ch{Mi}{}\ch{La}{Ferma} su di me i tuoi occhi}
\rigo{\ch{Do\#m}{}\ch{La}{La} tua mano stendi e donami la vita}
\strofa
\begin{ritornello}
\rigo{\ch{Mi}{}\ch{La}{Abbracciami} Dio dell'eternità}
\rigo{\ch{Do\#m}{}\ch{La}{Rifugio} dell'anima grazie che opera}
\rigo{\ch{Mi}{}\ch{La}{Riscaldami} fuoco che libera}
\rigo{\ch{Do\#m}{}\ch{La}{Manda} il tuo Spirito Maranathà Gesù}
\strofa
\end{ritornello}
\rigo{\ch{Mi}{}\ch{La}{Gesù} asciuga il nostro pianto}
\rigo{\ch{Do\#m}{}\ch{La}{Leone} vincitore della tribù di Giuda}
\rigo{\ch{Mi}{}\ch{La}{}\ch{Do\#m}{Vieni} nella tua potenza questo cuore sciogli}
\rigo{\ch{La}{Con} ogni sua paura}
\strofa
\begin{ritornello}
\rigo{\ch{Mi}{}\ch{La}{Abbracciami} Dio dell'eternità}
\rigo{\ch{Do\#m}{}\ch{La}{Rifugio} dell'anima grazie che opera}
\rigo{\ch{Mi}{}\ch{La}{Riscaldami} fuoco che libera}
\rigo{\ch{Do\#m}{}\ch{La}{Manda} il tuo Spirito Maranathà Gesù}
\strofa
\end{ritornello}
\rigo{\ch{Mi}{}\ch{Si}{}\ch{La}{}\ch{La}{Per} sempre io canterò la tua immensa fedeltà}
\rigo{\ch{Mi}{}\ch{Si}{}\ch{La}{}\ch{La}{Il} tuo spirito in me in eterno ti loderà}
\rigo{\ch{Mi}{}\ch{Si}{}\ch{La}{}\ch{La}{Per} sempre io canterò la tua immensa fedeltà}
\rigo{\ch{Mi}{}\ch{Si}{}\ch{La}{}\ch{La}{Il} tuo spirito in me in eterno ti loderà}
\strofa
\begin{ritornello}
\rigo{\ch{Mi}{}\ch{La}{Abbracciami} Dio dell'eternità}
\rigo{\ch{Do\#m}{}\ch{La}{Rifugio} dell'anima grazie che opera}
\rigo{\ch{Mi}{}\ch{La}{Riscaldami} fuoco che libera}
\rigo{\ch{Do\#m}{}\ch{La}{Manda} il tuo Spirito Maranathà Gesù}
\end{ritornello}
\par\endgroup

\begingroup
\cantotesta{202}{Andate per le strade}{}{}{488pt}
\rigo{\ch{Sim}{Andate} per \ch{Re}{le} str\ch{La7}{ade} in t\ch{Re}{utto} il mondo}
\rigo{c\ch{Sim}{hiam}\ch{Fa\#}{ate} \ch{Sim}{i} mi\ch{La}{ei} \ch{Re}{amici} per far festa}
\rigo{c'è u\ch{Sim}{n} posto\ch{Sol}{} per\ch{La}{} ciascuno a\ch{Sim}{lla} mia mensa.}
\strofa
\rigo{Nel v\ch{Sol}{ostro} c\ch{La}{amm}\ch{Re}{ino} annunciate il Vangelo,}
\rigo{d\ch{Sim}{icen}\ch{Mim}{do:} \ch{Fa\#}{"E'} \ch{Sim}{vicino} il Regno dei cieli".}
\rigo{G\ch{Re7}{uarite} \ch{Sol}{i} malat\ch{La}{i,} mon\ch{Re}{date} i lebbrosi,}
\rigo{r\ch{Sim}{ende}\ch{Fa\#m}{te} la\ch{Do\#7}{} vita\ch{Fa\#}{} a chi l'ha perduta.}
\strofa
\rigo{Vi è stato donato con amore gratuito}
\rigo{ugualmente donate con gioia per amore.}
\rigo{Con voi non prendete nè oro nè argento,}
\rigo{perché l'operaio ha diritto al suo cibo.}
\strofa
\rigo{Entrando in una casa donatele la pace.}
\rigo{Se c'è chi vi rifiuta e non accoglie il dono,}
\rigo{la pace torni a voi ed uscite dalla casa}
\rigo{scuotendo la polvere dai vostri calzari.}
\strofa
\rigo{Ecco io vi mando come agnelli in mezzo ai lupi,}
\rigo{siate dunque avveduti come sono i serpenti,}
\rigo{ma liberi e chiari come le colombe:}
\rigo{dovrete sopportare prigioni e tribunali.}
\strofa
\rigo{Nessuno è più grande del proprio maestro}
\rigo{né il servo è più importante del suo padrone.}
\rigo{Se hanno odiato me, odieranno anche voi}
\rigo{ma voi non temete io non vi lascio soli.}
\par\endgroup

\begingroup\senzaaccordi
\cantotesta{203}{Con te faremo cose grandi}{}{}{469pt}
\rigo{Con te faremo cose grandi}
\rigo{il cammino che percorreremo insieme}
\rigo{di te riempiranno sguardi}
\rigo{la speranza che risplenderà nei volti}
\rigo{tu la voce che rischiara}
\rigo{tu la voce che ci chiama}
\rigo{tu la gioia che da vita ai nostri sogni}
\rigo{Parlaci Signore come sai}
\rigo{sei presente nel mistero in mezzo a noi}
\rigo{chiamaci col nome che vorrai}
\rigo{e sia fatto il tuo segno su di noi}
\rigo{tu la voce che rischiara}
\rigo{tu la voce che ci chiama}
\rigo{tu la gioia che da vita ai nostri sogni.}
\strofa
\rigo{Con te faremo cose grandi}
\rigo{il cammino che percorreremo insieme}
\rigo{di te riempiranno sguardi}
\rigo{la speranza che risplenderà nei volti}
\rigo{tu l'amore che dà vita}
\rigo{tu il sorriso che ci allieta}
\rigo{tu la forza che raduna i nostri giorni.}
\strofa
\rigo{Guidaci Signore come sai}
\rigo{da chi soffre chi è più piccolo di noi}
\rigo{strumenti di quel regno che tu fai}
\rigo{di quel regno che ora vive in mezzo a noi}
\rigo{tu l'amore che dà vita}
\rigo{tu il sorriso che ci allieta}
\rigo{tu la forza che raduna i nostri giorni.}
\strofa
\rigo{Con te faremo cose grandi}
\rigo{il cammino che percorreremo insieme}
\rigo{di te si riempiranno sguardi}
\rigo{la speranza che risplenderà nei volti.}
\par\endgroup

\begingroup
\cantotesta{204}{Eccomi per sempre}{}{\bLC}{893pt}
\rigo{\ch{Mi}{}\ch{Si4}{}\ch{La4/2}{}\ch{Si4}{}\ch{Mi}{}}
\strofa
\rigo{L'avventura \ch{La2}{della} vita parte con \ch{Mi}{semplicità,}}
\rigo{nel sentiero \ch{La2}{che} dal prato alla montagna \ch{Mi}{arriverà.}}
\rigo{Con la gioia \ch{Do}{di} giocare, l'entusiasmo di \ch{Sol}{scoprire:}}
\rigo{la natura \ch{Re}{l'amicizia} e l'armonia.\ch{Resus4}{}}
\rigo{Poi l'impegno \ch{La2}{di} uno stile con fiducia e \ch{Mi}{lealtà.}}
\rigo{Sempre pronta \ch{La2}{per} accogliere ogni sfida che \ch{Mi}{verrà.}}
\rigo{Con il senso \ch{Do}{del} dovere, della fede, \ch{Sol}{dell'onore,}}
\rigo{del rispetto, \ch{Re}{della} generosità.}
\strofa
\begin{ritornello}
\rigo{\ch{Do}{Eccomi,} per \ch{Re}{sempre,} dirò \ch{Sol}{“eccomi”.}}
\rigo{\ch{Do}{Per} servire in ogni \ch{Re}{attimo,} per dire \ch{Sol}{sempre} sì}
\rigo{\ch{Do}{Eccomi,} per \ch{Re}{sempre,} dirò \ch{Sol}{“eccomi”.}}
\rigo{\ch{Do}{Oggi} e \ch{Sol}{domani,} oltre i \ch{Re}{confini,} \ch{La}{lungo}\ch{Mi}{}}
\rigo{una strada \ch{Si4}{che} mai \ch{La2/4}{finirà.}\ch{Mi}{}}
\strofa
\end{ritornello}
\rigo{Sentinella della vita, dono per l'umanità.}
\rigo{Prima o poi la tua partenza verso il mondo arriverà.}
\rigo{Donna scelta dal Signore, per poter testimoniare,}
\rigo{che la scelta dell'amore è libertà.}
\rigo{Lungo un secolo di storia, nel futuro che verrà,}
\rigo{la promessa di ogni guida sempre si rinnoverà.}
\rigo{Un messaggio per sperare, un impegno da portare,}
\rigo{nel futuro della nostra civiltà.}
\strofa
\begin{ritornello}
\rigo{\ch{Do}{Eccomi,} per \ch{Re}{sempre,} dirò \ch{Sol}{“eccomi”.}}
\rigo{\ch{Do}{Per} servire in ogni \ch{Re}{attimo,} per dire \ch{Sol}{sempre} sì}
\rigo{\ch{Do}{Eccomi,} per \ch{Re}{sempre,} dirò \ch{Sol}{“eccomi”.}}
\rigo{\ch{Do}{Oggi} e \ch{Sol}{domani,} oltre i \ch{Re}{confini,} \ch{La}{lungo}\ch{Mi}{}}
\rigo{una strada \ch{Si4}{che} mai \ch{La2/4}{finirà.}\ch{Mi}{}}
\strofa
\end{ritornello}
\rigo{Anche quando la \ch{Lam7}{fatica} ti fa chiedere “perché”\ch{Mim}{}}
\rigo{Anche quando la \ch{Lam7}{certezza} lascia il posto a tanti \ch{Mim}{“se”}}
\rigo{Segui sempre \ch{Do}{il} tuo sentiero, tante mani \ch{Sol}{stringerai}}
\rigo{Con la luce \ch{Re}{della} Fede la tua strada \ch{Resus4}{troverai}}
\strofa
\begin{ritornello}
\rigo{\ch{Do}{Eccomi,} per \ch{Re}{sempre,} dirò \ch{Sol}{“eccomi”.}}
\rigo{\ch{Do}{Per} servire in ogni \ch{Re}{attimo,} per dire \ch{Sol}{sempre} sì}
\rigo{\ch{Do}{Eccomi,} per \ch{Re}{sempre,} dirò \ch{Sol}{“eccomi”.}}
\rigo{\ch{Do}{Oggi} e \ch{Sol}{domani,} oltre i \ch{Re}{confini,} \ch{La}{lungo}\ch{Mi}{}}
\rigo{una strada \ch{Si4}{che} mai \ch{La2/4}{finirà.}\ch{Mi}{}}
\end{ritornello}
\par\endgroup

\begingroup
\cantotesta{205}{Fratello sole e sorella luna}{}{}{382pt}
\rigo{\ch{Re}{Dolce} sent\ch{Fa\#m}{ire} come\ch{Sol}{} nel\ch{La7}{} mio cuore \ch{Fa\#m}{}}
\rigo{o\ch{Sim}{ra,} \ch{La}{umi}\ch{Sol}{lmen}\ch{Fa\#m}{te,} s\ch{Sim}{ta} n\ch{Mim}{ascendo} amo\ch{Mi7}{re.}         \ch{La7}{}}
\rigo{D\ch{Re}{olce} capir\ch{Fa\#m}{e} che no\ch{Sol}{n} son pi\ch{La7}{ù} solo     \ch{Fa\#m}{}}
\rigo{m\ch{Sim}{a} che s\ch{La}{on} par\ch{Sol}{te} d\ch{Fa\#m}{i} un'\ch{Sim}{imme}\ch{Mim}{nsa} vita,   \ch{La}{}      \ch{Re}{}}
\rigo{c\ch{Sim}{he} gene\ch{Solm}{rosa} \ch{Re}{ris}\ch{Sol}{plende} \ch{Solm}{intorno} \ch{La}{a} me:      \ch{Re}{}}
\rigo{d\ch{Sim}{ono} \ch{Solm}{di} Lu\ch{Re}{i,} \ch{Sol}{del} \ch{Solm}{Suo} immenso \ch{La}{amor.}       \ch{Re}{}}
\strofa
\rigo{Ci ha dato il cielo e le chiare stelle,}
\rigo{fratello Sole e sorella Luna;}
\rigo{la madre Terra con frutti, prati e fiori}
\rigo{il fuoco, il vento, l'aria e l'acqua pura,}
\rigo{fonte di vita per le Sue creature:}
\rigo{dono di Lui, del Suo immenso amor,}
\rigo{dono di Lui, del Suo immenso amor.}
\strofa
\rigo{Sia laudato nostro Signore}
\rigo{che ha creato l'universo intero.}
\rigo{Sia laudato nostro Signore}
\rigo{noi tutti siamo Sue creature:}
\rigo{dono di Lui, del Suo immenso amor,}
\rigo{beato chi Lo serve in umiltà.}
\par\endgroup

\begingroup
\cantotesta{206}{Laudato sii, Signore mio}{}{}{823pt}
\commento{Rit.}
\rigo{L\ch{Sol}{audato} sii, Signore mio \ch{Lam}{laudato} \ch{Re}{sii,} Signore mio.      \ch{Mim}{}}
\rigo{L\ch{Do}{audato} sii, Signore mio lau\ch{Re}{dato} si\ch{Lam}{i,} Signore \ch{Re}{mio.}           \ch{Sol}{}}
\strofa
\rigo{\ch{Sol}{Per} il sole d'ogni g\ch{Lam}{iorno}}
\rigo{c\ch{Re}{he} riscalda e dona vita,   \ch{Mim}{}}
\rigo{e\ch{Do}{gli} illumina il cammino\ch{Re}{}}
\rigo{d\ch{Lam}{i} chi cerca Te, Signore.\ch{Re}{}}
\strofa
\commento{Rit.}
\strofa
\rigo{Per la luna e per le stelle}
\rigo{io le sento mie sorelle,}
\rigo{le hai formate su nel cielo}
\rigo{e le doni a chi è nel buio.}
\strofa
\commento{Rit.}
\strofa
\rigo{Per la nostra madre terra}
\rigo{che ci dona fiori ed erba,}
\rigo{su di lei noi fatichiamo}
\rigo{per il pane d'ogni giorno.}
\strofa
\commento{Rit.}
\strofa
\rigo{Per chi soffre con coraggio}
\rigo{e perdona nel Tuo amore}
\rigo{Tu gli dai la pace Tua}
\rigo{alla sera della vita.}
\strofa
\commento{Rit.}
\strofa
\rigo{Per la morte che è di tutti}
\rigo{io la sento ogni istante,}
\rigo{ma se vivo nel Tuo amore}
\rigo{dona un senso alla mia vita.}
\strofa
\commento{Rit.}
\strofa
\rigo{Per l'amore che è nel mondo}
\rigo{tra una donna e l'uomo suo,}
\rigo{per la vita dei bambini}
\rigo{che il mio mondo fanno nuovo.}
\strofa
\commento{Rit.}
\strofa
\rigo{Io Ti canto, mio Signore}
\rigo{e con me la creazione}
\rigo{Ti ringrazia umilmente}
\rigo{perché Tu sei il Signore.}
\strofa
\commento{Rit.}
\par\endgroup

\begingroup
\cantotesta{207}{Mani}{}{}{835pt}
\rigo{\ch{Re}{Vorrei} che le parol\ch{Sol}{e} mutassero \ch{La}{in} preghiera       \ch{Re}{}}
\rigo{e riv\ch{Sol}{ederti} \ch{Sim}{o} Padre che dipi\ch{Sol}{ngevi} il cie\ch{La}{lo}}
\rigo{Sapes\ch{Re}{si} quante volte\ch{Sol}{} guardando q\ch{La}{uesto} mondo        \ch{Re}{}}
\rigo{v\ch{Sol}{orrei} che tu tornass\ch{Sim}{i} a rito\ch{Sol}{ccar}\ch{La}{e} il cuore\ch{Re}{.}}
\rigo{V\ch{Sim}{orrei} che le mie mani av\ch{Fa\#m}{essero} la\ch{Sol}{} forza     \ch{Re}{}}
\rigo{per s\ch{Sol}{ostenere} ch\ch{Mi}{i} non può camminare    \ch{La}{}}
\rigo{V\ch{Sim}{orrei} che questo cuore c\ch{Fa\#m}{he} esplode in sen\ch{Sol}{timenti}         \ch{Re}{}}
\rigo{d\ch{Sol}{iventasse} c\ch{Sim}{ulla} per chi\ch{Sol}{} non ha più madr\ch{La}{e...}}
\strofa
\commento{Rit.}
\rigo{M\ch{Re}{ani,} prendi queste mie mani,       \ch{La}{}}
\rigo{fanne vita, fanne amore  \ch{Sol}{}}
\rigo{braccia aperte per riceve\ch{Sim}{re...Chi}\ch{La}{} è solo.}
\strofa
\rigo{\ch{Re}{Cuore,} prendi questo mio cuore,    \ch{La}{}}
\rigo{fa che si spalanchi al mondo \ch{Sol}{}}
\rigo{germogliando per quegli occhi    \ch{Sim}{}}
\rigo{LA                 SOL     (LA)}
\rigo{che non sanno pianger più.}
\strofa
\rigo{Sei tu lo spazio che desidero da sempre,}
\rigo{so che mi stringerai e mi terrai la mano.}
\rigo{Fa che le mie strade si perdano nel buio}
\rigo{ed io cammini dove cammineresti Tu.}
\rigo{Tu soffio della vita prendi la mia giovinezza}
\rigo{con le contraddizioni e le falsità}
\rigo{Strumento fa che sia per annunciare il Regno}
\rigo{a chi per queste vie Tu chiami Beati...}
\strofa
\commento{Rit.}
\strofa
\rigo{Noi giovani di un mondo che cancella i sentimenti}
\rigo{e inscatola le forze nell'asfalto di città}
\rigo{Siamo stanchi di guardare siamo stanchi di gridare}
\rigo{ci hai chiamati siamo Tuoi cammineremo insieme...}
\strofa
\commento{Rit.}
\rigo{M\ch{Mi}{ani,} prendi queste nostre mani         \ch{Si}{}}
\rigo{fanne vita fanne amore,  \ch{La}{}}
\rigo{braccia aperte per riceve\ch{Do\#m}{re...}\ch{Si}{Chi} è solo.}
\rigo{C\ch{Mi}{uori,} prendi questi nostri cuori   \ch{Si}{}}
\rigo{fa che siano testimon\ch{La}{i}}
\rigo{che Tu chiami ogni uomo a\ch{Do\#m}{} far festa co\ch{Si}{n} Dio.      \ch{La}{}}
\par\endgroup

\begingroup\senzaaccordi
\cantotesta{208}{Popoli tutti}{}{}{70pt}
\rigo{Popoli tutti lodate il Signore, amen,}
\rigo{cantiamo amici, gloria al Signore, amen.}
\par\endgroup

\begingroup
\cantotesta{209}{Ti ringrazio mio Signore}{}{}{716pt}
\rigo{\ch{Do}{}\ch{Lam}{}\ch{Fa}{}\ch{Sol}{Amatevi} l'un l'altro come Lui ha amato voi}
\rigo{\ch{Do}{}\ch{Lam}{}\ch{Fa}{}\ch{Sol}{e} siate per sempre Suoi amici}
\rigo{\ch{Fa}{}\ch{Sol}{}\ch{Do}{}\ch{Mim}{}\ch{Lam}{}\ch{Lam7}{e} quello che farete al più piccolo tra voi,}
\rigo{\ch{Fa}{}\ch{Sol}{}\ch{Do}{}\ch{Fa}{}\ch{Do}{credete} l'avete fatto a Lui.}
\strofa
\begin{ritornello}
\rigo{\ch{Do}{}\ch{Lam}{}\ch{Fa}{}\ch{Sol}{}\ch{Do}{}\ch{Lam}{Ti} ringrazio mio Signore non ho più paura,}
\rigo{\ch{Fa}{}\ch{Sol}{}\ch{Do}{}\ch{Lam}{}\ch{Fa}{}\ch{Sol}{}\ch{Do}{}\ch{Do7}{perché} con la mia mano nella mano degli amici miei}
\rigo{\ch{Fa}{}\ch{Sol}{}\ch{Do}{}\ch{Mim}{}\ch{Lam}{cammino} tra la gente della mia città}
\rigo{\ch{Fa}{}\ch{Sol}{}\ch{Do}{}\ch{Do7}{e} non mi sento più solo}
\rigo{\ch{Fa}{}\ch{Sol}{}\ch{Do}{}\ch{Mim}{}\ch{Lam}{}\ch{Lam7}{non} sento la stanchezza, guardo dritto avanti a me}
\rigo{\ch{Fa}{}\ch{Sol}{}\ch{Do}{}\ch{Fa}{}\ch{Do}{perché} sulla mia strada ci sei Tu...}
\strofa
\end{ritornello}
\rigo{Se amate veramente perdonatevi tra voi}
\rigo{nel cuore di ognuno ci sia pace:}
\rigo{il Padre che è nei cieli vede tutti i figli suoi}
\rigo{con gioia a voi perdonerà.}
\strofa
\begin{ritornello}
\rigo{\ch{Do}{}\ch{Lam}{}\ch{Fa}{}\ch{Sol}{}\ch{Do}{}\ch{Lam}{Ti} ringrazio mio Signore non ho più paura,}
\rigo{\ch{Fa}{}\ch{Sol}{}\ch{Do}{}\ch{Lam}{}\ch{Fa}{}\ch{Sol}{}\ch{Do}{}\ch{Do7}{perché} con la mia mano nella mano degli amici miei}
\rigo{\ch{Fa}{}\ch{Sol}{}\ch{Do}{}\ch{Mim}{}\ch{Lam}{cammino} tra la gente della mia città}
\rigo{\ch{Fa}{}\ch{Sol}{}\ch{Do}{}\ch{Do7}{e} non mi sento più solo}
\rigo{\ch{Fa}{}\ch{Sol}{}\ch{Do}{}\ch{Mim}{}\ch{Lam}{}\ch{Lam7}{non} sento la stanchezza, guardo dritto avanti a me}
\rigo{\ch{Fa}{}\ch{Sol}{}\ch{Do}{}\ch{Fa}{}\ch{Do}{perché} sulla mia strada ci sei Tu...}
\strofa
\end{ritornello}
\rigo{Sarete suoi amici se vi amate tra di voi}
\rigo{e questo è tutto il Suo Vangelo.}
\rigo{L'amore non ha prezzo, non misura ciò che da}
\rigo{l'amore confini non ne ha.}
\strofa
\begin{ritornello}
\rigo{\ch{Do}{}\ch{Lam}{}\ch{Fa}{}\ch{Sol}{}\ch{Do}{}\ch{Lam}{Ti} ringrazio mio Signore non ho più paura,}
\rigo{\ch{Fa}{}\ch{Sol}{}\ch{Do}{}\ch{Lam}{}\ch{Fa}{}\ch{Sol}{}\ch{Do}{}\ch{Do7}{perché} con la mia mano nella mano degli amici miei}
\rigo{\ch{Fa}{}\ch{Sol}{}\ch{Do}{}\ch{Mim}{}\ch{Lam}{cammino} tra la gente della mia città}
\rigo{\ch{Fa}{}\ch{Sol}{}\ch{Do}{}\ch{Do7}{e} non mi sento più solo}
\rigo{\ch{Fa}{}\ch{Sol}{}\ch{Do}{}\ch{Mim}{}\ch{Lam}{}\ch{Lam7}{non} sento la stanchezza, guardo dritto avanti a me}
\rigo{\ch{Fa}{}\ch{Sol}{}\ch{Do}{}\ch{Fa}{}\ch{Do}{perché} sulla mia strada ci sei Tu...}
\end{ritornello}
\par\endgroup

\begingroup
\cantotesta{210}{Voi siete di Dio}{}{}{727pt}
\rigo{\ch{Sol}{}\ch{Re}{}\ch{Do9}{}\ch{Mim7}{}\ch{Do9}{}\ch{Re}{}}
\strofa
\rigo{\ch{Sol}{}\ch{Re}{}\ch{Sol}{Tutte} le stelle della notte,}
\rigo{\ch{Sol}{}\ch{Re}{}\ch{Mim}{le} nebulose, le comete,}
\rigo{\ch{Sol}{}\ch{Re}{}\ch{Sol}{il} sole su una ragnatela,}
\rigo{\ch{Do}{}\ch{Sol}{}\ch{Re}{}\ch{Do}{è} tutto vostro e voi siete di Dio.}
\strofa
\rigo{\ch{Mi}{}\ch{Si}{}\ch{Mi}{Tutte} le rose della vita,}
\rigo{\ch{Mi}{}\ch{Si}{}\ch{Do\#m}{il} grano, i prati, i fili d'erba,}
\rigo{\ch{Mi}{}\ch{Si}{}\ch{Mi}{il} mare, i fiumi, le montagne,}
\rigo{\ch{La}{}\ch{Mi}{}\ch{Si}{}\ch{Mi}{è} tutto vostro e voi siete di Dio.}
\strofa
\rigo{\ch{Sol}{}\ch{Re}{}\ch{Sol}{Tutte} le musiche e le danze,}
\rigo{\ch{Sol}{}\ch{Re}{}\ch{Mim}{i} grattacieli, le astronavi,}
\rigo{\ch{Sol}{}\ch{Re}{}\ch{Sol}{i} quadri, i libri, le culture,}
\rigo{\ch{Do}{}\ch{Sol}{}\ch{Re}{}\ch{Sol}{è} tutto vostro e voi siete di Dio.}
\strofa
\rigo{\ch{Mi}{}\ch{Si}{}\ch{Mi}{Tutte} le volte che perdono,}
\rigo{\ch{Mi}{}\ch{Si}{}\ch{Do\#m}{quando} sorrido e quando piango,}
\rigo{\ch{Mi}{}\ch{Si}{}\ch{Mi}{quando} mi accorgo di chi sono,}
\rigo{\ch{La}{}\ch{Do\#m}{}\ch{Si}{}\ch{La}{è} tutto vostro e voi siete di Dio,}
\rigo{\ch{La}{}\ch{Mi}{}\ch{Si}{}\ch{Mi}{è} tutto nostro e noi siamo di Dio.}
\strofa
\rigo{\ch{Fa}{}\ch{Do}{}\ch{Sib9}{}\ch{Rem7}{}\ch{Sib9}{}\ch{Do}{}}
\strofa
\rigo{\ch{Fa}{}\ch{Do}{}\ch{Fa}{Tutte} le stelle della notte,}
\rigo{\ch{Fa}{}\ch{Do}{}\ch{Rem}{le} nebulose, le comete,}
\rigo{\ch{Fa}{}\ch{Do}{}\ch{Fa}{il} sole su una ragnatela,}
\rigo{\ch{Sib}{}\ch{Fa}{}\ch{Do}{}\ch{Sib}{è} tutto vostro e voi siete di Dio.}
\strofa
\rigo{\ch{Re}{}\ch{La}{}\ch{Re}{Tutte} le rose della vita,}
\rigo{\ch{Re}{}\ch{La}{}\ch{Sim}{il} grano, i prati, i fili d'erba,}
\rigo{\ch{Re}{}\ch{La}{}\ch{Re}{il} mare, i fiumi, le montagne,}
\rigo{\ch{Sol}{}\ch{Re}{}\ch{La}{}\ch{Re}{è} tutto vostro e voi siete di Dio.}
\strofa
\rigo{\ch{Fa}{}\ch{Do}{}\ch{Fa}{Tutte} le musiche e le danze,}
\rigo{\ch{Fa}{}\ch{Do}{}\ch{Rem}{i} grattacieli, le astronavi,}
\rigo{\ch{Fa}{}\ch{Do}{}\ch{Fa}{i} quadri, i libri, le culture,}
\rigo{\ch{Sib}{}\ch{Fa}{}\ch{Do}{}\ch{Fa}{è} tutto vostro e voi siete di Dio.}
\strofa
\rigo{\ch{Re}{}\ch{La}{}\ch{Re}{Tutte} le volte che perdono,}
\rigo{\ch{Re}{}\ch{La}{}\ch{Sim}{quando} sorrido e quando piango,}
\rigo{\ch{Re}{}\ch{La}{}\ch{Re}{quando} mi accorgo di chi sono,}
\rigo{\ch{Sol}{}\ch{Sim}{}\ch{La}{}\ch{Sol}{è} tutto vostro e voi siete di Dio,}
\rigo{\ch{Sol}{}\ch{Re}{}\ch{La}{}\ch{Re}{è} tutto nostro e noi siamo di Dio.}
\par\endgroup

\sottosezione{Mariani}

\begingroup\senzaaccordi
\cantotesta{211}{Come Maria}{}{}{444pt}
\rigo{Vogliamo vivere, Signore,}
\rigo{offrendo a Te la nostra vita;}
\rigo{con questo pane e questo vino accetta quello che noi siamo.}
\rigo{Vogliamo vivere, Signore,}
\rigo{abbandonati alla tua voce,}
\rigo{staccati dalle cose vano,}
\rigo{fissati nella vita vera.}
\strofa
\rigo{Vogliamo vivere come Maria,}
\rigo{l'irragiungibile,}
\rigo{la Madre amata,}
\rigo{che vince il mondo con l'Amore}
\rigo{e offrire sempre la Tua vita}
\rigo{che viene dal Cielo.}
\strofa
\rigo{Accetta dalle nostre mani}
\rigo{come un offerta a Te gradita,}
\rigo{i desideri di ogni cuore}
\rigo{le ansie della nostra vita.}
\rigo{Vogliamo vivere, Signore,}
\rigo{accesi dalle tue parole,}
\rigo{per ricordare ad ogni uomo}
\rigo{la fiamma viva del tuo Amore.}
\strofa
\rigo{Vogliamo vivere come Maria,}
\rigo{l'irragiungibile,}
\rigo{la madre amata}
\rigo{che vince il mondo con l'Amore}
\rigo{e offrire sempre la Tua vita}
\rigo{che viene dal Cielo.}
\rigo{e offrire sempre la Tua vita}
\rigo{che viene da Cielo.}
\par\endgroup

\begingroup
\cantotesta{212}{Ecco il nostro sì}{}{}{723pt}
\rigo{\ch{Re}{}\ch{Mim}{}\ch{Sol}{}\ch{La4}{}\ch{La}{}}
\strofa
\rigo{\ch{Re}{}\ch{Fa\#m}{}\ch{Sol}{}\ch{La}{Fra} tutte le donne scelta in Nazareth,}
\rigo{\ch{Sim}{}\ch{Sol}{}\ch{Mim}{}\ch{La}{}\ch{Re}{sul} tuo volto risplende il coraggio di quando hai detto “Sì”.}
\rigo{\ch{Fa\#m}{}\ch{Sol}{}\ch{La}{Insegna} a questo cuore l'umiltà,}
\rigo{\ch{Sim}{}\ch{Sol}{}\ch{Mim}{}\ch{La}{il} silenzio d'amore, la Speranza nel figlio tuo Gesù}
\strofa
\begin{ritornello}
\rigo{\ch{Re}{}\ch{Mim}{}\ch{Sol}{}\ch{La}{Ecco} il nostro Sì, nuova luce che rischiara il giorno,}
\rigo{\ch{Fa\#m}{}\ch{Sol}{}\ch{Mim}{}\ch{La}{è} bellissimo regalare al mondo la Speranza.}
\rigo{\ch{Re}{}\ch{Mim}{}\ch{Sol}{}\ch{La}{Ecco} il nostro Sì, camminiamo insieme a te Maria,}
\rigo{\ch{Fa\#m}{}\ch{Sol}{}\ch{Mim}{}\ch{La}{Madre} di Gesù, madre dell'umanità}
\strofa
\end{ritornello}
\rigo{\ch{Re}{}\ch{Fa\#m}{}\ch{Sol}{}\ch{La}{Nella} tua casa il verbo si rivelò}
\rigo{\ch{Sim}{}\ch{Sol}{}\ch{Mim}{}\ch{La}{}\ch{Re}{nel} segreto del cuore il respiro del figlio Emmanuel.}
\rigo{\ch{Fa\#m}{}\ch{Sol}{}\ch{La}{Insegna} a queste mani la fedeltà,}
\rigo{\ch{Sim}{}\ch{Sol}{}\ch{Mim}{}\ch{La}{a} costruire la pace, una casa comune insieme a te.}
\strofa
\begin{ritornello}
\rigo{\ch{Re}{}\ch{Mim}{}\ch{Sol}{}\ch{La}{Ecco} il nostro Sì, nuova luce che rischiara il giorno,}
\rigo{\ch{Fa\#m}{}\ch{Sol}{}\ch{Mim}{}\ch{La}{è} bellissimo regalare al mondo la Speranza.}
\rigo{\ch{Re}{}\ch{Mim}{}\ch{Sol}{}\ch{La}{Ecco} il nostro Sì, camminiamo insieme a te Maria,}
\rigo{\ch{Fa\#m}{}\ch{Sol}{}\ch{Mim}{}\ch{La}{Madre} di Gesù, madre dell'umanità}
\strofa
\end{ritornello}
\rigo{\ch{Re}{}\ch{Fa\#m}{}\ch{Sol}{}\ch{La}{Donna} dei nostri giorni sostienici,}
\rigo{\ch{Sim}{}\ch{Sol}{}\ch{Mim}{}\ch{La}{}\ch{Re}{guida} il nostro cammino con la forza di quando hai detto “Sì”.}
\rigo{\ch{Fa\#m}{}\ch{Sol}{}\ch{La}{Insegnaci} ad accogliere Gesù,}
\rigo{\ch{Sim}{}\ch{Sol}{}\ch{Mim}{}\ch{La}{noi} saremo Dimora, la più bella poesia dell'anima}
\strofa
\begin{ritornello}
\rigo{\ch{Re}{}\ch{Mim}{}\ch{Sol}{}\ch{La}{Ecco} il nostro Sì, nuova luce che rischiara il giorno,}
\rigo{\ch{Fa\#m}{}\ch{Sol}{}\ch{Mim}{}\ch{La}{è} bellissimo regalare al mondo la Speranza.}
\rigo{\ch{Re}{}\ch{Mim}{}\ch{Sol}{}\ch{La}{Ecco} il nostro Sì, camminiamo insieme a te Maria,}
\rigo{\ch{Fa\#m}{}\ch{Sol}{}\ch{Mim}{}\ch{La}{Madre} di Gesù, madre dell'umanità}
\end{ritornello}
\par\endgroup

\begingroup\senzaaccordi
\cantotesta{213}{Madonna degli scout}{A. Ghetti - V. Ghetti}{}{217pt}
\rigo{Madonna degli scouts, ascolta! T'invochiam:}
\rigo{concedi un forte cuore a noi, che ora partiam.}
\rigo{La strada è tanto lunga, il freddo già ci assal:}
\rigo{respingi tu, Regina, lo spirito del mal.}
\strofa
\rigo{E il ritmo dei passi ci accompagnerà:}
\rigo{là, verso gli orizzonti lontani si va.}
\strofa
\rigo{E lungo quella strada non ci lasciare tu:}
\rigo{nel volto di chi soffre saprem trovar Gesù.}
\rigo{Allor ci fermeremo le piaghe a medicar}
\rigo{e il pianto di chi è solo sapremo consolar.}
\strofa
\rigo{E il ritmo dei passi ci accompagnerà...}
\par\endgroup

\begingroup\senzaaccordi
\cantotesta{214}{Mamma del cielo}{}{\bLC}{217pt}
\rigo{Mamma del Cielo, Vergine Santa}
\rigo{a te d’intorno il branco canta}
\rigo{Fa’ che ci amiamo come fratelli}
\rigo{mentre la pista noi percorriam.}
\strofa
\rigo{Mamma del Cielo, rendici tu}
\rigo{cortesi e buoni come Gesù.}
\strofa
\rigo{Fa che io ami ogni vicino}
\rigo{come se fosse Gesù bambino.}
\rigo{Del mio meglio fare prometto}
\rigo{per diventare un Tuo lupetto}
\strofa
\rigo{Al cuor divino stringimi tu}
\rigo{come stringevi a te Gesù.}
\par\endgroup

\sottosezione{Preghiere scout}

\begingroup\senzaaccordi
\cantotesta{215}{Preghiera del lupetto}{}{\bLC}{196pt}
\rigo{Buon Gesù}
\rigo{io ti amo con tutto il mio cuore,}
\rigo{perché mai troverò migliore amico di Te.}
\rigo{Per farti piacere io voglio,}
\rigo{con l’aiuto di Maria nostra mamma del cielo,}
\rigo{fare del mio meglio,}
\rigo{come Tu facevi ogni giorno a Nazareth.}
\rigo{Insegnami dunque ad ascoltare non me stesso}
\rigo{ma la voce,}
\rigo{quella dei miei genitori,}
\rigo{dei miei superiori,}
\rigo{dei miei capi. Così sia.}
\par\endgroup

\begingroup\senzaaccordi
\cantotesta{216}{Preghiera della sera}{}{\bLC}{1175pt}
\rigo{Al cader della giornata}
\rigo{noi leviamo i cuori a Te,}
\rigo{Tu l'avevi a noi donata}
\rigo{bene spesa fu per Te.}
\strofa
\rigo{Te nel bosco, nel ruscello}
\rigo{Te nel monte Te nel mar,}
\rigo{Te nel cuore del fratello,}
\rigo{Te nel mio cercai d'amar.}
\strofa
\rigo{La strofa segreta}
\strofa
\rigo{"...e se un bel dì l'ASCI risorgerà tutti compatti ci troverà!" La guerra era finita da pochi mesi e il sogno della canzone delle Aquile randagie si avverava. Man mano che le città venivano liberate, immediatamente si ricostituivano ad opera dei vecchi Capi le prime unità e tornavano a sventolare fiamme e guidoni della risorta ASCI.}
\rigo{Così, per il S. Giorgio del 1946 il commissariato regionale lombardo fu in grado di organizzare ad Affori un grandioso raduno che attirò in un'atmosfera di entusiasmo e di fraternità molti Gruppi di altre regioni e persino dalla Svizzera. Fraterno amico, da tanti anni, dell'indimenticabile Baden e di Arnaldo Basini, fui invitato anch'io col Riparto che avevo fondato a Parma nel '45; e sebbene le condizioni dei trasporti fossero ancora molto precarie e difficili, accettai con gioia. Lo stile scout esigeva (non so se l'usanza sopravviva ancora) che chi si reca in visita a un campo porti sempre con sé qualcosa di utile da offrire: sia pure un semplice ramo per il fuoco.}
\strofa
\rigo{Così pensai di portare al campo dei fratelli milanesi, insieme con le mie quattro squadriglie, qualcosa che durasse più di un po' di legna: portai in dono le parole che avevo scritto poche settimane prima, da cantare sulla melodia della nota canzone friulana: "Ai preat la biele stele". In una notte di veglia forzata, infatti, mi erano venute spontanee e rapide, una dopo l'altra, le tre strofe di quella "Preghiera della sera" che quasi tutti gli scout, almeno parzialmente, cantano al fuoco di bivacco. Dico 'almeno parzialmente' poiché quasi tutte le raccolte di canti scout oggi non riportano la seconda, che è la più bella delle tre strofe di cui è composta la canzone "Al cader della giornata", che Baden gradì moltissimo e che volle inserita subito nel canzoniere "Canti di mezzanotte". La volete riportare sul vostro quaderno di caccia questa povera strofa derelitta?... Eccovela:}
\strofa
\rigo{Se non sempre la mia mente in Te pura si affissò;}
\rigo{se talora stoltamente da Te lungi si attardò...}
\strofa
\rigo{io Signor, ne son dolente: Te ne chiedo, o Dio, mercé!}
\rigo{el mio meglio lietamente, io farò doman per Te.}
\strofa
\rigo{Poi vengono i cieli che sembran prati e il bivacco delle stelle; ma adesso capite, se non la conoscevate ancora, perché questa strofa, a me che la scrissi per voi, sembri più bella delle altre. Voi che ne pensate?}
\strofa
\rigo{d. Tar (Aquila Azzurra)}
\rigo{da «scout-proposta educativa»}
\strofa
\rigo{Se non sempre la mia mente}
\rigo{in Te pura s'affissò;}
\rigo{se talora stoltamente}
\rigo{da Te lungi s'attardò,}
\rigo{mio Signore ne son dolente}
\rigo{Te ne chieggo, o Dio, mercè!}
\rigo{Del mio meglio lietamente}
\rigo{io farò doman per Te.}
\strofa
\rigo{I tuoi cieli sembran prati}
\rigo{e le stelle tanti fior...}
\strofa
\rigo{Son bivacchi dei beati}
\rigo{stretti in cerchio al lor Signor.}
\strofa
\rigo{Quante stelle quante stelle}
\rigo{dimmi Tu la mia qual'è'?}
\rigo{Non ambisco la più bella}
\rigo{basta sia vicino a Te!..}
\par\endgroup

\sottosezione{Preghiere a tavola}

\begingroup
\cantotesta{217}{Dalla nascita del sol}{tradizionale}{}{178pt}
\rigo{\ch{Re}{Dalla} nascita del sol, fino al suo tramonto vorrò}
\rigo{\ch{Mim}{Lodare} il \ch{La7}{Signor,} con \ch{Re}{tutto} \ch{Re7}{il} cuor. (bis)}
\strofa
\begin{ritornello}
\rigo{\ch{Sol}{Loda} il \ch{Re}{Signor,} \ch{Fa\#7}{lodatelo} voi}
\rigo{\ch{Sol}{servitori} \ch{La7}{suoi,} \ch{Re}{glorificate} \ch{Re7}{il} creator.}
\rigo{\ch{Sol}{Benedetto} \ch{Re}{è} il nome \ch{Fa\#7}{Suo}}
\rigo{\ch{Sol}{Da} ora in \ch{La7}{poi,} e per \ch{Re}{l'eternità.}}
\end{ritornello}
\par\endgroup

\begingroup\senzaaccordi
\cantotesta{218}{E prima di mangiar}{tradizionale}{}{128pt}
\rigo{Per i sentieri da non smarrir}
\rigo{per il presente e l'avvenir}
\rigo{per questo pan che noi mangiam}
\rigo{per questi amici Ti lodiam!}
\rigo{E prima di mangiar cantiamo tutti in cor}
\rigo{gloria gloria gloria gloria gloria al Signor!}
\par\endgroup

\begingroup\senzaaccordi
\cantotesta{219}{Grazie}{tradizionale}{}{284pt}
\rigo{Grazie, signore per il tuo amore testimoniato da una croce.}
\rigo{Grazie per il sole che ci riscalda e ci fa vivere.}
\rigo{Grazie per l’amore di papà e mamma che richiama il tuo amore.}
\rigo{Grazie per la maestà delle montagne, che richiamano la tua grandezza.}
\rigo{Grazie per i piccoli fiori del prato che mi richiamano la tua fedeltà}
\rigo{Grazie per questo passero che vola nel cielo che mi richiamala libertà}
\rigo{Grazie per lo Spirito che mi illumina ogni giorno.}
\rigo{Grazie per gli amici che camminano con me.}
\rigo{Grazie per la Chiesa, madre di tutti i viventi.}
\rigo{Grazie per ogni uomo di buona volontà che cammina nelle tue vie.}
\par\endgroup

\begingroup\senzaaccordi
\cantotesta{220}{Per questo pan}{tradizionale}{}{146pt}
\rigo{Per questo pan, per ogni don,}
\rigo{noi ti lodiam Signor}
\rigo{Pour ce repas, pour toute joie}
\rigo{nous Te louons, Seigneur!}
\rigo{Für Hilf und Kraft, für taglich Brot}
\rigo{wir danken Dir, mein Gott!}
\rigo{For bread strength, for ev'ry gift}
\rigo{we sing praise you Lord!}
\par\endgroup

\end{multicols}

\sezione{Canti del risveglio}{}

\begin{multicols}{2}\raggedcolumns

\begingroup
\cantotesta{221}{Al chiaror del mattin}{}{\bEG}{519pt}
\rigo{Al c\ch{Sol}{hiaror} del mattin, al chiaror del mattin}
\rigo{ci sv\ch{Re7}{eglia} una canzon, ci sveglia\ch{Do}{} una canzon.   \ch{Sol}{}}
\rigo{Al chiaror del mattin, al chiaror del mattin}
\rigo{ci sv\ch{Re7}{eglia} una canzon, è sorto il sol.   \ch{Sol}{}}
\strofa
\commento{Rit.}
\rigo{Su leviam, su leviam,}
\rigo{c'è un magnifico \ch{Re7}{ciel}\ch{Do}{o} limpido;\ch{Sol}{}}
\rigo{Su leviam, su leviam,}
\rigo{non è tem\ch{Re7}{po} di esitar\ch{Do}{.}     \ch{Sol}{}}
\strofa
\rigo{Se la pioggia cadrà, se la pioggia cadrà}
\rigo{e tutto bagnerà, e tutto bagnerà.}
\rigo{Se la pioggia cadrà, se la pioggia cadrà}
\rigo{più divertente ancor sarà.}
\strofa
\commento{Rit.}
\strofa
\rigo{Se il vento verra', se il vento verrà}
\rigo{e tutto muovera' e tutto bagnerà.}
\rigo{Se il vento verra', se il vento verrà}
\rigo{piu' divertente ancor sara'.}
\strofa
\commento{Rit.}
\strofa
\rigo{Se la neve verra', se la neve verrà}
\rigo{e tutto coprira', e tutto coprirà.}
\rigo{Se la neve verra', se la neve verrà}
\rigo{piu' divertente ancor sara'.}
\strofa
\commento{Rit.}
\par\endgroup

\begingroup
\cantotesta{222}{Suona già il risveglio}{}{\bEG}{403pt}
\begin{ritornello}
\rigo{\ch{Fa}{Suona} già il risveglio,}
\rigo{spunta il \ch{Do}{dì,} il mattino è \ch{Fa}{bello;}}
\rigo{su, su, lesti, via saltiamo in \ch{Do}{piè.}}
\rigo{\ch{Rem}{L'aria} \ch{Do}{}\ch{Fa}{pura} al \ch{Do}{mattin}}
\rigo{va \ch{Fa}{nei} polmoni e ci \ch{Do}{fa} \ch{Fa}{cantar:}}
\strofa
\end{ritornello}
\rigo{\ch{Do}{perchè} c'è il \ch{Fa}{sol,} \ch{Do}{ch'è} tanto \ch{Fa}{bel,}}
\rigo{\ch{Do}{quando} i \ch{Fa}{V.P.} non l'han \ch{Do}{veduto} ancor.}
\strofa
\begin{ritornello}
\rigo{\ch{Fa}{Suona} già il risveglio,}
\rigo{spunta il \ch{Do}{dì,} il mattino è \ch{Fa}{bello;}}
\rigo{su, su, lesti, via saltiamo in \ch{Do}{piè.}}
\rigo{\ch{Rem}{L'aria} \ch{Do}{}\ch{Fa}{pura} al \ch{Do}{mattin}}
\rigo{va \ch{Fa}{nei} polmoni e ci \ch{Do}{fa} \ch{Fa}{cantar:}}
\strofa
\end{ritornello}
\rigo{\ch{Do}{se} non c'è il \ch{Fa}{sol} ch'è \ch{Do}{tanto} \ch{Fa}{bel,}}
\rigo{\ch{Do}{col} nostro \ch{Fa}{canto} tornerà \ch{Do}{seren.}}
\strofa
\begin{ritornello}
\rigo{\ch{Fa}{Suona} già il risveglio,}
\rigo{spunta il \ch{Do}{dì,} il mattino è \ch{Fa}{bello;}}
\rigo{su, su, lesti, via saltiamo in \ch{Do}{piè.}}
\rigo{\ch{Rem}{L'aria} \ch{Do}{}\ch{Fa}{pura} al \ch{Do}{mattin}}
\rigo{va \ch{Fa}{nei} polmoni e ci \ch{Do}{fa} \ch{Fa}{cantar.}}
\end{ritornello}
\par\endgroup

\end{multicols}

% ---- Indici
\titolosezione{Indice alfabetico}
\begingroup\small
\begin{multicols}{2}
\lettera{A}
\voce{14}{A caccia andiamo}
\voce{201}{Abbracciami}
\voce{170}{Accogli i nostri doni}
\voce{153}{Acqua siamo noi}
\voce{5}{Addio lupetto}
\voce{98}{Ah io vorrei tornare}
\voce{126}{Airone}
\voce{221}{Al chiaror del mattin}
\voce{70}{Al passo del guidon}
\voce{44}{Alla fiera dell'est}
\voce{165}{Alleluia la tua festa}
\voce{166}{Alleluia! È il tempo della gioia}
\voce{167}{Alleluia, ed oggi ancora}
\voce{168}{Alleluia, la parola tua}
\voce{169}{Alleluia, passeranno i cieli}
\voce{202}{Andate per le strade}
\voce{41}{Andiamo al falò}
\voce{15}{Antica traccia}
\voce{171}{Antica, eterna danza}
\voce{45}{Arturo e Gertrude}
\voce{16}{Attorno alla rupe}
\lettera{B}
\voce{46}{Banana cocco baobab}
\voce{180}{Beati voi}
\voce{47}{Bella però sta macchina}
\voce{172}{Benedetto sei tu}
\voce{181}{Benedetto Signore (cerco solo te)}
\voce{173}{Benedici, o Signore}
\voce{134}{Bersagliere ha cento penne}
\voce{78}{Bevi bevi}
\voce{99}{Bocca di rosa}
\voce{17}{Branco di Akela}
\voce{18}{Buon lupetto}
\voce{19}{Buona caccia}
\lettera{C}
\voce{182}{C'è ancora un sole}
\voce{20}{Caccia al sambhur}
\voce{154}{Camminerò}
\voce{155}{Cantiamo a te}
\voce{100}{Canto del ricordo}
\voce{21}{Canto del risveglio}
\voce{7}{Canto della promessa}
\voce{6}{Canto della promessa del lupetto}
\voce{22}{Canto della sera}
\voce{101}{Canto delle Aquile Randagie (La luna delle vette)}
\voce{10}{Canto delle Scolte d'Assisi}
\voce{3}{Canto di gruppo}
\voce{102}{Canto di Montegemoli}
\voce{183}{Canto di San Damiano}
\voce{103}{Canto Eurojam 94}
\voce{4}{Canto Riparto TV7}
\voce{79}{Canzone dell'amicizia}
\voce{80}{Caramba}
\voce{104}{Carnet di marcia}
\voce{156}{Che gioia ci hai dato}
\voce{105}{Col cappellone}
\voce{48}{Colla dei francobolli}
\voce{106}{Colore del sole}
\voce{184}{Come fuoco vivo}
\voce{211}{Come Maria}
\voce{185}{Come ti ama Dio}
\voce{186}{Come tu mi vuoi}
\voce{203}{Con te faremo cose grandi}
\lettera{D}
\voce{157}{Dall'aurora al tramonto}
\voce{217}{Dalla nascita del sol}
\voce{23}{Danza dei cani rossi}
\voce{24}{Danza del fiore rosso}
\voce{25}{Danza della tregua dell'acqua}
\voce{187}{Dimmelo tu}
\voce{135}{Dove sei stato mio bell'alpino}
\voce{81}{Drizza la tenda}
\lettera{E}
\voce{107}{E mi la dona bionda}
\voce{218}{E prima di mangiar}
\voce{188}{E sono solo un uomo}
\voce{212}{Ecco il nostro sì}
\voce{174}{Ecco quel che abbiamo}
\voce{204}{Eccomi per sempre}
\voce{71}{El pescador da Ciosa}
\voce{82}{El tram de Opcina}
\voce{158}{Emmanuel}
\voce{136}{Era una notte che pioveva}
\voce{8}{Estote parati}
\lettera{F}
\voce{72}{Fare well}
\voce{108}{Faria faria oh}
\voce{109}{Fiume rosso}
\voce{49}{Fli flai}
\voce{26}{Fratelli alla candida luna}
\voce{205}{Fratello sole e sorella luna}
\voce{175}{Frutto della nostra terra}
\voce{189}{Fuoco per il mondo}
\lettera{G}
\voce{127}{Geordie}
\voce{137}{Gerusalemme}
\voce{190}{Gesù sei qui con me}
\voce{159}{Gioia che invade l'anima}
\voce{27}{Gli allegri folletti}
\voce{160}{Grandi cose}
\voce{219}{Grazie}
\voce{110}{Guarda lontano}
\voce{111}{Guardati intorno scout}
\lettera{I}
\voce{83}{I do gobeti}
\voce{191}{I.N.R.I. (io non ritorno indietro)}
\voce{28}{Il canto del cuculo}
\voce{192}{Il canto dell'amore}
\voce{50}{Il duca di Barnabò}
\voce{128}{Il falco}
\voce{29}{Il lupo ferito}
\voce{84}{Il mio galletto}
\voce{112}{Il pescatore}
\voce{51}{Il ranocchio ciclista}
\voce{30}{Il risveglio della giungla}
\voce{129}{Il testamento del capitano}
\voce{52}{Il tortellino}
\voce{130}{Il vascello fantasma}
\voce{1}{Inno d'Italia}
\voce{2}{Inno FSE}
\voce{193}{Insieme}
\voce{161}{Invochiamo la tua presenza}
\voce{194}{Io sarò con te}
\lettera{J}
\voce{53}{Jack è in cucina}
\voce{85}{John Brown}
\voce{54}{Jona nella balena}
\voce{113}{Joska la rossa}
\voce{55}{Jump down}
\lettera{K}
\voce{114}{Kamaludu}
\lettera{L}
\voce{56}{L'albero piantato}
\voce{31}{L'ultimo canto di Akela}
\voce{86}{L'uomo di Cromagnon}
\voce{57}{La battaglia di Magenta}
\voce{87}{La bella la va al fosso}
\voce{58}{La bella polenta}
\voce{88}{La bomba imbriaga}
\voce{32}{La cassa misteriosa}
\voce{115}{La ceseta de Transaqua}
\voce{89}{La damigiana}
\voce{90}{La Gigia l'è malada}
\voce{116}{La guadalupa}
\voce{117}{La guerra di Piero}
\voce{42}{La leggenda del fuoco}
\voce{138}{La leggenda del Piave}
\voce{59}{La macchina del capo}
\voce{118}{La montanara}
\voce{119}{La mula de Parenzo}
\voce{131}{La piroga}
\voce{139}{La tradotta}
\voce{140}{La vien giù dalle montagne}
\voce{206}{Laudato sii, Signore mio}
\voce{91}{Lazy boy}
\voce{33}{Lupo salta su}
\lettera{M}
\voce{213}{Madonna degli scout}
\voce{214}{Mamma del cielo}
\voce{207}{Mani}
\voce{73}{Masafacia'}
\voce{92}{Me compare Giacometo}
\voce{74}{Milan brucia}
\voce{141}{Monte Canino}
\voce{142}{Monte Pasubio}
\lettera{N}
\voce{60}{Napuletana mia}
\voce{195}{Nel tuo volto la mia strada}
\voce{34}{Nella giungla misteriosa}
\voce{93}{Non si va in cielo}
\lettera{O}
\voce{176}{Ogni mia parola}
\voce{94}{Oleanna}
\voce{61}{Onci onci}
\voce{62}{One banana}
\lettera{P}
\voce{196}{Pane del cielo}
\voce{220}{Per questo pan}
\voce{120}{Pim pam}
\voce{208}{Popoli tutti}
\voce{215}{Preghiera del lupetto}
\voce{216}{Preghiera della sera}
\voce{197}{Preghiera semplice}
\lettera{Q}
\voce{63}{Quattro pirati}
\voce{143}{Quel mazzolin di fiori}
\lettera{S}
\voce{43}{Salite o fiamme}
\voce{35}{Salta o lupo}
\voce{64}{San Michele}
\voce{179}{Santo (Dio del cielo)}
\voce{132}{Scende la notte (esploratori)}
\voce{11}{Se incontrerò}
\voce{177}{Se m'accogli}
\voce{144}{Se non ci conoscete}
\voce{198}{Sei fuoco e vento}
\voce{12}{Servire}
\voce{178}{Servo per amore}
\voce{95}{Si va}
\voce{36}{Siam nella vasta giungla}
\voce{65}{Siamo andati alla caccia del leon}
\voce{133}{Signor, fra le tende schierati}
\voce{145}{Signore delle cime}
\voce{146}{Stelutis alpinis}
\voce{13}{Strade delle mie routes}
\voce{199}{Su ali d'aquila}
\voce{147}{Sul cappello}
\voce{96}{Sul cappello un bel fior}
\voce{75}{Sul colle}
\voce{66}{Sul lago Tanganika}
\voce{148}{Sul Monte Nero}
\voce{121}{Sul pajon}
\voce{149}{Sul ponte di Bassano}
\voce{150}{Sul ponte di Perati}
\voce{222}{Suona già il risveglio}
\voce{200}{Symbolum 77}
\lettera{T}
\voce{151}{Ta-pum}
\voce{122}{Terra della betulla}
\voce{209}{Ti ringrazio mio Signore}
\voce{76}{Tinchete Ritinchete}
\lettera{U}
\voce{37}{Ula, ula, ula}
\voce{38}{Un bravo lupo}
\voce{67}{Un ciodo de fero vecio}
\voce{68}{Un elefante}
\voce{9}{Un mondo, una promessa}
\voce{69}{Uomini della filibusta}
\lettera{V}
\voce{152}{Valsugana}
\voce{77}{Vecchio scarpone}
\voce{123}{Vento della sera}
\voce{124}{Vento sottile}
\voce{162}{Vieni spirito d'amore}
\voce{97}{Vinassa, vinassa}
\voce{163}{Vivere la vita}
\voce{164}{Vocazione}
\voce{125}{Voglio girare il mondo}
\voce{210}{Voi siete di Dio}
\lettera{Z}
\voce{39}{Zampe agili e silenti}
\voce{40}{Zanna bianca}
\end{multicols}
\endgroup
\titolosezione{Indice per categoria}
\begingroup\small
\begin{multicols}{2}
\categoriaindice{Canti associativi}
\voce{3}{Canto di gruppo}
\voce{4}{Canto Riparto TV7}
\voce{1}{Inno d'Italia}
\voce{2}{Inno FSE}
\categoriaindice{Canti cerimoniali}
\voce{5}{Addio lupetto}
\voce{7}{Canto della promessa}
\voce{6}{Canto della promessa del lupetto}
\voce{10}{Canto delle Scolte d'Assisi}
\voce{8}{Estote parati}
\voce{11}{Se incontrerò}
\voce{12}{Servire}
\voce{13}{Strade delle mie routes}
\voce{9}{Un mondo, una promessa}
\categoriaindice{Canti di lupettismo}
\voce{14}{A caccia andiamo}
\voce{15}{Antica traccia}
\voce{16}{Attorno alla rupe}
\voce{17}{Branco di Akela}
\voce{18}{Buon lupetto}
\voce{19}{Buona caccia}
\voce{20}{Caccia al sambhur}
\voce{21}{Canto del risveglio}
\voce{22}{Canto della sera}
\voce{23}{Danza dei cani rossi}
\voce{24}{Danza del fiore rosso}
\voce{25}{Danza della tregua dell'acqua}
\voce{26}{Fratelli alla candida luna}
\voce{27}{Gli allegri folletti}
\voce{28}{Il canto del cuculo}
\voce{29}{Il lupo ferito}
\voce{30}{Il risveglio della giungla}
\voce{31}{L'ultimo canto di Akela}
\voce{32}{La cassa misteriosa}
\voce{33}{Lupo salta su}
\voce{34}{Nella giungla misteriosa}
\voce{35}{Salta o lupo}
\voce{36}{Siam nella vasta giungla}
\voce{37}{Ula, ula, ula}
\voce{38}{Un bravo lupo}
\voce{39}{Zampe agili e silenti}
\voce{40}{Zanna bianca}
\categoriaindice{Fuoco di bivacco}
\sottocategoriaindice{Apertura}
\voce{41}{Andiamo al falò}
\voce{42}{La leggenda del fuoco}
\voce{43}{Salite o fiamme}
\sottocategoriaindice{Bans}
\voce{44}{Alla fiera dell'est}
\voce{45}{Arturo e Gertrude}
\voce{46}{Banana cocco baobab}
\voce{47}{Bella però sta macchina}
\voce{48}{Colla dei francobolli}
\voce{49}{Fli flai}
\voce{50}{Il duca di Barnabò}
\voce{51}{Il ranocchio ciclista}
\voce{52}{Il tortellino}
\voce{53}{Jack è in cucina}
\voce{54}{Jona nella balena}
\voce{55}{Jump down}
\voce{56}{L'albero piantato}
\voce{57}{La battaglia di Magenta}
\voce{58}{La bella polenta}
\voce{59}{La macchina del capo}
\voce{60}{Napuletana mia}
\voce{61}{Onci onci}
\voce{62}{One banana}
\voce{63}{Quattro pirati}
\voce{64}{San Michele}
\voce{65}{Siamo andati alla caccia del leon}
\voce{66}{Sul lago Tanganika}
\voce{67}{Un ciodo de fero vecio}
\voce{68}{Un elefante}
\voce{69}{Uomini della filibusta}
\sottocategoriaindice{Crescendo}
\voce{70}{Al passo del guidon}
\voce{71}{El pescador da Ciosa}
\voce{72}{Fare well}
\voce{73}{Masafacia'}
\voce{74}{Milan brucia}
\voce{75}{Sul colle}
\voce{76}{Tinchete Ritinchete}
\voce{77}{Vecchio scarpone}
\sottocategoriaindice{Picco}
\voce{78}{Bevi bevi}
\voce{79}{Canzone dell'amicizia}
\voce{80}{Caramba}
\voce{81}{Drizza la tenda}
\voce{82}{El tram de Opcina}
\voce{83}{I do gobeti}
\voce{84}{Il mio galletto}
\voce{85}{John Brown}
\voce{86}{L'uomo di Cromagnon}
\voce{87}{La bella la va al fosso}
\voce{88}{La bomba imbriaga}
\voce{89}{La damigiana}
\voce{90}{La Gigia l'è malada}
\voce{91}{Lazy boy}
\voce{92}{Me compare Giacometo}
\voce{93}{Non si va in cielo}
\voce{94}{Oleanna}
\voce{95}{Si va}
\voce{96}{Sul cappello un bel fior}
\voce{97}{Vinassa, vinassa}
\sottocategoriaindice{Calando}
\voce{98}{Ah io vorrei tornare}
\voce{99}{Bocca di rosa}
\voce{100}{Canto del ricordo}
\voce{101}{Canto delle Aquile Randagie (La luna delle vette)}
\voce{102}{Canto di Montegemoli}
\voce{103}{Canto Eurojam 94}
\voce{104}{Carnet di marcia}
\voce{105}{Col cappellone}
\voce{106}{Colore del sole}
\voce{107}{E mi la dona bionda}
\voce{108}{Faria faria oh}
\voce{109}{Fiume rosso}
\voce{110}{Guarda lontano}
\voce{111}{Guardati intorno scout}
\voce{112}{Il pescatore}
\voce{113}{Joska la rossa}
\voce{114}{Kamaludu}
\voce{115}{La ceseta de Transaqua}
\voce{116}{La guadalupa}
\voce{117}{La guerra di Piero}
\voce{118}{La montanara}
\voce{119}{La mula de Parenzo}
\voce{120}{Pim pam}
\voce{121}{Sul pajon}
\voce{122}{Terra della betulla}
\voce{123}{Vento della sera}
\voce{124}{Vento sottile}
\voce{125}{Voglio girare il mondo}
\sottocategoriaindice{Chiusura}
\voce{126}{Airone}
\voce{127}{Geordie}
\voce{128}{Il falco}
\voce{129}{Il testamento del capitano}
\voce{130}{Il vascello fantasma}
\voce{131}{La piroga}
\sottocategoriaindice{Ultimi}
\voce{132}{Scende la notte (esploratori)}
\voce{133}{Signor, fra le tende schierati}
\categoriaindice{Canti alpini}
\voce{134}{Bersagliere ha cento penne}
\voce{135}{Dove sei stato mio bell'alpino}
\voce{136}{Era una notte che pioveva}
\voce{137}{Gerusalemme}
\voce{138}{La leggenda del Piave}
\voce{139}{La tradotta}
\voce{140}{La vien giù dalle montagne}
\voce{141}{Monte Canino}
\voce{142}{Monte Pasubio}
\voce{143}{Quel mazzolin di fiori}
\voce{144}{Se non ci conoscete}
\voce{145}{Signore delle cime}
\voce{146}{Stelutis alpinis}
\voce{147}{Sul cappello}
\voce{148}{Sul Monte Nero}
\voce{149}{Sul ponte di Bassano}
\voce{150}{Sul ponte di Perati}
\voce{151}{Ta-pum}
\voce{152}{Valsugana}
\categoriaindice{Canti liturgici}
\sottocategoriaindice{Ingresso}
\voce{153}{Acqua siamo noi}
\voce{154}{Camminerò}
\voce{155}{Cantiamo a te}
\voce{156}{Che gioia ci hai dato}
\voce{157}{Dall'aurora al tramonto}
\voce{158}{Emmanuel}
\voce{159}{Gioia che invade l'anima}
\voce{160}{Grandi cose}
\voce{161}{Invochiamo la tua presenza}
\voce{162}{Vieni spirito d'amore}
\voce{163}{Vivere la vita}
\voce{164}{Vocazione}
\sottocategoriaindice{Alleluia}
\voce{165}{Alleluia la tua festa}
\voce{166}{Alleluia! È il tempo della gioia}
\voce{167}{Alleluia, ed oggi ancora}
\voce{168}{Alleluia, la parola tua}
\voce{169}{Alleluia, passeranno i cieli}
\sottocategoriaindice{Offertorio}
\voce{170}{Accogli i nostri doni}
\voce{171}{Antica, eterna danza}
\voce{172}{Benedetto sei tu}
\voce{173}{Benedici, o Signore}
\voce{174}{Ecco quel che abbiamo}
\voce{175}{Frutto della nostra terra}
\voce{176}{Ogni mia parola}
\voce{177}{Se m'accogli}
\voce{178}{Servo per amore}
\sottocategoriaindice{Santo}
\voce{179}{Santo (Dio del cielo)}
\sottocategoriaindice{Comunione}
\voce{180}{Beati voi}
\voce{181}{Benedetto Signore (cerco solo te)}
\voce{182}{C'è ancora un sole}
\voce{183}{Canto di San Damiano}
\voce{184}{Come fuoco vivo}
\voce{185}{Come ti ama Dio}
\voce{186}{Come tu mi vuoi}
\voce{187}{Dimmelo tu}
\voce{188}{E sono solo un uomo}
\voce{189}{Fuoco per il mondo}
\voce{190}{Gesù sei qui con me}
\voce{191}{I.N.R.I. (io non ritorno indietro)}
\voce{192}{Il canto dell'amore}
\voce{193}{Insieme}
\voce{194}{Io sarò con te}
\voce{195}{Nel tuo volto la mia strada}
\voce{196}{Pane del cielo}
\voce{197}{Preghiera semplice}
\voce{198}{Sei fuoco e vento}
\voce{199}{Su ali d'aquila}
\voce{200}{Symbolum 77}
\sottocategoriaindice{Finale}
\voce{201}{Abbracciami}
\voce{202}{Andate per le strade}
\voce{203}{Con te faremo cose grandi}
\voce{204}{Eccomi per sempre}
\voce{205}{Fratello sole e sorella luna}
\voce{206}{Laudato sii, Signore mio}
\voce{207}{Mani}
\voce{208}{Popoli tutti}
\voce{209}{Ti ringrazio mio Signore}
\voce{210}{Voi siete di Dio}
\sottocategoriaindice{Mariani}
\voce{211}{Come Maria}
\voce{212}{Ecco il nostro sì}
\voce{213}{Madonna degli scout}
\voce{214}{Mamma del cielo}
\sottocategoriaindice{Preghiere scout}
\voce{215}{Preghiera del lupetto}
\voce{216}{Preghiera della sera}
\sottocategoriaindice{Preghiere a tavola}
\voce{217}{Dalla nascita del sol}
\voce{218}{E prima di mangiar}
\voce{219}{Grazie}
\voce{220}{Per questo pan}
\categoriaindice{Canti del risveglio}
\voce{221}{Al chiaror del mattin}
\voce{222}{Suona già il risveglio}
\end{multicols}
\endgroup

% ---- Ultima pagina: QR code verso la web app
\clearpage
\thispagestyle{empty}
\noindent
\begin{minipage}[c]{0.38\textwidth}
  \centering
  \includegraphics[width=44mm]{qr.pdf}
\end{minipage}\hfill
\begin{minipage}[c]{0.58\textwidth}
  {\fontsize{16}{19}\selectfont\bfseries Il canzoniere sul telefono\par}
  \vspace{3mm}
  Inquadra il codice con la fotocamera del telefono: si apre il canzoniere
  online, con gli spartiti in PDF (dove disponibili) e la modalit\`a notte
  per il bivacco.\par
  \vspace{4mm}
  {\small\textbf{In caso di problemi con il QR\ldots}\par
  \begin{enumerate}\itemsep=0pt\topsep=1mm
    \item apri il browser del telefono (Chrome, Safari\ldots);
    \item scrivi nella barra degli indirizzi:\\
          {\small\url{https://lc-ncltt.github.io/canzoniere_scout/}}
    \item cerca il canto per numero (quello stampato accanto al titolo);
    \item per riaprirlo senza connessione, aggiungilo alla schermata Home.
  \end{enumerate}}
\end{minipage}

\end{document}
